\chapter{引力实验基础}%引力实验与理论检验

\section{概述}\label{sec:eq-prin}

引力理论的检验标准当然是科学的基础——实验，为此自然需要一个关于引力实验本身的理论，这样就可从实验结果依据实验理论来评判引力理论。Dicke 从 1960s 开始所从事的实验研究使人们对等效原理的理解逐渐深刻，并意识到应把等效原理摆在引力理论的约束这一地位上。

Galileo 性 $m_I=m_G$ 可准确地称为\textbf{弱}(weak)\textbf{等效原理}，简称 WEP。Einstein 将力学实验推至一切物理实验的等效原理则称 \textbf{Einstein 等效原理}，简称 EEP。亦有学者称之\textbf{强}(strong)\textbf{等效原理}，简称 SEP。
不过，Will\cite{Will18} 选择区分 EEP 和 SEP。
他将 EEP 表述为：
\begin{itemize}
    \item WEP 成立；
    \item \textbf{局域 Lorentz 不变性}：\emph{任何局域非引力实验的结果与观者 4-速无关}；
    \item \textbf{局域位置不变性}：\emph{任何局域非引力实验的结果与其时空点本身无关}。
\end{itemize}
注意局域一词仍指空间点。之后又在 SEP 的表述中强调\textbf{自引力系统}\footnote{自引力系统(self-gravitating system)指各部分由引力聚集或能产生引力场的系统，大到恒星、行星，小到日常所见的物体皆于考虑之列。}与测试粒子一样遵守 WEP，且将局域 Lorentz 和位置不变性中的非引力实验换为任意实验。可见这种要求似乎更强，它不仅考虑系统所处的外引力场，还考虑系统内物体所激发的自引力场，即引力的主被动方面皆考虑。
它们对引力理论选择的具体影响如何？目前任何引力理论都满足 WEP，因为验证很早就已开始且精度较高，这件近乎铁打的事实使人很难寻找一个不符 WEP 的理论。可论证若 EEP 成立，则引力一定是时空几何效应，严格满足一个度规理论的基本性质：
\begin{itemize}
    \item \emph{时空上能定义度规}；
    \item \emph{自由的测试粒子的世界线满足该度规下的测地线方程}；
    \item \emph{局域惯性系中非引力物理定律满足狭义相对论}。
\end{itemize}
前文已从 WEP 窥探到 2 的成立。\textbf{Schiff 猜想}(conjecture)认为，凡完备且自洽的满足 WEP 的理论皆必须满足 EEP。这一推论在 $TH\epsilon\mu$ 形式中可证，更一般的耦合则有反例。当然，引力的度规理论是很多的，除了广义相对论，还有 Brans-Dicke 理论。

由于目前已知的度规理论除广义相对论均不满足 SEP，故有人提出：若 SEP 成立，广义相对论可能是唯一选择。关于这一命题的讨论目前还不够严格，故只能算猜想。
目前学术界对等效原理的分类未达成共识，Ohanian\cite{Ohanian} 就批评用自引力系统强行区分 WEP 和 EEP。
对此，论证细节无法说得更通俗了，只有在我们学会后续知识、具体实验后才能有更深理解。
本章讲的涉及等效原理的实验，仍沿用 Will 的划分方法。

设在自由落体电梯中以相同初速释放两个材质不同的重物。略去潮汐，比较两者的加速度是否相同，检验的是 WEP。若在电梯中并排放置两种原子钟，改变电梯的初速或换一处重复实验，看两钟的频率比是否改变，便可检验 EEP 对局域非引力实验的要求。若改用扭秤测量两球间的引力，同样比较在不同速度、不同地点测得的 $G$ 是否相同，检验的便是 SEP 对局域引力实验的要求。

WEP 目前最精确的检验是 21 世纪于太空中完成的，得到钛合金与铂合金\footnote{MICROSCOPE 于 2022 年公布的最终结果为 $\eta(\mathrm{Ti},\mathrm{Pt})=(-1.5\pm2.3\pm1.5)\times10^{-15}$，两项不确定度依次为统计不确定度（$1\sigma$）和系统不确定度。见 P.~Touboul 等，\href{https://doi.org/10.1103/PhysRevLett.129.121102}{Phys. Rev. Lett. \textbf{129}, 121102 (2022)}。}间 $m_G/m_I$ 的差异小于 $10^{-14}$，且数据在地球上空各处基本一致；关于 EEP，对量子非引力实验的描述需用波函数，而波函数至少散布在一个区域内，这样强引力场的潮汐效应会显现在波函数中，但可选择对弱引力场中的小区域波函数进行实验，已有数据显示这种量子系统可满足；SEP 目前最精确的检验利用了地月测距(lunar-laser-ranging)，测得两天体间 $m_G/m_I$ 的差异不大于 $5.5\times 10^{-13}$，引力结合能的贡献差异不大于 $1.3\times 10^{-3}$。

原则上不同材质、不同内部结构的物体，$m_I,m_G$ 有可能不同，这是因为不同物体中质子、中子、电子数目各不相同．即使假定三者等基本粒子的 $m_I,m_G$ 相等，它们结合在一起，还会释放出结合能．而结合能分为化学、核子、电磁的结合能等，对自引力系统还有引力的结合能．基本上非质点物体都存在自引力．目前没有更底层的原理来保证这些结合能对 $m_I,m_G$ 的贡献也相同．这只能用实验检验，特别是检验不同材质、不同内部结构的物体是否满足 $m_I=m_G$．的确有人提出过，为兼容其它相互作用，可适当地放弃 $m_I=m_G$，这基本等同于大幅度修改经典引力理论．目前没有引力理论违背 $m_I=m_G$．

\section{EEP}\label{sec:exp-eep}

以下实验的分类、后 Newton 记号和历史测值主要依据概述所引 Will 专著第 2 至 12 章。各实验的精度对应文中所述观测年代。

回看实验的次序，容易理解为什么本章先讲钟尺，后讲场方程。1959 年的金星雷达回波使行星的距离原则上可以用地面一只钟测定。1960 年前后的塔内 $\gamma$ 射线实验，则能把十几至几十米的高度差化作实验室内的频率比较。这两种精密计时有不同用处：前者沿着穿过太阳系的光路积分，后者拿一个已知势差比较本地的能级。1961 年发表的 Brans-Dicke 理论与广义相对论对后一实验有相同预言，对前一实验却未必相同。仅说某一相对论效应“已经看见”，并不能由一项观测分辨这两个理论。此后从行星回波到航天器相干转发，改进的是光行时中可以分出的微小项。1974 年发现双星脉冲星，又使研究对象从绕日运动的弱自引力天体转向以中子星作钟的系统。2015 年首次直接记录双黑洞并合波形，精密计时的对象成了时空本身的变化。各实验所能回答的问题不同，故须从其实际读数出发。

\subsection{实验预先假定了什么？}

Dicke 在 1964 年的 Les Houches 讲义中从一个比度规理论更弱的起点谈实验。先承认事件可用四维可微流形上的点表示，容许任意坐标来标记它们。起初既不指定哪两点相距多少，也不指定怎样把不同点的矢量作平行比较。前者需要度规，后者需要联络。若一开始就要求一切钟读度规的固有时、一切自由质点走该度规的测地线，再以这些假定处理钟和落体的读数，便已经排除了实验本可看见的一类偏离。

坐标可任意选取，是不是说理论中不能有优选参考系？设某个模型先指定一个宇宙时间 $T$。即使最初只在特别的坐标中写出模型，只需把 $T$ 当作标量场，各同时面写成 $T=\text{常数}$，它在换坐标后仍有完全确定的表示。又如一份平直背景度规 $\eta$，也可用处处满足 $\operatorname{Riem}[\eta]=0$ 的张量场表示。故方程形式上的一般协变性无法决定这样一个时间、这样一份背景究竟存在与否。物理问题是：拿什么仪器可以指出 $T$ 的方向，又或者能否测出相对于 $\eta$ 的运动？如果所用的场能与仪器发生额外耦合，便可能有办法。后面的优选系实验就是在寻找这种读数。

还应区分由场方程决定的量与事先给定的量。做一次太阳附近的实验，给场方程以太阳的质量分布作为源，还须给出边界条件。有些场会随所放的源改变，有些背景结构在理论中却只能取预先指定的形式。写出一个能重现太阳外度规的公式，本身并未说明把太阳换成两颗有内部压强的星时怎样计算，也未说明用另一种钟时该读哪一种时间。要和这一整组实验比较，理论就须同时给出源、引力场和测量器件的动力学。另一方面，同一个理论若既从物质方程预言质心的一种运动，又从场方程和边界条件推出另一种运动，它已无法给实验一个确定答案。这是要求理论完备、自洽的实际含义。

可以从一个对比体会这些要求。设一粒子原先静止在实验室，给定其静止平均寿命 $\tau_0$，再加速至 $v$ 并让它自由飞行。若狭义相对论适用，实验室中的平均飞行距离为 $\gamma v\tau_0$，其中 $\gamma=(1-v^2/c^2)^{-1/2}$。在此处 $\gamma$ 指 Lorentz 因子。大气中产生的高速 $\mu$ 子若没有钟慢效应，按 $\tau_0\simeq2.2\,\mu\mathrm s$ 估计，即便以光速飞行也只有约 $660$ 米的平均程。实验可以把不同高度、不同能量处的通量作比较，检验的仍是物质钟的运动学。反过来，若要求一种引力理论在弱场、低速时给出轨道和扭秤的 Newton 极限，检验的则是引力场怎样由质量产生。在后者中可用轨道标定一次 $GM$，却不能再逐一调整每一种钟的走时率来符合前者。因此要把实验逐步分开。先比较物体和钟对同一外场的响应，再比较在不同物质分布下外场本身如何建立。

一个背景模型首先应能通过哪些无引力检验？Lorentz 变换见 \S\ref{sec:PoincareLorentz}。实验还须分别问：光速是否随发光体速度改变？运动钟是否按同一个 Lorentz 因子减慢？这两件事不是一次静置的光速测量所能同时决定的。设一种发射理论预言沿视线看到的传播速度为 $c+kv_{\rm em}$，$v_{\rm em}$ 是光源沿视线速度。两束由不同运动状态发出的光传播同样距离 $D$，一阶传播时间差约为 $-kD\Delta v_{\rm em}/c^2$。对轨道速度已知的双星，这项差会把观测相位调制成貌似异常的轨道形状。若光已被星际介质吸收、重新发射，却可能失去原来光源的信息。1977 年利用三个银河系内的 X 射线双星，在源速约为 $10^{-3}c$ 的条件下给出 $|k|<2\times10^{-9}$。另一个独立办法是使高速 $\pi^0$ 在已知位置衰变成两个光子，再将选准方向的光子沿地面已知距离作飞行计时。1964 年这类实验以约 $30\,\mathrm m$ 的光路、接近 $c$ 的光源速度，得 $|k|<10^{-4}$。后一界限数值较松，但直接控制了源与路程，不依赖双星轨道或星际介质。

运动钟可以是衰变粒子。若静止平均寿命为 $\tau_0$，恒速样品在实验室时刻 $t$ 剩下的粒子数应为 $N(t)=N(0)\exp[-t/(\gamma\tau_0)]$。求它的海平面通量不能只拿一个“产生高度”除以速度，而要知道取样粒子的能谱和能量损失。Rossi 和 Hall 比较不同高度处选定能量范围的 $\mu$ 子，以此检验时间膨胀。$\mu$ 子的静止平均寿命约 $2.2\,\mu\mathrm s$，相应的 $c\tau_0$ 约 $660\,\mathrm m$。储存环中可用同一批粒子作受控的能量选择。1966 年实验以 $v\simeq0.997c$ 的 $\mu$ 子得到约为静止值十二倍的寿命，与当时的预言在百分之几内相合。取这组数据所检验的是衰变规律按轨道积累的固有时，不是把粒子的向心加速度代入一条静态重力红移公式。电磁谱线、寿命和光行时用了完全不同的发生、接收机制，正因如此，才可共同检查在局域自由落体系中使用同一狭义相对论的依据。

\subsubsection{从观测差额到引力动力学}

水星轨道、日食星像和谱线红移常并称为三大检验。三者在理论选择上实际并不等价。水星是先由历表知道有差额，再问一份场方程能否在已知 $GM_\odot$ 下给出同一差额。太阳附近的光传播是先给预言，再取得变形过的星场同平日星场比较。而一个给定势差内的首阶引力红移只需标准钟与自由落体实验室的局域规律，能够由好几种相互不等价的场方程共同预言。相应的几何计算参见 \hyperlink{subsection.2.2.1}{\S\,2.2.1} 至 \hyperlink{subsection.2.2.3}{\S\,2.2.3}。因此增加同类红移观测的数目，虽可大幅提高钟的普适性检验，却不会自动补上对空间曲率或场的自相互作用所欠的检验。

这一区别在早期观测中颇为实际。太阳光球有上升、下沉的气体和压强引起的谱线移动，其量级可超过待测的引力位移。1917 年太阳光谱中未能可靠分离红移，一度被理解为对理论不利；改用不同谱线或不同日面位置时，气体运动却也随之改变。白矮星的紧致度大得多，但若是亮星的暗伴星，采到的光谱可混入主星的散射光，预言又需知道白矮星的质量和半径。实验室内的核共振源所给的势差极小，反而可使能级、吸收体运动与校准误差成为受控量。正因如此，观测上把一个数做大与把一个理论检验做干净是两件不同的事情。

类似地，早期由膨胀率所得的宇宙年龄曾同当时估计的恒星年龄不合。这里参与对比的有星系距离标尺、恒星演化模型、宇宙的物质模型以及引力场方程；修订天体距离或恒星年龄就能改变差额。发展替代理论时因而不能只给一条背景膨胀曲线，还须回答同一理论怎样解释引力束缚的天体、局域钟尺和光的频率。20 世纪中叶的标量--张量理论使这一要求变得具体：一个理论可以完整给出太阳外度规和宇宙背景，满足 Einstein 等效原理，又与广义相对论预言不同的地月相对运动。Dicke 的实验框架同后来的 PPN 计算，分别解决的是“哪一类物理假设正在被检验”和“这个读数具体依赖哪些系数”两个问题。

还要分清强场与动态。对一个束缚源可用紧致度 $\mathcal C=GM/(Rc^2)$ 估计相对于 Newton 引力的偏离，用几何改变的时间尺度 $T$ 与光穿过系统的时间 $R/c$ 相比，判断它是否来得及明显演化。一个静止黑洞在视界附近 $\mathcal C$ 为一的量级，但若没有扰动，仪器可以只测其定常几何。双中子星在轨道间距上的 $Gm/(ac^2)$ 可很小，星面紧致度却约为 $0.1$，脉冲计时同时测到强自引力和缓慢轨道运动。到了并合，轨道尺度处的场与变化速率都增大，才可在同一信号中检验强场动力学。潮汐仪器真正测到的则是曲率而非 $\mathcal C$ 本身；两个质量极不相同的黑洞在各自相同倍数的引力半径处有相近的 $\mathcal C$，曲率量级却按质量的平方反比改变。说明测了哪一种“强”，比只说一个实验处于强引力场更能判明它是否探到了新效应。

\subsection{物体的组成与自由落体}

若整个实验室连同仪器一同下落，怎样觉察两种材料的 $m_G/m_I$ 略有不同？设 $a_1,a_2$ 是两个试验体沿同一外场方向的加速度，定义有符号的\textbf{Eötvös 参数} $\eta_{12}=2(a_1-a_2)/(a_1+a_2)$。只要两物体处于同一外场，即使不知道场强也能比较这一比值。更严格地说，须把两条轨道化到同一初始事件、同一速度再比。例如在地面释放先后错开的两物体，靠近地面的一者自会落得略快。若初速不同，地球自转还会给出不同的 Coriolis 加速度。这些差别用同材料的试验体也会发生。

设第 $i$ 个物体惯性质量为 $m_i$，内能中属于某种相互作用的一项占其静能的份额 $f_i^A=E_i^A/(m_ic^2)$，其中能量及其正负号按同一约定计算。若该项对被动引力质量的贡献是对惯性质量贡献的 $1+\delta_A$ 倍，则 $a_i/g=1+\sum_A\delta_A f_i^A$。只保留异常系数的一次项，得 $\eta_{12}=\sum_A \delta_A(f_1^A-f_2^A)$。故不能把 $|\eta_{12}|$ 的上限直接当作任一种结合能的异常上限。两材料中占比相同的成分在差中消失。各种占比虽不同，其贡献也可能恰好相消。选取第三种材料能否分开两种贡献，要看新增的一组占比差与原来那组是否成比例。

以静电能为例。半径 $R$ 的均匀带电球电场能正比于 $Z^2/R$，其中 $Z$ 为质子数。原子核的半径近似正比于质量数 $A$ 的立方根，去掉质子的自作用，再除以 $A$ 个核子的静能，即得核静电能份额近似为 $f^{\rm ES}\simeq7.5\times10^{-4}Z(Z-1)/A^{4/3}$。铍和铝的典型核种分别有 $(Z,A)=(4,9),(13,27)$，代入得 $f^{\rm ES}$ 约为 $4.8\times10^{-4}$ 和 $1.4\times10^{-3}$。若假定只有静电能耦合异常，测出 $|\eta|<10^{-13}$ 意味着 $|\delta_{\rm ES}|$ 在 $10^{-10}$ 左右或更小。同一个落体误差经不同的内能份额换算，限度可差许多数量级。核内弱相互作用能的占比很小，又在材料之间相差甚微，因而很精确的落体数据也未必给出同样精确的弱相互作用能界限。

如何在地面取得这样小的加速度差？对单摆，两种材料制成的摆锤即使长度 $l$ 相同，若 $m_G/m_I$ 有别，周期 $P=2\pi\sqrt{lm_I/(gm_G)}$ 也应有别。同样的周期差却也可由悬点移动、摆长变化产生。扭秤把问题改成了找一个本应消失的力矩。Eötvös 把不同材料的重物装在水平杆两端，用细丝悬住杆。若引力恰沿悬丝，即使两物体落得不一样快，也不能使杆绕丝转动。地球自转使铅垂线沿地球引力与离心加速度的合方向，除在赤道与两极等处外，它和地球引力的方向有小偏角。于是地球引力在垂直悬丝的平面里留下一分量 $g_\perp$。交换两种材料的位置，这一分量产生的差动力矩应反号。

设两端惯性质量同为 $m$，臂长各为 $l$，杆与场的横向投影成角 $\theta$。一阶近似下绕丝的力矩大小为 $ml\,|\eta_{12}g_\perp\sin\theta|$。如图~\ref{fig:exp-torsion}，慢慢转动仪器，扭角就有一个随方位绕行一周而重复的分量。若转得不能视为静态，还须用悬丝的转动惯量、阻尼和扭转常数从角度反算力矩。真正关心的是在已知转频上、相位指向场源且随材料对调反号的那个量。Dicke 等人改用太阳作引力源，由地球自转来改变杆相对太阳的方位，便在一个太阳日的周期上找信号。太阳同时加速地球和扭秤。对随地球运动的仪器，两者共同的加速度相消，可能留下的仍是不同材料的响应差。

\begin{figure}[htbp]
    \centering
    \begin{tikzpicture}[x=1cm,y=1cm,>=stealth,font=\small,line cap=round]
        \begin{scope}
            \node at (0.25,1.42) {扭秤俯视};
            \draw[line width=1.1pt] (-1.22,0.79) -- (1.22,-0.79);
            \filldraw[fill=white,line width=.75pt] (-1.22,0.79) circle (.105);
            \fill (1.22,-0.79) circle (.105);
            \node[above left,inner sep=2pt] at (-1.22,0.79) {材料 2};
            \node[below right,inner sep=2pt] at (1.22,-0.79) {材料 1};
            \node[above right,inner sep=1pt] at (-.84,.55) {$l$};
            \draw[->] (0,0) -- (1.90,0) node[right,inner sep=2pt] {$g_\perp$};
            \draw[->] (-33:.68) arc[start angle=-33,end angle=0,radius=.68];
            \node at (.91,-.27) {$\theta$};
            \filldraw[fill=white,line width=.65pt] (0,0) circle (.095);
            \fill (0,0) circle (.028);
            \draw[thin] (-.07,-.08) -- (-.42,-.46) -- (-.71,-.46);
            \node[left,inner sep=2pt] at (-.71,-.46) {悬丝};
        \end{scope}
        \begin{scope}[shift={(3.12,0)}]
            \node at (1.44,1.42) {绕悬丝的力矩};
            \draw[->] (0,-.82) -- (0,.91) node[above,inner sep=1pt] {$\tau_z$};
            \draw[->] (-.12,0) -- (3.30,0) node[right,inner sep=2pt] {$\theta$};
            \draw[line width=.85pt,domain=0:360,samples=72,smooth,variable=\t]
                plot ({3*\t/360},{.62*sin(\t)});
            \draw[dashed,line width=.65pt,domain=0:360,samples=72,smooth,variable=\t]
                plot ({3*\t/360},{-.62*sin(\t)});
            \node[below left,inner sep=1pt] at (0,0) {$0$};
            \node[below,fill=white,inner sep=.5pt] at (1.5,-.16) {$\pi$};
            \node[below,inner sep=1pt] at (3.04,-.08) {$2\pi$};
            \draw[line width=.85pt] (.18,-1.07) -- (.65,-1.07);
            \node[right,inner sep=2pt] at (.65,-1.07) {原位};
            \draw[dashed,line width=.65pt] (1.72,-1.07) -- (2.19,-1.07);
            \node[right,inner sep=2pt] at (2.19,-1.07) {对调};
        \end{scope}
    \end{tikzpicture}
    \caption{Eötvös 扭秤。两端惯性质量相同。由原来指向材料 1 的杆臂逆时针转至场的横向投影，所转的角为 $\theta$。取 $\eta_{12}>0$，交换材料后仍用原杆臂计角。}
    \label{fig:exp-torsion}
\end{figure}

附近的墙、山体，乃至移动的器材，都会在秤的两端产生不一样大的普通引力。以杆长乘以场的梯度就给出这种差加速度的最低阶估计。简单交换两块质量相同而形状不同的物体，未必能使它消失。故精密扭秤由多块试验体组成，安放位置使整个悬摆对外场梯度敏感的低阶非球对称质量多极矩尽量小，同时保留由材料分布给出的差动力矩。也可把附加的已知质量移到扭秤附近，故意放大某一梯度，以测定剩余的耦合。华盛顿大学的 Eöt-Wash 实验以这些办法比较铍、钛等材料，把 $|\eta|$ 压到约 $2\times10^{-13}$ 的量级。报告这个数时还应说明加速度朝向哪一场源：地球、太阳和人为移动的场源，既有不同的组成，也与仪器相隔不同的距离。

卫星把自由落体的时间从几秒延长到许多轨道周期。MICROSCOPE 用静电力使一对同轴圆筒相对卫星保持静止。两圆筒在测量轴上的运动相同，扣除干扰后，静电控制加速度之差恰与所求的落体加速度之差反号。一筒用铂铑合金，另一筒用钛铝钒合金，另备一对同为铂铑合金的圆筒检查干扰。若前一对的中心相差 $\Delta x^j$，即使 $\eta=0$，外场梯度也会给出 $\partial_jg^i\,\Delta x^j$ 的差加速度。它随卫星姿态的变化可用来估计两中心的偏离。这和差加速度直接随地心方向的投影改变是可区别的。在轨高度约 $720\,\mathrm{km}$，当地重力加速度仍约 $8\,\mathrm{m/s^2}$。以 $10^{-15}$ 的相对量级作比较，所求的差加速度只有约 $10^{-14}\,\mathrm{m/s^2}$。2017 年首批异质材料数据给出 $10^{-14}$ 量级的限度，2022 年最终结果及其统计、系统不确定度见本章概述。

还可用原子替代宏观试验体。让短激光脉冲分开、反射和重新合并自由下落的原子波包，三个脉冲等间隔发出，间隔为 $T$。在均匀加速度 $a$ 下，三个时刻的激光相位组合含有位置的二阶差分 $z(0)-2z(T)+z(2T)=aT^2$，故干涉相位中的加速度项是 $k_{\rm eff}aT^2$，其中 $k_{\rm eff}$ 是沿测量轴的有效波数。比较两种同位素，两套波数、脉冲间隔与初始位置均须分别校准。共用激光反射镜有助于抵消镜面振动，却不能自动消去原子云初始位置、速度不同所致的梯度项。这里比较的是从干涉条纹读出的自由落体加速度。如欲作钟速实验，须另指明哪两个内态充作钟、怎样读出它们的相对相位。

\subsubsection{从材料组合反推耦合}

再问一步，表中所谓“强相互作用能的上限”是怎样来的？惯性质量可以直接称量，核内每一项能量的期望值却须依靠核模型。半经验核质量公式里，同一种单核子静能对不同材料的贡献几乎都是一，故灵敏的部分多半来自结合能。短程核力在核内饱和，体能项近似随 $A$ 增长，除以质量后几乎不随材料变。表面、同位旋不对称和核子配对等修正的 $A,Z$ 依赖各不相同，正是它们使核能的占比差不为零。静磁能又与静电能不同。满壳层的核子电流可以互相抵消，不能用带电均匀球的静电公式代替它。比如用壳模型估得铝、铂中核静磁能占总静能分别约 $4.1\times10^{-7},2.4\times10^{-7}$。若只靠这对试样寻找静磁能异常，有效杠杆便是两数之差 $1.7\times10^{-7}$。这种核结构估计的不确定性不在扭角测量的统计误差内。

强、弱相互作用还有一个易忽略的区别。对具有确定宇称的核基态，弱相互作用中宇称为奇的部分，一阶能量修正是本征态中一个奇算符的期望值，因而为零。弱作用中保持宇称的部分仍可贡献静能。用低能弱相互作用的估计，它在总静能中的份额只有约 $10^{-8}$ 或更小，且随中子、质子比例改变得缓慢。故 Eöt-Wash 的铍--钛数据在每次只容许一个系数非零的假设下，给出的强相互作用、静电、超精细和弱相互作用异常界限依次约为 $4.9\times10^{-11},1.3\times10^{-10},2.3\times10^{-7},2.0\times10^{-2}$。测量材料换成铍--铝，静电项的界限变为 $2.1\times10^{-11}$，弱作用项却仅改至 $1.5\times10^{-2}$。同样是重测一次落体，并不能使各类结合能的界限都按同一倍数改善。

若希望同时检验两个耦合，而不假定其中一个已知为零，应怎样处理？把第 $k$ 对材料的测值写作 $y_k=D_{k1}\delta_1+D_{k2}\delta_2$，其中 $D_{kA}=f^A_{k,1}-f^A_{k,2}$。两行数据若满足 $D_{11}D_{22}=D_{12}D_{21}$，则无论误差多小，仍只确定一个线性组合。其比值十分接近时，耦合本身的误差可能远大于各 $y_k$ 的测量误差所暗示的值。还可利用另一种办法：同一对材料向不同场源下落，若异常来自可区分中子数和质子数的长程场，场源的组成也进入该矩阵。在地面以人为移动的近源取得的零结果，与让同样试样朝向太阳取得的零结果，便有不同的解释。类似地，地月测距给的是整体天体的加速度差。用成分相似的地面模型先限制普通结合能的贡献，才可进一步用地月测距检查自引力。

原子干涉实验所称的 $a$ 也不是从一幅静止条纹直接读出的。若反射镜加速度为 $a_m$，自由演化时激光脉冲在原子处印下的是镜、原子相对位置的相位，一阶有 $\varphi=k_{\rm eff}(a-a_m)T^2$。对两种原子，先分别除以 $k_{{\rm eff},i}T_i^2$，再作差才是差加速度。只减相位而不管比例因子，会使本应共同的镜面振动泄入结果。在线性梯度 $a(z)=a_0+\Gamma z$ 中，各波包中心若起初相差 $\Delta z_0,\Delta v_0$，则差结果还会包含与 $\Gamma\Delta z_0,\Gamma T\Delta v_0$ 成正比的项。这与宏观圆筒中心不重合的误差有同一来源。改变脉冲间隔可分辨某些项，但若 $\Delta z_0$ 恒定，它与均匀场中恒定的材料异常具有同样的领先 $T^2$ 依赖，还须独立确定初始位置或改变梯度。

\subsection{有限力程与第五种力}

同一份差分数据还可用来寻找新相互作用。设质量为 $M$ 的点源带有一种荷 $Q$，试验体质量为 $m_i$、荷为 $q_i$。把两荷分别归一化为无量纲的单位质量荷 $s_S,s_i$，在 Newton 势能上附加 Yukawa 型一项，可写为 $V_i=-GMm_i[1+\alpha_5s_Ss_i e^{-r/\lambda}]/r$。这里 $\lambda$ 是力程，$\alpha_5$ 的定义随荷的归一化而定。求径向导数，以指向场源为加速度的正向，则一阶下 $\eta_{12}(r)\simeq \alpha_5 s_S(s_1-s_2) (1+r/\lambda)e^{-r/\lambda}$。若新力按重子数耦合，原子质量因结合能而不严格正比于重子数，两种电中性材料仍可有不同的 $s_i$。同一材料对着成分不同的场源也会有不同响应。若反过来取 $s_i$ 对一切材料相等，差分落体便完全失去灵敏度。但在 $r\sim\lambda$ 处，合力对距离的依赖仍偏离平方反比律，两种实验所问的事情便分开了。

在 $r\ll\lambda$ 处，几何因子 $(1+r/\lambda)e^{-r/\lambda}$ 近于一。若荷质比又是普适常数，单在这一尺度上测引力，只会把新作用计入测得的 $G$。在 $r\gg\lambda$ 处，信号按指数衰减。真实场源还有形状，例如一段厚度远大于 $\lambda$ 的山壁，只有面向仪器的一薄层对新力有可观贡献，不能把整座山的荷都集中在其质心来算。为特定力程解释实验时，应将上述势对场源、试验体两者的分布积分，并算出沿测量轴或绕悬丝的分量。

短距离实验常把相同形状的孔阵列分别做在扭摆和旋转的吸引盘上。若一圈有 $N$ 个等间隔孔，吸引盘每转一周，孔阵列给出的力矩重复 $N$ 次。再在盘下加一层错开的孔阵列，使它的 Newton 力矩抵消上层的。若新力的力程比两层厚度小，下层对新力的贡献弱，抵消便不完全。改变盘与摆的间距，可寻找该转频谐波是否按有限力程衰减。盘和摆之间还须有静电屏蔽，并检查表面电势和残余磁性。单次在小距离测到一个吸力，不能凭它断定有额外引力。

\subsubsection{把引力势同实验的距离尺度相比较}

对已知几何形状，有限力程的积分可很简单。设一层横向远大于 $\lambda$ 的均匀物质平板占 $-d<z'<0$，其密度为 $\rho_S$、单位质量荷为 $s_S$。试验体在 $z>0$。先对每一薄层的柱面坐标半径 $b$ 积分，由 $R^2=b^2+(z-z')^2$ 有 $b\,\d b=R\,\d R$，故 $\int_0^\infty e^{-R/\lambda}b\,\d b/R=\lambda e^{-(z-z')/\lambda}$。再积厚度，得到每单位试验质量的额外势和指向平板的加速度大小
\[
 \delta\phi(z)=-2\pi G\alpha_5\rho_S s_Ss_i\lambda^2
     e^{-z/\lambda}(1-e^{-d/\lambda}),\qquad
 |\delta a_z|=2\pi G|\alpha_5\rho_S s_Ss_i|\lambda
     e^{-z/\lambda}(1-e^{-d/\lambda}).
\]
若 $d\ll\lambda$，信号近似随面密度 $\rho_Sd$ 增加。若 $d\gg\lambda$，增厚远离试验体的一面几乎无效。无限平板的 Newton 引力在外部不依赖 $z$，短程力却随间距指数衰减，故扫动间距比仅增大场源质量更能分辨两者。实际屏蔽膜虽能抑制静电力，仍有自身的引力，且占去了两测试面之间的一段距离。当 $\lambda$ 比这段距离还小时，力程灵敏度会很快丧失。

平方反比律还可在比实验室大得多的尺度检查。在理想球形星球外部，若相互作用严格为 $1/r$ 势，则 $g(r)r^2$ 与半径无关。比较地面、矿井和高塔的重力时却不能照搬这个等式：仪器可能在物质内部，山体、地下密度及地表水的分布都会改变 $g$。因此在某处测得与简单地球模型不合的重力垂向剖面，不足以给第五种力指定力程。轨道方法有另一种退化。若只测近圆卫星的一条轨道，引力偏离可被吸入用 $n^2r^3$ 定出的 $GM$。以同一颗中心天体在不同半径处的轨道或独立测距资料作比，才能不把同一个 $GM$ 当作每条轨道各自的自由量。20 世纪后期中程第五种力搜索在约 $1$ 至 $10^4\,\mathrm m$ 力程范围内给出的相对强度上限，随实验与模型约从 $10^{-3}$ 到 $10^{-6}$，这些数并非任意力程通用的一个界限。

航天器导航里也曾出现过很像新长程力的信号。Pioneer 10、11 远离太阳时，实测 Doppler 频移与当时的轨道及仪器模型有缓慢的差，解释成沿视线加速度是 $10^{-9}\,\mathrm{m/s^2}$ 量级。倘若异常真来自普适的太阳引力，它还应作用于其它在相近轨道的物体。另一方面，航天器电源和舱室放热不对称，净辐射功率 $P_{\rm an}$ 携出的动量每秒约为 $P_{\rm an}/c$，质量 $m$ 的飞行器自会受到 $P_{\rm an}/(mc)$ 的反冲。对 $250\,\mathrm{kg}$ 的器材，数十瓦的定向热辐射已足以给出上述量级。后来的遥测、热设计和导航资料支持这种热反冲解释。这个例子说明，偏离一套参考历表的残差还不是对某一引力参数的测量，须先问残差是共用于所有试验体，还是只随一台设备的姿态、功率而变。

\subsection{局域 Lorentz 不变性}

在局域实验室，把光学谐振腔的轴转过直角，或保持轴向而改变整个实验室的速度，若能测出差异，就可能从内部辨出相对某一优选系的运动。这里比较的对象须明确。给定一根杆只说它变短多少，尚须指明拿什么和它比较。说某束光变慢，也得交代以何种尺、钟测量。能够从实验读出的是无量纲的频率比、相位或到达时间与参考周期之比。

先看腔。真空腔长 $L$，忽略镜面相移，阶数 $n$ 的纵模满足 $\nu_n=n\bar c/(2L)$，这里 $\bar c$ 是沿该轴往返的平均光速。可使两台激光分别锁到两个相互垂直的腔，再使它们干涉，以拍频比较两个共振频率。若 $\bar c$ 沿平面内某方向有微小的二阶各向异性，可写 $\bar c(\theta)=c[1+b\cos 2(\theta-\theta_0)]$。两个相同长度的腔一减，一阶下便有 $(\nu_1-\nu_2)/\nu\simeq2b\cos2(\theta-\theta_0)$。腔转半周时两镜虽已互换，往返光路却与原来相同，所以应找两倍转频。若腔长也随方位改变，测到的系数还要减去相应的 $\Delta L/L$。因此把拍频完全解释成光速各向异性，本身就含有对造尺物质的假设。

可用一个简单模型把这项假设展示出来。取物质粒子极限速度 $c_0$ 为速度单位，在假想优选系中令真空光速为 $C$，速度均以 $c_0$ 为单位。假设杆的平行长度按物质的 Lorentz 因子收缩，静长为 $L$ 的杆以速度 $w\ll1$ 平移时，平行长度为 $L(1-w^2/2)$。在优选系中，光走过平行臂再返回所需的时间为 $t_\parallel=L(1-w^2/2)[(C-w)^{-1}+(C+w)^{-1}]$。在横臂中，光的横向分速度为 $\sqrt{C^2-w^2}$，故 $t_\perp=2L/\sqrt{C^2-w^2}$。于是 $t_\parallel-t_\perp= \frac{L}{C}(C^{-2}-1)w^2+O(w^4)$。换算到随仪器走的钟，因右端已是 $w^2$ 小量，在此阶不改结果。若 $C=1$，两项抵消，这就是 Michelson--Morley 零结果在该模型中的意义。若两臂本来不等长，即使方向不变，干涉条纹还可随实验室公转速度改变，这提供了另一种比较。以上全程只用往返光行时。若声称测量单程光速，就须另说明远处的钟如何同步。

原子核也能作方向的探针。Hughes 等人与 Drever 分别用核磁共振辨认很小的能级各向异性。以锂核为例，外磁场将 $^7\mathrm{Li}$ 自旋为 $3/2$ 的基态分为 $m=-3/2,-1/2,1/2,3/2$ 四级。普通 Zeeman 能量正比于 $mB$，三个相邻跃迁的频率应相等。若另有相对于星空固定的四极扰动，沿磁场量子化轴的能移可含 $[3m^2-J(J+1)]$ 一项。于是 $|m|=3/2$ 的两级一起移动，$|m|=1/2$ 的两级向另一方向移动，原来的一条共振线便可分开。地球自转使量子化轴与假想优选方向之间的夹角改变，从而还能用恒星日的周期认出它。早期实验对额外分裂给出的上限为 $0.04\,\mathrm{Hz}$，相当于 $h\Delta\nu\simeq1.7\times10^{-16}\,\mathrm{eV}$。比较两种共处一地的原子、自旋或腔钟还有好处：它们若都同样响应磁场漂移，可用两种频率的适当组合消去这项干扰，同时保留对不同粒子成分的响应。

怎样将腔实验与核实验放在同一理论中？\textbf{标准模型扩展}(Standard Model Extension, SME)在标准模型各物质场和规范场的作用量中加入可与背景张量缩并的项。以平直背景中保持 CPT 的电动力学为例，取 $c=\hbar=1$，光子部分可在 $-F_{\mu\nu}F^{\mu\nu}/4$ 上加
$-(k_F)^{\mu\nu\alpha\beta}F_{\mu\nu}F_{\alpha\beta}/4$，电子部分则可在动能项中加入与 $c^{\mu\nu}\bar\psi\gamma_\mu \mathrm{i}D_\nu\psi$ 成正比的项及其 厄米共轭。前者改变场的传播，后者也能经由原子间距、跃迁能改变尺和钟。张量 $k_F$ 可取 Riemann 张量的指标对称性，并去掉一项只会重定义普通 Maxwell 动能的双迹，留有 $19$ 个独立分量。此为光子这一组算符的数目，不能把 $19$ 个未知数都用一个拍频上限分别限制。

改变坐标时，场和背景张量一起变换，上述作用量仍是标量。实际旋转仪器时，旋转的是实验物质，相对于外部背景的方向改变。这一区别使得把一切方程写成张量形式并不能保证 Lorentz 实验有零结果。可把各背景分量先给在一套相对于恒星近似固定的坐标中，再用已知的台站纬度、自转相位和转台角度把它们投影到实验轴上。对一对互相垂直的腔，最低阶的差响应会挑出形如 $K_{ij}(n_1^in_1^j-n_2^in_2^j)$ 的组合，其中 $\bm n_1,\bm n_2$ 为两轴方向。另一个纬度、不同材料的腔或不同粒子的跃迁通常挑出另一组合。先列出这些组合，再谈各系数可达到何种上限，才不会把单一仪器的零结果解释得过强。

\subsubsection{同一实验在优选系和仪器系中的描述}

取物质极限速度为一，在假想优选系中沿 $x$ 轴放置实验室的速度 $w$。按由物质尺钟给出的 Lorentz 变换，光线四维波矢的分量变为 $\omega'=\gamma_w(\omega-wk_x)$、$k'_x=\gamma_w(k_x-w\omega)$。优选系中 $\omega=C|\bm k|$，并不意味着对 $C\ne1$ 在实验室也还有 $\omega'=C|\bm k'|$。譬如正反向光脉冲在仪器系的坐标速率大小分别为 $(C-w)/(1-wC)$ 与 $(C+w)/(1+wC)$。一阶展开可见两速率各有与 $(C^2-1)w$ 成正比而反号的改正。在同一路径上相继往返，对这类改正的一次项恰好不灵敏，于是前述干涉仪的主要效应落在二次速度上。若以非标准方式规定分立时钟的同时性，单程速度的表达会改变。在同一分束器重合的两波相位差却不随这种规定改变。

有些实验用运动中的原子束或转子代替一根干涉臂。在理想横向 Doppler 实验中，发射体速度与视线正交，一阶纵向项消失，所检验的是运动钟速的二阶项。交替用相反方向的光照射粒子束并选取同一类粒子，也可用频率的乘积消去一阶速度依赖。以普通狭义相对论为例，一条本征频率为 $\nu_0$ 的谱线，对严格共线的相反两束光有 $\nu_\pm=\nu_0/[\gamma_v(1\mp v/c)]$，故 $\nu_+\nu_-=\nu_0^2$。如果两束光各自选到速度不同的原子，乘积就不再消去一阶项。若以优选速度模型解释这些测量，必须在同一模型中计算吸收原子和实验室参考钟的频率，不能只修改传播中的光。比较不同实验前，还须说明给出的限度是对仪器系中的直接各向异性，还是除去了实验室相对优选系速度因子的基本理论系数。

Hughes--Drever 型谱线限度尤其容易在这里被误读。设一种结合能 $E_A$ 在核内的各向异性部分为 $\delta_A E_A$，可测的能移当然与核结构有关。若各向异性来自一个在宇宙静系中原本各向同性的背景，实验室以约 $370\,\mathrm{km/s}$ 运动才产生一个方向，则二阶张量的无迹部分正比于 $w_iw_j-\delta_{ij}w^2/3$，$w^2/c^2$ 只有约 $1.5\times10^{-6}$。核能异向的限制与对优选系基本常数的限制可以差这一个因子。实验还可用冷原子或离子陷阱延长自由进动时间，减小碰撞使谱线展宽的影响。但长期观测更换了限制因素，地磁场的日变化可能与待测周期接近。只有按恒星日而不仅仅按太阳日积累相位，并用另一种对磁场敏感性已知的钟作共同参照，这个时间优势才用得上。

一个单程例子是以光纤连接的两台氢微波激射钟。两端相距约 $21$ 公里，在初始时刻校定相对相位，再随地球自转比较多日的相位变化。若 $\bm n,\bm n_0$ 分别为测量时和初始时的传播方向，实验室相对于宇宙背景的速度为 $\bm w$，则在本章所用的 $c^2$ 模型和该链路的响应下，最低阶的异常有 $\Delta\varphi/\varphi\simeq(4/3)\Upsilon_0(\bm w/c)\cdot(\bm n-\bm n_0)$，$\varphi=2\pi\nu L/c$ 为基准传播相位。故静止的时钟零点差不会模仿这个方向变化，而两钟的日变化、链路温度和仪器延迟必须单独监测。当时 JPL 实验的日变上限折为 $|\Upsilon_0|<3\times10^{-4}$。这是对有钟、有传输媒介的整个实验模型作的解释，不能把 $4/3$ 搬给一切在真空中传播的光子。与之互补的 Mössbauer 转子实验把发射体置于轮缘、吸收体置于轮心，改变轮缘速度与外部速度的相对方位，在很窄的核共振线上找钟慢的方向依赖；双光子吸收则以窄的原子跃迁线读出速度组合。它们不需同一条长光纤，所要扣除的载体形变和频移亦不同。

\subsubsection{SME 中哪些系数组合可以测？}

把 $k_F$ 的时间、空间指标分开，电磁拉氏量可等价地写成
\[
 \mathcal L_\gamma=\tfrac12\bm E\cdot(1+\kappa_{DE})\bm E
 -\tfrac12\bm B\cdot(1+\kappa_{HB})\bm B+\bm E\cdot\kappa_{DB}\bm B .
\]
单位取 $c=1$。前两个矩阵各自对称，双迹条件使 $\operatorname{tr}(\kappa_{DE}+\kappa_{HB})=0$，混合矩阵 $\kappa_{DB}$ 无迹。可以不必记四指标张量的每个分量，直接对这个拉氏量分别按 $\bm E,\bm B$ 求偏导，得到置换场 $\bm D=(1+\kappa_{DE})\bm E+\kappa_{DB}\bm B$ 与磁场强度 $\bm H=(1+\kappa_{HB})\bm B-\kappa_{DB}^{\mathsf T}\bm E$。这两式只借用介质电动力学的记号，物质空间仍是真空。用普通无源 Maxwell 方程消去 $\bm B$，对给定波矢 $\bm k$，电磁波两个横偏振组成一个二阶本征值问题。若两本征频率不同，同一条光路就能积累可测的偏振相位差。这与两垂直腔比较不同的传播方向是两种观测。

令 $\kappa_{HE}=-\kappa_{DB}^{\mathsf T}$。通常的分解是将 $\kappa_{DE},\kappa_{HB}$ 的半和、半差分别归入宇称偶的 $\tilde\kappa_{e+},\tilde\kappa_{e-}$，后者还须扣去迹 $\tilde\kappa_{\rm tr}=\operatorname{tr}(\kappa_{DE})/3$。将 $\kappa_{DB},\kappa_{HE}$ 的半和、半差分别归入宇称奇的 $\tilde\kappa_{o+},\tilde\kappa_{o-}$。于是 $e+,e-,o-$ 各有五个独立分量，$o+$ 反对称而只有三个，最后加上一个标量，合为十九。线性近似下，使两偏振传播速度分开的系数在 $e+,o-$ 两组中，而一对普通往返谐振腔对 $e-$ 等不引起双折射的方向项仍很灵敏。二十世纪末至本世纪初的腔实验对一部分 $e-$ 系数组合给出约 $10^{-14}$ 的量级，对一部分 $o+$ 组合约 $10^{-10}$。两数的差别部分来自实验运动速度所带的抑制。若强令一个未知系数以外的其余分量全为零，可报出单个系数的上限。容许它们一起变化时，实验只能给出实际拟合得到的组合。

费米子部门也有类似的选择。取参考电子的各向同性 $c_{00}$ 及光子部门 $\tilde\kappa_{\rm tr}$ 为例，若只看自由电子和光子，匀常地重选时空单位可以把电子动能恢复成普通 Dirac 形式，却会同时改变光子部分的 $\bm E^2,\bm B^2$ 系数。在电子的极限速度取一的单位中，由两个系数的比算出光速 $C$，一阶有
$1-C^{-2}=-2(\tilde\kappa_{\rm tr}-4c_{00}/3)$。所以观测到电子、光子具有同样的极限速度，只给这一个组合的约束。原子核含质子、中子，体系内的费米子系数也未必等于电子系数。比较两种原子的自旋进动，若把假想的自旋依赖能移写作 $-\bm b\cdot\bm\sigma$，沿外磁场的两个自旋投影给出的裂距正比于 $2\bm b\cdot\hat{\bm B}$。换一种核种，可同时改变质子和中子自旋的权重。若这两个权重靠核模型给出，则应与谱线误差分开报告。某些钟比较实验曾把这类能量组合限制到 $10^{-27}$ 至 $10^{-32}\,\mathrm{GeV}$ 量级。这是带能量量纲的自旋耦合限度，不能与腔实验中无量纲的 $k_F$ 数值直接作大小比较。

天文光的传播还能检验能量依赖。设只多一个色散项 $E^2=p^2c^2[1+f_n(pc/E_{\rm Pl})^{n-2}]$，其中 $E_{\rm Pl}$ 为 Planck 能量。群速 $\mathrm dE/\mathrm dp\simeq c[1+(n-1)f_n(E/E_{\rm Pl})^{n-2}/2]$，要寻找两种能量脉冲之间的相对传播延迟，须有 $n\ne2$ 且相应系数不使群速修正消失。若爆发源先发低能光再发高能光，发射时差也会进入同一个观测量。可比较不同距离的多次事件以利用距离依赖，又可改看同一源的偏振：两圆偏振的波数若差 $\Delta k$，线偏振面的转角便为 $\Delta kD/2$。若该角在探测器带宽内已扫过许多周期，收到的带宽平均线偏振度应下降。观察到高偏振度的遥远源因而能约束双折射，却无法用这项资料排除两偏振同速但同时有色散的模型。

\subsection{两种引力红移实验}

局域位置不变性还涉及钟在何处、何时走得一样快。怎样比较相距很远的两只钟？让静止的发射者连续发出由钟规定的波峰，静止的接收者也以本地同种钟计数。在静态场中，相继两个波峰的坐标传播时间相等，故收发两端相邻波峰的坐标时差相等，两端测得的频率之比由 $\d\tau/\d t$ 之比给出。沿用 \eqref{eq:h-g-phi} 的 Newton 势 $\phi$（即下文的 $-U$），弱场下静止钟的 $\d\tau/\d t\simeq1+\phi/c^2$。若把钟 $A$ 固定在塔底，以固有频率 $\nu_0$ 发光，塔顶用同型钟比较收到的光，写 $\Delta\phi=\phi_{\rm rec}-\phi_{\rm em}>0$，则可用
$\nu_{\rm rec}/\nu_0-1=-(1+\alpha_A)\Delta\phi/c^2$ 表示实测频移。\hyperlink{subsection.2.2.1}{\S\,2.2.1} 算出的红移相当于 $\alpha_A=0$，一切理想钟都应给出此值。这一实验只需比较钟速，并不要求收发两端事先对准钟面上的时刻。

塔高几十米，$\Delta\phi/c^2\simeq gh/c^2$ 不过 $10^{-15}$ 量级，怎样分辨出这样小的频移？对自由原子产生的窄谱线，发射时的反冲及原子的热运动都会妨碍相同跃迁的共振吸收。Pound 和 Rebka 采用 $^{57}\mathrm{Fe}$ 的 $\gamma$ 射线，使发射、吸收有一部分在晶格不激发声子的情况下进行。动量可由整块晶体承担，这部分射线的反冲能可忽略，因而保有足够窄的共振线。若上下两端固定，抵达时的频率已略偏离吸收线中心。给源或吸收体一个可控的视向速度 $v$，使 $v/c$ 抵消引力频移。由恢复共振所需的速度便知红移。交换传播方向，引力项随之反号，还可区别某些只随源、吸收体制备条件改变的线心差。在此后一系列改进实验中，红移系数与一相符至约百分之一。

增加高差可以增大引力频移，却会同时引入运动频移。1976 年 Gravity Probe A 用火箭把氢微波激射钟送至约 $10^4\,\mathrm{km}$ 高空，再与地面同种钟作比。只听星上发来的一个频率是否足够？令 $R(t)$ 为站星距离，约定远离时 $\dot R>0$。下行信号的一阶相对频移中有 $-\dot R/c$，远大于待测的引力项。不妨另由地面发出载波，由星上按已知的固定频率比相干转发。用地面的同一时标比较回波时，一阶 Doppler 项近于 $-2\dot R/c$，不含独立星钟的钟速差。把转发比归一化后，用单程频移减去一半往返频移，就消去了这一大项。实际比较还要使上下行的测量时刻对应，火箭在一段光行时间内并非匀速，不能对任意两个同时收到的频率机械相减。

若暂略去传播介质和仪器误差，运动钟的钟速为 $\d\tau/\d t\simeq1+\phi/c^2-v^2/(2c^2)$。因而去掉一阶 Doppler 项后，仍须由跟踪资料求星上与地面钟速度平方之差，并分别算出两端的势。引力异常改变的是随势差变化的那个系数。火箭上升、到达顶点再下降，在一次飞行中就扫过许多不同势差。造钟时留下的常频偏可同这个随高度变化的项分开。该实验所得 $|\alpha_{\rm H}|<2\times10^{-4}$，下标 H 特指氢钟。若假定所有钟种的额外漂移都一样，单靠两只同型钟和以同一物理过程制成的尺，并不能分开所有可能的变化。对它的解释中还含有外场势差由何种测量确定这一条件。

另一种计时从一开始便让钟分路，最后回到同一地点直接对读。对低速、弱场的两条路径 $A,B$，相对于同一坐标时 $t$，有
\[
 \tau_A-\tau_B\simeq\frac{1}{c^2}\int
  [\phi_A-\phi_B-\tfrac12(v_A^2-v_B^2)]\,\mathrm dt .
\]
这里 $v$ 应在同一参考系中取，例如都用不随地球转动的地心系。Hafele 和 Keating 用铯钟随民航客机分别向东、向西绕地球，返回后与地面钟比较。对于地速相同、方向相反的两次飞行，势项随飞行高度而变，速度平方中的交叉项却随飞行方向变号。若地面随自转有速度 $v_E$，客机相对地面沿纬圈有速度 $\pm v_f$，则 $(v_E\pm v_f)^2-v_E^2=v_f^2\pm2v_Ev_f$。因此只把空速平方代入而遗漏地面本身的运动，甚至可能错判钟差的符号。原实验的精度约为十分之一量级，但它的读数是一对在相同起点和终点直接比较的累积钟差，和塔上持续比较两个静止频率标准不同。采用这种办法不用规定远隔钟的初始同步，代价是要知道整个飞行过程的路线、速度和高度。

另有一种\textbf{零红移实验}，将工作原理不同的钟 $A,B$ 并排放置，两者一起随地球运动。在静态弱场中，若把同种钟红移公式的异常系数归给钟在给定坐标时内的走时率，则两钟本地频率比的变化可写为 $\Delta\ln(\nu_A/\nu_B)=(\alpha_A-\alpha_B)\Delta\phi/c^2$。共同的引力红移在相除时消去，实验只需查两钟频率比是否随外场有规律地变化。地球轨道稍椭圆，在近日点和远日点之间移动，太阳引力势有约 $3\times10^{-10}c^2$ 的峰峰变化。若两类钟的 $\alpha$ 相差 $10^{-6}$，该比值的年变化也只有 $3\times10^{-16}$ 左右。以原子钟和谐振腔钟作的早期实验对两者异常系数之差给出百分之一量级的限制。后来把不同原子跃迁互比，可把某些组合限制到百万分之一量级。设若所有 $\alpha$ 都等于同一个非零值，这个实验仍会得到零，这正是还需测绝对红移系数的缘故。

还有一个容易弄混的例子。在地球正午和午夜两处各放一只同种钟，向阳的一只比另一只离太阳近约两个地球半径。能否径取这两点太阳势差，预言它们之间有很大的日夜频差？在地心随地球绕日自由落下的坐标中，太阳外场在地心的一阶梯度已与坐标系的加速抵消。展开太阳势，扣去常数和这一次项，才得到与 $x^ix^j\partial_i\partial_j\phi_\odot$ 成正比的潮汐项。把地球当作自由落体，忽略其自转和形变时，不能测到前一种偶极的日夜频差。零红移实验却是把两种不同的钟放在一起，检验它们内部的频标会否随太阳外场而异，两个问题的设备和信号都不相同。

钟的比较还可检验非引力常数的变化。若跃迁频率依赖精细结构常数 $\alpha_{\rm em}$、电子质子质量比等，可算出每种跃迁对无量纲常数 $X$ 的灵敏系数 $K_A^X=\partial\ln\nu_A/\partial\ln X$。两只钟的读数给出 $\Delta\ln(\nu_A/\nu_B)=\sum_X(K_A^X-K_B^X)\Delta\ln X$。此处单个频率系数使用同一参考频标，真正可测的是两系数之差。两种同受一个普适频率尺度支配的跃迁，未必是好的比较对象。原子光学跃迁与超精细跃迁的差别较大，因而可分辨的组合也不同。以若干年的实验室资料约束的是现时变化率。以遥远吸收云的几条原子谱线作比较，则要把各谱线共有的宇宙学红移与不同跃迁间的相对位移分开，所得是两个历元之间的变化。若吸收线来自速度略有不同的气团，谱线的相对移动也能模仿常数的变化。

\subsubsection{从太阳系的频标到不同时代的物理常数}

前述零实验中的外部势按时间改变，常数本身也可能在长时间尺度上漂移。欲区分二者，可将实测钟比写成 $\Delta\ln(\nu_A/\nu_B)=b\,\Delta\phi_\odot/c^2+d\,(t-t_0)$。其中第一项在近日点附近有确定相位，第二项是单调漂移。在一小段近似直线的轨道弧上两项会混淆，覆盖整季乃至整年方能分开。若外部环境自身也随季节改变，如室温或磁场有同相变化，增加拟合年项不足以消除系统误差，还须用仪器的温度、磁场灵敏度校正。对弱相互作用、引力等其它常数也应先选择无量纲组合。只声明某个以秒的幂或以米表示的有量纲常数改变，而不给时钟和长度基准，并不是单独可复现实验。

表~\ref{tab:exp-historical-constants} 按观测手段给出 2018 年参考书汇集的一组历史界限。这里 $X$ 均为无量纲量。$\alpha_{\rm em}=e^2/(\hbar c)$ 采用 Gaussian 电磁单位，$\alpha_{\rm W}=G_Fm_p^2c/\hbar^3$ 以 Fermi 常数和质子质量定义，质量比取 $m_e/m_p$。实验室钟的 $z=0$ 表示近代直接测得的漂移。其余各行把不同时代的差异折成每年变化率，一阶相当于 $|\ln[X(t_0)/X(t_z)]|/(t_0-t_z)$。只有再假定单调、近似线性漂移，它才可当作同一时段每年的导数。地质样品的 $z$ 是按宇宙年龄换算的历元，并非测得的光谱红移。
\begin{table}[htbp]
\centering\small
\caption{由钟、地质样品与远方谱线得到的历史限度。末列为相对变化率上限。}
\label{tab:exp-historical-constants}
\setlength{\tabcolsep}{5pt}
\begin{tabular}{@{}llcc@{}}
\toprule
观测手段 & $X$ & 历元 $z$ & 上限（$\mathrm{yr}^{-1}$）\\
\midrule
实验室钟的频率比 & $\alpha_{\rm em}$ & $0$ & $1.3\times10^{-16}$\\
实验室钟的频率比 & $m_e/m_p$ & $0$ & $3.3\times10^{-15}$\\
Oklo 天然反应堆 & $\alpha_{\rm em}$ & $0.15$ & $5\times10^{-17}$\\
Oklo 天然反应堆 & $\alpha_{\rm W}$ & $0.15$ & $1\times10^{-11}$\\
陨石中的 $^{187}\mathrm{Re}$ 衰变 & $\alpha_{\rm em}$ & $0.45$ & $3.4\times10^{-16}$\\
类星体吸收谱 & $\alpha_{\rm em}$ & $0.4$--$2.3$ & $1.2\times10^{-16}$\\
类星体吸收谱 & $m_e/m_p$ & $2.6$--$3.0$ & $3\times10^{-15}$\\
原初轻元素丰度 & $\alpha_{\rm W}$ & $\sim10^9$ & $5\times10^{-12}$\\
\bottomrule
\end{tabular}
\end{table}

例如 Leefer 等（2013）比较镝的两种同位素，近简并的能级对 $\alpha_{\rm em}$ 变化有较大而不同的频率响应，给出表中实验室的第一项。若改用两种超精细跃迁，核磁矩、质量比与 $\alpha_{\rm em}$ 都可影响读数，不能把钟比的一条直线拟合同时当成每个基本常数的独立界。$m_e/m_p$ 的两行采用 2018 年汇编值。在分子吸收谱中，要用对质量比有不同灵敏度的谱线彼此比较，同时拟合同一片吸收云的共同速度。对金属吸收谱，Srianand 等（2004）的分析给出 $\dot\alpha_{\rm em}/\alpha_{\rm em}=(-0.6\pm0.6)\times10^{-16}\,\mathrm{yr}^{-1}$，汇编以 $1.2\times10^{-16}\,\mathrm{yr}^{-1}$ 列作限度。另一些样本曾报告非零结果，由此不能将这些不同假设、不同统计处理的数合成一个普适置信区间。

约十八亿年前，加蓬 Oklo 一带含铀矿体在地下水慢化中子时曾间歇维持链式裂变。$^{149}\mathrm{Sm}$ 的热中子俘获截面有一个靠近 $0.1\,\mathrm{eV}$ 的共振。今日由矿样中的 $^{149}\mathrm{Sm}/{}^{147}\mathrm{Sm}$ 及反应堆运行期间的中子积分通量，可反推当时 $^{149}\mathrm{Sm}$ 的有效截面，再限制共振能的变化。Shlyakhter 于 1976 年提出这种比较，对精细结构常数的上限参考了 Damour 与 Dyson（1996）、Fujii（2004）及 Petrov 等（2006）的重分析。要从核能级反推 $\alpha_{\rm em}$，须用静电能对共振位置的灵敏度。改问弱耦合时，则须用另一种核响应。表中 Oklo 对 $\alpha_{\rm W}$ 的一行取 2018 年汇编所列值，不能仅由对 $\alpha_{\rm em}$ 的限度换单位得到。若同时容许强相互作用尺度、静电和弱作用都漂移，观测仍只约束实际核结构选出的组合。

在约四十五亿年前形成的陨石中，$^{187}\mathrm{Re}$ 的低能 $\beta$ 衰变是另一类时钟。母、女核丰度记载的是衰变率的时间积分。衰变 $Q$ 值对静电能的灵敏度须另由核模型换算，表中一项出自 Olive 等（2004）。原初核合成所含的弱作用信息又不同。在自然单位中，中子、质子互变的热反应率约正比于 $G_F^2 T^5$。膨胀率 $H$ 与反应率相当时，丰度比逐渐偏离热平衡，其后还有自由中子的衰变。到开始合成氦时留下多少中子，同时取决于 $G_F$、中子质子质量差和早期膨胀史。参考书引述 Malaney 与 Mathews（1993）、Reeves（1994）等研究，固定相应背景模型后才列出表中的弱作用界限。同一古样品、吸收云或丰度读数若容许多个常数联动，不能逐一反演表中每个 $X$。

\subsection{Schiff 猜想与电磁束缚物}

为什么称 Eötvös 实验、钟的红移与 Lorentz 实验为同一原理的三个方面？钟的跃迁能也计入物体的质量，故把钟提到不同高度所作的功，应当与在那里交换的光子能量有关。可以不指定完整场方程，先构造一个弱场、静态、无方向依赖的能量模型来说明。取地面某处为 $\phi=0$，设一原子的基态、激发态质量为 $m_1,m_2$，$\Delta m\,c^2=(m_2-m_1)c^2=h\nu_0$。若两态的被动引力质量为 $m_1^P,m_2^P$，在固定的参考时标下，它们静止在势 $\phi$ 处的能量可记为 $E_a(\phi)=m_ac^2+m_a^P\phi$，只保留势的一次项。这里的能量以同一个参考钟计量。当地固有钟的读数还要作相应的换算。

先在低处激发原子，带到高处使之退激，再把基态原子带回原点。设高低两处势差为 $\Delta\phi$，提升与放下相抵后净作功 $(m_2^P-m_1^P)\Delta\phi$。原子回到原来的状态，借入的激发能和退激时给出的能量之差，只能等于所作的净功。换言之，按同一时标计量，跃迁能之差为 $\Delta(E_2-E_1)=(m_2^P-m_1^P)\Delta\phi$。若被动引力质量与惯性质量相同，其相对于 $h\nu_0$ 的变化正是普通红移所需的 $\Delta\phi/c^2$。把偏离量写进钟的系数，便得 $\alpha_{21}=\frac{m_2^P-m_1^P}{m_2-m_1}-1$。因此即使两态各自的 $m_a^P/m_a-1$ 很小，跃迁能与总静能之比也很小，作差后得到的红移异常未必被同样强地抑制。反过来，若在这个模型里一切物体、一切内态均有同一个 $m_a^P/m_a$，把公共因子收入测得的外场强度后，任何跃迁都只给普通红移。循环论证用到了可定义守恒能量、运送物体不引起其它内态变化、源与钟共享一种能量计量等条件。仅从两块材料恰有相同的加速度，无法推出对一切钟都成立的结论。

若还容许方向依赖，惯性质量本身也可能要用张量描述。比如复合体的质心动能中出现 $\tfrac12(m\delta_{ij}+\delta m^I_{ij})v^iv^j$，其总动量与 $\bm v$ 就不再严格平行。在普通的各向同性外场里，若受力仍沿 $\bm g$，反解动量变化的方程也会给出依方向而异的加速度。反之若异常来自物体内部结合能随外场位置的变化，作用的是被动引力质量。这两类效应在单轴落体实验中能落入同一读数。转动物体、比较不同自旋态，再把结果与谱线的方位实验合用，才有希望认出是哪一类系数。Schiff 猜想认为，在一个足够完备且自洽的理论中，这些异常不能任意彼此分开。怎样把这种想法变成一个可证的命题？

\subsubsection{用能量和动量同时描述钟与落体}

把以上想法写成一个能逐项核算的模型。以已选参考物质的自由落体定义正势 $U$，使 $g_i=\partial_iU$。为容许各向异性，引入 $U^{jk}=G\int\rho'(x-x')^j(x-x')^k\,\d^3x'/|\bm x-\bm x'|^3$，其迹为 $U$。这个张量不是潮汐张量：它与势同量纲，对于一个点源即为 $Un^jn^k$，而潮汐还须对标量势求两次空间导数。把复合体在标准环境中的静质量写成 $M=M_0-E_B/c^2$，$E_B>0$ 是将它拆散所需的结合能。设守恒能量在外场和质心低速的一阶非平凡近似为
\[
 E=Mc^2+\tfrac12(M\delta_{jk}+\delta m^I_{jk})V^jV^k
              -(M\delta_{jk}+\delta m^P_{jk})U^{jk}.
\]
两种质量异常先作为不同量处理。它们本身依赖内部构型，可以为两个各向同性标量，也可有反映朝向的对称无迹部分。若 $E$ 是同一保守拉氏量导出的能量，则 $P_j=(M\delta_{jk}+\delta m^I_{jk})V^k$。在瞬时由静止释放、内部构型不变时，取力方程到异常的一次即可得到
\[
 A_i=g_i+\frac{\delta m^P_{jk}}{M}\partial_iU^{jk}
             -\frac{\delta m^I_{ij}}{M}g_j .
\]
在球形源的场中，即使异常惯性张量是常数，也可使加速度出现垂直于普通引力的分量。若两张量的各向同性系数分别为 $p_0,i_0$，则只余 $(p_0-i_0)g_i/M$。故一类物体在一种场源方向上的加速度为零差，并不等价于两个张量的每个分量都为零。

将一个复合体从 $a$ 态跃迁到 $b$ 态，改测光子。令 $\Delta$ 表示初态减末态，$h\nu_0=\Delta Mc^2$，以 $\alpha^I_{jk}=c^2\Delta(\delta m^I_{jk})/(h\nu_0)$、$\alpha^P_{jk}=c^2\Delta(\delta m^P_{jk})/(h\nu_0)$ 记跃迁系数。忽略比例为 $h\nu_0/(Mc^2)$ 的普通反冲修正，但保留质心与发射光方向 $\hat{\bm n}$ 的 Doppler 项。同时满足 $\Delta E=h\nu$、$\Delta\bm P=(h\nu/c)\hat{\bm n}$，消去发射前后那一微小速度差，得
\[
 \frac{\nu}{\nu_0}=
 \frac{1-(\delta_{jk}+\alpha^I_{jk})V^jV^k/(2c^2)
         -(\delta_{jk}+\alpha^P_{jk})U^{jk}/c^2}
      {1-\bm V\cdot\hat{\bm n}/c}
\]
至所保留阶次。直接相减两态在同一速度下的能量，会把含 $V^2$ 的项取成错误的正号。原因是初态和末态并没有相同的动量，反冲虽小，普通质心动能的一阶变化仍须由动量守恒求出。若同时观测两种跃迁，位置、速度和接收方向相同，分母与通常的钟速项在比值中相消。因而钟比实际挑出的是两条谱线的 $\Delta(\delta m)/\Delta M$ 之差，落体实验却问整个物体的 $\delta m/M$。二者可受同一异常耦合支配，却决不是同一个测值。若怀疑发出的光子本身有色散，还要改变 $\Delta\bm P$ 与 $\Delta E$ 之间的关系，这一式不能原封不动使用。

再想象完全不用光的循环。在高处以几个静止的粒子组成束缚物，把放出的能量另存，将束缚物及零散参考物都落到低处，分别把终态速度降为零，最后拆散束缚物并把各粒子缓慢运回。一个循环后粒子、仪器及储能器都回到原来的状态。束缚物下降的动能取决于惯性张量，低处拆散所需的功则取决于结合能的外场依赖，两者之差恰好补足把自由粒子升回所作的功。这是用守恒能量从另一方向得到同一加速度公式的办法。若模型另有会吸收能量的自由度，在循环末态也须使它恢复原状。只规定仪器表针回零而不管其它场的终态，并不构成这样的证明。

\subsubsection{静态球对称外场的四个系数}

Lightman 和 Lee 取静态球对称外场，选空间各向同性的坐标来描述一群有电荷的点粒子及其电磁束缚场。在本段取 $c=1$ 并用 Gaussian 电磁单位。设四个函数 $T,H,\epsilon,\mu$ 都只依赖外场势，远处用常规尺、钟和电荷单位归一化。假设粒子的自由运动由线元 $\d s^2=-T\d t^2+H\d\bm x^2$ 决定，但暂让电磁场有独立的系数，即取
\[
\begin{split}
 S_{\rm ng}={}&-\sum_a m_a\int\sqrt{T-Hv_a^2}\,\d t
             +\sum_a q_a\int A_\mu\d x_a^\mu\\
 &+\frac{1}{8\pi}\int\bigl(\epsilon\bm E^2-\mu^{-1}\bm B^2\bigr)\,\d t\,\d^3x .
\end{split}
\]
这里 $\bm E,\bm B$ 是由坐标势得出的分量。$\epsilon,\mu$ 描述外场对它们的作用，实验未在真空中放入介质。这个模型故意保留了两种可能不同的响应。一切无结构、无电荷的点粒子早已按同一 $T,H$ 运动，因此仅把点粒子放手，不会给 $\epsilon,\mu$ 任何条件。有约束的是内部由电磁场结合在一起的中性物体。

先看在事件 $P$ 附近能不能辨别速度。把四函数在这一点的值记作带下标 $0$。由粒子项开方不能为负可见粒子的极限坐标速度为 $c_0=\sqrt{T_0/H_0}$。在波长远小于四函数变化尺度时，Maxwell 方程的主部给出的电磁波坐标速度为 $c_{\rm em}=1/\sqrt{\epsilon_0\mu_0}$。用当地尺钟同时测这两个速度，其比值不因重新标定时间或长度而变。记 $\Upsilon_0=1-(c_0/c_{\rm em})^2=1-T_0\epsilon_0\mu_0/H_0$，则前面往返光时模型中的 $C=c_{\rm em}/c_0$ 满足 $C^2=(1-\Upsilon_0)^{-1}$。若 $\Upsilon_0\ne0$，粒子与电磁波便没有共同的不变锥面。只检验光速的绝对数值不足以比较两个锥面，必须把传播与物质构成的尺、钟联系起来。

怎样在上述坐标中实际构造自由落体实验室？取实验室中心瞬时静止于 $P$，以 $e^{\hat a}{}_{\mu}=\operatorname{diag}(\sqrt{T_0},\sqrt{H_0},\sqrt{H_0},\sqrt{H_0})$ 校准中心的时间和空间单位。若只作这个线性变换，$P$ 处的度规成为 Minkowski 形式，但质点方程仍含一次梯度项。到坐标差 $\Delta x^\mu$ 的二次项可取
\[
 X^{\hat a}=e^{\hat a}{}_{\mu}
 \left[\Delta x^\mu+\tfrac12\Gamma^\mu{}_{\nu\lambda}(P)
                          \Delta x^\nu\Delta x^\lambda\right].
\]
例如 $\Gamma^i{}_{00}=\partial_iT/(2H)$，一个由静止释放的质点有 $\mathrm d^2x^i/\mathrm dt^2=-\partial_iT/(2H)$。上式在 $X^{\hat i}$ 中增加一个与 $t^2$ 成正比的项，恰使这个质点在 $P$ 的局域加速度为零。时间坐标还要用到 $\Gamma^0{}_{0i}=\partial_iT/(2T)$，空间部分还含 $\Gamma^i{}_{jk}=(\delta^i_j\partial_kH+\delta^i_k\partial_jH-\delta_{jk}\partial_iH)/(2H)$。把 Maxwell 场也作为同一坐标变换下的二形式来处理，不能在变换质点坐标后把 $\bm E,\bm B$ 的旧分量留在原式中。

在 $P$，由 $F_{i0}=\sqrt{H_0T_0}F_{\hat i\hat0}$ 可定电场，磁场有 $F_{ij}=H_0F_{\hat i\hat j}$，四维坐标体积除以 $\sqrt{T_0}H_0^{3/2}$。所以局域电磁作用量中的两个项为 $\epsilon_0\sqrt{T_0/H_0}\,\hat{\bm E}^2$ 和 $-[\mu_0\sqrt{T_0/H_0}]^{-1}\hat{\bm B}^2$，这就是下面两个有效系数的由来。离开 $P$ 时，若这些系数依赖于地点，静电能受到与 $\hat X^i\partial_i\ln(\epsilon\sqrt{T/H})$ 成正比的调制，静磁能亦然。随加速坐标而来的局域推动还会混合 $\bm E$ 和 $\bm B$。电磁锥与粒子锥不同时，场作用量里可有与该相对速度点乘 $\hat{\bm E}\times\hat{\bm B}$ 成正比的项。这些都是一次外场梯度所留下的非普适效应，其尺度依赖与二次度规导数产生的普通潮汐不同。令两个有效系数互为常数倒数，则这种一阶局域实验重新取得普通电动力学的形式。

即使处处有 $\Upsilon_0=0$，位置不变性也尚未得到。用局域粒子单位把电磁作用量改写，决定静电与静磁强度的是 $e_*=\epsilon\sqrt{T/H}$ 和 $m_*=\mu\sqrt{T/H}$ 这两个函数。此处的 $e_*,m_*$ 都是系数，与粒子的荷、质量无关。两种速度相等只要求它们的积为一。一者随位置增大而另一者等比减小，仍能满足。若在每处各自改变电荷单位，可把 $e_*$ 吸收到某一处的 Coulomb 定律中，却不能使所有地点以同一种实验单位测得的精细结构常数都相同。为计量势的这种作用，选在 $P$ 邻域以 $\bm g=\bm\nabla U$ 归一化的正 Newton 势 $U$，并用原坐标的速度单位，定义
\[
    \Gamma_0=-c_0^2\left.\frac{\d\ln e_*}{\d U}\right|_P,\qquad
    \Lambda_0=-c_0^2\left.\frac{\d\ln m_*}{\d U}\right|_P .
\]
$\Upsilon_0$ 比较同一位置的两个速度，$\Gamma_0,\Lambda_0$ 则比较不同位置的两类耦合。只在一个点调单位，把某几个系数的数值调为一，并不会使它们的导数为零。

\subsubsection{用原子尺度与晶格尺度实现钟}

可以先不解完整的 Dirac 本征问题，找出哪几个系数进入每种钟。在本小节粒子局域单位中，局域电子质量 $m_e$ 固定，静电势为 $-e^2/(e_*r)$。对半径约 $a$ 的束缚态，动能约 $\hbar^2/(2m_ea^2)$，与库仑势能之和取极值得 $a\propto e_*\hbar^2/(m_ee^2)$。更精确地令 $r=ar'$，则整个无量纲 Schrödinger 方程只余量子数，主能级具有共同尺度 $m_e e^4/(\hbar^2 e_*^2)$。一个给定束缚态的归一化波函数在原点的密度因此正比于 $a^{-3}$。这些伸缩关系均由固有尺、钟测量。若写回静态坐标的能量，要另乘 $\sqrt T$，否则会把正常的引力红移也误算成位置不变性的破坏。

精细结构可由相对论动能修正 $-p^4/(8m_e^3)$、自旋--轨道项和 Darwin 项作一阶微扰。在 $p\sim\hbar/a$ 的原子上，动能的四次修正比主能量再小一个有效精细结构常数的平方，自旋--轨道项把库仑电场乘上动量与 $1/r$，恰有相同的标度。考虑质子磁矩与电子自旋在 $s$ 态中的接触相互作用时，磁偶极场的能量还带局域磁响应 $m_*$，故劈裂正比于 $m_*a^{-3}$。这一步假设核磁矩的组成没有其它额外位置依赖。若另有非电磁耦合改变质子磁矩，氢超精细跃迁所测的组合便相应增加一项，不能归入一个普通电子精细结构常数。以局域量表述 Dirac 方程还有一个好处：自旋矩阵的代数仍使用共同的粒子正交标架，不必因电磁相速有所改变而另给电子一套光子度规。有关标架和旋量联络的几何构造见 \S\ref{sec:tetrad}。

还可直接核对上述依赖。令局域 Dirac 粒子的荷为带符号的 $q$，$\bm\pi=\bm p-q\bm A$ 为机械动量。对缓变静电外场，把正负能分量分开，到精细结构所需精度，其正能 Hamilton 量除常见的 $m_e+\bm\pi^2/(2m_e)+q\phi_q$ 外，包含
\[
 -\frac{\bm\pi^4}{8m_e^3}
 -\frac{q\hbar}{2m_e}\bm\sigma\cdot\bm B
 -\frac{q\hbar^2}{8m_e^2}\bm\nabla\cdot\bm E
 -\frac{q\hbar}{8m_e^2}
       \bm\sigma\cdot(\bm E\times\bm\pi-\bm\pi\times\bm E).
\]
写的是 $c=1$ 的局域单位，省去了时间依赖外场和更高次磁场项。对点核，$\bm\nabla\cdot\bm E$ 有位于核上的分布项，说明 Darwin 能移只取样束缚态在核处的概率密度。对罕有的非点状核或重原子，核的有限尺度与相对论电子波函数都可能改变这一取样。自旋--轨道项来自电场在移动电子静系中表现出的磁效应及其运动学修正，不能把它的全部引力依赖都算在模拟真空磁导率的 $m_*$ 上。质子磁矩提供的外加偶极场则确由磁响应决定。两种看似都和自旋有关的效应在这一点上正好不同。

再把同一种固体制成电磁谐振腔。若晶格长度主要由 Coulomb 相互作用决定，设 $N$ 个晶格间距给出腔长，则 $L\propto Na\propto e_*$。光在腔中以 $(e_*m_*)^{-1/2}$ 的局域速度传播。忽略镜面相移、残余气体和材料色散后，其共振频率 $\nu_{\rm cav}\propto e_*^{-3/2}m_*^{-1/2}$。因此即使腔钟与氢钟都只靠电磁作用，两个频标的幂次也相当不同。若把腔长人为当作固定的坐标长度，就会遗漏 $e_*^{-1}$，预测错误的红移系数。实际材料的不同跃迁、固化方式、磁性和机械支撑可改变对基本常数的灵敏性，不应将这个简化结果直接推广到所有仪器。

这些标度用哪一种时间读出？在静止原子的局域粒子标架中记时间为 $\hat t$，在静态球对称坐标中记为 $t$。于固定空间点有 $\mathrm d\hat t=\sqrt{T}\,\mathrm dt$。例如以原子长度制成的腔，单位 $\hat t$ 内振荡次数正比于 $e_*^{-3/2}m_*^{-1/2}$。同一振荡的相位 $\mathrm d\theta=2\pi\nu_{\rm cav}^{\rm loc}\,\mathrm d\hat t=2\pi\nu_{\rm cav}^{\rm coord}\,\mathrm dt$，故
\[
 \nu_{\rm cav}^{\rm coord}\ \propto\ \sqrt{T}\,e_*^{-3/2}m_*^{-1/2}
   =\frac{H}{\sqrt T\,\epsilon^{3/2}\mu^{1/2}} .
\]
从坐标量也可得到同一结果：腔的坐标长度为其固有长度除以 $\sqrt H$，波的坐标传播速度为 $(\epsilon\mu)^{-1/2}$，两者之比即为坐标频率的尺度。若为度规耦合，$\epsilon=\mu=\sqrt{H/T}$，上式恰好只余 $\sqrt T$，与氢钟相同。在同一点用频率比读数时这个公共因子消失，在两个静止地点用同种钟作红移实验时它却正给出正常红移。现在依次把各钟的局域标度对外场求导，并恢复公共的 $\sqrt{T}$ 因子，依次以 p、f、hf 和 cav 表示主能级、精细结构、超精细结构和腔，得到
\[
 \alpha_{\rm p}=-2\Gamma_0,\quad
 \alpha_{\rm f}=-4\Gamma_0,\quad
 \alpha_{\rm hf}=-3\Gamma_0+\Lambda_0,\quad
 \alpha_{\rm cav}=-\tfrac12(3\Gamma_0+\Lambda_0).
\]
例如氢激射钟单独测红移，约束的是 $3\Gamma_0-\Lambda_0$。氢钟与腔钟并排比较，选出的是与 $\Gamma_0-\Lambda_0$ 成正比的另一个组合。这说明多一种钟不仅可能提高精度，也会增加可以区分的物理量。以上系数属于氢和特定腔的电磁模型。以铯、铷或重离子为频标时，原子的相对论性电子结构不同，须另算灵敏系数。

\subsubsection{由内部维里关系求异常质量}\label{sec:exp-binding-mass}

取两点电荷 $q_1,q_2$，以 $\mu_r$ 表示约化质量，$r=|\bm x_1-\bm x_2|$、$\bm n=(\bm x_1-\bm x_2)/r$。在局域粒子单位中写 $v_C(r)=q_1q_2/(e_*r)$。先只考虑异号电荷的稳定 Coulomb 轨道，并在长于内部运动的时段上平均。相对动量 $p_i$ 满足 $\dot p_i=q_1q_2 n_i/(e_*r^2)$。由 $\langle\mathrm d(r^jp_i)/\mathrm dt\rangle=0$ 得
\[
    \left\langle\frac{p_ip_j}{\mu_r}\right\rangle
       =-\langle v_C n_i n_j\rangle,qquad
    E_{Bij}:=-\tfrac12\langle v_C n_i n_j\rangle .
\]
取迹得 $E_B=-\langle v_C\rangle/2>0$，因为动能平均值为 $-\langle v_C\rangle/2$，而总内能为 $\langle v_C\rangle/2$。非球形的定向轨道可有 $E_{Bij}\ne E_B\delta_{ij}/3$。场的电应力有方向，因而同一个静电结合能在两个体系中虽可取相同的总数，也未必有相同的二阶矩。在多电荷体系，可对所有不重复的电荷对求和。无穷大的点电荷自能先并入粒子的质量参数，留下的是粒子间可随构型改变的相互作用能。

为什么内部静电结构会进入质心惯性？从全局作用量变分，场满足 $\bm\nabla\cdot(\epsilon\bm E)=4\pi\rho_q$、$\bm\nabla\times(\mu^{-1}\bm B)=4\pi\bm j+\partial_t(\epsilon\bm E)$，另外两式是 $F=\mathrm dA$ 的恒等式。固定外场缓慢运动电荷产生的场，其中 $\bm B$ 是内部速度的一阶，回代粒子方程的磁力则是速度二阶。若使整个体系以速度 $\bm V$ 平移，在这两阶中既有 $V^2 v_C$，也有 $(\bm V\cdot\bm n)^2v_C$。前者是动能中的标量修正，后者留下束缚能张量。即使内部每一电荷所受的首阶 Coulomb 作用在求和时相消，电磁场的能量、动量仍参与复合体质心的惯性。只把各裸粒子的 $m_a\bm v_a$ 相加，会漏掉这部分。

例如考虑静电场解中究竟哪些梯度会留下。改用通常的静电势 $\phi_q$ 使 $\bm E=-\bm\nabla\phi_q$，则 $\bm\nabla\cdot(\epsilon\bm\nabla\phi_q)=-4\pi\rho_q$。令 $b=\sqrt\epsilon$、$\phi_q=f/b$，左端恰为 $b\nabla^2 f-f\nabla^2 b$。若外场在体系大小上仅保留一次梯度，$\nabla^2b$ 便可忽略。因此带电粒子间的瞬时势为
\[
 \phi_q(\bm x)=\sum_a
  \frac{q_a}{|\bm x-\bm x_a|\sqrt{\epsilon(\bm x)\epsilon(\bm x_a)}} .
\]
这说明同一个电荷对的作用依赖场点与源点两处的 $\epsilon$，并非仅把普通 Coulomb 势的电荷替换成某一个质心位置的有效荷。求其对质心的力时，两端共同平移，两个平方根都要参与。较高阶的外场变化会使 $\nabla^2 b$ 进入方程，此处的近似核就不再足够。

再让电荷在外场大致均匀的小区域内缓慢运动，选 Lorenz 型规范，局域波动算符为 $\nabla^2-c_{\rm em}^{-2}\partial_t^2$。推迟标势在常数 $\epsilon_0,\mu_0$ 下可按内部速度展开。设 $s=|\bm x-\bm x'|$，则到第一次非静态修正，
\[
 \phi_q=\frac1{\epsilon_0}\int\frac{\rho'_q(t)}s\,\d^3x'
 +\frac1{2\epsilon_0 c_{\rm em}^2}\partial_t^2
        \int s\rho'_q(t)\,\d^3x'+\cdots,\qquad
 \bm A=\mu_0\int\frac{\bm j'(t)}s\,\d^3x'+\cdots .
\]
中间可能有的 $1/c_{\rm em}$ 项正比于总电荷的时间导数，已由电荷守恒消掉。矢势的最低项与内部电流同阶，取 $-\partial_t\bm A$ 或 $\bm v\times(\bm\nabla\times\bm A)$ 时应与标势中含二次时间导数的项一并保留。换另一套电磁规范可在这两个势之间调换一部分推迟项，最后进入粒子力方程的电场、磁场不变。一方面从这里来的速度平方项，另一方面从 $\sqrt{T-Hv^2}$ 展开至 $v^4$ 所得的粒子动能，在普通度规情形恰能合为一个标准的质心动能。若粒子的 $c_0$ 与场的 $c_{\rm em}$ 不同，剩余系数便正比于 $1-c_0^2/c_{\rm em}^2$。

保留首阶 Coulomb 结合能和最低的内速度修正，场的梯度只取一次，内部运动按上述维里关系消去。记 $M=M_0-E_B/c_0^2$，在 $P$ 处把局域表达换回定义 $c_0$ 的坐标单位。以质心总正则动量 $\bm P$ 为变量，内部平均后的 Hamilton 量含有
\[
 \langle H\rangle =
 Mc_0^2-MU+\frac{P^2}{2M}
 -\frac{2\Gamma_0 E_B}{c_0^2}U
 -\frac{\Upsilon_0}{M^2c_0^2}
       (E_B\delta_{ij}+E_{Bij})P^iP^j .
\]
这里只保留辨别该束缚能响应所需的项，不列普通质心动力学的高阶修正。势能修正也可由固定量子数的束缚能给出，它正比于 $e_*^{-2}$，所以对位置求导产生 $2\Gamma_0$。动能中 $P_iP_j$ 项的符号还要留心：在给定动量而非给定速度时，惯性增加会使动能减少。对 $P_i$ 求导得 $V_i$，再反解 $P_i$，在上面的守恒能量模型中读出
\[
 \delta m^P_{ij}=\frac{2\Gamma_0E_B}{c_0^2}\delta_{ij},\qquad
 \delta m^I_{ij}=\frac{2\Upsilon_0}{c_0^2}
                    (E_B\delta_{ij}+E_{Bij}).
\]
这些式子采用正的束缚能 $E_B$，若换成负的系统内能，表示系数的符号也须同时改换。它们说明即使局域自由点粒子看不见优选系，带内部电磁动量的复合体仍可能给出一个各向异性的质心惯性。

把后一式取球平均，$E_{Bij}=E_B\delta_{ij}/3$，上面的质心方程给出 $\bm A=\bm g\{1+[2\Gamma_0-8\Upsilon_0/3]E_B/(Mc_0^2)\}$。因此仅用各种宏观、无取向的静电束缚物作落体，就算它们全给出普通加速度，也只说明 $2\Gamma_0-8\Upsilon_0/3$ 为零。钟或定向体系提供了不同的方程。比如把同一体系的束缚张量相对 $\bm g$ 旋转，标量 $E_B$ 不变，而 $E_{Bij}g_j$ 可改变大小或出现横分量。要求对这些体系 WEP 仍成立才能分出 $\Upsilon_0$。磁相互作用能比静电能再高两个内部速度的幂次，在那一阶 $\Lambda_0$ 也可直接出现在落体公式中。每种具体实验的“验证”都有它所比较的材料、取向和速度阶次，不能同要求对任意束缚系统成立的理论命题互换。

因而在这类模型中，对一切体系及取向要求 WEP 给出 $\Gamma_0=\Upsilon_0=0$。要求此条件对任意外场位置都成立，则 $e_*$ 在各处恒定，且 $e_*m_*=1$。可一次性重选电荷单位使二者均为一，回到 $\epsilon=\mu=\sqrt{H/T}$。反过来将此式代进作用量，电磁部分正好就是用上述线元缩并 $F_{\mu\nu}$ 得到的 Maxwell 作用量。至此，在这个静态球对称、经典带电粒子与电磁场的模型内，确已从普适自由落体推出 EEP。超出这一范围仍有别的耦合。例如 Ni 容许赝标量乘上 $F_{\mu\nu}\,{}^*F^{\mu\nu}$。若该赝标量不是常数，偏振或电磁束缚系统的力矩可有异常，而质心加速度仍可与成分无关。因此若 WEP 的检验只涉及试验体质心的轨道，便不能以 Schiff 猜想替代谱线、钟和偏振实验。

\section{引力作为几何效应}\label{sec:exp-geometric}

\subsection{由钟、光线和质点确定度规}

把两只工作原理不同的钟放在一起校准，随后让它们沿同一世界线走到另一事件。若 EEP 成立，用各自的固定频率把振荡次数折合为时间，两钟应当给出同样的读数。因此固有时可归于世界线本身，写成 $c^2\d\tau^2=-g_{\mu\nu}\d x^\mu\d x^\nu$。一切非引力物质都以同样方式感受同一度规，称为\textbf{普适耦合}(universal coupling)。这里要紧的是“同一”。倘若氢钟和腔钟各自都有一套计算固有时的二次型，哪怕每套单看都与时空几何相仿，它们沿同一世界线的累积读数仍可能不同。

这种区别在前节的四函数模型里可直接看出。在静态坐标的同一地点，用远处已经对准的时标计数，静止钟在 $\d t$ 内的振荡次数应为 $\d N_A=\nu_{A0}\sqrt{T}\,\d t$，其中 $\nu_{A0}$ 由钟种决定。若 $\epsilon=\mu=\sqrt{H/T}$，氢的各类能级间距和腔钟频率均只拾取同一个 $\sqrt T$ 因子，因而所有理想钟测得 $\d\tau=\sqrt T\,\d t$。反过来设两钟各多出 $F_A(U),F_B(U)$ 两个因子。只要二者之比随位置改变，沿途每段记录的钟速比就会改变，一个 $g_{00}$ 已不足以说明全部读数。若二者及一切其它钟都带有完全相同的因子，单凭钟的比较不能把它与度规的选取分开。

光和钟各自提供的信息也不同。将度规改为 $\tilde g_{\mu\nu}=\Omega^2(x)g_{\mu\nu}$，任一类光切矢仍为类光，真空中的无参数光路亦保持。只描出同一地点向各空间方向发出的光线，只能确定光锥，不能辨出这个正函数 $\Omega$。类时钟沿一段曲线测得的线元却变为 $\d\tilde\tau=\Omega\d\tau$。物质钟和沿途的频率比较给长度、时间的标度以操作意义。再让有质量的自由试验体通过，检验它们的运动是否也符合同一度规。因而要将 EEP 用于几何，不能把“所有物体都落得一样快”单独当作对尺、钟和电磁规律的一切假设。

在任意事件均可选自由落体坐标使 $g_{\mu\nu}=\eta_{\mu\nu}$ 且 $\partial_\lambda g_{\mu\nu}=0$。沿一条无加速、无自转的参考世界线还能用平行移动的标架作相同安排。但在离开参考线一小段距离 $X^i$ 后，邻近自由质点的相对加速度出现与 $R^i{}_{0j0}X^j$ 成正比的项，见 \S\ref{sec:tidal}。若仪器特征长度为 $L$，在其范围内改用平直度规所忽略的无量纲量为 $|R|L^2$ 量级。比较两个相隔很高的钟当然要经过中间有引力的区域。EEP 所说的局域，是指每只钟读频率、每段光线在局域惯性系中传播时可以略去潮汐，并不宣称整次收发实验可在一个惯性系中完成。给定仪器大小、灵敏度和目标效应，才有实际的局域近似。

\subsubsection{有限仪器怎样赋予坐标量以长度、时间？}

相距很远的两条世界线没有一个与观者无关的“同时距离”。局域对钟和正交空间的构造见\hyperlink{subsection.1.1.1}{\S\,1.1.1}及 \S\ref{sec:sync-clock}，这里直接用其结果。对相对坐标静止的观者，邻近两条静止坐标线之间的固有长度平方为 $(g_{ij}-g_{0i}g_{0j}/g_{00})\d x^i\d x^j$。而地面台站给远处反射器报出的雷达距离为 $c(\tau_2-\tau_1)/2$，其中 $\tau_1,\tau_2$ 是同一只钟的收发时刻。邻近静止端点的极限下两者相合。对一条很长且随时变化的光路，后者包含传播时间及端点运动，不能直接解释为某一空间切片的弧长。

还有一个区别涉及“最小耦合”的范围。只知道无引力时拉氏量是什么，并不能仅凭记号替换排除一切曲率耦合。比如在带电标量的自由作用量中加 $-\lambda_R\sqrt{-g}\,R|\varphi|^2$，它在平直时空处处消失，因此再多平直时空的电磁测试也测不出 $\lambda_R$。到了曲率非零处，它却表现如依地点改变的质量平方。若本地非引力钟不同程度地依赖这种标量的质量，就可能从钟速中辨出曲率环境。实验对有限波包本身的寻常潮汐，与拉氏量里主动加入的点上曲率耦合并不是一回事。前者可从已给定的共同度规和包络尺度求，后者的系数则要另行确定。在本章把 EEP 化为普适耦合的低能模型时，取的是当地非引力定律回到狭义相对论的情形，不借助一般协变性强行证明所有此类额外算符都不存在。

\subsection{从非引力作用量读出实验量}

一个直接构造物质规律的办法，是在局域惯性系中写下狭义相对论的物质作用量，再用 $g_{\mu\nu}$、相应的体元与协变导数写成任意坐标下的形式。以下暂用 $c=1$。带电粒子的作用量为 $S_p=-m\int\d\tau+q\int A_\mu\d x^\mu$。变分第二项并分部积分，沿轨道偏离量前的系数为 $q(\partial_\mu A_\nu-\partial_\nu A_\mu)U^\nu$。于是 $mU^\nu\cd_\nu U^\mu=qF^\mu{}_\nu U^\nu$。当 $q=0$，质量从运动方程中消失，给定初始点和 4-速即可确定一条度规的类时测地线。若用带电试验体比较自由落体，则还须屏蔽电场或控制 $q/m$。一个很小的剩余电荷就会把右端变成材料依赖的加速度。

将真空电磁作用量写作 $S_{\rm em}=-(16\pi)^{-1}\int \sqrt{-g}\,F_{\mu\nu}F^{\mu\nu}\,\d^4x$。由于 $F=\d A$，反对称化时联络消失，仍有 $\cd_{[\lambda}F_{\mu\nu]}=0$。度规进入的是指标的缩并和体元。对 $A_\mu$ 变分得真空方程 $\cd_\mu F^{\mu\nu}=0$。不能简单把在平直背景中已交换过的高阶导数逐一换成协变导数。比如由这里的方程取 Lorenz 规范，可得矢势的波动方程，交换两次协变导数会出现与 Ricci 张量乘矢势成正比的项。因此判断一个曲率项是预设的非普适耦合，还是协变化既有方程的结果，应退回到作用量或未交换导数的一阶方程看。

在真空短波近似下，可写 $A_\mu=\operatorname{Re}\{(a_\mu+\cdots)\exp(\mathrm{i}\theta/\varepsilon)\}$，$\varepsilon$ 计数短波阶次，$k_\mu=\partial_\mu\theta$。Lorenz 条件在最高阶给出 $k^\mu a_\mu=0$，波动方程在 $\varepsilon^{-2}$ 阶给出 $k^\mu k_\mu=0$。又因 $\cd_\mu k_\nu=\cd_\nu k_\mu$，故 $k^\nu\cd_\nu k_\mu=\tfrac12\cd_\mu(k^\nu k_\nu)=0$，光线沿这同一度规的类光测地线走。下一阶给出 $2k^\nu\cd_\nu a^\mu+a^\mu\cd_\nu k^\nu=0$。以 $a^\mu=\mathcal A e^\mu$ 分开振幅和归一偏振，便有 $\cd_\mu(\mathcal A^2 k^\mu)=0$、$k^\nu\cd_\nu e^\mu=0$（偏振矢量只取其物理横向部分）。故光束越会聚，振幅通常越大。若只改变偏振仪相对于射线的方位，则须同时说明它的两根横轴沿途怎样搬运。场方程在一点按狭义相对论写成，传播一段有限距离却已含有曲率所造成的会聚和偏振基的改变。

频率的读出只需用\hyperlink{subsection.2.2.1}{\S\,2.2.1}的 $\omega=-u\cdot k$。静止端点可用 Killing 守恒量比较，运动端点则直接用各自的 4-速。这也说明在星地实验中把钟速与传播效应分开的便利虽依赖坐标，最终的收、发频率比却不依赖。

\subsubsection{电荷守恒来自哪一项？}

可以把按世界线写的粒子--场耦合改为一个时空积分。若 $z_a^\mu(\tau)$ 是带电粒子的世界线，定义
\[
 j^\mu(x)=\frac{1}{\sqrt{-g(x)}}\sum_a q_a
    \int u_a^\mu\,\delta^{(4)}(x-z_a(\tau))\,\d\tau .
\]
这里的 $\delta^{(4)}$ 是关于坐标测度 $\d^4x$ 的分布，除以 $\sqrt{-g}$ 才使积分给出不依赖坐标的电荷。于是 $\int\sqrt{-g}\,j^\mu A_\mu\,\d^4x=\sum_aq_a\int A_\mu\d z_a^\mu$。对 $A_\mu\mapsto A_\mu+\partial_\mu\chi$，每段世界线只贡献两端的 $q_a\chi$ 差。只要在时空内部的反应点上，进入和离开的总电荷相同，那么对任何有紧支撑的 $\chi$，所有端点差之和都为零。分部积分即给出 $\partial_\mu(\sqrt{-g}j^\mu)=0$，它既适用于粒子连续穿越区域，也适用于粒子数因成对产生而变化的反应。Gauss 定理和电荷的规范对称性已见 \S\ref{sec:Gauss} 与 \S\ref{sec:ED}，这里要分清守恒的是通量而非某个坐标密度。

取一张类空切片，单位未来法矢为 $n^\mu$，切片度量行列式为 $h$。切片上一群正电荷的总量为 $Q=-\int n_\mu j^\mu\sqrt h\,\d^3x$。同一世界管若无电流穿过侧面，便可用另一张切片来算同一个 $Q$。若改用实验者 $u$ 测电荷密度，他测的是 $-u_\mu j^\mu$，与上式中的 $-n_\mu j^\mu$ 可能不同。两位相对运动的观者量出不同的电荷密度，正如同一根带电棒被量出不同的长度。把电流与相应的有向三维面元一起积分，两者才可以比较。写成流体语言，这正是后 Newton 近似中守恒坐标密度与共动密度不可混作一个源的同一件事。

依本节的 $F_{\mu\nu}=\partial_\mu A_\nu-\partial_\nu A_\mu$ 及正号的 $j^\mu A_\mu$ 耦合，Maxwell 方程带源时为 $\nabla_\nu F^{\mu\nu}=4\pi j^\mu$。因为 $F^{\mu\nu}$ 反对称，还可单从这条场方程推出源电流散度为零。这是一项相容性条件，不能首先把电荷依赖地点假定为任意函数，然后还要求原来的场方程逐字不动。譬如强令一根自由粒子世界线的荷等于 $q(z)$，则规范变换后 $\int q\,\mathrm d\chi$ 除端点项外还剩 $-\int\chi\,\mathrm d q$。要使此模型成立，必须指明以何种新场或新电流补回这笔通量。前节 $e_*$ 随位置而变时，变化的是单位电荷之间的静电作用强弱，上述守恒电荷的定义并未被破坏。

\subsubsection{原子与基本场怎样测得同一光锥？}

不妨把量子测试物换成一束窄的带电粒子波包。仍取 $c=1$。设相位写成 $\exp(\mathrm{i}S/\hbar)$，定义 $p_\mu=\partial_\mu S-qA_\mu$。若物质作用量只与同一度规最小耦合，则标量场的带电 Klein--Gordon 方程在短波首阶给出 $g^{\mu\nu}p_\mu p_\nu=-m^2$。Dirac 场也一样。以度规号差 $(-+++)$ 为例，可选 $\{\gamma^\mu,\gamma^\nu\}=-2g^{\mu\nu}$ 的矩阵约定及普通 $\mathrm i\hbar\gamma^\mu D_\mu-m$ 算符。将短波形式代入，保留最高相位导数，得到作用在局域旋量振幅上的代数方程。再乘上反质量号的共轭因子，因 $p_\mu p_\nu$ 对称，矩阵乘积只余反对易子，便是同一条质量壳。旋量联络对振幅、偏振的输运在下一阶才出现。改变标架、曲率耦合的自旋效应及有限波包尺寸不在这个最低阶质心结果之内。

由质量壳对位置求导，利用 $\nabla_\nu p_\mu-\nabla_\mu p_\nu=qF_{\mu\nu}$，可进一步推出 $p^\nu\nabla_\nu p_\mu=qF_{\mu\nu}p^\nu$。当 $p^\mu=mu^\mu$ 时，便是按世界线作用量得到的 Lorentz 力方程。这里把短波条件用在相位相干的有限波束上，要求外场和包络的特征尺度均远大于波长。若检测的是一个束缚态跃迁，必须先求内部能谱，不能把整个原子当作一个具有单一自由粒子相位的波包。所以不同同位素的干涉仪检验波包质心的加速度，原子钟检验两个内态的能量差，尽管两者可用同一个实验室中的原子实现。

光子无质量且电中性，真空 Maxwell 方程的首阶特征式才是 $g^{\mu\nu}k_\mu k_\nu=0$。在物质介质中，电子也依同一度规耦合，但其集体响应可使电磁脉冲的群速度小于真空极限速度，甚至依频率、偏振而异。不能以玻璃的折射率或星际等离子体的色散直接否定普适光锥。对冷而稀薄、可略磁场的等离子体，局域色散关系是 $\omega^2=c^2k^2+\omega_p^2$，忽略曲率的微段上有 $v_g\simeq c[1-\omega_p^2/(2\omega^2)]$。因此沿路会积累正比于电子柱密度且近似反比于观测频率平方的时延。用多个频带先辨出这一频率依赖，正是以无线电光行时检验引力所需的实验操作之一。共同的局域非引力物理决定介质修正，比较理论时不能把它也归给真空光传播。

\subsection{连续物质与能动量的守恒}

把物质及其非引力场合记作 $\psi$。普适耦合可表为物质作用量 $S_{\rm m}[\psi,g]$ 的自变量中只有一个引力场，即物理度规 $g$。在物质场满足自身方程时，$S_{\rm m}$ 对任意具有紧支撑的微小坐标变换的变化均为零。按 $\delta S_{\rm m}=\tfrac12\int\sqrt{-g}\,T^{\mu\nu}\delta g_{\mu\nu}\,\d^4x$ 定义能动张量，用 $\delta g_{\mu\nu}=2\cd_{(\mu}\xi_{\nu)}$ 分部积分，即得 $\cd_\mu T^{\mu\nu}=0$，详细的 Noether 论证见 \S\ref{sec:Noe}。这里的 $T^{\mu\nu}$ 包含正参与交换能量、动量的一切非引力场。若只取带电物质、把电磁场留在外面，右端便有 Lorentz 力密度。带电物质失去的动量可由场带走，将二者相加才有上式。

书写成协变散度为零也不意味着有限区域里用各处坐标时间相加的能量必定守恒。将一指标降下并展开联络，有
\[
   \partial_\mu(\sqrt{-g}\,T^\mu{}_{\nu})
       =\tfrac12\sqrt{-g}\,T^{\alpha\beta}\partial_\nu g_{\alpha\beta}.
\]
若 $\partial_t g_{\alpha\beta}=0$，右端在 $\nu=0$ 的分量为零，才可把左端在一空间体积上积分，将内部的能量变化化为边界通量。一般情形应先找给定的对称性：若有 Killing 场 $\xi$，则 $J^\mu=T^{\mu\nu}\xi_\nu$ 的散度为零，给出相应的守恒荷。没有这种对称性，局域能动量平衡仍严格成立。这些关系由物质方程得到，换一套产生同一个度规的引力场方程，也不会改变其形式。

对连续物质，沿流线及其正交方向投影的推导见\hyperlink{subsection.1.2.3}{\S\,1.2.3}。若 $n$ 是共动粒子数密度、$\varepsilon$ 是固有能量密度，以 $u^2=-1$ 归一化，沿流线有 $\mathrm d(\varepsilon/n)=-p\,\mathrm d(1/n)$。横向方程给出的 4-加速度是 $a^\mu=-(g^{\mu\nu}+u^\mu u^\nu)\partial_\nu p/(\varepsilon+p)$。所以地面上静置的一柱流体即使坐标速度处处为零，仍有由压强梯度维持的非零 4-加速度。

对按静态 Killing 场静止的流体，取 $N=\sqrt{-g_{00}}$，则支撑加速度的空间协变分量为 $a_i=\partial_i\ln N$。上式给 $\partial_ip=-(\varepsilon+p)\partial_i\ln N$。恢复 $c$，在 $p\ll\rho c^2$、$\varepsilon\simeq\rho c^2$ 的弱场极限，有 $\partial_i p\simeq-\rho\partial_i\phi$，正是流体静力学。因而在有限高度内测压强差也能测当地重力，但压强与温度、状态方程都进入读数。若先把压强对能动张量的贡献略去，再去问它为什么支撑住一颗相对论性恒星，就已经在前提中排除了所求的效应。

\subsubsection{有限仪器中的场能与材料应力}

局域平衡式允许把场、吸收体和支架一起计算。仍用 Gaussian 单位，电磁能动张量为 $T_{\rm em}^{\mu\nu}=(F^{\mu\lambda}F^\nu{}_{\lambda}-g^{\mu\nu}F^{\alpha\beta}F_{\alpha\beta}/4)/(4\pi)$。利用场方程与 $\nabla_{[\alpha}F_{\beta\gamma]}=0$，其散度为 $\nabla_\mu T_{\rm em}^{\mu\nu}=-F^\nu{}_{\lambda}j^\lambda$。这一项的相反数正是带电物质的力密度。电荷束通过电容器时，仅对电荷积分会读出一个向外的动量改变。只对场积分则见动量减少。欲用实验检验总动量守恒，积分体积至少须囊括电场、电极、支撑它们的物体和侧面逸出的辐射。将 Maxwell 应力张量在边界上积分，所算的正是场传给或带出仪器的机械动量，不是另加一种电磁质量。

由此还可操作地解释为何统一的度规能致使不同内部组成的孤立小物体整体随同一局域重力下落。选一套以质心近似静止的局域 Lorentz 坐标来考察它，取无边界动量通量的平衡系统。在零阶平直背景下 $\partial_k T^{kj}=0$，乘 $x^i$ 对空间积分，表面积分消失即有 $\int T^{ij}\,\mathrm d^3x=0$。对于随内部周期运动而整体稳定的物体，这条关系取时间平均后同样成立。内部的电磁场或辐射气体单独通常并不满足它，例如各向同性辐射有 $T_{\rm rad}^{ij}=u_{\rm rad}\delta^{ij}/3$，须有容器壁的应力补回。一个对完整 $T^{\mu\nu}$ 的恒定小推动，积分所得总空间动量才等于总能量乘以速度再除以 $c^2$。若只取辐射场，动量响应看起来多了一个压强项，把它称为惯性质量异常就漏掉了墙的动量。

在引力中的局域平衡与这一积分论证并不冲突。若外场在系统大小 $L$ 上几乎恒定，用下落标架可把共同加速度从局域物质方程中消去，首个外场修正含潮汐与内部多极矩。若外场在 $L$ 上有可测曲率，或者侧面有热辐射、质量或内压流失，前段积分中的边界项就不能扔掉。受热的一块不对称辐射的金属板在真空中会反冲，一腔光气释放后质心也会变化。这些现象都可从能量、动量的流出量计算。比较普通结合能与引力结合能的分界就在此处：电磁、核和壁的应力可以归入 $S_{\rm m}$ 的总 $T^{\mu\nu}$，而自引力的能量并没有同样的普适局域物质张量。由前段论证把电荷容器换成强自引力的星，并不能自动证明它的全部结合能也有普适响应。

\subsection{边界条件与强等效原理}

非引力规律至此都由 $g$ 确定，却仍未回答同样的物质如何产生 $g$。在有作用量的理论中可把总作用量写成 $S=S_{\rm m}[\psi,g]+S_{\rm g}[g,\varphi,\ldots]$。若 $S_{\rm g}$ 除度规外还有标量 $\varphi$，它在物质的 Euler--Lagrange 方程中不必出现，可经度规的场方程影响 $g$。因此说度规理论的物质只与一个物理度规相耦合，并没有要求度规是唯一的动力学引力场。广义相对论与 Jordan 表述下的 Brans-Dicke 理论恰好在第一部分结构相同，第二部分却有差别。

也可在同一标量-张量理论中把度规重定义为 $\tilde g_{\mu\nu}=F(\varphi)g_{\mu\nu}$，令引力部分的曲率项取得较简单的形式。此时物质部分成为 $S_{\rm m}[\psi,F^{-1}\tilde g]$。例如质量为 $m$ 的粒子作用量可写为 $-m\int F^{-1/2}\d\tilde\tau$，在新变量下看起来像质量依赖于标量。两种表述只要同时变换尺、钟和耦合，就预言同一频率比与轨道。不能只更换度规变量，又默许钟仍直接读 $\tilde\tau$。判断 EEP 是否成立，应看能否找到一个度规，使一切非引力装置都以它标定读数。

现在在下落的封闭实验室里放两块可移动的引力源，做一次 Cavendish 实验。对实验自身的引力场，实验室必须大到可以包住源、试验体和它们之间的重要场区。对外部引力场，实验室又须小到潮汐可略。设仪器间距为 $d$，外场的变化尺度为 $\mathcal L$，可择一半径 $b$ 使 $d\ll b\ll\mathcal L$，在这个中间尺度给内场指定边界条件。这称为\textbf{准局域 Lorentz 系}。若先把一切物体都缩成同一个事件，仪器自身的引力也随之从问题中删去，便无从说明局域引力实验。

计算分两步。先略去实验室造成的小扰动，求出外部解在边界上给定的各场。再把两块源放入，在同一边界资料下解内区场方程。取随实验室下落的坐标，可令边界上的外部度规近似为 Minkowski 度规，所略去的是随潮汐而来的高阶不均匀性。广义相对论的外部条件至此已定，若实验室中有同样的质量、长度和时间标准，就会得到同样的局域引力规律。靠近恒星时固然可能受到相当大的外部潮汐，只要把它算入误差，或缩小设备便能减小。这与即使潮汐可忽略时局域引力本身仍依赖外界，是两种不同的效应。

在标量-张量理论中，还要给定 $\varphi$ 在边界上的值。标量的数值不能靠改变自由落体坐标消去。若它随宇宙演化或位置变化，室内两球在同一间距下的吸引力便可能不同。只要这个背景在实验期间及区域内可视为常数，单一标量值本身无方向，不能从中选出优选速度。标量若有不可忽略的梯度，则梯度已是另一个有方向的边界资料，以上取常数的近似须重新检查。若基本引力场还含有一个非零的类时矢量，它的边界值本身就选出一个静止系。实验室相对它的速度虽不会影响理想的局域非引力钟，却可能影响当地引力的强度和方向。后文 PPN 对优选系的检验，便是要从这类可能性中挑出可测的量。

有些理论还预先给定平直的背景度规或宇宙时间函数，它们不受局域物质与物理度规的反作用。纵使背景量处处有张量变换律，把物理度规化为 Minkowski 形式的同一套局域坐标一般也不能同时把两个度规化为同一标准形式。因此判断某个理论是否含有可观测的优选位置或方向，要看这些边界量怎样回到内场方程。若辅助场完全脱耦，方程中允许写出它，并不等于仪器一定测得到它。

从原则上说，用作用量表述还可解释为什么辅助引力场与物质的间接耦合不会破坏刚才的物质守恒律。取 $g_{\mu\nu}$ 和标量 $\varphi$ 为独立的场变量，若 $S_{\rm m}$ 没有显式依赖 $\varphi$，则在求引力场方程时物质通过 $\delta S_{\rm m}/\delta g_{\mu\nu}$ 出现在度规方程，标量方程只有引力部分的变分。消去度规方程之后，标量仍可间接地由物质迹 $T$ 激发。当写成别一组变量时，标量方程会显式有源，物质应力也可能相对于新联络显得不守恒。须把表观的附加力和换过的尺钟单位一起换回，才有可比的两种答案。并非所有候选理论都有一个完备的局域引力作用量。有时引力场方程单独给定，但物质仍由统一物理度规耦合。此前从 $S_{\rm m}$ 及非引力场方程得到的局域守恒律并不因此失效，孤立引力系统能否另有全局守恒量则需单查其引力动力学。

最后把实验室的试验体换成一颗有可观自引力的星。组成星体的粒子可以逐个遵守普通的 WEP。然而星的总质量中还有把它束缚起来的引力能。若当地的有效引力常数改变，保持重子数不变而重新求平衡，半径、结合能和总质量都可能改变。用这样两颗内部结构不同的星比较自由落体，比用两种金属比较又多检验了一项。SEP 同时要求这种自引力物体的普适自由落体，以及准局域引力实验对地点、时刻和速度的不变性。前者可在地月轨道中查，后者还可由重力仪、行星历表和双星轨道查。要把这些读数同各理论对比，下一步须把它们在弱场、低速下的差异写进一套共同的度规。



\subsubsection{引力结合能的循环与唯一性问题}

能否像讨论电磁束缚物一样，从大物体的自由落体推知局域引力实验的普适性？试取两个小球，由一根应力已计入的支架把它们保持在给定固有距离 $r_{\rm p}$。若在某处的局域 Cavendish 常数为 $G_{\rm L}$，两球从相距无穷远缓慢靠近释放的引力能在最低阶为 $G_{\rm L}m_1m_2/r_{\rm p}$。假设两物质的非引力质量均遵守 EEP，只有 $G_{\rm L}$ 随外场改变。将结合系统缓慢运往另一处，然后拆散，再将零散部件各自运回。闭合循环中，另一处拆散所作的功与原处结合放出的功之差，必须在全部运输过程中补回。如果结合态的运输和平移因其自引力而有不同的被动质量，则这项差异有去处。若一方面令所有自引力程度不同的物体的惯性质量与被动质量的比例都完全一样，另一方面又任意指定随位置变的结合能，通常不能同时保持循环能量守恒。

为量出这个效应，可令 $G_{\rm L}$ 在实验室位置附近缓变。取各部件静质量与固有尺寸暂固定，给定外势变动 $\delta U$ 会令 $\delta E_{\rm bind}=(m_1m_2/r_{\rm p})\delta G_{\rm L}$。因而对整体有效被动质量的改变量带有 $E_{\rm bind}\,\partial\ln G_{\rm L}/\partial U$，对两个结合能占比不同的系统会不一样。真实的星没有刚性支架，给定重子数后须求出它在不同背景场中的新平衡，再按总能量的变化定义响应。若背景还选出方向，则 $G_{\rm L}$ 可同时依赖仪器轴与方向的夹角。旋转装置也要计入异常力矩作的功。在一般 PPN 框架中，后文会求出这些可辨的方向项。这个循环为自引力的等效原理给出同类的联系，却仍假定保守动力学和可解的平衡族，不能作为对任意引力理论的无条件定理。

如果恰好只有度规传递长程引力，又能不能由 SEP 单独导出 Einstein 方程？在四维的一个较窄问题中，可要求场方程的几何端仅含度规及至多二阶的导数，满足对称性、协变性和恒等协变无散，且具有适当的局域性、正则性。此时的唯一性结果才使 Einstein 张量和 $g_{\mu\nu}$ 的常系数组合受到偏爱。其比例还要由 Newton 极限定出，另有宇宙学常数待测。若允许高阶导数或其它动力学场，这一推理的前提已变。另一条路线从平直背景上一种无质量二自旋场的线性耦合出发，要求它普适地耦合于守恒应力，进而把场自身贡献也逐次纳入源。这种自耦合构造同样引入平直极限、局域作用量等假定，详见 \S\ref{sec:Einstein} 对场方程的讨论。因此是实验先问谁与尺钟耦合，再用场动力学和弱场参数比较候选理论，不能倒过来因某种理论实现了 SEP 就先免除这些比较。

\section{后 Newton 理论}\label{sec:exp-pn}

\hyperlink{subsection.2.2.2}{\S\,2.2.2} 从 Schwarzschild 解求得水星的额外进动。若把其它行星的引力也算进去，场既不球对称，也不静止。所幸在水星轨道上 $GM_\odot/(rc^2)\simeq2.5\times10^{-8}$，行星的速度也远小于光速。可从 Newton 多体解出发，把场和运动同时修正。这称为\textbf{后 Newton 近似}(post-Newtonian approximation, PN)。以下沿用式 \eqref{eq:h-g-phi} 的符号约定，另记正值势 $U=-\phi$，取 $x^0=ct$。
太阳既自转也脉动，其外场还含有质量多极矩。自转轴与黄道面法线约成 $7.25^\circ$，与银盘面法线约成 $67.23^\circ$。光球并非刚体。相对于远星，赤道的自转周期约 $25.05$ 日，近两极约 $34.4$ 日，赤道线速度约 $7.19\times10^3\,\mathrm{km/h}$。因自转而生的质量四极矩已会改变 Newton 轨道。同一自转的质量流还会产生相对论性的拖曳。定轨时须分开两项。

% 黄 §3.12.1，式 (3.12.1)–(3.12.5)、表 3.12.1。Weinberg §9.1。
\subsection{近区与计数}

取源的尺度为 $R$，质量为 $M$，典型速度为 $v$。束缚系统的维里估计给出 $v^2\sim GM/R$。圆轨道上就是 $v^2=GM/R$。记 $\epsilon\sim v^2/c^2\sim U/c^2\ll1$，Newton 阶记为 0PN，相对大小为 $O(\epsilon^n)$ 的修正称为 $n$PN。对非束缚粒子须分别检查速度与场强。源演化的时间约为 $R/v$，故 $\p_i=O(1/R)$，$\p_0=c^{-1}\p_t=O(\sqrt\epsilon/R)$。静力平衡给出 $p/\rho\sim GM/R$，普通星体的比内能 $\Pi$ 通常亦为 $O(v^2)$。相对论性气体即使质心静止，也不能这样计数。

记 $q^\mu=\d x^\mu/\d t$，取 $q^0=c$。以 $\lambda$ 为仿射参数，在式 \eqref{eq:general-coordinate} 中用时间分量消去 $\d^2t/\d\lambda^2$，空间方程化为
\begin{equation}\label{eq:pn-coordinate-geo}
 \frac{\d q^i}{\d t}=-\left(\Gamma^i{}_{\mu\nu}
                  -\frac{q^i}{c}\Gamma^0{}_{\mu\nu}\right)q^\mu q^\nu .
\end{equation}
其中沿速度的一项来自重选参数。慢质点的坐标加速度须由 $v^2/R$ 算到 $v^4/(Rc^2)$。在固有时根号中，$g_{00}$、$g_{0i}$、$g_{ij}$ 各乘 $1$、$2v^i/c$、$v^iv^j/c^2$，故相应需要四、三、二阶度规。这里以 $v/c$ 的幂数说阶次。保守无自旋源在时间反演下 $g_{00},g_{ij}$ 为偶，$g_{0i}$ 为奇。静质量不产生一阶 $g_{0i}$，质量流首次在三阶出现。四阶 $g_{00}$ 中度规一次扰动的平方与动能修正同阶，必须保留。按 $G$ 展开的\textbf{后 Minkowski 近似}则不要求物质低速。\S\ref{sec:Einstein} 的线性弱场结果可在弱场慢运动极限与 PN 展开比较。

场点距源 $d$ 处的光行时约为 $d/c$。若源的特征角频率为 $\Omega$，则\textbf{近区}要求 $R,d\ll c/\Omega$，方可展开推迟时刻。孤立系统另有一个条件。随质心自由落体可消去外场的共同加速度，外部潮汐在尺度 $R$ 上留下的相对加速度 $|\p_i\p_jU_{\rm ext}|R$ 若与所求修正相当，仍须计入。对太阳系，可在源势已小、宇宙演化又可忽略处取不转动的渐近惯性轴，归一化钟和尺。求波时还须在远处指定无入射辐射。这项边界条件选出推迟解，决定耗散项的符号。近、远区的匹配见 \S\ref{sec:pn-matching}。

例如平直空间切片的膨胀宇宙，令 $a(t_0)=1$、$H_0=(\dot a/a)_0$、$q_0=-(a\ddot a/\dot a^2)_0$。宇宙空间度规中的 $2H_0(t-t_0)$ 可由局域坐标变换消去。按 \S\ref{sec:tidal} 的约定，随质心自由落下的正交轴上有 $R_{\hat0 i\hat0 j}=q_0(H_0/c)^2\delta_{ij}$、$R_{\hat0 ijk}=0$。$g_{00}$ 和 $g_{ij}$ 的首项背景修正分别为 $-R_{\hat0 i\hat0 j}X^iX^j$ 和 $-R_{ikjl}X^kX^l/3$。当地可测量的是 $O((H_0d/c)^2)$ 的潮汐。若远处未计的单极项约为 $\epsilon GM/(c^2d)$，令两种误差相当，得 $d_0^3\sim\epsilon GM/H_0^2$。取 $M=M_\odot$、$\epsilon=10^{-6}$、$H_0\simeq2.3\times10^{-18}\ {\rm s}^{-1}$，有 $d_0\sim3\times10^{13}\ {\rm km}$，度规误差约 $5\times10^{-20}$。邻星和银河的潮汐须按各自的梯度另估。

用静止球体检验度规的阶次。由\hyperlink{section.2.1}{\S\,2.1} 的 Schwarzschild 解，令 $m=GM/c^2$、$u_R=m/R$，$R$ 为面积半径，则 $g_{00}=-1+2u_R$，$g_{RR}=1+2u_R+4u_R^2+\cdots$。在 $x^i=Rn^i$ 中，后一修正只乘 $\d R^2=(n_i\d x^i)^2$。改用各向同性半径 $\varrho$，记 $u=m/\varrho$，则
\[
 g_{00}=-1+2u-2u^2+O(u^3),\qquad
 g_{ij}=\big(1+2u+\tfrac32u^2\big)\delta_{ij}+O(u^3).
\]
两式中 $g_{00}$ 的二次项不同，源于 $R=\varrho(1+u/2)^2$。空间度规中的二次项虽也能写出，计算慢体的 1PN 轨道只需其中的一次项。

% 黄式 (3.12.6)–(3.12.31)。Weinberg §9.1，式 (9.1.4)–(9.1.56)。
\subsection{谐和坐标中的场方程}\label{sec:pn-components}

写 $g_{00}=-1+h+k+\cdots$、$g_{0i}=b_i+\cdots$、$g_{ij}=\delta_{ij}+s_{ij}+\cdots$，其中 $h,s_{ij}$ 为二阶，$b_i$ 为三阶，$k$ 为四阶。对于物质，写 $T^{00}=\rho_0c^2+e+\cdots$、$T^{0i}=cj_i+\cdots$、$T^{ij}=\sigma_{ij}+\cdots$。$\rho_0$ 是最低阶质量密度，$j_i$ 是质量流，$\sigma_{ij}$ 是最低阶动量通量。相对于 $\rho_0c^2$，$cj_i$ 为一阶，$e,\sigma_{ij}$ 为二阶。

由 $g^{\mu\lambda}g_{\lambda\nu}=\delta^\mu_\nu$ 逐阶求逆，得 $g^{00}=-1-h-k-h^2+O(c^{-6})$、$g^{0i}=b_i+O(c^{-5})$、$g^{ij}=\delta^{ij}-s_{ij}+O(c^{-4})$。联络和 Ricci 分量同样按时间指标的奇偶分阶。

以逗号表示偏导，以方括号中的数表示 $v/c$ 的幂次。将上述逆度规代入联络公式 \eqref{eq:metric-christ}，需要的分量为
\begin{align*}
 (\Gamma^i{}_{00})_{[2]}&=(\Gamma^0{}_{0i})_{[2]}=-\tfrac12h_{,i},&
 (\Gamma^0{}_{00})_{[3]}&=-\tfrac12h_{,0},\\
 (\Gamma^i{}_{jk})_{[2]}&=\tfrac12(s_{ij,k}+s_{ik,j}-s_{jk,i}),&
 (\Gamma^0{}_{ij})_{[3]}&=\tfrac12(s_{ij,0}-b_{i,j}-b_{j,i}),\\
 (\Gamma^i{}_{0j})_{[3]}&=\tfrac12(s_{ij,0}+b_{i,j}-b_{j,i}),&
 (\Gamma^i{}_{00})_{[4]}&=b_{i,0}-\tfrac12 k_{,i}+\tfrac12s_{ij}h_{,j}.
\end{align*}
例如最后一项来自逆空间度规乘低阶时间势。在 \S\ref{sec:Einstein} 的谐和坐标条件中，保留最先出现的各项，得
\[
 C_0:=\tfrac12(h+s_{jj})_{,0}-b_{j,j}=0,\qquad
 C_i:=s_{ij,j}+\tfrac12(h-s_{jj})_{,i}=0 .
\]
写 $\Delta=\p_i\p_i$。直接缩并曲率，$2(R_{ij})_{[2]}=-\Delta s_{ij}+C_{i,j}+C_{j,i}$，$2(R_{0i})_{[3]}=-\Delta b_i+C_{i,0}-C_{0,i}$。四阶的 $R_{00}$ 还含二阶联络的乘积。消去含 $C_\mu$ 及其导数的项后，有 $2(R_{00})_{[2]}=-\Delta h$、$2(R_{ij})_{[2]}=-\Delta s_{ij}$、$2(R_{0i})_{[3]}=-\Delta b_i$。四阶时间分量为 $2(R_{00})_{[4]}=-\Delta k+h_{,00}+s_{ij}h_{,ij}-h_{,i}h_{,i}$。例如 $2b_{i,i0}-s_{ii,00}=h_{,00}$ 用到了 $C_{0,0}=0$。$s_{ij,i}=(s_{ii}-h)_{,j}/2$ 则消去了 $s_{ii,j}h_{,j}$。这样保留下来的梯度平方与 $\Delta k$ 同阶。

场方程右边是 $S_{\mu\nu}=T_{\mu\nu}-g_{\mu\nu}T/2$。用上指标的物质变量，记 $A=\rho_0c^2$。则 $T_{00}=A+e-2hA+\cdots$，$T=-A-e+hA+\sigma_{jj}+\cdots$。一个坐标静止的观者有单位切矢 $n^0=(-g_{00})^{-1/2}$、$n^i=0$，所测能量密度为 $T_{00}/(-g_{00})=A+e-hA+\cdots$。换成观者的钟后，留下的势能项与单看 $T_{00}$ 时不同。逐分量取迹反转得
\begin{align*}
 S_{00}&=\tfrac12(A+e+\sigma_{kk}-2hA)+\cdots,\\
 S_{0i}&=-cj_i-T^{0i}_{(3)}+c(h\delta_{ij}-s_{ij})j_j-\tfrac12Ab_i+\cdots,\\
 S_{ij}&=\tfrac12 A\delta_{ij}+\sigma_{ij}+\tfrac12As_{ij}
           +\tfrac12(e-hA-\sigma_{kk})\delta_{ij}+\cdots .
\end{align*}
这里 $T^{0i}_{(3)}$ 是 $T^{0i}$ 的三阶项。在 $R_{0i}$、$R_{ij}$ 的场方程中，此次只需 $S_{0i}$、$S_{ij}$ 的首项。

% 黄式 (3.12.32)–(3.12.45)。Weinberg §9.1 后半及 §9.2。
\subsection{三个势与运动的钟}\label{sec:pn-three}

现在解反迹场方程。记 $s=|\bm x-\bm x'|$，对任意源密度 $f$ 定义 $\mathcal I[f]=G\int f(t,\bm x')\,\d^3x'/s$。最低两个势为 $U=\mathcal I[\rho_0]$、$V_i=\mathcal I[j_i]$。因 $\nabla^2\mathcal I[f]=-4\pi G f$，在 $R_{\mu\nu}=8\pi G S_{\mu\nu}/c^4$ 中分别取 $00,ij$ 的二阶及 $0i$ 的三阶，衰减解即为 $h=2U/c^2$、$s_{ij}=h\delta_{ij}$、$b_i=-4V_i/c^3$。此时空间谐和条件自动满足，时间条件化为 $\dot U+\bm\nabla\cdot\bm V=0$。对它取 Laplace 算符，所得正是 $-4\pi G(\dot\rho_0+\p_i j_i)=0$。孤立边界排除了附加的调和函数。

四阶时间方程还含有 $|\bm\nabla U|^2$。用 $2|\bm\nabla U|^2=\nabla^2(U^2)-2U\nabla^2 U$，可将非线性部分抽出，置 $k=2(Q-U^2)/c^4$。余下的方程及其近区解为
\begin{align}
 \nabla^2Q&=\ddot U-4\pi G(e+\sigma_{ii}),&
 Q&=\mathcal I[e+\sigma_{ii}]+\tfrac12\ddot X,\qquad
 X=G\int\rho_0(t,\bm x')s\,\d^3x'.\label{eq:pn-Q}
\end{align}
这是因为 $\nabla^2s=2/s$，故 $\nabla^2X=2U$。还可直接展开推迟核。$\mathcal I[\rho_0]$ 的一阶推迟改正为 $-(G/c)\d(\int\rho_0\,\d^3x')/\d t$，孤立系统的总质量在此阶守恒，故为零。二阶改正为 $\ddot X/(2c^2)$。$X$ 本身在远处随 $r$ 增长，进入近区度规的却是它的时间导数。至此，任意慢源的谐和度规可写为
\begin{equation}\label{eq:pn-harmonic}
 \begin{aligned}
 g_{00}&=-1+\frac{2U}{c^2}+\frac{2(Q-U^2)}{c^4}+O(c^{-6}),\\
 g_{0i}&=-\frac{4V_i}{c^3}+O(c^{-5}),\qquad
 g_{ij}=\left(1+\frac{2U}{c^2}\right)\delta_{ij}+O(c^{-4}).
 \end{aligned}
\end{equation}
若用 Weinberg 的几何单位，这三个势分别是 $\phi=-U$、$\bm\zeta=-4\bm V$、$\psi=-Q$。下节的 $\psi$ 则专指标准 PPN 坐标中的标量势。静止球体外 $\bm V=0$，但星内 $e+\sigma_{ii}$ 仍可给 $Q$ 一个 $1/r$ 项。用外部轨道测得的质量吸收此项后，$g_{00}$ 的显式二次项才只剩 $-2U^2/c^4$。均匀流体球的压强、内能与结合能怎样合成这一质量，见 \S\ref{sec:pn-binding}。

考虑在场中运动的钟。令 $\bm v=\d\bm x/\d t$，将式 \eqref{eq:pn-harmonic} 代入线元，按 $c^{-1}$ 展开。扣去静能后的单位质量拉氏量 $\ell=c^2(1-\d\tau/\d t)$ 为
\begin{equation}\label{eq:pn-test-L}
 \ell=\frac{v^2}{2}+U+\frac{1}{c^2}
   \left(\frac{v^4}{8}+\frac32Uv^2-\frac12 U^2+Q-4\bm v\cdot\bm V\right)
   +O(c^{-4}).
\end{equation}
地面钟和卫星钟相比较，势与速度在 $\d\tau/\d t$ 中同阶，地面处 $U$ 较大，卫星的 $v^2/2$ 通常较大，两项对钟速差的贡献可有相反的符号。若用电磁信号比钟，还须另算光路。对自由下落的钟，时间测地线还给 $\d\ln(\d t/\d\tau)/\d t=(\p_t U+2\bm v\cdot\bm\nabla U)/c^2+O(c^{-4})$。其中一个 $\bm v\cdot\bm\nabla U$ 来自沿轨道求势的导数，另一个来自动能随 Newton 加速度的改变。

对质点用式 \eqref{eq:pn-test-L} 作变分，$D_t=\p_t+\bm v\cdot\bm\nabla$ 沿轨道作用于势，保留至 1PN 得
\begin{align}
 \dot{\bm v}={}&\bm\nabla U+\frac{1}{c^2}\big\{
 (v^2-4U)\bm\nabla U+\bm\nabla Q+4\p_t\bm V
 -4\bm v\times(\bm\nabla\times\bm V)\nonumber\\[-2pt]
 &\hspace{37mm}-\bm v(3\dot U+4\bm v\cdot\bm\nabla U)\big\}+O(c^{-4}).
 \label{eq:pn-test-a}
\end{align}
式中 $\dot U=\p_tU$ 表示场点固定。只留最低阶，$U=GM/r$ 给出 $\dot{\bm v}=-GM\bm n/r^2$，$\bm n=\bm x/r$，方向指向场源。$\bm V$ 的时间导数作用于瞬时静止的粒子，旋度则乘以粒子速度。后文式 \eqref{eq:pn-general-motion} 将这两种作用推广到一般 PPN 度规。

光子的速度在零阶就是 $c$，须回到零测地线。记 $\bm u=\d\bm x/\d t$、$\bm n=\bm u/|\bm u|$。零线元给出 $|\bm u|/c=1-2U/c^2+O(\epsilon^{3/2})$。在式 \eqref{eq:pn-coordinate-geo} 中只取场的二阶项，得横纵两向的坐标加速度 $\dot{\bm u}=2\bm\nabla U-4\bm n(\bm n\cdot\bm\nabla U)$。其纵向分量恰为坐标光速沿程的改变量。以 $\ell_{\rm ray}$ 表示坐标弧长，横向分量给 $\d\bm n/\d\ell_{\rm ray}=2[\bm\nabla U-\bm n(\bm n\cdot\bm\nabla U)]/c^2$。两个固定端点的积分及与时延的关系见 \S\ref{sec:pn-light}，球对称场的精确测地线见 \hyperlink{subsection.2.2.3}{\S\,2.2.3}。

% Weinberg §§9.3、9.8。热流体所需的额外联络见式 (9.1.73)–(9.1.75)。
\subsection{物质的守恒与逐次解法}\label{sec:pn-matter}

场作用于源，源又决定场。为看清求解的次序，先不指定物质模型。将 \S\ref{sec:Einstein} 的 $\cd_\mu T^{\mu\nu}=0$ 按阶分开，最低两式为 $\dot\rho_0+\p_i j_i=0$、$\dot j_i+\p_j\sigma_{ij}=\rho_0 U_{,i}$。前式已保证谐和条件，后式是 Newton 引力下的动量平衡。写 $T^{0i}=cj_i+F_i/c+\cdots$、$T^{ij}=\sigma_{ij}+\Sigma_{ij}/c^2+\cdots$，将同一守恒方程再算一阶，得
\begin{align}
 \dot e+\p_iF_i={}&-\rho_0\dot U,\nonumber\\
 \dot F_i+\p_j\Sigma_{ij}={}&
 \rho_0[\p_i(Q-2U^2)+4\dot V_i]+(e+\sigma_{kk})U_{,i}
 -4\sigma_{ij}U_{,j}\nonumber\\[-2pt]
 &+4j_j(\p_jV_i-\p_iV_j)-4j_i\dot U .\label{eq:pn-matter-balance}
\end{align}
可用时间方程核对符号。$\sqrt{-g}=1+2U/c^2+O(c^{-4})$，体积因子对能流给出的 $2j_iU_{,i}$ 被 $\Gamma^0{}_{0i}$ 项抵消，余下 $-\rho_0\dot U$。这里守恒的是物质的协变能动量。在时变的外场中，物质的坐标能量本身可以改变。逐阶求系数时，$F_i,\Sigma_{ij}$ 也包括低阶源随轨道修正产生的改变量。若用未展开的密度和速度作未知量，应将相邻两阶的守恒方程相加。两种写法都只需上一步的场。

以只有引力相互作用的粒子为例。把 \S\ref{sec:Einstein} 的点粒子能动张量在 $t=\text{常数}$ 面上展开，令 $\delta_A=\delta^3(\bm x-\bm x_A)$，则
\[
 \rho_0=\sum_A m_A\delta_A,\quad j_i=\sum_A m_A v_A^i\delta_A,\quad
 \sigma_{ij}=\sum_A m_A v_A^iv_A^j\delta_A,\quad
 e=\sum_A m_A(\tfrac12v_A^2-U_A)\delta_A .
\]
最后一式的 $-U_A$ 来自 $(\d t/\d\tau_A)/\sqrt{-g}$，$U_A$ 只计其它天体在 $A$ 处产生的势。由 $\p_t\delta_A=-v_A^i\p_i\delta_A$，质量方程恒满足，动量方程则逐体给 $\dot{\bm v}_A=\bm\nabla U_{\rm ext}(\bm x_A)$。局域弹性碰撞可用碰撞前后的总能动量匹配。无碰撞时，先解 Newton 多体问题，用所得位置和速度构成 $j_i,e,\sigma_{ij}$，再由式 \eqref{eq:pn-Q} 求 $Q$，将 $Q,\bm V$ 代回式 \eqref{eq:pn-test-a} 更新轨道。点源的自场只在取外场时排除。每个天体自身的结合能已计入它的质量。对有限体先积分再取单极极限，见 \S\ref{sec:pn-finite}。

若以流体描述有大小的天体，月球的 Newton 潮汐即可超过地球场的 1PN 修正。沿用 \hyperlink{subsection.1.2.3}{\S\,1.2.3} 的理想流体张量与粒子数流，$\rho$ 是固有静质量密度，$\mathcal E=\rho(c^2+\Pi)$ 是共动能量密度。取 $u^\mu=\d x^\mu/\d\tau$、$D_t=\p_t+v^j\p_j$，则 $u^i=v^iu^0/c$。\textbf{守恒密度} $\rho_*=\rho\sqrt{-g}\,u^0/c$ 满足 $\dot\rho_*+\bm\nabla\cdot(\rho_*\bm v)=0$，随流体元运动的 $\int\rho_*\,\d^3x$ 恰为重子质量。若以 $\rho_0=\rho_*$ 定义 $U$，由钟速及体积因子有 $\rho_*/\rho=1+(v^2/2+3U)/c^2+O(c^{-4})$。于是
\[
 j_i=\rho_*v_i,\qquad \sigma_{ij}=\rho_*v_iv_j+p\delta_{ij},\qquad
 e=\rho_*(\Pi+v^2/2-U).
\]
这使式 \eqref{eq:pn-Q} 的瞬时源成为 $\rho_*(\Pi+3v^2/2-U)+3p$。若改以固有密度定义势，$2U/c^2$ 自身要补上一项四阶改正，两种密度不能只换其一。理想流体的全部 $F_i,\Sigma_{ij}$ 由 \S\ref{sec:pn-source-stress} 的分量式取 $\gamma=1$ 即得。

冷的简并流体可取单参数状态方程 $p=p(\mathcal E)$。先以连续性方程和 $\rho D_t\bm v=\rho\bm\nabla U-\bm\nabla p$ 求 Newton 解，按该状态方程更新各势的源。具体令 $d=\mathcal E/c^2$、$W=d+[p+d(v^2+2U)]/c^2$。此时 $T^{00}/c^2=W-p/c^2$、$T^{0i}/c=Wv^i$，动量通量为 $Wv^iv^j+p(1-2U/c^2)\delta_{ij}$，各留到其首个修正。将能量的两阶相加，便有 $\p_t(W-p/c^2)+\bm\nabla\cdot(W\bm v)=-d\dot U/c^2$。将同一动量通量代入式 \eqref{eq:pn-matter-balance} 的空间方程并加上 Newton 项，即可推进 $\bm v$。等价的沿流线写法见式 \eqref{eq:exp-fluid-p}，其中取 $\gamma=1$、$B_i=-4V_i$、$K=2(Q-U^2)$，即回到这里的谐和坐标。

温度若是独立变量，$p=p(\rho,\Pi)$，还需求热能。为从时间方程中扣去重子静能，置 $E_{\rm th}=T^{00}-c\rho u^0$、$F^i_{\rm th}=cT^{i0}-c^2\rho u^i$。前者含有扣除静能后的动能与引力项，并非单指 $\rho\Pi$。将能量守恒减去 $c^2$ 倍的重子守恒，得
\[
 \p_t(\sqrt{-g}\,E_{\rm th})+\p_i(\sqrt{-g}\,F^i_{\rm th})
 =-c\sqrt{-g}\,\Gamma^0{}_{\alpha\beta}T^{\alpha\beta}.
\]
热能相对静能小两阶。故在 Newton 步，连续、动量、去静能的时间方程分别留到速度的一、二、三次，在 1PN 步分别留到三、四、五次。为求右端，除前面的联络外只须补上
\[
 (\Gamma^0{}_{00})_{[5]}=-\tfrac12(k_{,0}+h h_{,0}+b_i h_{,i}),\qquad
 (\Gamma^0{}_{0i})_{[4]}=-\tfrac12(k_{,i}+h h_{,i}).
\]
$\Gamma^0{}_{ij}$ 乘以本已二阶的应力，取三阶即够。配上状态方程及重子守恒，各步的未知量与独立方程数相合。无黏绝热时，沿流线的热方程也可用 $D_t\Pi=pD_t\rho/\rho^2$ 表示。

% Weinberg §9.4。
\subsection{有限源的外场}\label{sec:pn-source-moments}

设物质限于半径 $b$ 内，场点在源外近区 $b\ll r\ll c/\Omega$。令 $\bm n=\bm x/r$，则 $1/|\bm x-\bm x'|=1/r+\bm n\cdot\bm x'/r^2+O(b^2/r^3)$。任一紧支集密度 $s$ 的势，前两项只取决于单极 $M[s]=\int s\,\d^3x$ 与偶极 $\bm D[s]=\int\bm x s\,\d^3x$。例如 $U$ 的两矩由 $\rho_0$ 给出。矢量势各分量也可这样展开，
\[
 V_i=\frac{GP_i}{r}+\frac{G n_j}{r^2}\int x'^j j_i'\,\d^3x'+O(r^{-3}),
 \qquad \bm P=\int\bm j\,\d^3x .
\]
连续性方程给出 $\dot M[\rho_0]=0$、$\dot{\bm D}[\rho_0]=\bm P$。源定常时 $\bm P=0$，对 $\int x^ix^j\rho_0\,\d^3x$ 求导又知流的一次矩是反对称的。它的三个分量构成自转角动量 $\bm J=\int\bm x\times\bm j\,\d^3x$，故 $\bm V=G\bm J\times\bm x/(2r^3)+O(r^{-3})$。质量偶极取决于原点，质心静系中的自转偶极却不能用平移消去。

在 $g_{00}$ 中可将 $U+Q/c^2$ 一起作多极展开，因为 $U^2$ 中再作这一替换已超出 1PN。其紧支集源为 $s=\rho_0+(e+\sigma_{kk})/c^2$。若源定常，保留到偶极阶，标量外场可写为 $G M[s]/|\bm x-\bm D[s]/M[s]|$。故外部质量含能量与应力的修正，将原点移至该有效质量的中心即可消去标量偶极。源随时间改变时，还须加上 $\ddot X/(2c^2)$。此项可直接将 $|\bm x-\bm x'|=r-\bm n\cdot\bm x'+[x'^2-(\bm n\cdot\bm x')^2]/(2r)+\cdots$ 代入 $X$ 后求导。孤立源的总质量守恒、质心匀速，故头两项的二阶时间导数为零。记 $z=\bm n\cdot\bm x'$，则 $X$ 在 $1/r$、$1/r^2$ 阶所取的矩分别为 $\int\rho_0'(x'^2-z^2)\,\d^3x'$ 和 $\int\rho_0' z(x'^2-z^2)\,\d^3x'$，系数各为 $G/2$。因此即使只读 $1/r$ 项，随时间变化的二阶矩也会出现。

对球对称的标量源，角平均 $\langle1/|\bm x-\bm x'|\rangle_{\bm n'}=1/\max(r,r')$。因而定常球体的标量外场直至表面都只有单极。谐和半径 $r$ 与面积半径在 Schwarzschild 外场中相差 $GM/c^2$，据此可用\hyperlink{section.2.1}{\S\,2.1} 的精确解核对 $g_{00}$ 中的平方项。即使球体径向脉动，其球对称真空外部仍由该节的 Birkhoff 定理决定。

设球体各层以角速度 $\bm\omega(r')$ 缓慢转动，保留转速的一次项，则 $\bm j'=\rho_0(r')\bm\omega(r')\times\bm x'$。用 $\langle\bm x'/|\bm x-\bm x'|\rangle_{\bm n'}=\bm x r'^2/[3\max(r,r')^3]$ 积分，有
\[
 \bm V(\bm x)=\frac{G}{3}\!\int\!
 \frac{r'^2\,\bm\omega(r')\times\bm x}{\max(r,r')^3}\,\d m',
 \qquad
 \bm J=\frac23\int r'^2\bm\omega(r')\,\d m',\qquad
 \d m'=4\pi\rho_0(r')r'^2\d r'.
\]
所以即使表面与内层的转速不同，整个球外区域的 $\bm V$ 都由一个 $\bm J$ 给定。若物质只在壳层上，空腔内则有 $\bm V=(G/3)[\int(\bm\omega/r')\,\d m']\times\bm x$，其旋度处处相同。

% Weinberg §9.5。Icarus 数值例据黄书第 3 章习题 8。
\subsection{二体运动与近心点}\label{sec:pn-peri}

两体质量为 $m_1,m_2$，记 $m=m_1+m_2$、$\mu=m_1m_2/m$、$\eta=\mu/m$。以 $\bm x=\bm x_1-\bm x_2$ 为相对位置，令 $r=|\bm x|$、$\bm n=\bm x/r$、$\bm v=\dot{\bm x}$、$\dot r=\bm n\cdot\bm v$。Einstein、Infeld 与 Hoffmann (EIH) 原从天体周围的真空场方程反求运动。这种边界积分与流体算法的关系见 \S\ref{sec:pn-matching}。在质心谐和坐标中可直接写相对拉氏量。将式 \eqref{eq:exp-nbody-lag} 取广义相对论的参数并去掉质心运动，得
\begin{equation}\label{eq:exp-rel-L}
 \frac{L}{\mu}=\frac{v^2}{2}+\frac{Gm}{r}
 +\frac{1}{c^2}\left\{\frac{1-3\eta}{8}v^4+
 \frac{Gm}{2r}\left[(3+\eta)v^2+\eta\dot r^2-\frac{Gm}{r}\right]\right\}
 +O(c^{-4}).
\end{equation}
变分后，式 \eqref{eq:g3-relative-ppn} 中乘 $Gm/r,v^2,\dot r^2,\dot r\bm v$ 的四个系数依其记号为 $A=4+2\eta$、$B=1+3\eta$、$C=3\eta/2$、$D=4-2\eta$。$\eta\to0$ 还原单粒子方程。对圆轨道，1PN 项中可用 $v^2=Gm/r$，故 $\Omega^2=(Gm/r^3)[1-(3-\eta)Gm/(rc^2)]$。由时间平移守恒，将此式代入 $E=\bm v\cdot\p L/\p\bm v-L$，用 $x=(Gm\Omega/c^3)^{2/3}$ 消去 $r$，所得结合能的首项为 $-\mu c^2 x/2$，其 1PN 相对修正为 $-(3/4+\eta/12)x$。半径随坐标选择而变，$E(\Omega)$ 可直接与波的能流联用。高一阶的结果见 \S\ref{sec:pn-wave-energy}。

偏心轨道可用 \textbf{Runge--Lenz 矢量}的方向标记近心点。令 $\bm h=\bm x\times\bm v$、$\bm e=\bm v\times\bm h/(Gm)-\bm n$。在 Newton 椭圆上两者均为常矢量。若 $\bm A$ 为扰动加速度，则 $Gm\dot{\bm e}=\bm A\times\bm h+\bm v\times(\bm x\times\bm A)$。令 $p=a(1-e^2)$、$h=\sqrt{Gmp}$，在未扰动椭圆上以真近点角 $f$ 作参数，有 $r=p/(1+e\cos f)$、$\dot r=he\sin f/p$、$r^2\dot f=h$。将 $\dot{\bm e}$ 投影到面内垂直于 $\bm e$ 的方向，再沿一周积分，即得
\begin{equation}\label{eq:pn-apsis-integral}
 \Delta\varpi=\frac{1}{eGm}\int_0^{2\pi}r^2
 \left[-A_R\cos f+A_T(1+r/p)\sin f\right]\d f .
\end{equation}
其中 $A_R,A_T$ 是 $\bm A$ 的径向与顺行分量。由二体方程，
\begin{align*}
 \frac{r^2c^2}{Gm}A_R&=(4+2\eta)\frac{Gm}{r}-(1+3\eta)v^2
                              +(4-\eta/2)\dot r^2,\\
 \frac{r^2c^2}{Gm}A_T&=(4-2\eta)\dot r\,h/r .
\end{align*}
两项给出的进动分别为 $-(2-4\eta)\pi Gm/(pc^2)$ 和 $(8-4\eta)\pi Gm/(pc^2)$，相加后质量比消去。故 \hyperlink{subsection.2.2.2}{\S\,2.2.2} 的每周进动 $6\pi Gm/(pc^2)$ 对质量可比的双星仍成立。作一数值检验。小行星 Icarus 的 $a\simeq1.61\times10^8\ {\rm km}$、$e\simeq0.827$，每世纪约绕日 $89$ 周，得每世纪约 $10.0''$，可与早期测值 $(9.8\pm0.8)''$ 比较。严格圆轨道没有近心点方向，其小偏心率极限应理解为径向振动与绕转频率之差。

Runge--Lenz 矢量在式 \eqref{eq:pn-apsis-integral} 中的一阶变化对扰动加速度是线性的。若干微扰的一周进动可分别计算再相加。以太阳角动量 $\bm J_\odot$ 为例，定常矢量势给试验粒子的附加加速度为 $\bm A_J=\frac{2G}{c^2r^3}\,\bm v\times[\bm J_\odot-3\bm n(\bm n\cdot\bm J_\odot)]$。取顺行赤道轨道，$\bm J_\odot$ 沿轨道法线。此时 $A_R=2GJ_\odot h/(c^2r^4)$，$A_T=-2GJ_\odot\dot r/(c^2r^3)$。代入同一个进动积分，得每周 $\Delta\varpi_J=-8\pi GJ_\odot/[c^2\sqrt{GM_\odot}\,p^{3/2}]$，符号与球单极进动相反。取早期估值 $J_\odot=1.7\times10^{41}\,\mathrm{kg\,m^2\,s^{-1}}$、水星半通径 $5.55\times10^7\,\mathrm{km}$，得每周 $-2.1\times10^{-11}$ 弧度，约合每世纪 $-0.0018$ 角秒。即使将这项角动量放大 $25$ 倍，也只有每世纪约 $0.04$ 角秒。还有太阳扁率、转动动能对 $Q$ 的各向异性贡献，以及其余行星摄动的 1PN 修正。实际历表怎样分离它们，见 \S\ref{sec:pn-ephemeris}。

% Weinberg §§9.6–9.7。
\subsection{陀螺的轴与旋转球壳}\label{sec:pn-transport-shell}

如何将\hyperlink{subsection.1.3.3}{\S\,1.3.3} 的输运律化成相对于远星的角速度？Schiff 沿用 Papapetrou 和 Fock 的方法，先求有体积陀螺的运动再取小体极限。若忽略自旋对轨道的反作用，亦可直接将陀螺视为点状探针。以 $S_\mu$ 表示自旋协矢，其三个坐标分量记为 $\bm s=(S_i)$。正交条件给 $S_0=-\bm v\cdot\bm s/c$。于是平行移动所保持的平方长度，至相对 $c^{-2}$ 阶为 $(1-2U/c^2)s^2-(\bm v\cdot\bm s)^2/c^2$。引入 $\bm{\mathcal S}=\left(1-\frac{U}{c^2}\right)\bm s-\frac{\bm v(\bm v\cdot\bm s)}{2c^2}$，其欧氏长度正是陀螺静系中的自旋长度。对上式求导，已乘 $c^{-2}$ 的加速度用 $\dot{\bm v}=\bm\nabla U$。将 $\d S_i/\d t=\Gamma^\alpha{}_{i\nu} S_\alpha\,\d x^\nu/\d t$ 代入，$\bm s\,\d U/\d t$ 项相消，余下的反对称部分为
\[
 \frac{\d\bm{\mathcal S}}{\d t}=\bm\Omega\times\bm{\mathcal S},
 \qquad \bm\Omega=\frac{1}{c^2}
 \left(\frac32\bm v\times\bm\nabla U+2\bm\nabla\times\bm V\right).
\]
第一项是测地进动，第二项是中心天体自旋与陀螺自旋的耦合，形似原子超精细相互作用。若空间度规的系数改成 $2\gamma U/c^2$，恒长矢量第一项相应改成 $(1-\gamma U/c^2)\bm s$，测地进动系数便由 $3/2$ 变为 $\gamma+1/2$。将球外的 $\bm V$ 代入并沿椭圆平均，结果及两种角速度的方向见 \S\ref{sec:pn-gyro-observe}。

Pugh 与 Schiff 建议用卫星携带陀螺。用望远镜让导星的像落在随轴运动的光电探测器上，实际量的是入射光与转轴的夹角。令 $u$ 为陀螺的单位切矢（$u\cdot u=-1$），$k$ 为光的未来指向波矢，则 $\cos\theta=(S\cdot k)/[\sqrt{S\cdot S}(-u\cdot k)]$，因为在陀螺本身的正交标架中 $S^{\hat0}=0$。在地心系设远方光线的传播方向为 $\bm n_\infty$，在接收点静止的正交标架中其方向为 $\bm n_\infty+\delta\bm n$。其中 $\delta\bm n\perp\bm n_\infty$ 由光线方程给出，量级为 $U/c^2$。令 $\bm b=\bm v/c$、$z=\bm b\cdot\bm n_\infty$。在这个正交标架与陀螺静系间作瞬时 Lorentz 推进，将上面的不变量展开得
\[
 \cos\theta=\hat{\bm{\mathcal S}}\cdot\!
 \left[\bm n_\infty+\delta\bm n
       +z\bm n_\infty-\bm b+
       (z^2-\tfrac12 b^2)\bm n_\infty-\tfrac12 z\bm b\right]
 + O\bigl((v/c)^3\bigr).
\]
这里光的局域方向已归一化。若从坐标波矢出发，静止标架的空间尺度与频率都须换算，才能得到同一个 $\delta\bm n$。括号内关于 $\bm b$ 的一次、二次项是相应阶的光行差。地球近旁，单周内的偏折、二阶光行差与测地进动均为 $O(v^2/c^2)$，逐周分辨它们需约 $10^{-10}$ 弧度的角精度。同一轨道位置观测 $N$ 周后，周期修正仍有界，轴的进动却累积至 $O(Nv^2/c^2)$。扣除光行差后，夹角余弦的长期改变即为 $\bm n_\infty\cdot\Delta\hat{\bm{\mathcal S}}$。

取半径 $r$ 的近地圆极轨道，地球单极给出的测地进动约为每年 $8.4(R_\oplus/r)^{5/2}$ 角秒。若初始转轴沿轨道法线，这一项不使转轴偏移，而地球自转给出每年约 $0.055(R_\oplus/r)^3$ 角秒的偏转。在 $r\simeq R_\oplus$ 处二者之比约 $6.5\times10^{-3}$。这是一种只显出自转效应的设想，Gravity Probe B 用导星分辨两种进动的办法见 \S\ref{sec:pn-gyro-guide}。

陀螺在自身携带的惯性标架中有固定方向，故上面所得也是该标架相对于远星的转速，与陀螺本身转得多快无关。按 Mach 的想法，周围物质的运动应参与确定惯性系。旋转球壳给这一想法一个可算的例子。空腔内 $U$ 为常数，只有质量流的旋度留下。若薄球壳以同一 $\bm\omega$ 转动，$U=GM/R$，则 $\bm\Omega_{\rm in}=4GM\bm\omega/(3Rc^2)$，整个空腔内的惯性轴均沿壳的转向拖动。若 $GM/(Rc^2)$ 不再很小，还能沿用此式吗？

对于旋转外场，可将\hyperlink{subsection.2.5.2}{\S\,2.5.2} 的 Kerr 解改用渐近直角坐标，与这里的偶极直接比较。令 $m=GM/c^2$、$\bm a=\bm J/(Mc)$、$r=|\bm x|$，取在远处满足 $\varrho/r\to1$ 的根 $\varrho^4-(r^2-a^2)\varrho^2-(\bm a\cdot\bm x)^2=0$。在入射 Kerr--Schild 坐标下，线元为
\begin{align*}
 \d s^2&=-c^2\d t^2+\d\bm x^2+
  \frac{2m\varrho^3}{\varrho^4+(\bm a\cdot\bm x)^2}
  (c\,\d t+\bm\ell\cdot\d\bm x)^2,\\
 \bm\ell&=\frac{\varrho\bm x-\bm a\times\bm x+
 (\bm a\cdot\bm x)\bm a/\varrho}{\varrho^2+a^2}.
\end{align*}
上面的四次方程保证 $\bm\ell^2=1$。在远处 $g_{00}=-1+2m/r+O(r^{-2})$，$g_{ij}=\delta_{ij}+2mn_i n_j/r+O(r^{-2})$。混合分量有两项，$g_{0i}=2mn_i/r-2m(\bm a\times\bm n)_i/r^2+O(r^{-3})$。前项为 $2m\p_i\ln r$，来自时间坐标，旋度为零。后项恰是谐和度规的自转偶极。由此读出系统在静系的总动量为零、能量为 $Mc^2$、角动量为 $Mc\bm a$。仅给外部真空解还不能指定壳内的惯性轴，须在壳面匹配内外解。

Brill 与 Cohen 对自转角速度只留一次项，对质量则保持所有阶。设 $R$ 为壳的各向同性坐标半径，$q=\max(r,R)$，$A=(1-m/2q)/(1+m/2q)$、$B=(1+m/2q)^2$，则
\[
 \d s^2=-A^2c^2\d t^2+B^2\bigl[\d r^2+r^2\d\theta^2+
 r^2\sin^2\theta(\d\phi-\Omega(r)\d t)^2\bigr]+O(\omega^2).
\]
在 $r<R$，$A,B,\Omega$ 均为常数。取 $t'=At$、$r'=Br$、$\phi'=\phi-\Omega_{\rm in} t$，即得平直度规。规定无穷远处 $\Omega=0$，外部真空解有 $\Omega(r)\propto[rB(r)]^{-3}$，比例常数由壳面匹配决定。若 $\chi=m/(2R)$，支撑静态薄壳的切向应力与面能量密度之比为 $b_s=\chi/[2(1-\chi)]$。将角动量方程跨越壳面作积分，得
\[
 \frac{\Omega_{\rm in}}{\omega}
 =\left[1+\frac{3(R-m/2)}{4m(1+b_s)}\right]^{-1}
 =\frac{4\chi(2-\chi)}{(1+\chi)(3-\chi)} .
\]
弱场极限给出 $4m/(3R)$。当 $R\to m/2$ 时，壳的面积半径 $R(1+\chi)^2\to2m$，拖曳率趋于壳的转率。这个极限须连同维系薄壳的应力一起求，不能从弱场式外推。\footnote{D.~R.~Brill 与 J.~M.~Cohen，Phys. Rev. \textbf{143} (1966) 1011。这里区分各向同性半径和面积半径。}

% Weinberg §9.9。
\subsection{Brans--Dicke 理论中的球体外场}\label{sec:pn-bd}

取 \S\ref{sec:scalar} 的 Brans--Dicke 场方程，令耦合 $\omega$ 为常数，标量无势。在孤立系统的边界令 $\phi\to\phi_0$，写 $\phi=\phi_0(1+\Psi)$，以 $q=1/(\omega+2)$ 记标量的强度。将 \eqref{eq:BDtrace} 代回度规方程，$R_{\mu\nu}$ 的物质源正比于 $T_{\mu\nu}-(\omega+1)g_{\mu\nu}T/(2\omega+3)$，另有 $\cd_\mu\cd_\nu\Psi/(1+\Psi)$ 与 $\omega\p_\mu\Psi\p_\nu\Psi/(1+\Psi)^2$。迹 $T$ 只取物质的贡献。比较标量及 $00$ 方程的最低阶 Poisson 源，即有 $\Psi=q u+O(c^{-4})$，其中 $u=U/c^2$，势中的 $G$ 由慢质点的加速度或钟速差标定。它与 $1/\phi_0$ 的换算见 \S\ref{sec:pn-bd-parameters}。

现在取\emph{物理度规}的谐和坐标。记 $g_{00}=-1+2u+k$、$g_{ij}=\delta_{ij}+s_{ij}$、$g_{0i}=b_i$，各式只保留 1PN 所需的阶次。沿用 \S\ref{sec:pn-components} 的上指标物质源。令 $K=1-q/2$，逐阶得
\begin{align*}
 \nabla^2s_{ij}&=2(1-q)\delta_{ij}\nabla^2u-2q\,u_{,ij},\\
 \nabla^2 b_i&=\frac{16\pi G}{c^3}Kj_i-2q\,u_{,i0},\\
 \nabla^2k&=2(1-q)u_{,00}+2s_{ij}u_{,ij}-2(2+q)|\bm\nabla u|^2
 \\[-2pt]
 &\quad +\frac{8\pi G}{c^4}\big[(4+q)\rho_0c^2u-e-(1-q)\sigma_{ii}\big].
\end{align*}
其中 $u_{,0}^2$ 已是六阶。与 \S\ref{sec:pn-bd-aux} 对照，同一源在辅助度规的谐和坐标中有各向同性的 $g_{ij}$，在这里则多了标量的 Hessian。

设源静止且球对称。取 $Y=X/c^2$，则 $\nabla^2Y=2u$，正则性给出 $Y'(r)=2r^{-2}\int_0^r s^2u(s)\,\d s$。故在球内外皆有 $s_{ij}=2(1-q)u\delta_{ij}-qY_{,ij}$。置 $m=GM/c^2$，球外 $u=m/r$，并以 $mR_e^2=\int_0^\infty[m/s-u(s)]s^2\,\d s$ 定义一个质量矩。被积式只在球内非零。令 $n_i=x_i/r$，在球外积分即得
\begin{align*}
 s_{ij}&=u[(2-3q)\delta_{ij}+q n_i n_j]
       +\frac{2q mR_e^2}{r^3}(\delta_{ij}-3n_i n_j),\\
 k&=-(2-q)u^2-\frac{2q m^2R_e^2}{r^4}+\frac{D}{r} .
\end{align*}
长度 $D$ 由球心正则解在表面的匹配确定。时间度规中的 $1/r$ 项将远处测得的质量长度改为 $m+D/2$。再作 $x'^i=x^i(1+qmR_e^2/r^3)$，空间度规由平直项产生的改正抵消 $r^{-3}$ 项，而 $2m/r$ 改用 $r'$ 表示时正好抵消 $k$ 中的 $r^{-4}$ 项。因此这颗弱自引力球在所求阶次的外部度规只余实测质量。与谐和半径下的球对称展开比较，$n_i n_j$ 的系数是 $(1-\gamma)u$，时间度规二次项是 $-(\gamma-1+2\beta)u^2$，即 $\gamma=1-q$、$\beta=1$。

若球体缓慢转动而质量流定常，$u_{,i0}=0$，矢量势便是广义相对论结果的 $K=(2\omega+3)/(2\omega+4)$ 倍。因此给定相同的实测 $GM$ 和自转角动量，自转引起的陀螺和轨道近心点进动均按此因子改变。正 $\omega$ 时 $K<1$，两种自转效应均较弱。另在外势 $U_{\rm e}$ 中作局域 Cavendish 实验，实验室虽可自由下落，标量的背景值仍是 $\phi_0(1+qU_{\rm e}/c^2)$。固定 $\omega$ 的局域耦合与该值成反比，故两地点测值之差为 $\Delta G_{\rm L}/G=-q\,\Delta U_{\rm e}/c^2$，与 \S\ref{sec:pn-local-g} 的固有距离与支撑力计算相合。取 Weinberg 当时举例所用的 $\omega=6$，仅计地球的势，从远高轨道移到地面，$U_\oplus/c^2\simeq6.9\times10^{-10}$，相对变化约为 $-8.6\times10^{-11}$。同一 $q$ 也乘在弱自引力体的负结合能分数上，给出其被动质量与惯性质量之差，见 \S\ref{sec:pn-sep}。

\subsection{精度与适用范围}

对静止球体，以轨道标定 $GM$ 后，还可分别问单位质量产生多少空间曲率、时间度规中的二次项有多大。标准坐标用 $\gamma,\beta$ 表示这两个系数，广义相对论中均为一。式 \eqref{eq:exp-ppn-all} 给出有质量流及优选系等效应时的推广。近日点进动测 $2+2\gamma-\beta$，星光偏折与 Shapiro 时延测 $1+\gamma$。偏折角和总进动的归一化分别见 \S\ref{sec:pn-light}、\S\ref{sec:pn-ephemeris}。区分各参数须合用不同观测。

非自旋二体的保守动力学先有 1PN、2PN 修正，辐射带走能量与角动量，在相对 Newton 加速度的 2.5PN 阶首次产生\textbf{辐射反作用}。对自旋天体固定 $\chi_A=cS_A/(Gm_A^2)$ 作计数，自旋--轨道和自旋--自旋耦合分别起于 1.5PN、2PN。前者相对 Newton 束缚能为 $O(\chi_A x^{3/2})$，慢转行星的 $\chi_A$ 须按实际自转另估。波的尾项也可使可观测量出现半整数阶，例见 \S\ref{sec:pn-wave-energy}。黑洞或中子星的内部即使不满足弱场条件，其早期双星轨道仍可满足 $x\ll1$。在广义相对论的低阶轨道方程中，星体结构主要归入质量与自旋，有限大小的潮汐效应可用更高阶的矩描述。

误差亦可用精确解估计。在 \hyperlink{subsection.2.2.2}{\S\,2.2.2} 的有效势中取圆轨道条件，得 $E/(\mu c^2)=(1-2x)/\sqrt{1-3x}-1$。其二次项正是前面 $\eta\to0$ 时的 $3x^2/8$。但到最内稳定圆轨道 $x=1/6$，精确结合能的绝对值为 $1-\sqrt{8/9}$，Newton 截断却给 $1/12$。偏心轨道也应在近心点检查场强与速度。接近并合时须借助重求和、黑洞微扰或数值相对论，视质量比及所需精度而定。

\section{PPN 形式}\label{sec:exp-ppn}

设已由行星轨道标定太阳的 $GM$，再用一束掠日光测量空间曲率。\S\ref{sec:exp-pn} 的静止球体展开说明，为何二者对同一质量的响应仍由独立系数 $\gamma$ 计量。$\beta$ 的定义还须指明径向坐标。若在面积半径下直接读取 $g_{00}$ 的二次项，与各向同性半径下不会得到同一数，见 \hyperlink{section.2.1}{\S\,2.1}。

当质量分布会运动、形变时，可把度规写成一组共同的势，而将理论间的差别留在各势的系数中，称\textbf{参数化后 Newton 形式}(parametrized post-Newtonian formalism, PPN)。如何选势？先按阶次，再按空间旋转下的性质选取。譬如 $g_{00}$ 在 $c^{-4}$ 阶需要量纲为速度四次方的标量，$g_{0i}$ 在 $c^{-3}$ 阶则需要矢量。沿用 \S\ref{sec:pn-three} 的势积分，简记 $P[f]=\mathcal I[\rho_* f]$，则 $U=P[1]$、$V_i=P[v_i]$。若 $\Pi$ 是单位静质量的内能，$p$ 是压强，便有四个标量势 $\Phi_1=P[v^2]$、$\Phi_2=P[U]$、$\Phi_3=P[\Pi]$、$\Phi_4=P[p/\rho_*]$。式 \eqref{eq:pn-Q} 的超势也采用 $\rho_0=\rho_*$。于是 $\p_i\dot X$ 也是所需阶次的矢量。

先取只有 $\gamma,\beta$ 两个自由参数的一族理论。在标准 PPN 规范中，运动流体产生的度规为
\begin{align*}
 g_{00}&=-1+\frac{2U}{c^2}+\frac{2}{c^4}(\psi-\beta U^2)+O(c^{-6}),\\
 \psi&=\tfrac12(2\gamma+1)\Phi_1-(2\beta-1)\Phi_2+\Phi_3+3\gamma\Phi_4,\\
 g_{0i}&=-\frac{1}{c^3}\left[2(1+\gamma)V_i+\tfrac12\p_i\dot X\right]+O(c^{-5}),\\
 g_{ij}&=(1+2\gamma U/c^2)\delta_{ij}+O(c^{-4}).
\end{align*}
取 $\gamma=\beta=1$，$g_{0i}$ 仍与上节的谐和坐标相差一项 $\p_i\dot X$。这一项来自时间坐标的变换。至于 $\psi$ 的系数，则与我们选用守恒密度 $\rho_*$ 有关。若将 $U$ 改用固有密度 $\rho$ 定义，连 $2U/c^2$ 的一部分也要重新分配到 $c^{-4}$ 阶去。

\subsection{怎样选势和坐标}

设想直接让度规的十个分量各带一个待测系数，会有什么困难？转动坐标轴，原来的分量就混合了。即使张量本身不变，换一个径向坐标也会改变 $U^2$ 的系数。故要先问所选源量在空间转动下如何变换，又要问由此造出的势中哪些可以用坐标变换消去。取背景宇宙各向同性时的静系为起点，场点只可通过 $\bm s=\bm x-\bm x'$ 进入积分，否则仅把坐标原点平移就会改变预言。再要求进入度规的势及其导数组合在远离源时衰减，按阶次截断，就得到下面有限的几类。这个构造假定物质源是理想流体，加速度、密度的高阶导数或一种新的作用程不任意作为源量加入。

先看 $g_{ij}$。一个静止质量元可以用 $\delta_{ij}$ 乘 $1/s$ 构成二阶张量，也可用 $s_i s_j/s^3$。记后一种势为 $U_{ij}=G\int \rho_*'s_i s_j\,\d[3]{x'}/s^3$。由于 $\p_i\p_j s=\delta_{ij}/s-s_i s_j/s^3$，有 $U_{ij}=\delta_{ij}U-X_{,ij}$。这里及下文逗号后面的拉丁指标表示普通空间导数。在最低非平凡阶作 $x^i\mapsto x^i+k\p_iX/c^2$，度规扰动就改变 $-2kX_{,ij}/c^2$。故可消去这个 Hessian，余下各向同性的空间度规。剩下的 $\gamma$ 才是可以在同一规范中比较的系数。

再看 $g_{0i}$。把同一质量元的速度沿 $\bm s$ 作投影便得到第二个矢量势 $W_i=G\int\rho_*'(\bm v'\cdot\bm s)s_i\,\d[3]{x'}/s^3$。对 $X$ 求时间导数，把 $\p_t\rho_*'$ 换成 $-\bm\nabla'\cdot(\rho_*'\bm v')$ 作分部积分，再对场点求空间导数，恰得 $\p_i\dot X=W_i-V_i$。可见 $V_i,W_i,\p_i\dot X$ 不能作为三个独立势。类似地，作 $t\mapsto t+k'\dot X/c^4$ 能在 $g_{00}$ 的 $\ddot X$ 项与 $g_{0i}$ 的 $\p_i\dot X$ 项之间搬移系数。标准规范选成 $g_{00}$ 不显含 $\ddot X$。谐和规范便于解方程，标准规范便于比较实验系数，两者都不改变用固有时和角度作的比较。

$g_{00}$ 的选择最多。把静质量、动能、引力能、内能和压强分别作源得到的势已见上文。除动能 $v'^2$ 外，速度在源点到场点连线上的投影平方也是标量。置 $\hat{\bm s}=\bm s/s$，相应的势为 $\Phi_6=P[(\bm v'\cdot\hat{\bm s})^2]$，即在 $U$ 的被积式中插入该投影平方。还可由两个不同的源元造势。$U^2$ 是场点分别与两源元相连，$\Phi_2=P[U]$ 是一个源元通过另一个源元处的势影响场点。第三种连接还取决于三点组成的三角形。令 $\bm x_0=\bm x$ 为场点，$\bm x_a,\bm x_b$ 为积分点，$\bm r_{ab}=\bm x_a-\bm x_b$，$r_{ab}=|\bm r_{ab}|$。定义 Whitehead 势
\begin{equation}
 \Phi_W=G^2\!\iint\rho_{*a}\rho_{*b}\,
 \frac{\bm r_{0a}}{r_{0a}^3}\cdot
 \left(\frac{\bm r_{ab}}{r_{0b}}-\frac{\bm r_{0b}}{r_{ab}}\right)
 \d[3]{x_a}\d[3]{x_b} .\label{eq:exp-whitehead}
\end{equation}
这里 $\rho_{*a}=\rho_*(t,\bm x_a)$。这三个二次势都有速度四次方的量纲，通常彼此独立。若只求球对称单体的外场，若干势会只给质量参数加一个常数修正；若要比较实验室相对银河中心的取向，就不能事先舍去带方向的势。实验的几何构型决定哪些系数实际可测。

\subsection{一般源的十参数度规}

将以上自由度收集起来，便可给运动的物质分布写出一般的标准 PPN 度规。令 $\bm w$ 为坐标系相对优选静系的速度，$X_{,ij}=\p_i\p_jX$。取 $G$ 为边界处静系内实测的引力常数。保留慢测试粒子达到 1PN 所需的项，有
\begin{align}
 g_{00}={}&-1+\frac{2U}{c^2}
        +\frac{2\psi-2\beta U^2+\Phi_{\rm F}}{c^4}+O(c^{-6}),\nonumber\\
 g_{0i}={}&\frac{1}{c^3}\left\{-\left[2(1+\gamma)+\frac{\alpha_1}{2}\right]V_i
       -\frac{1+\alpha_2-\zeta_1+2\xi}{2}\p_i\dot X
       +\Phi_{{\rm F}i}\right\}+O(c^{-5}),\nonumber\\
 g_{ij}={}&(1+2\gamma U/c^2)\delta_{ij}+O(c^{-4}),\label{eq:exp-ppn-all}\\[-2pt]
 \psi={}&\tfrac12(2\gamma+1+\alpha_3+\zeta_1-2\xi)\Phi_1
              -(2\beta-1-\zeta_2-\xi)\Phi_2 \nonumber\\
       &+(1+\zeta_3)\Phi_3+(3\gamma+3\zeta_4-2\xi)\Phi_4
              -\tfrac12(\zeta_1-2\xi)\Phi_6-\xi\Phi_W .\nonumber
\end{align}
标记 F 的两项只收集在 $\bm w=0$ 时消失的贡献，其余优选系参数在 $\psi$ 和 $g_{0i}$ 中仍可出现。具体为
\begin{align*}
 \Phi_{\rm F}&=(\alpha_3-\alpha_1)w^2U
           +(2\alpha_3-\alpha_1)w^i V_i+\alpha_2w^iw^jX_{,ij},\\
 \Phi_{{\rm F}i}&=-\tfrac12\alpha_1 w_iU+\alpha_2w^jX_{,ij}.
\end{align*}
可用两个极限检验此式。令所有优选系参数为零，则两个 F 势消失，任意匀速平移不能引入可测的绝对速度。再令 $\xi,\zeta_i$ 为零，就回到上面的两参数度规。反之，只让 $\bm w=0$ 而保留 $\alpha_1$，静系中的质量流 $V_i$ 仍产生异常的 $g_{0i}$。所以即使不知道优选系在哪个方向，也可拿绕旋转天体的陀螺去限制若干参数组合。

为何不让每个含 $\bm w$ 的项再各带一个待测参数？先在首选静系求出度规，再把整个解包括物质速度和积分所取的同时面一同变换到速度为 $\bm w$ 的系，两个表达式必须描述同一个实验。最低阶已有 $\bm v_{\rm rest}=\bm v+\bm w$，故 $\Phi_1$ 的变换自然产生 $2\bm w\cdot\bm V+w^2U$。较高阶还须变换时间，否则只作 Galilei 速度相加，就会遗漏 Lorentz 变换对 $g_{00}$ 的贡献。把这些运动学项和度规分量本身的变换一并收集，再作小的坐标调整恢复标准规范，余下的恰是上面含 $\alpha_i$ 的组合。给两个参考系各配一套独立参数，就失去了拿同一理论比较不同运动实验的意义。

标准 PPN 共有十个无量纲参数。余下八个的实验含义如下。
\begin{itemize}
 \item $\xi$ 容许局域引力实验随外界物质分布而改变。例如相对于银河中心转动扭秤，两质量间的引力强度可能有方向差异，称为\textbf{优选位置效应}。
 \item $\alpha_1,\alpha_2,\alpha_3$ 容许\textbf{优选参考系效应}。若引力场选定一个静系，实验室相对它的速度 $\bm w$ 也可进入度规，由 $w^2U$、$\bm w\cdot\bm V$ 等项影响轨道或自转。
 \item $\alpha_3$ 又与 $\zeta_1,\zeta_2,\zeta_3,\zeta_4$ 一同检验孤立系统的总动量守恒。此处的总动量包括引力场的贡献，仅凭物质场的局域方程 $\cd_\nu T^{\mu\nu}=0$ 还不能保证它守恒。
\end{itemize}
例如从优选系换到相对它以 $\bm w$ 运动的系，$g_{0i}$ 将多出 $[-\tfrac12\alpha_1 w_i U+\alpha_2 w^j\p_i\p_j X]/c^3$。这给出了随实验室运动改变的两种方向依赖，单测静止球体的 $g_{00}$ 无从分辨它们。

广义相对论中这八个参数全为零。一个理论若在此阶有总能量、总动量守恒律，就须有 $\alpha_3=\zeta_1=\zeta_2=\zeta_3=\zeta_4=0$，称\textbf{半守恒理论}。若进一步要求在任一渐近惯性系中总角动量守恒、质心作匀速直线运动，则还须有 $\alpha_1=\alpha_2=0$，称\textbf{全守恒理论}。因此即使已用光线和行星轨道测定 $\gamma,\beta$，也还有相互独立的实验可做。

\subsection{同一组势在两套坐标中的系数}

还少了一种可能的源。加速度也有正确的量纲，比如 $\bm s\cdot\dot{\bm v}'$ 与 $v'^2$ 同量纲，若直接逐项枚举会多出一个势。这里可用物质运动在最低阶消去它。沿用 $P[f]$ 的记号，以 $\Phi_5=P[\bm s\cdot\bm\nabla' U']$ 表示由源点所处的引力梯度组成的势。对有界源及不显含时间的核 $F(\bm x,\bm x')$，连续性给
\[
 \p_t\int \rho_*'F\,\d[3]{x'}=
            \int \rho_*'\bm v'\cdot\bm\nabla'F\,\d[3]{x'} .
\]
取 $F=G/s$ 立即有 $\p_i V_i=-\dot U$。取 $F=Gs$ 再微分一次，$\bm v'$ 本身的导数换成 Euler 加速度 $\bm a'=\bm\nabla' U'-\bm\nabla' p'/\rho_*'$。来自 $\bm\nabla' p'$ 的一项作分部积分，用 $\bm\nabla'\cdot\hat{\bm s}=-2/s$，便得 $\ddot X=\Phi_1+2\Phi_4-\Phi_5-\Phi_6$。这里以 $\rho_*$ 代 $\rho$ 只用于一个本来已经四阶的项。因而在 $g_{00}$ 中保留 $\ddot X$ 和 $\Phi_5$ 二者之一就够了。

有另一个恒等式能看出空间坐标怎样改变看似固定的非线性系数。令 $P[\bm\nabla'X']$ 的每个分量都照 $P$ 作势，则从 \eqref{eq:exp-whitehead} 的双积分按 $\bm x_a-\bm x_b=(\bm x_a-\bm x_0)+(\bm x_0-\bm x_b)$ 分解，可得
\[
 \Phi_W=-U^2-\Phi_2-\bm\nabla U\cdot\bm\nabla X+
                         \bm\nabla\cdot P[\bm\nabla'X'] .
\]
因此对两源元的积分，既可直接算 $\Phi_W$，也可从两个单积分的势及导数求出它。后一条途径尤其适合检验数值程序的符号。它并不说 $\Phi_W$ 总是 $U^2$ 和 $\Phi_2$ 的线性组合，最后两项通常不可略。

具体地，以下在本小节的坐标变换中暂取 $G=c=1$。令新坐标 $\bar x^\mu=x^\mu+\epsilon^\mu$，用 $\epsilon_0=a\dot X$、$\epsilon_i=b X_{,i}$ 指定小位移，指标在这一阶以 $\eta$ 升降。源的守恒质量元 $\rho_*\,\d[3]{x}$ 在这次空间改标下不变，改变的是场点和源点的坐标距离。把距离的倒数展开，并用上式整理，得到旧坐标算出的 $U=\bar U+b(\bar U^2+\bar\Phi_2+\bar\Phi_W)$。若只变换度规的两个指标而仍把 $U$ 当作同一个数，就会遗漏这一组量。按新坐标的势统一表示后，同阶改变为
\begin{align*}
 \bar g_{ij}-g_{ij}&=-2bX_{,ij},\\
 \bar g_{0i}-g_{0i}&=-(a+b)\dot X_{,i},\\
 \bar g_{00}-g_{00}&=-2a\ddot X+2b(U^2+\Phi_2+\Phi_W).
\end{align*}
等号右边的量此时均用新坐标求值。先取 $b$ 消去空间度规的 $X_{,ij}$，再取 $a$ 消去时间度规中与 $\Phi_5$ 等价的项，便是标准 PPN 规范。换别的 $b$ 可以改变球外 $U^2$ 的系数，但会同时引入其余各项。用一个人的度规和另一个人的测地线、距离定义拼合来读 $\beta$，并不能算作两理论的比较。

对于求解场方程，另一种有用的选择是\textbf{广义谐和规范}。例如 Jordan 标量--张量理论对辅助度规施加谐和条件，未必等于对物理度规施加它。把这种情况统一写成 $\p_\nu[(1-a_{\rm H}U)\sqrt{-g}\,g^{\mu\nu}]=0$，只取条件所需的 PN 阶。普通谐和坐标相当于 $a_{\rm H}=0$。考虑半守恒理论，由标准坐标作 $\epsilon_0=\lambda_1\dot X+\lambda_3w^jX_{,j}$、$\epsilon_i=\lambda_2X_{,i}$。将 \eqref{eq:exp-ppn-all} 和行列式代入上述条件，$\bm\nabla U$、$\dot U$ 和 $\bm w\cdot\bm\nabla U$ 对一般源可独立改变。令其系数分别为零，有
\begin{align*}
 2\lambda_2&=1+a_{\rm H}-\gamma,\qquad
 2\lambda_1=a_{\rm H}-\gamma+\tfrac12\alpha_1-\alpha_2-2\xi,\\
 2\lambda_3&=2\alpha_2-\tfrac12\alpha_1 .
\end{align*}
两个规范的位移至此完全确定，不必重新解一遍物质方程。特别可取 $a_{\rm H}=\gamma-1$ 使空间部分仍各向同性。在这个坐标中令 $\psi$ 为标准规范的半守恒标量势，则
\begin{align*}
 g_{00}&=-1+2U+2(\psi-\beta U^2)
       +(1-\alpha_1/2+\alpha_2+2\xi)\ddot X+\Phi_{\rm H},\\
 g_{0i}&=-[2(1+\gamma)+\alpha_1/2]V_i-\tfrac14\alpha_1\dot X_{,i}
          -\tfrac12\alpha_1 w_i U+\tfrac14\alpha_1w^jX_{,ij},\\
 \Phi_{\rm H}&=-\alpha_1(w^2U+\bm w\cdot\bm V)+\alpha_2w^iw^jX_{,ij}
                   -(2\alpha_2-\alpha_1/2)w^j\dot X_{,j}.
\end{align*}
按本节约定略去更高阶，$g_{ij}=(1+2\gamma U)\delta_{ij}$。取广义相对论的参数，立即得前面遇见的 $+\ddot X$ 和 $-4V_i$。故在谐和坐标中当然也可测定 PPN 参数，只要连规范中产生的附加项都按同一套公式比较。以下仍以标准规范为准。

\subsection{运动参考系中的同时面}

为什么变换 Newton 势也会进入 1PN 度规？在优选静系和实验室分别给一颗流体元贴上标记。$\rho,p,\Pi$ 是该质元瞬时共动系中测得的标量，不会因改用实验室坐标而改变。但是 $U$ 的积分取静系的一张同时面，实验室的 $U$ 则取另一张同时面。原先位于同一时刻的各个质元，在新坐标中须沿各自的世界线挪到同一时刻。只把 $\bm v$ 改成 $\bm v+\bm w$ 会漏掉这个效应。

仍用几何制来数速度幂。令 $(t,\bm x)$ 为静系，$(\tau,\bm y)$ 为以 $\bm w$ 运动的坐标系。对场点与源点记 $\bm d=\bm y-\bm y'$、$\bm n=\bm d/|\bm d|$，$\bm v'$ 指源点在运动系的速度。由 Lorentz 变换，静系同时的那一对事件在运动系满足 $\Delta\tau=-\bm w\cdot\bm d$ 到最低阶。为了重建 $\tau$ 常数的积分，需把质元坐标按 $\bm v'\Delta\tau$ 展开。连同沿 $\bm w$ 方向的长度变换，最后有
\[
 \frac{1}{|\bm x-\bm x'|}
 =\frac1{|\bm d|}\left[1+\tfrac12(\bm w\cdot\bm n)^2
              +(\bm w\cdot\bm n)(\bm v'\cdot\bm n)\right]
 +O(\epsilon^2/|\bm d|).
\]
取固有标量时没有上述括号，取守恒质量元时沿同一组质元积分也没有额外的粒子数因子。假设撇在下列各势上专指运动系而非源点，于是有
\begin{align*}
 U&=U'+w^jW'_j+\tfrac12w^jw^kU'_{jk},&
 V_j&=V'_j+w_jU',\\
 \Phi_1&=\Phi'_1+2w^jV'_j+w^2U',&
 \Phi_6&=\Phi'_6+2w^jW'_j+w^jw^kU'_{jk},\\
 X_{,0j}&=X'_{,0j}-w^kX'_{,jk}.&&
\end{align*}
每式只保留在度规中有用的阶次。$U_{jk}=\delta_{jk}U-X_{,jk}$，故第一式可以完全改用前面的势表。$\Phi_2,\Phi_3,\Phi_4,\Phi_W$ 在此精度已无需修正。$\Phi_1$ 里的 $w^2U$ 容易从速度平方看出，$U$ 自身的改正却能在 $g_{00}$ 的二次阶抵消一部分这样的项。由此可知同一理论在两个系中余下的 $\bm w$ 系数为什么不是任意可选的。

将变换过的各势放回度规，又会在 $g_{00}$ 中遇到一项正比于 $w^j\dot X_{,j}$。这是第一个参考系的标准坐标条件未在第二个系中自动成立。再作只改变时间的小变换，$\epsilon_0$ 取为 $w^jX_{,j}$ 的常数倍，$g_{00}$ 中这项就能移去，相应的 $w^jX_{,ij}$ 转入 $g_{0i}$。最后留下的是 \eqref{eq:exp-ppn-all} 中两个 F 势。若 $\alpha_i=0$，作这一套手续后，任意速度 $\bm w$ 完全从实验方程中消失。若仅 $\alpha_2\ne0$，同一个仪器转过直角，$w^iw^jX_{,ij}$ 可改变，即使仪器相对静系的速率不变也能留下信号。后面的重力仪和自转轴实验正可区别这两种情形。以下在给出实验量时恢复 $c$。

\subsection{局域守恒与孤立系统的守恒量}

物质所满足的 $\nabla_\nu T^{\mu\nu}=0$ 与包含场在内的整体守恒律，区别见 \S\ref{sec:Einstein}、\S\ref{sec:Noe}。现在要求在渐近 Descartes 坐标中构造普通散度为零的密度 $\tau^{\mu\nu}$。这样 $P^\mu=c^{-1}\int\tau^{\mu0}\,\d[3]{x}$ 满足 $\d P^\mu/\d t=-\oint\tau^{\mu j}\d S_j$。问题在于给定的 PPN 系数能否使这样的积分存在，以及何时它还给出守恒的角动量。

角动量多一层要求。对 $x^i\tau^{j\nu}-x^j\tau^{i\nu}$ 求散度，除去已因动量守恒而为零的项，只余 $\tau^{ji}-\tau^{ij}$。若空间应力可取为对称，就得到空间角动量的守恒。再比较 $x^0\tau^{i\nu}-x^i\tau^{0\nu}$ 的通量，可见能流与动量密度之差还决定能量中心能否匀速运动。一个优选系理论即使在自己的静系里有对称的空间应力，换到运动的实验室后也未必保有这两条定律。全守恒理论正是在这一点上比半守恒理论多了约束。

一个具体的检验是把相同总能量以不同形式装进系统。$\Phi_3$ 以单位静质量内能 $\Pi$ 为源，$\Phi_4$ 以压强为源。若机械能化成热，$\int\rho_*\Pi\,\d[3]{x}$ 和 $\int p\,\d[3]{x}$ 一般都要变。假设它们在外部引力场产生的质量变化与内部能量守恒的计算不符，双体相互吸引时就可能有剩余合力。$\zeta_i$ 量度的正是这类不相容。压力容器本身也有应力，计算封闭物体的主动质量时须把器壁算进去；只算里面气体的 $3p$ 项，会把普通应力平衡误判成守恒律破坏。

\subsection{从余力求守恒条件}

如何从实际的 PPN 方程中认出上面五个非守恒参数？令 $G=c=1$，把 $\p_\nu T^{j\nu}$ 用场强、物质速度和压力梯度表示。每当出现 $\rho_* U_{,j}$，就利用 $\nabla^2 U=-4\pi\rho_*$ 作分部积分。譬如定义 $\mathcal T_{jk}[f,g]=f_{,(j}g_{,k)}-\delta_{jk}\bm\nabla f\cdot\bm\nabla g/2$，则 $\p_k\mathcal T_{jk}[f,g] =\tfrac12(f_{,j}\nabla^2g+g_{,j}\nabla^2f)$。故 $\rho_* U_{,j}=-(4\pi)^{-1}\p_k\mathcal T_{jk}[U,U]$。物质所受 Newton 引力被写成穿过边界的场应力，这一步对任意密度形状都成立。到 1PN 再做一次，可以把含 $\bm V,\Phi_1,\ldots$ 的大部分项也收进空间散度和时间导数。对 $f=g$ 得到的应力自动对称；不同势互乘时须保留指标的次序，稍后才能判断场所带角动量。

具体来说，任取一常数 $a$，从 $[1+(5\gamma+a-1)U]T^{\mu\nu}$ 出发另加场项。不同的 $a$ 相当于重新分配物质和场携带的动量，不应改变全空间积分的结果。消去纯散度以后，空间方程留下的余力密度可取为
\begin{align*}
 Q^j=U_{,j}\bigg[&\tfrac12(\alpha_3+\zeta_1)\rho_*v^2
 +\frac{\zeta_1}{8\pi}\nabla^2\Phi_6
 +\frac{\zeta_2}{8\pi}|\bm\nabla U|^2
 +\zeta_3\rho_*\Pi+3\zeta_4p
 +\alpha_3\rho_*(\bm w\cdot\bm v)\bigg].
\end{align*}
这个式子的用法不是逐一对六项找实验，而是问能否对\emph{任意}流体写出一个局域场动量，把 $\int Q^j\,\d[3]{x}$ 化成它的全时间导数。在某一个静止球体中，这个积分可以靠反射对称为零，即使各系数并不为零。若任意改变物质的内能分布，$\rho_*\Pi$ 的余项也必须消失；改压强分布、速度的方向以及实验室的 $\bm w$ 则依次给出其它条件。应用 Poisson 恒等式再把其中一项移入通量，只会在另一组独立源上留下余项。对标准 PPN 势族，要求这样的一般动量积分存在便给出 $\alpha_3=\zeta_1=\zeta_2=\zeta_3=\zeta_4=0$。这里的论证针对任意源，不能由一次孤立系统未见异常加速倒推出这五个数各自为零。

满足这五个条件时，可给出一个实际可积分的动量。略去高阶，在标准规范中能量（含静能）与动量的 $j$ 分量分别为
\begin{align*}
 P^0&=\int\rho_*\left(1+\tfrac12v^2-\tfrac12 U+\Pi\right)\d[3]{x},\\
 P^j&=\int\rho_*\bigg\{ v^j\left(1+\tfrac12v^2-\tfrac12 U+\Pi+p/\rho\right)
       -\tfrac12(1+\alpha_1)V_j-\tfrac12(1+\alpha_2)X_{,0j}\\
    &\hspace{31mm}-\tfrac12(\alpha_1 w_jU-\alpha_2w^kX_{,jk})\bigg\}\d[3]{x}.
\end{align*}
$P^0$ 中 $-\rho_*U/2$ 是总结合能的密度写法，它的积分与把密度写成 $-|\bm\nabla U|^2/(8\pi)$ 等价。证明只须用 Poisson 方程再将 $\p_i(U\p_iU)$ 移到远边界。一个闭合系统完全可能满足这个 Newton 阶的能量积分，却在下一阶有 $\int Q^j$。所谓“能量不守恒”总要附上所指的阶次和被积分的系统。

要求能量中心作匀速运动，还须有 $\alpha_1,\alpha_2$ 为零。定义能量中心 $\bm X=(P^0)^{-1}\int \bm x\,\rho_*(1+v^2/2-U/2+\Pi)\,\d[3]{x}$。用连续性和最低阶 Euler 方程对分子求导，内部压力所作的功化为机械能与内能的交换；利用交换双积分变量的恒等式，可消去 $\int\rho_*(V_j-Uv_j)\,\d[3]{x}$。在全守恒情形剩下的恰是 $P^j$，故 $\dot X^j=P^j/P^0$。恢复单位是 $\dot X^j=c^2P^j/E$。反之，允许 $\alpha_1,\alpha_2$ 后，场的应力和能流一般仍留反对称的部分。守恒张量可加一个恒等散度为零的项，却不能任意指定其所有分量。把以某一速度运动的系统同静系统相比较，就能分开这种真正的差异与密度重新分配造成的差异。

这个结论对检验孤立天体自加速尤其重要。若只找到物质之间一对不等的作用力，应再问场是否正储存或辐射所差的动量；若辐射穿过积分面，应把通量加回。当有限物质体与其近区场均作了积分、取至 1PN 且无远区辐射带走相应动量时，长期不为零的质心加速度才直接检验上述余项。太阳系内各天体的内部结构近于球对称，恰是许多余项因奇偶性不可见的情形，故实验还要借助偏心双星这一非球对称的“组合体”。

\section{度规理论与后 Newton 极限}

怎样从场方程读出 PPN 参数？先找出物质钟尺所测的度规，选取离系统足够远而仍可忽略宇宙演化的区域作边界。令此处度规趋于 Minkowski 度规，其它引力场则取各自的背景值。由静止测试粒子的加速度确定 $g_{00}$ 中 $1/r$ 项的系数，把它归入实测的 $G$。然后逐阶解出 $g_{ij},g_{0i},g_{00}$，转到同一 PPN 规范，逐项比较即可。

\subsection{边界归一化与场方程的逐级比较}\label{sec:pn-source-stress}

选用什么坐标能让物理度规在边界处取成 Minkowski 形式？若只有一份动力学度规，直接把坐标尺和钟校准即可。若背景还带有不能调节的平直张量，先把背景写成 $\eta_{\mu\nu}$ 后，物理度规的均匀部分通常是 $\operatorname{diag}(-C_0,C_1,C_1,C_1)$，而不是 $\eta_{\mu\nu}$。$C_0,C_1$ 由宇宙的匹配条件给定。令 $\bar t=\sqrt{C_0}\,t$、$\bar{\bm x}=\sqrt{C_1}\,\bm x$，物理度规在新坐标中归一，原来背景张量的分量却一起变了。若在变换后又把背景强定为旧数值，就修改了理论。假设原坐标给 $g_{00}=-C_0+2kU_0$，$U_0=\int\rho_*'\d[3]{x'}/s$ 暂不含 $G$；由于 $\d m_*=\rho_*'\d[3]{x'}$ 不变，$U_0=\sqrt{C_1}\,\bar U_0$。故运动方程 Newton 阶实际读出的常数是 $G_{\rm N}=k\sqrt{C_1}/C_0$。引力作用量中的系数和台秤或轨道测得的 $G$ 至此才可比较。

在这样归一以后的坐标中，先从 $00$ 和可能的辅助标量方程求出 $U$ 的系数，然后求 $ij$ 的二阶项，还不必知道压强或系统的质量流。其次用 $0i$ 方程及辅助矢量的空间方程求三阶项。最后返回 $00$ 方程，各低阶解的平方、一阶源的推迟、物质的热能和内应力都由此进入四阶项。若较低阶还有标量 $\phi-\phi_0$，耦合函数须在 $\phi_0$ 处也展开到足够阶，可能产生一个先前不存在的参数。最后才换到标准规范。这套顺序可用完美流体写成显式的源作检验。例如当 $g_{ij}=(1+2\gamma U/c^2)\delta_{ij}+\cdots$ 时，令 $H=\Pi+p/\rho_*$，有
\begin{align*}
 T^{00}/c^2&=\rho_*\!\left[1+\frac{\Pi+v^2/2+(2-3\gamma)U}{c^2}\right]+O(c^{-4}),\\
 T^{0i}/c&=\rho_*v^i\!\left[1+\frac{H+v^2/2+(2-3\gamma)U}{c^2}\right]+O(c^{-4}),\\
 T^{ij}&=\rho_*v^iv^j\!\left[1+\frac{H+v^2/2+(2-3\gamma)U}{c^2}\right]
          +p\delta^{ij}(1-2\gamma U/c^2)+O(c^{-4}).
\end{align*}
余项按各自最低的物质尺度理解。$T^{00}$ 的这一级没有独立的 $p$ 项，原因是理想流体张量中焓所含的 $p$ 被 $pg^{00}$ 的最低部分抵消；同样的抵消不适用于 $T^{ij}$ 的迹。这个小检验可发现将静质量密度、重子坐标密度替错或把物质压力重复计入的错误。流体应力张量的变分推导见 \hyperlink{subsection.1.2.3}{\S\,1.2.3}，这里从它出发，只需要将源排列到合适阶次。

如果有两组理论在逐项比较后得到相同的十个数，1PN 度规在标准规范中便相同；这个结论的前提是两种理论都属于同一势族，配以相同的物质分布、实测 $G$ 和匹配条件。这一步没有比较星体内部是否有非微扰分支，也没有比较波区的传播方程。反过来，某个理论若只有在一套特定的背景值下符合太阳系数据，观测本身通常给出的是背景值与裸耦合的联合允许区，而不能唯一确定作用量中的每一个常数。

\subsection{从场方程求出广义相对论的系数}

可以先把这套手续在熟悉的理论中做一遍。暂用几何制 $G=c=1$，采用前面定义的守恒密度，令 $\mathfrak g^{\mu\nu}=\sqrt{-g}\,g^{\mu\nu}$、$H^{\mu\nu}=\eta^{\mu\nu}-\mathfrak g^{\mu\nu}$。这样定义的 $H^{\mu\nu}$ 在一阶等于通常的迹反转度规扰动，在高阶则还含度规的乘积。谐和条件正好写成 $\p_\nu H^{\mu\nu}=0$。把 Einstein 方程的一阶波算符放在左边，其余项移到右边，就有
\[
 \square_\eta H^{\mu\nu}=-16\pi\tau^{\mu\nu},\qquad
 \tau^{\mu\nu}=(-g)T^{\mu\nu}+\frac{1}{16\pi}\Lambda^{\mu\nu}[H],
 \qquad \p_\nu\tau^{\mu\nu}=0 .
\]
$\Lambda^{\mu\nu}$ 从 $H\,\p\p H$、$\p H\,\p H$ 等二次项起步。左边用了平直算符，右边的场本身仍是未知的；选推迟解并迭代，便可求出前述的场能动量，而无须先把引力能分配到各流体元。

第一次迭代，源中只保留 $T^{00}=\rho_*+O(\rho_*\epsilon)$ 和 $T^{0i}=\rho_*v^i+O(\rho_*\epsilon^{3/2})$。所得为 $H^{00}=4U+O(\epsilon^2)$、$H^{0i}=4V_i+O(\epsilon^{5/2})$，而 $H^{ij}=O(\epsilon^2)$。要把结果还原成度规，必须将 $\sqrt{-g}$ 也展开。于是
\begin{align*}
 g_{00}&=-1+\tfrac12H^{00}-\tfrac38(H^{00})^2+\tfrac12H^{kk}
                         +O(\epsilon^3),\\
 g_{0i}&=-H^{0i}+O(\epsilon^{5/2}),\qquad
 g_{ij}=(1+\tfrac12H^{00})\delta_{ij}+O(\epsilon^2).
\end{align*}
所以第二次迭代中只需求 $H^{00}$ 和空间迹 $H^{kk}$，并不需要解出空间部分的全部六个分量。注意第一次求得的 $4U$ 此时已经通过 $(H^{00})^2$ 参与下一阶。

将第一次解带回 $\tau^{\mu\nu}$，保留给 1PN 度规作贡献的源项，有
\begin{align*}
 \square_\eta H^{00}&=-16\pi\rho_*(1+\tfrac12v^2+3U+\Pi)
                          +14|\bm\nabla U|^2,\\
 \square_\eta H^{kk}&=-16\pi(\rho_*v^2+3p)+2|\bm\nabla U|^2 .
\end{align*}
例如 $3p$ 来自 $T^{kk}$ 中三个正交方向的应力，$|\bm\nabla U|^2$ 却处处可有，甚至延伸到物质外。后一项看似需要另外定义一个没有紧支集的势。其实由 $\nabla^2U=-4\pi\rho_*$ 及 $\nabla^2\Phi_2=-4\pi\rho_*U$，可写 $|\bm\nabla U|^2=\tfrac12\nabla^2(U^2)-\nabla^2\Phi_2$。这样它在满足既定边界条件的 Poisson 解中只会产生 $U^2/2-\Phi_2$，不再是新未知量。最低阶 $\rho_*$ 的推迟修正则是式 \eqref{eq:pn-Q} 中的 $\ddot X/2$。合起来得
\begin{align*}
 H^{00}&=4U+7U^2+2\Phi_1-2\Phi_2+4\Phi_3+2\ddot X,\\
 H^{kk}&=U^2+4\Phi_1-2\Phi_2+12\Phi_4 ,
\end{align*}
两行只写到所需阶次。代回 $g_{00}$，$U^2$ 的总系数为 $7/2-6+1/2=-2$。这给出另一个不同于球对称解的检验，而且同时得到一般流体的系数 $\psi=\tfrac32\Phi_1-\Phi_2+\Phi_3+3\Phi_4$。这时 $g_{00}$ 还剩 $+\ddot X$。作 $t_{\rm PPN}=t_{\rm harm}-\dot X/2$，它便消失，而 $g_{0i}$ 成为 $-4V_i-\p_i\dot X/2$。与 \eqref{eq:exp-ppn-all} 逐项相比，$\gamma=\beta=1$，余下八个系数皆为零。从这个例子可见，用真空球解可以很快读出两个数，要证明其它各项全无异常，必须允许物质运动、有压强且不具球对称性。以下恢复 $G,c$。

\subsubsection{非线性源的构造与外部质量的核对}\label{sec:pn-binding}

若要在计算中自行产生这里用到的二阶源，不必为每种物质另写一个引力应力公式。用相同的几何制和谐和坐标，定义四指标量
$\mathcal H^{\alpha\mu\beta\nu}=\mathfrak g^{\alpha\beta}\mathfrak g^{\mu\nu}
-\mathfrak g^{\alpha\nu}\mathfrak g^{\beta\mu}$。它对 $\alpha,\mu$ 及 $\beta,\nu$ 各自反对称，对这两对指标的交换对称。因此 $\p_\beta\p_\mu\p_\nu\mathcal H^{\alpha\mu\beta\nu}=0$ 恒成立。Landau--Lifshitz 形式可以反过来作为可计算的定义，
\[
 16\pi(-g)t_{\rm LL}^{\alpha\beta}
 =\p_\mu\p_\nu\mathcal H^{\alpha\mu\beta\nu}
                       -2(-g)G^{\alpha\beta}.
\]
把右边的 $G^{\alpha\beta}$ 用同一个度规算出，二阶导的线性项相消，剩下的是度规密度一阶导的乘积。若要直接求谐和波动方程的源，一种等价而更便于程序逐阶截断的写法是 $\Lambda^{\alpha\beta}=-\square_\eta H^{\alpha\beta}-2(-g)G^{\alpha\beta}$。这里的等式应在求解前用度规和导数展开，消去线性项，再把所得的非线性多项式带到已知的低阶解上。若把未知的精确 Einstein 张量仍留在右边，则方程只是恒等式。这个构造还说明为何一次迭代就须使用\emph{密度} $\mathfrak g^{\mu\nu}$。若过早将它换成单纯的 $g^{\mu\nu}$，所得 $\Lambda$ 从线性阶就会不同。

取一颗不转动、静力平衡且近乎均匀的弱场星，作为有限大小的例子。恢复 $G,c$，记重子质量为 $M_b=\int\rho_*\,\d[3]{x}$，用 Newton 半径 $R$。在最低阶 $\rho_*=3M_b/(4\pi R^3)$，积分 Poisson 方程给内部正势 $U(r)=GM_b(3R^2-r^2)/(2R^3)$。由 $\d p/\d r=\rho_*\d U/\d r$ 及 $p(R)=0$，得 $p(r)=\rho_*GM_b(R^2-r^2)/(2R^3)$。此时 $\int\rho_*U\,\d[3]{x}=6GM_b^2/(5R)$、$\int p\,\d[3]{x}=GM_b^2/(5R)$，正好符合 $3\int p=-\Omega$；$\Omega=-3GM_b^2/(5R)$ 是负的结合势能。$R$ 与密度在这些四阶势源中均只需用 Newton 值。具有同一 $M_b$ 的星若半径不同，其结合能就不同。

由于物质和引力能都已计入，外部真空，并不意味着由内部积分产生的 $\Phi_2,\Phi_3,\Phi_4$ 到表面便消失。在此球星外三者都是 $1/r$ 型，代入广义相对论的 $\psi$，其组合恰为 $G(E_{\rm int}+\Omega)/r$，其中 $E_{\rm int}=\int\rho_*\Pi\,\d[3]{x}$。故由外部钟读出的单极质量为 $M=M_b+(E_{\rm int}+\Omega)/c^2$，改用它来表示 $g_{00}$ 时，其余显式二次项才简化为 $-2(GM/rc^2)^2$。这给出了从流体外场过渡到点质量度规的具体条件，也能核对 $3\Phi_4$ 的系数。若把压强引力孤立地解释为 $+3\int p/c^2$ 的额外质量，却仍另按 Newton 结合能全额更新 $M$，就不能得到上述一致的外部解。

同一个例子还可说明测定参数时要分清哪些量已经被 Kepler 定律吸收。若只改变 $\Phi_3$ 在 $g_{00}$ 中的系数，单个定常球体的 $1/r$ 项仍可通过重新拟合 $GM$ 而符合数据。若让系统内部的能量形式改变，或令不同结合能占比的物体在同一外场自由落体，同一个质量重定义便不能普遍消去差别。实际实验的灵敏度在很大程度上取决于有没有独立的第二种质量比较方法。

\subsection{从物理度规标定引力常数}\label{sec:pn-bd-parameters}

以 \S\ref{sec:scalar} 的标量-张量理论为例。忽略标量势，使标量场在太阳系尺度上是长程的。远离源处令 $\phi=\phi_0$，记 $\omega_0=\omega(\phi_0)$。由标量场方程，源除了激发度规扰动，还激发一个与质量成正比的标量扰动。依第一章作用量的归一化，Cavendish 实验所测的引力常数为 $G=(2\omega_0+4)/[(2\omega_0+3)\phi_0]$。以此值定义 $U$，才可把空间度规同上节比较。结果为
\[
 \gamma=\frac{1+\omega_0}{2+\omega_0},\qquad
 \beta-1=\frac{\phi_0\omega'_0}{(3+2\omega_0)(4+2\omega_0)^2},\qquad
 \omega'_0=\left.\frac{\d\omega}{\d\phi}\right|_{\phi_0}.
\]
其它八个参数为零。Brans-Dicke 理论的 $\omega$ 是常数，故 $\beta=1$，$\gamma-1=-1/(2+\omega)$。在 $2\omega+3>0$ 的通常分支上，若测得 $|\gamma-1|<\delta$，就有 $\omega>1/\delta-2$。若 $\omega$ 随 $\phi$ 改变，太阳系实验只测得背景 $\phi_0$ 附近的耦合，不能据此外推中子星内部的标量场。

\subsection{标量场的两种度规表示}

取 $c=1$，用 $\phi_0$ 将背景归一化，定义辅助度规 $\tilde g_{\mu\nu}=(\phi/\phi_0)g_{\mu\nu}$。令 $\phi=\phi_0(1+\Psi)$ 且 $|\Psi|\ll1$。可在 $\tilde g$ 的谐和坐标求解，把标量的二阶导数从度规的线性方程中分离出去。记 $\tilde G=1/\phi_0$，$\tilde U=\tilde G\int\rho_*'\d[3]{x'}/s$。标量的源是能动量迹，冷物质的 $T$ 约为 $-\rho_*$。忽略标量自势后，标量方程的 Newton 阶为 $\nabla^2\Psi=-8\pi\tilde G\rho_* /(3+2\omega_0)$，故 $\Psi=2\tilde U/(3+2\omega_0)$。辅助度规在同一阶有 $\tilde g_{00}=-1+2\tilde U$、$\tilde g_{ij}=(1+2\tilde U)\delta_{ij}$。回到物质所用的物理度规，则 $g_{00}=-1+2\tilde U+\Psi$、$g_{ij}=(1+2\tilde U-\Psi)\delta_{ij}$。前一个加号使实测势 $U=\tilde U+\Psi/2$ 增大，后一个减号使单位实测势产生的空间曲率变小。两式的扰动之比即给 $\gamma=(1+\omega_0)/(2+\omega_0)$。若直接以 $\tilde g$ 追踪原子钟，会错误地读到 $\gamma=1$。即使采用 Einstein 表象使真空度规方程变得简单，也须在物质作用量中保留相应的标量耦合。

为比较一般的耦合函数，可另给标量规定使动能项为 $-2(\tilde\nabla\varphi)^2$ 的归一化。令 $g_{\mu\nu}=A^2(\varphi)\tilde g_{\mu\nu}$，定义 $\alpha_0=(\d\ln A/\d\varphi)_0$、$b_0=(\d\alpha/\d\varphi)_0$，这里的 $\alpha_0$ 与优选系参数 $\alpha_i$ 无关。线性近似中，每个弱自引力物体的标量源正比于 $\alpha_0m$，另一个物体对标量的响应也正比于 $\alpha_0$。故交换标量相对于交换张量的强度为 $\alpha_0^2$，静力实验量得 $G=G_*A_0^2(1+\alpha_0^2)$，其中 $G_*$ 是辅助度规作用量里的常数。改写两个 PPN 参数，有
\[
 \gamma-1=-\frac{2\alpha_0^2}{1+\alpha_0^2},\qquad
 \beta-1=\frac{b_0\alpha_0^2}{2(1+\alpha_0^2)^2}.
\]
第一式与 $\alpha_0^2=(3+2\omega_0)^{-1}$ 相容。第二式说明，$\alpha_0^2$ 若很小，附近弱场实验对耦合函数斜率 $b_0$ 的约束也会减弱。用中子星作源时，内部结合能占比很大，不能简单将弱场荷 $\alpha_0m$ 代给整个星体。后文的敏感度正是把这个缺口补上。

由 \S\ref{sec:scalar} 的辅助场变分，在 $f^{\prime\prime}(q)\ne0$ 的分支度规 $f(R)$ 等价于 $\omega=0$ 且带势的标量理论，$\phi=f_R(q)$，$V=\phi q-f(q)$。

若该标量在实验尺度上确是无质量的，立即有 $\gamma=1/2$。然而 $V$ 在背景处的曲率一般不为零，只报这个数便可能错判实验。

例如 $f(R)=R+aR^2$ 且 $a>0$，在平直真空附近线性化，标量的逆作用程为 $m_s=1/\sqrt{6a}$。对孤立、弱自引力的点源，忽略非线性屏蔽，时间和空间扰动分别正比于 $(1+\e^{-m_s r}/3)/r$ 和 $(1-\e^{-m_s r}/3)/r$。在 $m_s r\gg1$ 时附加场被指数抑制，在 $m_s r\ll1$ 时两个系数之比才趋于 $1/2$。用行星运动定质量还要对势求导，$\e^{-m_s r}$ 会变成 $(1+m_s r)\e^{-m_s r}$。因此用某个尺度的动力学质量去解释另一个尺度的光路，须在同一距离依赖模型里同时拟合。此处的 $m_s$ 是逆长度，恢复 $c$ 后作用程仍为 $1/m_s$。

若新增的场是单位类时矢量，它的背景值就选出了一个时间方向。以 Einstein-以太理论为例，其 $\gamma=\beta=1$，但 $\alpha_1,\alpha_2$ 一般随矢量场的耦合常数而变。两个理论因而可以给出同样的太阳光线偏折，却给出不同的双星轨道面进动。

\subsection{从一次标量响应到非线性参数}\label{sec:pn-bd-aux}

为看清 $\omega'_0$ 究竟怎样被近日点实验测到，以下用 $G_{\rm N}=c=1$，取标量的背景值为 $\phi_0$、归一化物理标量为 $\phi/\phi_0=1+\Psi$。引入两种不同的量
\[
 z=\frac{1}{4+2\omega_0},\qquad
 L_s=\frac{\phi_0\omega'_0}{(3+2\omega_0)(4+2\omega_0)} .
\]
$z$ 决定标量源相对于实测 Newton 势的强度，$L_s$ 决定源响应如何随局域标量值继续改变。可在辅助度规的谐和坐标中取 $H_*^{\mu\nu}=\eta^{\mu\nu}-\sqrt{-g_*}g_*^{\mu\nu}$。第一次迭代给 $\Psi=2zU$、$H_*^{00}=4(1-z)U$。乘回物质度规中的共形因子，时间扰动的两个贡献 $2(1-z)U+2zU=2U$ 正好落实 $G_{\rm N}$ 的定义，空间扰动则为 $2(1-2z)U$。所以取实测 $G$ 之后，同一光线的空间曲率仍知道标量分担了多大份额。

第二次迭代时，标量本身的梯度平方是引力场的源，且 $\omega=\omega_0+\phi_0\omega'_0\Psi+\cdots$ 同时进入标量的传播方程。把它们与能动量迹一起展开，最后将答案用物理守恒密度的势表示。在辅助度规谐和、宇宙背景静止的坐标中，有
\begin{align*}
 g_{00}&=-1+2U+2\{\psi_s-(1+zL_s)U^2\}+\ddot X,\\
 \psi_s&=\tfrac12(3-4z)\Phi_1-(1+2zL_s)\Phi_2+\Phi_3+3(1-2z)\Phi_4,\\
 g_{0i}&=-4(1-z)V_i,\qquad g_{ij}=[1+2(1-2z)U]\delta_{ij}.
\end{align*}
到标准坐标移去 $\ddot X$，或直接在 $a_{\rm H}=\gamma-1=-2z$ 的广义谐和规范中比较，都给 $\gamma=1-2z$、$\beta=1+zL_s$。$\Phi_2$ 和 $U^2$ 是两个不同的源，但同一个 $L_s$ 同时出现在它们的系数中。故两个任意函数若在当前背景有相同的 $\omega_0,\omega'_0$，它们用迥异的形式写在作用量中也可以有相同的 1PN 预言。要区分更高导数，须用高一阶的实验或取样于明显不同的背景。纯 Brans--Dicke 的 $\omega'_0=0$，单凭 $\beta$ 一项不能限制它；而光传播和标量荷所携带的双星辐射仍可限制它。

再问有势的情形。采用物理度规变分时，由 $f_R R-2f+3\square_g f_R=8\pi G T/c^4$ 可知，标量的质量应在\emph{具体背景}上读出。若 $R_0$ 是常曲率真空，先解 $f_R(R_0)R_0-2f(R_0)=0$，再对迹方程线性化。在 $f_{RR}(R_0)\ne0$ 时，扰动的 Klein--Gordon 方程给逆作用程平方 $m_s^2=[f_R(R_0)-R_0f_{RR}(R_0)]/[3f_{RR}(R_0)]$。代入此前的二次例子才得到 $1/(6a)$。若 $m_s^2<0$，把它代进指数衰减的 Yukawa 解本来就没有根据，应先检查所取背景能否长期存在。若质量随环境变，实验室、行星和宇宙学的界限也不能仅按一种固定作用程作图。到这里的计算都假设引力源对标量响应微弱；对强自引力星体的另一分支将在本章后半单独求解。

\subsection{由附加场到优选方向}

类时单位矢量 $u^\mu$ 要在度规的每个切空间中满足 $g_{\mu\nu}u^\mu u^\nu=-1$。这项约束固定的是模，不固定方向，故 $u^\mu$ 的空间扰动仍要解场方程。Einstein-以太理论在两导数阶可由 $\cd_\mu u_\nu$ 的四种缩并给出四个无量纲常数 $c_1,c_2,c_3,c_4$。前三种分别可取 $\cd_\mu u_\nu\cd^\mu u^\nu$、$(\cd_\mu u^\mu)^2$、$\cd_\mu u_\nu\cd^\nu u^\mu$，第四种是矢量线汇加速度 $u^\nu\cd_\nu u^\mu$ 的平方。可用乘子实施归一化。与一开始就指定一个不随物质改变的背景矢量相比，在变分时保留 $u^\mu$ 的方程是重要的区别。

在边界取 $u^\mu=(1,\bm0)$ 的坐标，静止球体只激发球对称的组合。将 Newton 常数换成静力实测值后，$\beta=\gamma=1$。要看剩下的效应，须让物体相对于这个边界运动。按通常 $c_i$ 的符号约定，令 $c_{ij}=c_i+c_j$、$c_{123}=c_1+c_2+c_3$，在非奇异的参数分支有
\begin{align*}
 G_{\rm N}&=G_{\rm bare}(1-c_{14}/2)^{-1},\qquad
 \alpha_1=-\frac{8(c_3^2+c_1c_4)}{2c_1-c_1^2+c_3^2},\\
 \alpha_2&=\tfrac12\alpha_1-
 \frac{(2c_{13}-c_{14})(c_{13}+c_{14}+3c_2)}
 {c_{123}(2-c_{14})}.
\end{align*}
其余 PPN 参数取广义相对论的值。第一行说明，改变 $G_{\rm bare}$ 的归一化就可以让行星在 Newton 阶照常运行；它不会把 $\alpha_1$ 也消掉。例如取 $c_4=-c_3^2/c_1$ 可令 $\alpha_1=0$，还需另一条关系才使 $\alpha_2$ 也为零。即使做到这一步，额外自由度的波速和能流还要分别同信号传播和双星衰减相比。PPN 的十个数描述的是弱场近区运动，不能替代理论对波本身所作的预言。

Bekenstein 的张量-矢量-标量理论（TeVeS）还有两种度规须分清。引力作用量用一个辅助度规写成，钟尺所测的度规则由辅助度规、矢量场和标量场共同决定。取后 Newton 极限时，应展开后一个度规。理论另引入加速度尺度 $a_0$，使绕孤立质量 $M$ 的物体在极低加速度下满足 $a^2/a_0\simeq GM/r^2$。对圆轨道代入 $a=v^2/r$，便得 $v^4\simeq GMa_0$，可见远处的转速何以不再随半径下降。太阳系所检验的则是 $a\gg a_0$ 时的另一极限。

标准 PPN 假定势的空间形式已经给定，也有不宜硬套的情形。例如有质量的标量可产生 $e^{-r/\lambda}/r$ 型势，其中 $\lambda$ 是作用程。若观测正跨越 $r\sim\lambda$，就应比较不同距离的引力，只给一对常数 $\gamma,\beta$ 无法表出这种变化。

\subsection{受约束与不受约束的矢量场}

单位类时矢量的条件究竟排除了什么？在只有一份物理度规的普适耦合理论中，取引力作用量（本小节用几何制）
\[
 S_g=\frac{1}{16\pi G_*}\int\!\sqrt{-g}\left[
 R-\mathcal K^{\alpha\beta}{}_{\mu\nu}
       (\nabla_\alpha u^\mu)(\nabla_\beta u^\nu)
       +\lambda(u^\mu u_\mu+1)\right]\d[4]{x},
\]
其中 $\mathcal K^{\alpha\beta}{}_{\mu\nu}=c_1g^{\alpha\beta}g_{\mu\nu}
+c_2\delta^\alpha_\mu\delta^\beta_\nu+c_3\delta^\alpha_\nu\delta^\beta_\mu
-c_4u^\alpha u^\beta g_{\mu\nu}$。乘子 $\lambda$ 的方程消去模的扰动，而矢量方程沿 $u^\mu$ 的投影定出 $\lambda$。在静系内只剩三个独立的方向变化。空间矢量又分为相对于一个给定波矢的横向与纵向部分。即使慢物质所见的静态 $g_{00}$ 恰与广义相对论相同，这些变化也能参与波动或优选系运动。

前文已将 $\alpha_1$ 写成四个 $c_i$ 的函数。若从实验出发，先令 $\alpha_1=\alpha_2=0$，可问四个 $c_i$ 是否全被消掉。取 $c_1\ne0$ 及分母非零的支，由 $\alpha_1=0$ 得 $c_4=-c_3^2/c_1$。再取 $c_2=(c_1+c_3)(c_3-2c_1)/(3c_1)$，则 $\alpha_2$ 也为零。此时仍可用 $c_+=c_1+c_3$、$c_-=c_1-c_3$ 作两个自由参数，$c_{14}=2c_+c_-/(c_++c_-)$。若只在一个无自转太阳外测光线与行星，根本看不见这二维参数空间。另有方程本身的限制：若某模的速度平方为负，小扰动即随时间指数增长；若动能为负，向外传播的波也未必带正能量。这些条件和由数据得到的上限须分开施加。

一个独立限制来自运动很快的粒子。若有与物质耦合的引力模显著慢于光速，高能粒子相对于以太静系可超过该模的相速度，进而辐射类似 Čerenkov 辐射的引力波而失能。在以太理论常用的允许支上，张量模的速度平方（以光速为一）是 $1/(1-c_{13})$。脉冲定时或光与引力波的联合观测约束这一传播速度；它并不只依赖 $\alpha_1,\alpha_2$ 本身。反过来，一个矢量理论即使传播速度按观测调得接近一，若未约束其它系数，仍可能让双星轨道面对优选系方向进动。

若 $u_\mu$ 还处处正交于某一族同时面，则存在局域标量 $T$ 使 $u_\mu=-\p_\mu T/\sqrt{-g^{\alpha\beta}\p_\alpha T\p_\beta T}$。由 \S\ref{sec:timelike} 的线汇分解，此时旋度为零。以时间叶层为基本自由度的低能理论称作 khronometric 理论。以 $\alpha_K,\beta_K,\lambda_K$ 记它的三个动力学耦合，常用定义下有
\[
 \alpha_1=\frac{4(\alpha_K-2\beta_K)}{\beta_K-1},\qquad
 \alpha_2=\frac{\alpha_1}{2}+
 \frac{(\alpha_K-2\beta_K)(\alpha_K+\beta_K+3\lambda_K)}
 {(2-\alpha_K)(\beta_K+\lambda_K)} .
\]
这里 $\gamma=\beta=1$，其余异常参数为零。两个分母非零时，取 $\alpha_K=2\beta_K$ 便能同时让两种优选系实验得到零，叶层却并未从作用量中消失。可见对“找到了一个特殊时间方向”的物理解释不能只依据公式中是否出现一个 $T$；还要问该方向通过什么耦合进入正在读的仪器。

反之若无模长约束，真空中的矢量方程可以是齐次的。空间均匀、类时的解具有任意背景模长 $K$，它的值和引力作用量中的系数一起进入 PPN 结果。早期一般矢量--张量模型容许 $K^\mu K_\mu R$、$K^\mu K^\nu R_{\mu\nu}$、反对称导数 $F_{\mu\nu}F^{\mu\nu}$ 和 $\nabla_\mu K_\nu\nabla^\mu K^\nu$。由两曲率项和导数项作分部积分，若干不同的写法其实等价。在 Will--Nordtvedt 的无约束模型中，取相应参数 $c_1=-1$、$c_2=c_3=c_4=0$，有 $\gamma=\beta=1$ 而 $\alpha_2=K^2/(1+K^2/2)$，这里 $K^2=-K^\mu K_\mu$ 为背景正数。故轨道面和毫秒脉冲星给出的上限，在弱场框架下限制的是这个背景模长。Hellings--Nordtvedt 的选择则一般可使 $\gamma,\beta,\alpha_1,\alpha_2$ 同时变化。若令相同的 $K$ 在不同实验中任意重选，就不再是在同一宇宙匹配条件下检验同一理论。

\subsubsection{从曲率耦合换到矢量导数的耦合}

一些早期矢量理论用 $(\omega,\eta_K,\epsilon_K,\tau_K)$ 标记作用量。为与前面的 $c_i$ 对照，以下取 $x^0=ct$、度规号差 $(-+++)$，把引力作用量写作 $S_g=c^3(16\pi G_*)^{-1}\int\sqrt{-g}\,\mathcal L_g\,\d^4x$，物质项另加。$G_*$ 是作用量中的裸常数，$c_i$ 和上述四个系数均无量纲，矢量 $K^\mu$ 也无量纲。置 $K^2_{\rm L}=K^\mu K_\mu$（类时场为负），$F^K_{\mu\nu}=2\nabla_{[\mu}K_{\nu]}$。若无模长约束，一种常见的写法为
\[
 \mathcal L_g=(1+\omega K^2_{\rm L})R+
     \eta_K K^\mu K^\nu R_{\mu\nu}
     -\epsilon_K F^K_{\mu\nu}F_K^{\mu\nu}
     +\tau_K\nabla_\mu K_\nu\nabla^\mu K^\nu .
\]
$F^K$ 与 Maxwell 场强有相同的反对称形式，但其它项一般不允许 $K_\mu\mapsto K_\mu+\nabla_\mu f$，不能据此删去一个纵向自由度。

为何作用量中似乎比 $c_i$ 多了一种曲率缩并？记 $I_1=(\nabla_\mu K_\nu)(\nabla^\mu K^\nu)$、$I_2=(\nabla_\mu K^\mu)^2$、$I_3=(\nabla_\mu K_\nu)(\nabla^\nu K^\mu)$。有 $F_K^2=2(I_1-I_3)$，而交换两个协变导数又有
\[
 K^\mu K^\nu R_{\mu\nu}
 = I_2-I_3-\nabla_\mu\big(K^\mu\nabla_\nu K^\nu
                          -K^\nu\nabla_\nu K^\mu\big).
\]
于是将导数项写为 $-c_1 I_1-c_2 I_2-c_3 I_3$，对应 $c_1=2\epsilon_K-\tau_K$、$c_2=-\eta_K$、$c_3=\eta_K-2\epsilon_K$。特别有 $c_{123}=-\tau_K$。一般二导数作用量还可含 $c_4(K^\mu\nabla_\mu K_\nu)(K^\alpha\nabla_\alpha K^\nu)$，上列无约束子类取 $c_4=0$。若再加乘子令 $K^2_{\rm L}=-1$，则 $1+\omega K^2_{\rm L}=1-\omega$ 为常数。在它非零时可将整个引力部分重新归一化，同时重标 $G_*$ 和相对于 Ricci 项计量的导数耦合。若矢量模长仍可变，这样做就不再能消去 $\omega$。例如 Hellings--Nordtvedt 的子类取 $c_1=2$、$c_2=2\omega$、$c_3=-2(1+\omega)$，三个数的和为零，但这并不是单位以太理论中可以直接代入的特殊点。去掉模长条件后，宇宙边界处的 $-K^2_{\rm L}$ 也成为物理解的一份资料。

\subsection{低加速度规律与物理度规}

把银河中一个重子的运动规律从 $a=GM/r^2$ 改成 $a^2/a_0=GM/r^2$ 并未给出光线的运动规律。以星系团的弱透镜质量和星系内恒星运动互比，正需要二者。TeVeS 的做法是以辅助度规 $\tilde g_{\mu\nu}$ 和相对于它归一化的矢量 $K^\mu$、标量 $\phi$ 构造 $g_{\mu\nu}=\e^{-2\phi}\tilde g_{\mu\nu} -2\sinh(2\phi)K_\mu K_\nu$。在 $K$ 的局域静系里，第一项使时间和空间一同伸缩，第二项只作用于类时方向。当 $K^\mu=(1,\bm0)$ 且辅助度规局域为平直时，简单相加即有 $-g_{00}=\e^{2\phi}$、$g_{ij}=\e^{-2\phi}\delta_{ij}$。故标量不仅能改变钟速，也改变光锥相对于辅助度规的开口。若只使用一个共形因子，这两项会在零测地线计算中相消。

标量作用量取为某函数 $F(y)$ 的积分，其中 $y=k\ell^2(\tilde g^{\mu\nu}-K^\mu K^\nu)\p_\mu\phi\p_\nu\phi$，$k$ 为无量纲耦合，$\ell$ 为长度尺度。在静止弱场中，变分产生一个流 $\mu(y)\bm\nabla\phi$，$\mu=\d F/\d y$。对球对称质量分布，场方程一次积分便有 $r^2\mu(y)\,\d\phi/\d r\propto M(<r)$。若梯度较大处 $\mu\to1$，就得到普通的 $1/r^2$ 标量力；若小梯度处 $\mu$ 与 $|\bm\nabla\phi|$ 成正比，才得到 $|\bm\nabla\phi|\propto1/r$。两个条件规定的是插值函数的两端，中间形状要由横跨不同加速度的资料决定。附加尺度 $a_0$ 给出一个区分运动规律的加速度，而非一个对所有星系相同的转折半径。

在高加速度端，广义矢量作用量的常数也记作 $c_i$，$\kappa=k/(8\pi)$。把外部背景取为 $\phi_0$ 后，实验室量得
$G_{\rm N}=2G_*[1+\kappa(2-c_{14})]/(2-c_{14})$。这一极限有 $\gamma=\beta=1$，$\xi,\alpha_3,\zeta_i$ 都为零。

\subsubsection{两个优选系参数的共同归一化}

把高加速度解化到物理标准 PPN 坐标时，时间和空间的变换分别含 $\exp(\phi_0)$ 与 $\exp(-\phi_0)$，而质量势须用实测 $G_{\rm N}$ 定义。为使耦合和常数的量纲一目了然，令
$D=2c_1-c_1^2+c_3^2$、$d=2-c_{14}$、$q=2c_{13}-c_{14}$、$\kappa=k/(8\pi)$，
并记 $g_*=G_*/G_{\rm N}=d/[2(1+\kappa d)]$。以下只讨论给出通常 Newton 极限且 $D\,d\,c_{123}\ne0$ 的支。取 $s_4=\sinh(4\phi_0)$、$s_2=\sinh^2(2\phi_0)$。以前文只保留单位矢量时的结果为上标 $\mathrm{AE}$，则同一套归一化给出
\begin{align*}
 \alpha_1-\alpha_1^{\rm AE}
     &=-\frac{16g_*}{D}
            \left[\kappa c_1 d-c_3s_4+2(1-c_1)s_2\right],\\
 \alpha_2-\alpha_2^{\rm AE}
     &=-2g_*\left[A_1 \kappa-2A_2s_4-A_3s_2\right],\\
 A_1&=\frac{q^2}{c_{123}d}+\frac{4c_1d}{D}
                    -\frac{6(1+c_{13}-c_{14})}{d},\\
 A_2&=\frac{q^2}{c_{123}d^2}-\frac{4(1-c_1)}{D}
                +\frac{2(1-c_{13})}{d}\left(\frac{2}{c_{123}}+\frac{3}{d}\right),\\
 A_3&=\frac{q^2}{c_{123}d^2}+\frac{4c_3}{D}
                +\frac{2}{d}\left[\frac{3(1-c_{13})}{c_{123}}-\frac{q}{d}\right].
\end{align*}
在 $G_{\rm N}=c=1$ 的单位中，$g_*$ 的数值就是前式所写的 $G_*$。若保留 SI 单位，PPN 参数仍是无量纲数，前因子是裸常数与实测常数的比，不应留下孤立的 $G_*$，也不需给双曲函数中的 $\phi_0$ 添加 $c$ 的幂次。这个写法还可作两项直接检验。在 $\kappa=\phi_0=0$ 时，$g_*=d/2$，确有 $G_{\rm N}=G_*/(1-c_{14}/2)$，且两修正同时消失。在固定 $\kappa$、很小的 $\phi_0$ 下，背景的线性响应来自 $s_4\simeq4\phi_0$，第二种背景函数 $s_2\simeq4\phi_0^2$ 至下一阶才出现。调节一个 $c_i$ 消去 $\alpha_1$ 的线性响应，不保证 $A_2$ 也消失。故同一天体群给出两个优选系参数的限制，比只知道它们之一为零多了一条独立条件。

插值还有一个即使忽略优选速度也存在的后果。令系统外有近于常数的场 $\bm g_e$，在 $\bm g_e$ 上叠加内部小场 $\bm g_i$。线性化通量 $\mu(|\bm g|)\bm g$，得 $\delta(\mu g_j)=\mu_e(\delta_{jk}+L_e\hat g_{e j}\hat g_{e k})g_i^k$，其中 $L_e=(\d\ln\mu/\d\ln|\bm g|)_e$。解内部 Poisson 型方程时，平行与垂直外场的二阶导便可取不同系数。对太阳系，需要连接内部高加速度解与外部银河场占优势的区域，不能把上式的小内场近似一直用到日心。其回到内区的领先异常势可呈无迹四极形式 $\delta U=(Q/2)[(\hat{\bm g}_e\cdot\bm x)^2-r^2/3]$。相应相对加速度与 $r$ 成正比，沿银河方向的分量又不同于横向分量。这称为外场效应。土星测距限制的是按给定插值函数与银河边界条件算出的 $Q$，故即使在行星处 $a\gg a_0$，仍有实验可作。

\subsection{曲率修正、有质量的引力与双度规}

有没有一种修正，连 \eqref{eq:exp-ppn-all} 的势都不足以表达？考虑将一个标量与曲率平方耦合。四维中 $R_{\mu\nu\lambda\sigma}R^{\mu\nu\lambda\sigma}-4R_{\mu\nu}R^{\mu\nu}+R^2$ 的积分在常系数下给出 Gauss--Bonnet 拓扑项，对局域场方程无贡献。把常系数改成场的函数，分部积分会留下该函数的导数，才改变解。另一个例子是与 Pontryagin 密度 $\prescript{*}{}{R}^{\alpha\beta\mu\nu}R_{\alpha\beta\mu\nu}$ 耦合，其中星表示对 Riemann 张量的一对指标取对偶。此时旋转特别重要，因为单个无自转球体缺少能形成这一赝标量的方向。对一种取给定时间背景的 Chern--Simons 模型，弱场 $g_{0i}$ 会有一项正比于 $(\bm\nabla\times\bm V)_i$。标准 PPN 度规只把 $\bm V$ 和 $\bm\nabla\dot X$ 当作相应的矢量势。把十个标准系数全部设为广义相对论的值，也没有除掉这项。陀螺和旋转天体附近的轨道才直接取样这个度规分量。若标量是动力学场，其边界条件和方程也须同时求解，上述背景模型的系数不能照搬。

有质量自旋二场所遇到的是另一种困难。考虑线性 Fierz--Pauli 模型，用相同的两个守恒源 $T_{\mu\nu},T'_{\mu\nu}$ 比较引力交换。在几何制中，无质量自旋二场的交换振幅分子正比于 $T_{\mu\nu}T'^{\mu\nu}-TT'/2$，有质量场即使质量很小也给 $T_{\mu\nu}T'^{\mu\nu}-TT'/3$。对两个冷物质源，$T\simeq-\rho$，两者的静力系数分别为 $1/2$ 和 $2/3$。若用静力实验把两理论的 $G$ 校准为相同，有质量理论的裸耦合便只需无质量理论的 $3/4$。对光却有 $T'=0$，第二项本来就无贡献。因此线性理论在作用程远大于冲量参数时仍只给广义相对论 $3/4$ 的光偏折，也即 $\gamma=1/2$。这称为 van Dam--Veltman--Zakharov 不连续性，来自额外的纵向自由度。

在线性近似失效的区域，结论还可能改变。有些非线性有质量引力模型的纵向场在源附近被压低，称为 Vainshtein 屏蔽。例如在含三次 Galileon 自相互作用的模型中，Schwarzschild 半径 $r_s$ 与引力子的约化 Compton 波长 $\lambda_g$ 可组合出 $r_V\sim(r_s\lambda_g^2)^{1/3}$。在 $r\ll r_V$ 时须把非线性场方程的解与远区线性解匹配。由此可知，太阳系内的偏折、星系尺度的势以及波的飞行时限制，分别取样于解的不同部分。由波形限制色散关系时也须说明假定的发射波形，不能将色散上限直接当作全部有质量理论的静力上限。

最后比较几个历史例子。若物理度规始终只是平直度规乘同一个标量函数，$\d s^2=A^2(\phi)\eta_{\mu\nu}\d x^\mu\d x^\nu$，则类光曲线不随共形因子改变。这样的纯标量理论能给慢质点以加速度，却不产生太阳附近实测的光线偏折。Whitehead 理论则走了相反的一条路。它在孤立点质量外给出 Schwarzschild 几何，经典的单体检验均可符合。一旦放入外界物质，\eqref{eq:exp-whitehead} 中取向相关的项不再消失。在标准 PPN 规范中，这类准线性构造可把 $\xi$ 锁定在 $\beta$，Whitehead 的选择给 $\xi=\beta=1$。使 $\beta$ 接近近日点所要求的值，就预言了重力仪和月球轨道本应看见的异常；两项实验不能同时满足。

Rosen 双度规理论用一份平直度规规定引力场方程的背景，而物质钟尺只和另一份度规耦合。在平直背景取 Minkowski 分量的坐标中，物理度规的远处时间、空间尺度可分别为 $c_0,c_1$，即背景物理线元写成 $-c_0 c^2\d t^2+c_1\d\bm x^2$。把钟尺重新归一后留下的优选系参数为 $\alpha_2=c_0/c_1-1$。调节背景可使太阳系弱场效应很小，强自引力双星的运动和辐射却仍能与观测不符。于是用一个物体校准 Newton 常数，用第二个物体校准空间曲率，仍有第三种系统能检验同一套动力学方程。

\subsection{由无量纲参数转向有长度量纲的耦合}

对曲率平方理论，先要决定把它当作任意曲率下的精确作用量，还是在低能只保留首项的有效理论。在后一种解释中，取曲率尺度 $L_R$ 满足 $|R_{\mu\nu\alpha\beta}|\sim L_R^{-2}$，若相对于 $R$ 的新项为 $\ell^2 R_{\mu\nu\alpha\beta}R^{\mu\nu\alpha\beta}$，只要 $\ell^2/L_R^2\ll1$ 就可按耦合 $\ell^2$ 再作一层展开。先以 Newton 或 Einstein 解作零阶背景，算出标量源或修正应力，用其求度规的一阶偏离。若迭代中又保留了一部分同未知高阶算符一样大的项，那么从得到的精确根推断新分支就超出了这种有效理论的精度。

\subsubsection{二次曲率不变量的系数怎样比较}

仍以 $c^3/(16\pi G_*)$ 乘在整个引力作用量前，并选标量 $\vartheta$ 为无量纲。设 $b_i(\vartheta)$ 具有长度平方的量纲，$b_s$ 无量纲，$V_s(\vartheta)$ 的量纲为长度的负二次方。于是包括最低阶标量动能的一种统一写法是
\begin{align*}
 \mathcal L_g={}&R+b_1 R^2+b_2 R_{\mu\nu}R^{\mu\nu}
       +b_3 R_{\mu\nu\rho\sigma}R^{\mu\nu\rho\sigma}
       +b_4\,{}^*RR \\
       &-\tfrac12 b_s\left(g^{\mu\nu}\nabla_\mu\vartheta\nabla_\nu\vartheta
                         +2 V_s(\vartheta)\right),\qquad
 {}^*RR=\tfrac12\varepsilon^{\gamma\delta\rho\sigma}
             R^\alpha{}_{\beta\rho\sigma} R^\beta{}_{\alpha\gamma\delta}.
\end{align*}
$\varepsilon^{\gamma\delta\rho\sigma}$ 为随定向选定的体积张量。若改把 Ricci 项外的 $1/(16\pi G_*)$ 留在被积式内，文献中的二次曲率系数还须除以 Ricci 项系数，才是这里的 $b_i$。每项同 $R$ 都有相同的量纲。对标量变分得 $b_s(\square\vartheta-V_s')=-\sum_i b_i'\mathcal I_i$，其中 $\mathcal I_i$ 按上式四个曲率不变量的次序排列，撇对 $\vartheta$ 求导。由此可读出某一给定几何能否直接激发这个标量。

在四维把 $b_1=b_3=b(\vartheta)$、$b_2=-4b(\vartheta)$、$b_4=0$，就得到标量耦合的 Gauss--Bonnet 组合。若 $b$ 为常数，在给定拓扑且正确计入边界项时，这项不产生局域体场方程。若 $b'\ne0$，一般既会激发标量，也会通过 $b$ 的导数反馈到度规。若改令 $b_1=b_2=b_3=0$、$b_4=\ell_{\rm CS}^2\vartheta/4$，则为线性标量耦合的 Chern--Simons 例。取 $b_s\ne0$ 时标量有自身的方程，取 $b_s=0$ 而继续将它变分时则出现 $\,{}^*RR=0$ 的约束。若把 $\vartheta$ 直接指定为时空上的函数，应从度规方程、Bianchi 恒等式和物质守恒检验相容性。有非零背景梯度时也会受到 Pontryagin 约束。在该非动力学时间背景的标准弱场解中令 $f_{\rm CS}=4b_4=\ell_{\rm CS}^2\vartheta$，则以长度坐标 $x^0=ct$ 和本章含 $G_{\rm N}$ 的 $V_i$ 表示，附加项可写为 $\delta g_{0i}=2(\partial_0 f_{\rm CS})[\bm\nabla\times(\bm V/c^3)]_i$。系数和标量归一的选择已在 $f_{\rm CS}$ 中指定，量纲为长度的 $\partial_0 f_{\rm CS}$ 正好抵消旋度的逆长度。这一个与宇称有关的质量流项超出了标准 PPN 势表。

最后，只有 $b_1=a$ 为常数也包括 $f(R)=R+aR^2$ 的低曲率模型。Ricci 平坦的真空使它的额外源项消失，而上述 Gauss--Bonnet 组合在同一真空未必消失。此例的标量来自度规迹方程的附加自由度，原作用量中毋须另置一个有动能的 $\vartheta$。这和已预设一个独立标量再把它耦合到曲率平方，是两种不同的变分问题。

若只有 Gauss--Bonnet 组合与标量耦合，分部积分不能消去非恒定的耦合函数，标量的源正比于该曲率组合。孤立黑洞外的 Ricci 张量虽为零，Riemann 张量平方仍非零，故与 $R^2$ 单项发生作用的情形不同。这正好解释了为什么两种理论在普通弱场物体的 PPN 参数上可能不易分辨，但在双黑洞的偶极辐射或强场外部几何上却可能有不同的信号。比较不同文献中对某个长度的上限，要先看它乘哪一个曲率不变量，是否有长程传播的标量，以及被用作源的是普通星还是黑洞。

还可用作用量说明 Rosen 理论与 Whitehead 理论在何处走上不同的道路。在 Rosen 理论中以平直背景的导数 $|$ 构造 $\eta^{\mu\nu} g^{\alpha\beta}g^{\gamma\delta}
(g_{\alpha\gamma|\mu}g_{\beta\delta|\nu}-g_{\alpha\beta|\mu}g_{\gamma\delta|\nu}/2)$，引力部分体元是 $\sqrt{-\eta}\,\d[4]{x}$；物质部分仍取物理度规 $g$。实验室的钟尺吸收 $C_0,C_1$ 的均匀缩放，但两份度规的相对本征值可由 $\eta^{-1}g$ 的迹等不变量表出，$\alpha_2=C_0/C_1-1$ 正是这一残余。若二者在边界完全重合，太阳系的十参数可以相同，紧致双星在该理论下却可能产生负能流的偶极辐射，预言轨道周期增加，与计时方向相反。选定一个背景来通过一项检验，并不许可在同一理论的另一项检验中自行改变辐射的能量符号。

在 Whitehead 的构造中，反而是用背景 $\eta$ 的过去光锥决定场点收到哪一个物质元的信息，再用这些推迟物质元构造物理度规。早期还可另假设钟怎样走；把它作为度规理论来测时，必须规定钟、光、试验体全用 $g$。静止点源的光锥积分可得与 Schwarzschild 等价的度规，但一般多源的非线性叠加、有限天体自身尺寸的贡献、引力辐射并未由该相等性决定。具体地，模型的优选位置效应给出过大的异常地潮，自能项引起过大的地月轨道极化，依赖源半径的外场项还会在卫星近地点上累积；偏心双星可同时试验质心自加速和预言的轨道反阻尼。五种读数利用不同的物体和时标，却涉及同一组不能各自更改的场方程。因而“给出正确的静止球体解”至多是在这个特定边界问题上过了一关。

还有一些理论把物理度规在宇宙时标 $T$ 的等时面上固定成各向同性形式 $g=(B-A)\d T\otimes\d T+B\eta$。在 $T$ 的静系中不论源如何运动，度规都无独立的 $g_{0j}$ 质量流项。变到标准 PPN 规范可生成 $\dot X_{,j}$，但不能任意生成独立的 $V_j$。对照 \eqref{eq:exp-ppn-all}，须有 $\alpha_1=-4(1+\gamma)$；而星光和行星要求 $\gamma$ 近于一，于是这类固定层结构强迫出很大的优选系效应。这里还没用到某一颗星的精确轨道，只比较度规可容许的势和实测系数，就已能排除整个类别。

\section{PPN 中的运动方程}

\subsection{坐标时与光的传播}\label{sec:pn-light}

先看光线。以 $\d\ell=|\d\bm x|$ 表示坐标欧氏线元，$\bm n=\d\bm x/\d\ell$ 表示光线方向。对静止弱场，只保留 $U/c^2$ 的一次项。由 $\d s^2=0$ 得 $c\,\d t=[1+(1+\gamma)U/c^2]\d\ell$。再由零测地线方程，方向的变化为 $\d\bm n/\d\ell=(1+\gamma)\bm\nabla_\perp U/c^2$，其中 $\bm\nabla_\perp U=\bm\nabla U-\bm n(\bm n\cdot\bm\nabla U)$。沿未扰动光路积分，纵向的光行时和横向的偏折都只含 $1+\gamma$。前一式的 $\d\ell/\d t$ 是坐标速率，当地自由落体观者仍测得光速 $c$。

式 \eqref{eq:pn-coordinate-geo} 中沿 $\bm q$ 的一项来自重选参数。若只将仿射参数的记号直接改成 $t$，连光在静态势中的坐标速率也会算错。在前述各向同性空间度规中，令 $\bm q=|\bm q|\bm n$，则零线元先给出 $|\bm q|=c[1-(1+\gamma)U/c^2]$。把上式投影到与 $\bm n$ 正交的平面，所有平行于 $\bm q$ 的项都消去，便回到方向方程。故沿一条未扰动光路，定时取样的是 $U$ 的积分，量角取样的是横向梯度的积分。这里用了一次迭代。若求二阶偏折，还要把一阶修正过的位置带回梯度，不能只在原直线上再积分一次二阶度规。

进一步，给定发射点和接收点，比给定发射方向更接近测距实验。对静态、透明的引力场，定义 $T(\bm x_e,\bm x_r)$ 为连接两点的传播时间。最低非平凡阶有
\[
 T=\frac{D}{c}+\frac{1+\gamma}{c^3}\int_0^D
 U(\bm x_e+s\bm k)\,\d s,\qquad
 \bm k=(\bm x_r-\bm x_e)/D,\quad D=|\bm x_r-\bm x_e| .
\]
在这个积分中，弯曲带来的路程变化为二阶量，因为在固定端点的直线附近路长的一阶变分为零。若从同一振荡器连续发射两个相邻波峰，即有 $t_r-t_e=T(\bm x_e(t_e),\bm x_r(t_r))$。对发射时刻微分，得到
\[
 \frac{\d t_r}{\d t_e}
  =\frac{1+\bm v_e\cdot\bm\nabla_e T}
         {1-\bm v_r\cdot\bm\nabla_r T},\qquad
 \frac{\nu_r}{\nu_e}=
 \frac{(\d\tau/\d t)_e}{(\d\tau/\d t)_r}\frac{\d t_e}{\d t_r}.
\]
只代入 $D/c$，就给出通常的视向 Doppler 因子。代入余项，则光线即使没有改变频率标准，也会因端点扫过不同的光路而产生附加频移。这正是下节利用上合信号测 $\gamma$ 的根据。若度规本身随时刻显著改变，$T$ 还显含发射时刻，在微分时须另加这一项。

接收端的角也须用当地的标架。设两束光同时到达，波矢为 $k^\mu,l^\mu$，接收者的 4-速 $u^\mu$ 满足 $u^\mu u_\mu=-c^2$。把两个波矢投影到其静止空间，并以各自所测频率归一，夹角满足 $\cos\theta=1+c^2(k\cdot l)/[(u\cdot k)(u\cdot l)]$。这样写能区分三件事：光路穿过太阳势时的偏折、观者运动引起的光行差、所选坐标在接收端的局域度规。单把两个坐标方向作欧氏点积，只在已经换到同一正交标架时才等于望远镜所测的角。

还可把前面的光线积分实际做到接收点，而不假定光来自无穷远。给定发射点 $\bm x_e$ 及初始欧氏单位方向 $\bm k$，用零阶直线参数 $\ell=c(t-t_e)$。对静止源元 $\d m$，令 $\bm s_e=\bm x_e-\bm x'$、$z_e=\bm k\cdot\bm s_e$、$\bm b=\bm s_e-z_e\bm k$。沿直线 $z=z_e+\ell$、$s=\sqrt{z^2+b^2}$，$\bm b$ 不变。借助 $\d(z/s)/\d\ell=b^2/s^3$ 和 $\d\ln(s+z)/\d\ell=1/s$，可将求偏折和求传播时间的两个积分同时化为端点项。令 $\delta\bm n(0)=0$，即有
\begin{align*}
 \delta\bm n(\ell)&=-\frac{(1+\gamma)G}{c^2}\int
           \frac{\bm b}{b^2}\left(\frac{z}{s}-\frac{z_e}{s_e}\right)\d m,\\
 \delta\bm x_\perp(\ell)&=-\frac{(1+\gamma)G}{c^2}\int
           \frac{\bm b}{b^2}\left(s-s_e-\ell\,\frac{z_e}{s_e}\right)\d m,\\
 \delta x_\parallel(\ell)&=-\frac{(1+\gamma)G}{c^2}\int
           \ln\frac{s+z}{s_e+z_e}\,\d m .
\end{align*}
第二行是在给定初方向下的横移，第三行是在同一坐标时刻比较的纵向行程差。求一去一回的时延时，它们与固定端点的边界条件不能混用。若发、收两点既定，零阶方向应先取为端点之差的单位矢 $\bm k_0$。计算以上积分后，还要把发射方向加 $\delta\bm k=-\delta\bm x_\perp(D)/D$，这样射线才真正到达接收点，$D$ 为两点间零阶距离。在接收点看到的方向是一整条光路的方向变化再加上这点发射修正。对于遥远射电源它常可忽略，行星间短弧、双星近旁的有限距离光源则未必如此。

几何上还有一个数值细节。在接收者处 $z$ 很负时，$s+z$ 是两个大数的差。可用 $(s+z)(s-z)=b^2$，把对数改算为 $\ln[(s_e-z_e)/(s-z)]$；同一条直线的 $b$ 在分子分母相消。两个形式在穿过 $z=0$ 时连续。若积分中包含多个离光路较近的质量元，$\bm b,z_e$ 应对每个源元另求，不能拿整个太阳系质心的冲量参数乘总质量。这是从孤立点源计算进入实际光路模型时所需的第一步。

\subsection{从度规求联络：光与慢物体要留到哪一项}

有时只想从一份已求得的度规模型生成定轨方程。为避免在每一行都写十个参数，令 $B_i=c^3g_{0i}$ 和 $K=c^4(g_{00}+1-2U/c^2)$ 分别取其最低非平凡阶；它们须用 \eqref{eq:exp-ppn-all} 给出的完整势及优选系项。在 \S\ref{sec:pn-components} 的分量展开中取 $h=2U/c^2$、$s_{ij}=2\gamma U\delta_{ij}/c^2$、$b_i=B_i/c^3$、$k=K/c^4$，即得所需各联络。空间指标用 $\delta_{ij}$ 升降。值得核查的是 $\Gamma^i{}_{00}$ 的最后一项。即使 $g_{00}$ 中的 $K$ 已包括非线性，还须展开逆空间度规，$g^{ij}=\delta^{ij}(1-2\gamma U/c^2)$，才会有 $2\gamma UU_{,i}$。另一处是 $\Gamma^0{}_{ij}$ 与 $\Gamma^i{}_{0j}$ 中 $\p B$ 的对称性不同。若给两者都取旋度，重排时间坐标产生的梯度势将无法正确相消。

代入以坐标时写成的测地线方程，有质量的试验体满足
\begin{align}
 a_i={}&U_{,i}+\frac1{c^2}\big\{(\gamma v^2-2\gamma U)U_{,i}
          +\tfrac12K_{,i}-\dot B_i
          +v^j(\p_iB_j-\p_jB_i)\nonumber\\
       &\hspace{27mm}-v_i[(2\gamma+1)\dot U+
                          2(1+\gamma)v^jU_{,j}]\big\}+O(c^{-4}).\label{eq:pn-general-motion}
\end{align}
例如一团运动物质产生的引力磁力不仅有 $\bm v\times(\bm\nabla\times\bm B)$，即使把试验体瞬时停下仍会有 $-\dot{\bm B}$。若改用一只运动钟的固有时，最后一组沿 $\bm v$ 的系数也会重排。原则上可代 $K=2\psi-2\beta U^2+\Phi_{\rm F}$、$B_i$ 的含参表达式直接积分任意测试粒子。下文对有压强的流体和自引力显著的扩展物体还要增加内部项，不能把这里的试验体方程套给每个天体的质量中心。

光则另有 $|\bm v|$ 在零阶就是 $c$ 的计数。把 $|\bm v|=c$ 代进一个已按 $v/c\ll1$ 截断的高阶慢体方程，一般不能得到同等精度的光学公式，因为慢体中忽略的高次速度项可能已经降至同一阶。求光的首个弱场偏折，只需先用零线元求纵向速率，再在横向方程保留 $U/c^2$ 的一次项。若求运动太阳或行星的更高阶光时，便须增加移动源、矢量势及端点在同一坐标规范内的修正。

\subsection{由单体到多体}

有质量的测试粒子不同。把 \S\ref{sec:exp-ppn} 的静止球体度规代入前述拉氏量，除去常数并除以质量，有
\[
 \frac{L}{m}=\frac{v^2}{2}+U
   +\frac{1}{c^2}\left[\frac{v^4}{8}+
   \left(\gamma+\tfrac12\right)Uv^2+
   \left(\tfrac12-\beta\right)U^2\right]+O(c^{-4}).
\]
第一项修正 $v^4/(8c^2)$ 在平直时空就有，后两项才将 $\gamma,\beta$ 带入轨道方程。

行星本身也有引力，何时可把它当成一个质点？对小流体元，$\cd_\nu T^{\mu\nu}=0$ 的 Newton 极限是 Euler 方程 $\rho\,\d\bm v/\d t=\rho\bm\nabla U-\bm\nabla p$，其中 $\d/\d t=\p_t+\bm v\cdot\bm\nabla$。将此方程及其 1PN 修正对一个天体的体积积分，表面压强为零时 Newton 阶的压强合力消失，各质元间的 Newton 引力也成对抵消。到下一阶须把内能、内部运动和引力结合能一并计入。对广义相对论中相距很远、近球形的无自旋天体，这些内部量在此阶可归入各自的质量，因而质心运动可用点质量描述。一般度规理论未必有这一性质，地月在太阳外场中的运动便能检验它。

在孤立静态球场中，可将粒子拉氏量的速度导数写成 $\bm p/m=\bm v[1+(v^2/2+(2\gamma+1)U)/c^2]$。求 $\d\bm p/\d t$ 时，前面的括号和 $\bm v$ 都须微分；在已乘 $1/c^2$ 的项里则可用 $\dot{\bm v}=\bm\nabla U$。整理得
\[
 \dot{\bm v}=\bm\nabla U+\frac{1}{c^2}
 \left\{[\gamma v^2-2(\beta+\gamma)U]\bm\nabla U
          -2(1+\gamma)(\bm v\cdot\bm\nabla U)\bm v\right\}
 +O(c^{-4}).
\]
第一项保持 Newton 形式。后面可分出一个总是沿场强的分量，另一个则随 $\bm v$ 改变方向。在圆轨道上 $\bm v\cdot\bm\nabla U=0$，只量某一个半径处的平均运动，修正易与场源质量混同。加入离心率不为零的轨道、或把测量延至不同半径，才能看到它随轨道相位的变化。令 $\beta=\gamma=1$ 即与前述二体方程的测试质量极限相符。

对多体还会有什么新内容？暂只许 $\gamma,\beta$ 偏离广义相对论，忽略自引力引起的质量差、潮汐和自转。令 $r_{ab}=|\bm x_a-\bm x_b|$、$\bm n_{ab}=(\bm x_a-\bm x_b)/r_{ab}$、$U_a=\sum_{b\ne a}Gm_b/r_{ab}$。省去静能后，标准坐标下一种 1PN 拉氏量为
\begin{align}
 L={}&\sum_a m_a\left(\frac{v_a^2}{2}+\frac{v_a^4}{8c^2}\right)
       +\frac12\sum_{a}\sum_{b\ne a}\frac{Gm_am_b}{r_{ab}}
       -\frac{2\beta-1}{2c^2}\sum_a m_a U_a^2 \nonumber\\
 &+\frac{1}{4c^2}\sum_a\sum_{b\ne a}\frac{Gm_am_b}{r_{ab}}
 \left[(2\gamma+1)(v_a^2+v_b^2)-(4\gamma+3)\bm v_a\cdot\bm v_b
 \right.\nonumber\\[-1pt]
 &\hspace{51mm}\left.-(\bm n_{ab}\cdot\bm v_a)(\bm n_{ab}\cdot\bm v_b)\right].
 \label{eq:exp-nbody-lag}
\end{align}
$\sum_a\sum_{b\ne a}$ 把每对质量计算两次，故 Newton 势能项有系数 $1/2$。把两体坐标改为质心坐标和相对坐标、置 $\gamma=\beta=1$，在质心系中便还原为 \eqref{eq:exp-rel-L}。也可令一星质量趋于零、其余只剩静止中心体，单位测试质量的拉氏量就还原为上式以前的球场结果。两项检验用的是不同极限，尤其能核对 $\bm v_a\cdot\bm v_b$ 与 $U_a^2$ 前的系数。

$U_a^2$ 是多体问题与一组孤立二体问题的区别所在。把太阳、地球、月球分别记作 $S,E,M$，则 $m_EU_E^2$ 含有 $2G^2m_E m_S m_M/(r_{ES}r_{EM})$。对月球位置作变分时不能把其中属于太阳的势丢掉。即使精确知道每一对天体单独存在时的进动，也不能将三条这样求出的轨道彼此相加得到历表。在数值积分中，可由 $\bm p_a=\p L/\p\bm v_a$ 及 $\dot{\bm p}_a=\p L/\p\bm x_a$ 先写成隐式方程，再在所有已乘 $1/c^2$ 的加速度上作 Newton 阶代换。若混入部分高阶加速度，便无意中换了截断约定。时间、空间平移及空间转动分别给出总能量、总动量及 $\sum_a\bm x_a\times\bm p_a$ 的守恒，可用作独立的积分精度检查。

\subsection{半守恒理论的多体作用量}

前面的 \eqref{eq:exp-nbody-lag} 只留 $\gamma,\beta$。若有能量与总动量守恒，却容许实验依赖 $\bm w$ 或外部物质方向，还能否只从轨道作用量出发？忽略体积矩、自能差和自旋，取 $G=c=1$，答案是肯定的。令 $L_{\gamma\beta}$ 为 \eqref{eq:exp-nbody-lag} 在几何制中的形式。除一个全导数外，可取 $L=L_{\gamma\beta}+\delta L_{\rm F}+\delta L_\xi$，其中
\begin{align*}
 \delta L_{\rm F}={}&-\frac14\sum_a\sum_{b\ne a}\frac{m_am_b}{r_{ab}}
 \left[(\alpha_1-\alpha_2)(\bm w+\bm v_a)\cdot(\bm w+\bm v_b)\right.\\
 &\hspace{28mm}\left.+\alpha_2[(\bm w+\bm v_a)\cdot\bm n_{ab}]
                            [(\bm w+\bm v_b)\cdot\bm n_{ab}]\right],\\
 \delta L_\xi={}&-\frac{\xi}{2}\sum_a\sum_{b\ne a}
 \sum_{d\ne a,b}\frac{m_am_bm_d}{r_{ab}^2}\,
       \bm n_{ab}\cdot\left(\frac{\bm x_b-\bm x_d}{r_{ad}}
                       -\frac{\bm x_a-\bm x_d}{r_{bd}}\right).
\end{align*}
上述两个 $\delta L$ 均为 1PN 修正。第一式实际由一对天体即可取样。第二式显式需要一个第三源；对一颗遥远天体，要在三重和的所有排列相加后再取大距离极限，否则会漏掉与方向有关的外势项。这正是比较孤立双星与置于银河场中的双星时要分开的两件事。若考察外界源对球体主动、被动质量的改正，仍要用前面的有限体积分，不应将同一自能贡献再加到这个点粒子作用量上。

以 $\hat{\bm p}_a=\bm p_a/m_a$ 记单位质量的正则动量。写 $L=L_0+L_1$，对每颗星有 $\bm p_a=m_a\bm v_a+\p L_1/\p\bm v_a$。将 $\bm v_a=\hat{\bm p}_a-\p_{\bm v_a}L_1/m_a+\cdots$ 代入 $\sum_a\bm p_a\cdot\bm v_a-L$，由 $L_0$ 产生的速度修正同第一项中的修正正好抵消。因此 $H_1=-L_1(\bm x,\hat{\bm p})$，而不是把 $L_1$ 里的动能项逐个换一个预想的符号。举例说，在 $\bm w=0$ 的系中，两体速度相关的 Hamilton 量可按无序对写为
\begin{align*}
 H_{1,{\rm pair}}=-\sum_{a<b}\frac{m_am_b}{2r_{ab}}\big[
 &(2\gamma+1)(\hat p_a^2+\hat p_b^2)
 -(4\gamma+3+\alpha_1-\alpha_2)\hat{\bm p}_a\cdot\hat{\bm p}_b\\
 &-(1+\alpha_2)(\hat{\bm p}_a\cdot\bm n_{ab})
                             (\hat{\bm p}_b\cdot\bm n_{ab})\big].
\end{align*}
再加 $-\sum_a m_a\hat p_a^4/8$、$+(2\beta-1)\sum_a m_aU_a^2/2$ 和 $-\delta L_\xi$ 就得到本规范中的全部一阶保守修正。不同正则坐标可给不同的系数表，计算同一轨道周期或时延必须同时换用相应的坐标到读数的关系。恢复 $c$ 后，$H_1$ 和两个 $\delta L$ 均另带 $c^{-2}$，源质量和速度则仍用普通单位并在势中恢复 $G$。

\subsection{从坐标加速度到轨道要素}

天文计时给出的是一串观测历元。实际拟合中常以在某一历元与真实轨道位置、速度相同的 Kepler 椭圆作起点，使其六个要素缓慢变化。这称为\textbf{密切轨道}。轨道受扰后，其瞬时要素和用整个时间段平均得到的要素可以不同。令非 Newton 加速度在径向、顺行、轨道法向上的分量为 $A_R,A_T,A_N$。由 Newton 能量 $E_{\rm N}=v^2/2-Gm/r$ 及角动量 $\bm h=\bm r\times\bm v$，立刻有 $\dot E_{\rm N}=\dot r A_R+(h/r)A_T$、$\dot h=rA_T$。对 $E_{\rm N}=-Gm/(2a)$、$h^2=Gma(1-e^2)$ 微分，得
\begin{align*}
 \dot a&=\frac{2a^2}{h}\left[e\sin f\, A_R+\frac{p}{r}A_T\right],\\
 \dot e&=\frac1h\left[p\sin f\,A_R+
              \big((p+r)\cos f+re\big)A_T\right].
\end{align*}
例如轨道切向若有微弱阻力，$\dot E_{\rm N}<0$，累积下来半长轴变小。反之，一个时间反演对称的保守摄动可使 $a$ 在一周之内改变，却在平均后只有角要素出现长期项。是否要把某一项当作可测的长期漂移，应在对观测历元作平均之后判断。

令 $i$ 为倾角、$\Omega_{\rm node}$ 为升交点经度、$\omega$ 为近心点角，$u=\omega+f$。由 $\dot{\bm h}=\bm r\times\bm A$，法向力只改变角动量的方向，其结果是 $\dot i=r\cos u\,A_N/h$、$\dot\Omega_{\rm node}=r\sin u\,A_N/(h\sin i)$。在转动的轨道面内读近心点角，还须从面内近心点转速减去 $\cos i\,\dot\Omega_{\rm node}$。故高度倾斜的卫星既能用轨道面检验拖曳，也可用长期进动约束中心体的多极矩。若 $e=0$ 或 $i=0$，公式中 $1/e$、$1/\sin i$ 的奇性只说明近心点或交点的坐标没有定义。改为拟合离心率矢量和轨道法向量的两个独立分量，轨道本身的响应依然正则。

\subsection{流体方程与有限天体}\label{sec:pn-finite}

源若是理想流体，沿流线的导数记为 $D_t=\p_t+v^j\p_j$。沿用 \hyperlink{subsection.1.2.3}{\S\,1.2.3} 的状态方程及绝热关系 $D_t\Pi=(p/\rho^2)D_t\rho$，配上 $\rho_*$ 的连续性方程，尚需求 1PN 动量方程。

为写出动量方程，设 $u^\mu=\d x^\mu/\d\tau$、$u^t=\d t/\d\tau$，则 $u^0=cu^t$。引入无量纲的比焓 $h=1+H/c^2$，其中 $H=\Pi+p/\rho$。以混合指标写 $\nabla_\mu T^\mu{}_i=0$，可把联络项化成 $\tfrac12T^{\alpha\beta}\partial_i g_{\alpha\beta}$。利用重子数守恒及 $D_th=D_tp/(\rho c^2)$，即得
\[
  D_t(hu_i)=\frac{h}{2u^t}u^\alpha u^\beta\partial_i g_{\alpha\beta}
        -\frac{1}{\rho u^t}\partial_i p .
\]
左端是单位静质量的协变动量随流而变。右端的度规导数只对场点求导，不对质元的速度求导。这一点颇似以坐标时为参数的质点 Euler--Lagrange 方程，只是左端还乘以焓，右端还有压强梯度。

为便于代入 \eqref{eq:exp-ppn-all}，取 $B_i=c^3g_{0i}$ 的最低阶，即
$B_i=-[2(1+\gamma)+\alpha_1/2]V_i-\tfrac12(1+\alpha_2-\zeta_1+2\xi)\partial_i\dot X+\Phi_{{\rm F}i}$。以 $K=2\psi-2\beta U^2+\Phi_{\rm F}$ 记 $g_{00}$ 的四阶系数。由 $u^t=1+(v^2/2+U)/c^2+O(c^{-4})$ 和 $\rho_*=\sqrt{-g}\rho u^t$ 展开上式，有
\begin{align}
 P_i:=hu_i={}&v_i+\frac{1}{c^2}
  \left[\left(H+\tfrac12v^2+(2\gamma+1)U\right)v_i+B_i\right]+O(c^{-4}),\nonumber\\
 D_tP_i={}&\partial_i U-\frac{1}{\rho_*}\partial_i p\nonumber\\
 &+\frac{1}{c^2}\left\{
   [H+(\gamma+\tfrac12)v^2+U]\partial_iU
       +\tfrac12\partial_iK+v^j\partial_iB_j
       -(3\gamma-1)\frac{U}{\rho_*}\partial_i p\right\}
       +O(c^{-4}).\label{eq:exp-fluid-p}
\end{align}
这里 $P_i$ 是流体的动量变量，与用 $P[f]$ 表示的势积分无关。所有势均依照同一时刻的源量计算。这是一个完整的 1PN Euler 方程：先从第一式将 $P_i$ 与 $v_i$ 互换，再以第二式推进即可。若欲显式求坐标加速度，只须在 $D_t(P_i-v_i)$ 内代入 Newton 阶方程。举例说，度规矢量势在加速度中的贡献为 $[-\partial_t B_i+v^j(\partial_i B_j-\partial_jB_i)]/c^2$。其中 $B_i$ 含有的纯梯度会从反对称项中消去，却仍会在时间导数中留下 $\partial_i\ddot X$。若只知道球体外的静态度规，就会漏掉来自源的运动和局域系速度的这些项。令 $p,\Pi$ 同时消失，则方程退化为这里所用坐标下的类时测地线方程。

平直背景中的惯性系数应还原为 \hyperlink{subsection.1.2.3}{\S\,1.2.3} 的焓密度。

再回到有外场的坐标系，由 $\sqrt{-g}=1+(3\gamma-1)U/c^2+O(c^{-4})$ 得 $\rho_*/\rho=1+(3\gamma U+v^2/2)/c^2+O(c^{-4})$。其中 $3\gamma U$ 来自三维固有体积，$v^2/2$ 则是流体运动造成的改正。当物理体积固定而坐标体积随所取空间坐标改变时，同一团流体可以有不同的坐标密度，质量并未改变。

取边界随星体物质运动的区域。记其重子质量为 $M_{*A}=\int_A\rho_*\,\d^3x$，相应质心为 $\bm z_{*A}=M_{*A}^{-1}\int_A\rho_*\bm x\,\d^3x$。连续性给出一般恒等式 $\mathrm d(\int_A\rho_*f\,\d^3x)/\mathrm dt=\int_A\rho_*D_tf\,\d^3x$，故 $M_{*A}$ 守恒，$\ddot{\bm z}_{*A}$ 是质元加速度的静质量平均。将 Newton 阶 Euler 方程积分，压强项成为表面力，若星体表面 $p=0$ 则消失。内部引力项是对两质元的二重积分，交换两点时 $\bm x-\bm x'$ 变号，力的总和为零。这两项虽都能强烈影响星体的形状，却不产生 Newton 阶的质心合力。

外场则在 $\bm z_{*A}$ 附近展开。令 $y^i=x^i-z_{*A}^i$，取外部真空区，$\nabla^2 U_{\rm ext}=0$。由于 $\int_A\rho_*y^i\,\d^3x=0$，对准这一中心没有质量偶极项，故
\[
 \ddot z_{*A}^i=\partial_iU_{\rm ext}
   +\frac{Q_A^{jk}}{2M_{*A}}\partial_i\partial_j\partial_k U_{\rm ext}
   +\cdots,\qquad
 Q_A^{jk}=\int_A\rho_*\left(y^jy^k-\tfrac13\delta^{jk}y^2\right)\d^3x .
\]
右端在质心取值。四极项相对单极项通常为 $(R_A/r)^2$ 量级，$R_A$ 为天体半径，$r$ 为外源距离。若外源本身亦有质量四极矩，须把它的势计入 $U_{\rm ext}$。相对点质量轨道的偏离，在完全满足普适自由落体的理论中也可以由这些形状项产生。

到 1PN 阶，应问所取的中心究竟是哪种质量的中心。置 $\bar{\bm v}=\bm v-\dot{\bm z}_{*A}$，只对星内物质积分求得自势 $\bar U$，并记 $\Omega_A=-\tfrac12\int_A\rho_*\bar U\,\d^3x<0$。以前式的守恒静质量加上质心系下动能、内能与引力势能，便得到此阶惯性质量及一种相应的质心定义：
\[
 M_A=\int_A\sigma_A\,\d^3x,\qquad
 z_A^i=M_A^{-1}\int_A\sigma_A x^i\,\d^3x,\qquad
 \sigma_A=\rho_*\left[1+\frac{\bar v^2/2+\Pi-\bar U/2}{c^2}\right].
\]
将整体公转速度算进 $\bar v$ 会把轨道能误作一颗星的静能。势能项的 $1/2$ 则源于不可重复计数同一对质元。略去外界作潮汐功、物质损失和辐射时，$M_A$ 在此阶守恒，而其一阶矩与 $\rho_*$ 的一阶矩仍可在 1PN 阶略有差别。

对体积积分整理内部项时，不能因外部看起来是一颗圆星，就先把内部压强和运动设为零。由 Newton 阶内部 Euler 方程乘 $y^i$ 再积分，取迹有 $\tfrac12\mathrm d^2 I/\mathrm dt^2=2T_{\rm int}+\Omega_A+3\int_Ap\,\d^3x$，其中 $I=\int_A\rho_*y^2\,\d^3x$，$T_{\rm int}$ 是质心系中的动能。这里忽略外部潮汐和表面应力。平衡星的左端为零，对周期变化星也可在许多内部周期上取平均。在一般 PPN 场方程中，哪种内能出现在远场 $1/r$ 项（主动质量），哪种内能乘在外部均匀场的梯度前（被动质量），与哪种内能出现在 $M_A$ 中，是三个可分开计算的问题。上述维里关系可消去从同一内部平衡重复出现的项。广义相对论中，球对称、弱自引力且无自旋的星在所取精度下只余同一个 $M_A$。一般情形，被动质量对引力结合能的异常会在加速度中留下与 $\Omega_A/(M_Ac^2)$ 成正比的项。它与上面的几何四极项有一个可辨之处：即使外场在整颗星上均匀，前者仍不会消失。

为区分三种质量，可以设想分别称量同一颗星。先由它在已知外力下的加速度定出惯性质量 $M_A$，再放开它，借其它已标定的小试验体所测的外场强度求其被动质量 $M_A^{\rm P}$，最后在远大于星半径处由绕星的小探针轨道测主动质量 $M_A^{\rm act}$。这些是方程中不同角色的定义，天文观测往往只能确定它们的组合。例如在只保留两体单极引力的模型中，$A$ 到 $B$ 的力和 $B$ 到 $A$ 的力的大小分别正比于 $M_A^{\rm act}M_B^{\rm P}$ 与 $M_B^{\rm act}M_A^{\rm P}$。若主动与被动质量之比随物质种类而变，单把两物体的机械动量相加便可能不守恒。还须检查理论中场携带的动量，才可由此判断一个孤立系统是否真会产生质心自加速。

\subsection{内部能量的张量与三种质量}

对球形行星只给一个标量质量就能解释许多资料。不过要导出这个简化，不宜先拿球对称去除可能有方向的内部量。以下用几何制，将 $A$ 内部的质元质量记作 $\d m=\rho_*\,\d[3]{x}$，内部位置 $\bm y=\bm x-\bm z_A$、速度 $\bm u=\bm v-\bm v_A$。对同一星内两点以 $\bm s=\bm x-\bm x'$、源点速度 $\bm u'$ 表示。定义
\begin{align*}
 I^{ij}&=\int_A y^iy^j\,\d m,\qquad
 S^{ij}=\int_A(y^i u^j-y^j u^i)\,\d m,\qquad
 T^{ij}=\tfrac12\int_A u^i u^j\,\d m,\\
 \Omega^{ij}&=-\tfrac12\iint_A \frac{s^is^j}{s^3}\,\d m\,\d m',\qquad
 P_A=\int_Ap\,\d[3]{x},\qquad E_A=\int_A\Pi\,\d m .
\end{align*}
$\Omega=\Omega^{ii}<0$ 是 Newton 结合势能，$T=T^{ii}$ 是内部动能。$I^{ij}$ 的无迹部分给质心二阶质量矩，$S^{ij}$ 的空间对偶给自旋。还有两种不同的速度--结合能相关量
\[
 H^{ij}=\iint_A \frac{u'^i s^j}{s^3}\,\d m\,\d m',\qquad
 K^{ij}=\iint_A\frac{(\bm u'\cdot\bm s)s^is^j}{s^5}\,\d m\,\d m' .
\]
它们来自 $\dot V_i$ 及 $\dot X_{,ij}$ 的内部部分。在一颗静止但有随机运动的星中，$\bm v_A=0$ 并不表明 $T^{ij}=0$；若星还在自转，$H^{ij}$ 的反对称部分更不能简单扔掉。

对 $I^{ij}$ 的两个时间导数用 Newton 阶连续性和 Euler 方程，可得一组减少独立积分数的等式：
\begin{align*}
 \tfrac12\dot I^{ij}&=\int_A u^iy^j\,\d m+\tfrac12 S^{ij},\\
 \tfrac12\ddot I^{ij}&=2T^{ij}+\Omega^{ij}+\delta^{ij}P_A,\\
 \tfrac12\dddot I^{ij}&=4H^{(ij)}-3K^{ij}+\delta^{ij}\dot P_A-2L^{(ij)},\qquad
 L^{ij}=\int_A u^i\p_jp\,\d[3]{x}.
\end{align*}
第二行的引力项可由将 $\int y^{(i}\p_{j)}\bar U\,\d m$ 交换两点后取半和直接得到，压强项则由表面 $p=0$ 给出。第三行是对同一式再求导，没有引入比 1PN 更精细的内部场方程。若内部以短于轨道周期的时标达到平衡，诸式的左边为零或可作内部周期平均。这允许只保留低阶结合能张量和自旋。若有激发的长寿命模或强烈潮汐破坏，这些左边本身也是可观测的，便不能用定常维里关系代换。

用这些恒等式消去内部动能等重复项，一种方便的组织方式是在惯性质量、主动引力质量、被动引力质量三者上各放一个张量。在 $M$ 中先收入各向同性的内部动能、热能及 $\Omega$；以下右边各项均用同一颗星的内部变量。以 $\delta$ 记三维单位阵，$\Omega^{ij}$ 保持张量意义，则
\begin{align*}
 (\tilde m_I)^{ij}={}&M\delta^{ij}
 +(\alpha_1-\alpha_2+\zeta_1)\Omega\delta^{ij}
        +(\alpha_2-\zeta_1+\zeta_2)\Omega^{ij},\\
 (\tilde m_P)^{ij}={}&M\delta^{ij}
 +(4\beta-\gamma-3-3\xi)\Omega\delta^{ij}-\xi\Omega^{ij},\\
 (\tilde m_A)^{ij}={}&M\delta^{ij}
 +\left[(4\beta-\gamma-3-3\xi-\alpha_3/2-\zeta_1/2-2\zeta_2)\Omega
                     +\zeta_3E_A\right]\delta^{ij}\\
 &+(3\zeta_4-\zeta_1-3\alpha_3/2)P_A\delta^{ij}
                       +(\zeta_1/2-\xi)\Omega^{ij}.
\end{align*}
回到普通单位时，以上所有内部能量项均除以 $c^2$。对任意外源，先用其主动质量张量在给定方向产生外部场梯度，再用受试物体的被动质量求合力，最后以惯性质量张量解加速度。把三者都定义为与方向无关的张量，代价是源的势要带指标。取 $r_B=|\bm x-\bm z_B|$、$\bm n_B=(\bm x-\bm z_B)/r_B$，例如可写 $\mathcal U^{ij}=\sum_{B\ne A}(\tilde m_A)_B^{ik} n_B^k n_B^j/r_B$，此处下标 $A$ 在 $\tilde m_A$ 中表示主动质量。准 Newton 部分用 $(\tilde m_I)_A^{ij}a_A^j=(\tilde m_P)_A^{kl}\p_i\mathcal U^{kl}$ 表出，另行加入自加速和通常的 1PN 轨道项。这样的写法使交换两物体时各自扮演的角色显而易见。

近球形时 $\Omega^{ij}=\delta^{ij}\Omega/3$，上式大大简化。尤其在半守恒理论中，$(\tilde m_P)^{ij}=(\tilde m_A)^{ij}$，两体相互作用仍可成对抵消；但这不要求它们等于 $(\tilde m_I)^{ij}$。球体的相对被动响应为 $1+\eta_{\rm N}\Omega/M$，其中
$\eta_{\rm N}=4\beta-\gamma-3-10\xi/3-\alpha_1+2\alpha_2/3-2\zeta_1/3-\zeta_2/3$。
对 Brans--Dicke 理论，$\beta=1$ 且只有 $\gamma$ 偏离一，便有 $\eta_{\rm N}=1/(2+\omega_0)$。自能为负，所以在通常的正 $\omega_0$ 分支上，引力结合得更紧的物体相对另一物体在外场中落得略慢。若所有实验物体都是同一尺度、相同密度的材料块，则其 $\Omega/M$ 都极小；即使轨道方程有这一项，普通材料扭秤也难以测定它。

还有一种效应和上面的外场响应不同，它甚至不需要附近有另一颗星。非守恒部分含有诸如 $\iint u'^2s^j/s^3\,\d m\,\d m'$、$\iint \Pi' s^j/s^3\,\d m\,\d m'$ 和把 $\d m'$ 换作 $p'\d[3]{x'}$ 为压力源的矢量。把物体及其内部流反射对称地成对积分，这些极矢量可为零。优选速度 $\bm w+\bm v_A$ 却可同内部自转构成一个即使球星也不消失的矢量，其加速度为 $-\alpha_3(w+v_A)^kH^{kj}/M$。设星近于刚性转动，$\bm u=\bm\omega_A\times\bm y$，交换双积分中的质元后，有 $H^{kj}=\epsilon^{k\ell m}\omega_A^\ell\Omega^{jm}$。取球对称的迹，遂得
\[
 \bm a_{\rm self}=-\frac{\alpha_3}{3}\,
    \frac{\Omega_A}{M_Ac^2}\,
             (\bm w+\bm v_A)\times\bm\omega_A .
\]
在最后一式中恢复普通单位。这个异常既需引力结合能，又需快自转，方向还得投影到观测视线上。中子星同时放大前两项，但将弱场推导用于中子星时要用所检验理论的强场系数，不能在未作这一步时把上限解释为对任意理论的 $\alpha_3$ 上限。

\subsection{内部矢量与一般多体的轨道方程}

\subsubsection{反射不对称的物体}

上一节的张量维里关系是真正求内部积分的捷径。若天体没有关于质心的反射对称性，还有几项不能由动能、势能与压强的标量积分描述。仍用 $G=c=1$ 和 $\d m=\rho_*\d[3]{x}$，在同一物体 $A$ 内取两点 $\bm x,\bm x'$，置 $\bm s=\bm x-\bm x'$、$s=|\bm s|$、$\bm q=\bm s/s^3$。两点的内部速度分别为 $\bm u,\bm u'$。为简写，以下用尖括号表示\emph{不作质量归一化}的双重积分，$\langle F\rangle_A=\int_A\int_A F\,\d m\,\d m'$。速度二次项有四种不同的缩并，
\begin{align*}
 \bm t_A&=\langle u'^2\bm q\rangle_A,&
 \bm{\mathcal T}_A&=\langle\bm u'(\bm u'\cdot\bm q)\rangle_A,\\
 \bm{\mathcal T}_A^*&=\langle\bm u(\bm u'\cdot\bm q)\rangle_A,&
 \bm{\mathcal T}_A^{**}&=
       \left\langle\frac{(\bm u'\cdot\bm s)^2}{s^2}\bm q\right\rangle_A .
\end{align*}
星号只用于区别缩并的位置。若把第一式中的 $u'^2$ 换为 $u^2$，交换两点后 $\bm q$ 会变号，不能省掉这个撇。再记 $\bar U(\bm x')=\int_A\d m''/|\bm x'-\bm x''|$、$\bm g_{\rm int}(\bm x')=\bm\nabla'\bar U(\bm x')$，则结合能、内能与压强还给出
\begin{align*}
 \bm{\mathcal O}_A&=\langle\bar U(\bm x')\bm q\rangle_A,&
 \bm{\mathcal O}_A^*&=-\langle[\bm g_{\rm int}(\bm x')\cdot\bm s]\bm q\rangle_A,\\
 \bm{\mathcal E}_A&=\langle\Pi'\bm q\rangle_A,&
 \bm{\mathcal P}_A&=\langle(p'/\rho_*')\bm q\rangle_A .
\end{align*}
这里特用 $\bm{\mathcal O}_A$ 表示一个\emph{矢量}，以免同自引力能量张量 $\Omega_A^{ij}$ 混淆。在第二式中代入 $\bm g_{\rm int}(\bm x')=-\int_A (\bm x'-\bm x'')\d m''/|\bm x'-\bm x''|^3$，便看出核心是三点间的两个方向之内积。末式的 $\d m'$ 与分母的密度相消，实际压力源是 $p'\d[3]{x'}$；在真空边界处不必真的除以一个为零的密度。这些量都只需星体内的流量、压强、热能和质量分布，第三个积分点也不可跑到一个伴星内部去。

可用两条恒等式核对上述定义，亦可看出为何在最终加速度中不需要八个彼此独立的向量。以内部运动为源，定义 $\bar{\bm V}(\bm x)=\int_A \bm u'\d m'/s$、$\bar{\bm W}(\bm x)=\int_A (\bm u'\cdot\bm s)\bm q\,\d m'$。对同一批质元取时间导数，利用连续性方程及最低阶 Euler 方程，有
\begin{align*}
 \frac{\d}{\d t}\int_A\bar{\bm V}\,\d m
      &=\bm{\mathcal T}_A+\bm{\mathcal T}_A^*+
                         \bm{\mathcal O}_A+\bm{\mathcal P}_A,\\
 \frac{\d}{\d t}\int_A\bar{\bm W}\,\d m
      &=-\bm t_A-\bm{\mathcal T}_A+\bm{\mathcal T}_A^*
                 +3\bm{\mathcal T}_A^{**}
                 -\bm{\mathcal O}_A^*-\bm{\mathcal P}_A .
\end{align*}
譬如第一行，对 $1/s$ 求导得到 $-(\bm u-\bm u')\cdot\bm q$；把其中含 $\bm u' (\bm u\cdot\bm q)$ 的项交换两个积分点即出现 $\bm{\mathcal T}_A^*$。对 $\bm u'$ 求导所生压力梯度经分部积分，把导数移到 $1/s$ 上，即得 $+\bm{\mathcal P}_A$。这里已用自由表面上 $p=0$，并在质心自由落体系中舍弃外场的潮汐体积矩。若流体被器壁约束，则须将表面应力与器壁一同计入。第二行的 3 来自对 $s^{-3}$ 求导。若内部平稳，或可在长期结构几乎不变时平均许多个内部周期，两条等式的左边可取零；不能在任意一张流体构型的瞬时图像上都如此代换。

略去自旋对外轨道的偶极力和高阶体积矩，此时有一组完全来自内部的 1PN 自加速。以 $M_A$ 为上文包括内部能量的标量质量，结果可写成
\begin{align*}
 M_A(a_A^j)_{\rm self}={}&
    -\tfrac12(\alpha_3+\zeta_1)t_A^j
    -\tfrac12\zeta_1(2\mathcal T_A^j-3\mathcal T_A^{**j})
    -\zeta_2\mathcal O_A^j-\zeta_3\mathcal E_A^j-3\zeta_4\mathcal P_A^j\\
    &-\alpha_3(w+v_A)^k H_A^{kj}.
\end{align*}
$w+v_A$ 一项表示 $\bm w+\bm v_A$ 的分量。前一行的各项在所有内部空间矢量同时反射时都变号；若密度、内能、压强在一对镜像位置相等，且 $\bm u(-\bm y)=-\bm u(\bm y)$，它们逐对积分为零。后一行却是外给的速度乘一个偶宇称张量，故可在旋转的近球形天体中留下上文给出的 $\alpha_3$ 自力。一颗不对称物体若也出现这些极矢量，只有 $\alpha_3,\zeta_1,\ldots,\zeta_4$ 中至少一个非零才有此处的自加速。由压力、内能或自引力的一项积分不为零，不能分别断言其质心自行加快。

\subsubsection{把一般轨道项按两个或三个体标计算}

从一颗星内各场点积分完毕，源点若在另一颗星中，应换用外部的轨道变量。设 $\bm r_{AB}=\bm z_A-\bm z_B$，$r_{AB}=|\bm r_{AB}|$，$\bm n_{AB}=\bm r_{AB}/r_{AB}$。现将 $A$ 的加速度拆成刚刚求得的自加速、上文由三种质量张量求得的准 Newton 加速度，以及以下纯轨道的 1PN 项。仍用几何制。写 $\bm k_A=\bm w+\bm v_A$、$\bm v_{AB}=\bm v_A-\bm v_B$，在每一对 $(A,B)$ 上置 $r=r_{AB}$、$\bm n=\bm n_{AB}$。定义一个速度平方量纲的系数
\begin{align*}
 F_{AB}={}&\gamma v_A^2-2(\gamma+1)\bm v_A\cdot\bm v_B
          +(\gamma+1)v_B^2-\tfrac32(\bm n\cdot\bm v_B)^2\\
 &+\tfrac12(\alpha_2+\alpha_3) k_B^2
          -\tfrac12\alpha_1\bm k_A\cdot\bm k_B
          -\tfrac32\alpha_2(\bm n\cdot\bm k_B)^2\\
 &-(2\gamma+2\beta+1+\tfrac12\alpha_1-\zeta_2)\frac{M_A}{r}
          -(2\gamma+2\beta)\frac{M_B}{r}.
\end{align*}
至此可以先把作用在 $A$ 上、只用 $A,B$ 两个标签的轨道修正相加：
\begin{align*}
 \bm a_A^{(2)}=\sum_{B\ne A}\frac{M_B}{r^2}\bigg\{
   &-F_{AB}\bm n+
      \big[\bm n\cdot((2\gamma+2)\bm v_A-(2\gamma+1)\bm v_B)\big]\bm v_{AB}\\
   &-\big[\bm n\cdot(\alpha_2\bm k_B+\tfrac12\alpha_1\bm v_{AB})\big]\bm k_B\bigg\}.
\end{align*}
$(2)$ 在此表示二体\emph{标签}，并非 2PN；这不是一个可以由牛顿轨道拟合来代换的孤立双星总质量。特别是 $F_{AB}$ 最后一行有两个不同的质量系数。若所有异常参数为零，先取 $\bm v_A=\bm v_B=0$ 的瞬间，则 $A$ 每单位 Newton 场强所得的径向修正为 $(5M_A+4M_B)/r$，$B$ 所得的则须交换质量。作相对加速度后质量分子为 $4(M_A+M_B)^2+2M_AM_B$，恢复了广义相对论二体式中的 $4+2\eta$。若一开始就令两者的中心质量相等，会丢掉这个检查。

再取三个互异的标签 $(A,B,C)$，以 $\bm d=\bm n_{AC}-\bm n_{BC}$，定义
\begin{align*}
 Q_{ABC}={}&\frac{2\gamma+2\beta-2\xi}{r_{AC}}
      +\frac{2\beta-1-2\xi-\zeta_2}{r_{BC}}\\
 &-\tfrac12(1+2\xi+\alpha_2-\zeta_1)
     \frac{r_{AB}}{r_{BC}^{2}}(\bm n_{AB}\cdot\bm n_{BC})
      -\xi\frac{r_{BC}}{r_{AC}^{2}}(\bm n_{AC}\cdot\bm n_{BC}),\\
 C_0={}&\tfrac12(4\gamma+3-2\xi+\alpha_1-\alpha_2+\zeta_1).
\end{align*}
$C_0$ 与三角形无关，$Q_{ABC}$ 有长度倒数的量纲。三标签贡献可压成
\begin{align*}
 \bm a_A^{(3)}
   =\sum_{B\ne A}\sum_{C\ne A,B}\frac{M_BM_C}{r_{AB}^{2}}
      \bigg\{&Q_{ABC}\bm n_{AB}
           -C_0\frac{r_{AB}}{r_{BC}^{2}}\bm n_{BC}\\
       &-\frac{\xi}{r_{AB}}
         [\bm d-3\bm n_{AB}(\bm n_{AB}\cdot\bm d)]\bigg\}.
\end{align*}
固定 $A$ 后，$(B,C)$ 是有序对，$B$ 与 $C$ 交换后的三角形必须再算一次。最后一行是 $\delta^{jk}-3n_{AB}^jn_{AB}^k$ 作用在 $d^k$ 上，既有沿 $AB$ 的分量，也有垂直于 $AB$ 的分量；只计引力势的大小，不足以重建由 $\xi$ 引起的方向效应。$\zeta_2$ 不仅在 $A,B$ 两标签的 $M_A/r_{AB}$ 项中出现，也改变邻星作用于 $B$ 处的势。$\zeta_1$ 既出现在与 $\bm n_{AB}$ 平行的角投影，又出现在沿 $\bm n_{BC}$ 的交叉项。反之，$\zeta_3,\zeta_4$ 的源是内部能量和压强，在已经把天体分成自加速、准 Newton 质量和纯轨道三部分后，不再在这两个纯轨道和式中额外出现。

实际运算可对所有互异的点先算一次 $r_{AB},\bm n_{AB}$ 及 $\bm v_A\cdot\bm v_B$，再分别累加 $\bm a_A^{(2)},\bm a_A^{(3)}$。若原问题只有两星，三标签的和为空，却须保留 $F_{AB}$ 中 $M_A/r,M_B/r$ 两项：它们正包含解场过程中反复用到同一源的效应。若有第三颗遥远的星 $C$，$\bm n_{AC}$ 与 $\bm n_{BC}$ 会接近，不能在所有项相加以前就把两个方向设为相同；例如 $\bm d/r_{AB}$ 可有同 $1/r_{AC}$ 一阶的外场贡献。应先取两内星的加速度差，再按 $r_{AB}/r_{AC}$ 展开。由共同外势产生的各向同性常数改正可归入局域实测质量标度，依外源方向而变的投影却可能留下。后面对优选位置的检验，所比较的正是这样的项。

在这两个和式中任一个质量或速度已经乘了 1PN 因子，只需用其 Newton 值；用已含内能修正的 $M_A$ 也只差下一阶。若恢复物理单位，则对 $F_{AB}$ 中的每一个 $M/r$ 换作 $GM/r$，以 $GM_B/r^2$ 作双标签项的外因子，三标签项的 $M_BM_C$ 换作 $G^2M_BM_C$，最后两式都乘 $c^{-2}$。准 Newton 项及其质量修正须按各自原先的单位另算，不能把整个加速度再一律除以 $c^2$。总方程可用于动力学平衡、尺度远小于天体间距的弱自引力星体；自旋力和质量四极力若可分辨，应按各自体积矩补入。内部积分含的是连续密度生成的自场，外部单极和则明确排除了 $B=A$。既不让外部点源在 $r_{AA}=0$ 自作用，也不应借这种排除删掉原本由有限星体内积分得到的结合能。

\subsection{一般双星的相对运动}

有了三种质量，还须问观测中所称的“双星总质量”是哪一种。为看清此点，取两颗近于球形、各自关于质心反射对称的星，让其质心系相对于所选静系暂为静止。记 $m_1,m_2$ 为各自惯性质量，$(m_P)_a,(m_A)_a$ 分别为球形平均后的被动和主动质量。对第一颗星，Newton 型加速度大小为 $G[(m_P)_1/m_1](m_A)_2/r^2$，第二颗星同理。故相对运动的 Kepler 项实际含有 $m_{\rm K}=\frac{(m_P)_1}{m_1}(m_A)_2+ \frac{(m_P)_2}{m_2}(m_A)_1$。若两星自能在观测中没有变化，仅量相对位置和周期就只能测 $Gm_{\rm K}$。把这个常数同 $G(m_1+m_2)$ 的某个很小差值宣称为由同一条轨道测出的异常，是把一项资料用了两次。要分离自能响应，可引入一个外部场，或比较另一项依赖质量的观测。在已乘 $1/c^2$ 的速度项中仍可用 $m=m_1+m_2$，差别在此已是下一阶。

令 $\eta=m_1m_2/m^2$、$\bm x=\bm x_1-\bm x_2$、$\bm n=\bm x/r$、$\bm v=\dot{\bm x}$、$\dot r=\bm n\cdot\bm v$。为写得紧凑，定义四个由标准 PPN 参数确定的系数
\begin{align*}
 A&=2\gamma+2\beta+\eta(2+\alpha_1-2\zeta_2),&
 B&=\gamma+\tfrac12\eta(6+\alpha_1+\alpha_2+\alpha_3),\\
 C&=\tfrac32\eta(1+\alpha_2),&
 D&=2+2\gamma-\eta(2-\alpha_1+\alpha_2).
\end{align*}
忽略随外部环境定向的项和内部自加速，由两星方程相减可得
\begin{align}
 \ddot{\bm x}={}&-\frac{Gm_{\rm K}}{r^2}\bm n
   +\frac{Gm}{r^2c^2}\left[
       \left(A\frac{Gm}{r}-Bv^2+C\dot r^2\right)\bm n+D\dot r\,\bm v\right]
   +\bm a_Q .\label{eq:g3-relative-ppn}
\end{align}
这里 $\bm a_Q$ 允许星 2 有轴对称的质量四极矩，星 1 的四极矩暂略。若 $\bm e_s$ 为星 2 的对称轴、$s=\bm e_s\cdot\bm n$、$R_2$ 为其参考半径，则 $\bm a_Q=3GmJ_{2}(R_2)R_2^2[(5s^2-1)\bm n-2s\bm e_s]/(2r^4)$，$J_2(R_2)$ 是以该半径归一的无量纲矩。其来源是 Newton 势的二阶质量矩，故未显式带 $c^{-2}$，即使比 1PN 项小也不应当一概按 PN 阶数删去。对称天体的极轴与赤道转动惯量分别为 $I_\parallel,I_\perp$ 时，$J_2=(I_\parallel-I_\perp)/(m_2R_2^2)$。令所有异常参数消失、再置 $J_2=0$，便回到先前保留的广义相对论二体方程。

可以把这四个系数同已求过的轨道平均直接对应。令 $f$ 为真近点角、$p$ 为半通径，扰动力分解成径向与顺行两项时，$D\dot r\bm v$ 的径向投影是 $D\dot r^2$。用 $Gm/r=(Gm/p)(1+e\cos f)$、$v^2=(Gm/p)(1+2e\cos f+e^2)$，径向项对进动的贡献为 $\pi(-A+2B)Gm/(pc^2)$，切向给 $2\pi D Gm/(pc^2)$。相加后，球对称 PN 部分的每转进动为
\[
 \Delta\varpi_{\rm PN}=
 \frac{6\pi Gm}{pc^2}\left[
 \frac{2+2\gamma-\beta}{3}
   +\frac{\eta}{6}(2\alpha_1-\alpha_2+\alpha_3+2\zeta_2)\right].
\]
质量比必须乘在后一个组合上：在测试质量极限，运动源的两体质量流消失，式子应退回静止单体外场。对水星，这一项受质量比抑制；对质量相当的双星则不再有这种抑制。套用弱场系数到双中子星以前，还须检验这些星能否用上述弱自引力质量张量表示。

若要联合测定四极矩和 PN 参数，可对 \eqref{eq:g3-relative-ppn} 中的 $\bm a_Q$ 也作一次轨道平均。在以对称轴为 $Z$ 轴的坐标中，倾角为 $i$，则 $s=\sin i\sin(\omega+f)$。把它投影到径向、顺行和法向，所得每转的 $\Delta p,\Delta e,\Delta i$ 在这一阶均为零，但 $\Delta\Omega_{\rm node}=-3\pi J_2(R_2/p)^2\cos i$，面内近心点角变量为 $\Delta\omega_Q=6\pi J_2(R_2/p)^2(1-5\sin^2i/4)$。这里及下式 $J_2$ 均按 $R_2$ 归一。观测者相对于固定方向得到的经度还要加 $\cos i\,\Delta\Omega_{\rm node}$。当 $i\to0$，两项合成 $3\pi J_2(R_2/p)^2$，不能只取 $\Delta\omega_Q$。故只看一颗近赤道行星的近心点，经常分不清中心体扁率与非线性引力；加入一条大倾角轨道的交点，就多了一个 PN 球对称场不给的量。

六个密切要素的前五个说明椭圆的大小、形状和朝向，第六个才告诉我们钟在何时经过曲线上的某点。如果选 $f$ 而非近心点历元作为独立角，尚有 $\dot f=h/r^2-\dot\omega-\cos i\,\dot\Omega_{\rm node}$。这是一条运动学恒等式，表达轨道面自身转动时，从固定观者看来“沿椭圆走过的角”须作哪些扣除。若把进动全部加在 $\omega$ 上却仍取 $\dot f=h/r^2$，在逐圈累计到达时刻时就会把同一部分进动算两遍。对圆轨道则宜采用平均经度与离心率矢量，避免在不存在近心点时给 $\omega$ 一个可测的初值。

\subsection{Cavendish 实验中的局域引力常数}\label{sec:pn-local-g}

设球形场源在外部引力场中自由下落，在离它一小段固有距离 $r_{\rm p}$ 处放一试验体。以非引力装置维持这一距离不变，并用测力装置读出单位试验质量所需的径向支撑加速度。将其随 $r_{\rm p}$ 改变而按 $M/r_{\rm p}^2$ 变化的部分取出，其系数就是\textbf{局域引力常数} $G_{\rm L}$。试验体仅为测力探针，源的 $M$ 用同一实验室的惯性质量单位标定。不能用 PPN 坐标距离和坐标加速度相乘便当作仪器读数，否则仅仅重新标度时间、长度也会给出一个假的 $G$ 变化。

取实验室与外部物质的相对运动在一次测量间变化很小，外部势 $U_{\rm e}$ 在源与试验体之间近似为常数。若 $r$ 是相应的坐标距离，沿同时线测长给出 $r_{\rm p}\simeq r(1+\gamma U_{\rm e}/c^2)$，钟速尚有 $\d\tau/\d t\simeq1-U_{\rm e}/c^2$ 的改正。这里只写由外场常数部分产生的项。局部源自身使度规偏离平直所给的二体修正可以另算。设 $\bm e$ 自场源指向试验体，试验体的 4-加速度 $A^\mu$ 投影在随室运动的径向单位标架上，才是加速度计要读出的数。还要减掉外场对两者共同的加速。两质点若都放手，外场所给的相对径向加速度为 $r_{\rm p}e^ie^j\partial_i\partial_j U_{\rm e}$。要保持距离，支撑力须提供相反的一项。它正比于 $r_{\rm p}$，不会混作 $1/r_{\rm p}^2$ 的系数。

怎样得到这一系数？把 $U=U_{\rm s}+U_{\rm e}$ 代进势和试验粒子的运动方程，$U_{\rm s}=GM/r$。对 $U^2$ 只留 $2U_{\rm s}U_{\rm e}$，对 $\Phi_2,\Phi_W$ 等二源势也只留一次含局部源、一次含外部源的贡献。另需保留 $w^2U_{\rm s}$ 及方向投影，其中 $\bm w$ 是仪器相对于所设优选系的速度。把保持 $r_{\rm p}$ 不变所需的力投影到固有标架，换成固有时间和距离，最后才对 $M/r_{\rm p}^2$ 读系数。设外部物体可视为质点，以 $R_a$ 表示第 $a$ 个外源到实验室的距离，$\bm n_a$ 为该连线单位矢量。置 $U_{\rm e}=\sum_a GM_a/R_a$、$U_{\rm e}^{ij}=\sum_a (GM_a/R_a)n_a^in_a^j$。取局部场源的半径相对 $r_{\rm p}$ 可忽略，则
\begin{align*}
 \frac{G_{\rm L}}{G}=1
 &- (4\beta-\gamma-3-\zeta_2-3\xi)\frac{U_{\rm e}}{c^2}
 -\frac{\alpha_1-\alpha_2-\alpha_3}{2}\frac{w^2}{c^2}\\
 &-\frac{\alpha_2}{2}\frac{(\bm w\cdot\bm e)^2}{c^2}
 +\xi\frac{U_{\rm e}^{ij}e^ie^j}{c^2}+O(c^{-4}).
\end{align*}
此处 $G$ 为前文在远处优选静系内归一化的常数。若外源只有一个，末项即正比于 $U_{\rm e}(\bm n\cdot\bm e)^2$，实验室即使自由下落、内部潮汐也已扣除，转动仪器仍可辨认这个方向。$\alpha_2$ 项则对准 $\bm w$。反转测量轴不能改变这两个平方项，转轴与所选方向夹角变化才会改变。还有两项不随仪器朝向改变，单在一处、以相同速度重复测量通常会将它们吸收到 $G_{\rm L}$ 的标定值中。在广义相对论这些修正全部消失，即使 $U_{\rm e}$ 很大，只要弱场近似仍成立，外部潮汐与局部 $G$ 仍是不同的量。

若实验室的场源是半径不可忽略的球，把局部 $X$ 一律写作 $GMr$ 还会漏下一项。取 $I_{\rm s}=\int\rho_*y^2\,\d^3y$，在球外对 $|\bm r-\bm y|$ 作角平均，精确得 $r+y^2/(3r)$，故 $X_{\rm s}=GMr+GI_{\rm s}/(3r)$。可见即使球外 $U_{\rm s}$ 完全等于点源的势，$X_{\rm s}$ 的二阶导数仍会把 $\chi=I_{\rm s}/(M r_{\rm p}^2)$ 带入计算。具体而言，$U_{\rm e}$ 的系数中 $3\xi$ 改为 $(3+\chi)\xi$，$w^2$ 的系数中 $\alpha_2$ 改为 $(1-\chi)\alpha_2$，最后两个方向项则同乘 $1-3\chi$。在普通 Newton 场中熟知的球壳定理，不能直接搬到所有 PPN 势上。

\subsubsection{径向支撑力的换算}

上述 $G_{\rm L}$ 的各向同性项有一个可手算的核对。设局部源和探针瞬时同静，忽略源的平动和自转，在几何制内只取一次含局部势 $U_s$、一次含外势 $U_e$ 的交叉项；局部源的纯二体 PN 项可由改变探针距离另行定出。令 $\bm e$ 从源指向探针，$r$ 为坐标距离，$r_p$ 为固有距离。先用 $r_p=r(1+\gamma U_e)$ 将 $r$ 换掉，再以 $E_r$ 指向源，$E_r^j=-e^j(1-\gamma U_e)$ 作固有投影。对不随局部装置变动的外场，上述两个换算和 $A^\mu$ 从固有时换成坐标时所带的 $\d\tau/\d t$ 因子分别都要作用到 Newton 项。若只在度规势的梯度上作一次修正，而拿坐标半径和钟速读 $G_{\rm L}$，所得结果连广义相对论中 $G_{\rm L}=G$ 的检查都未必通过。

具体地，对一般理论，在投影与距离换算以前，径向支撑项沿 $\bm e$ 含有
\begin{align*}
 -(\bm e\cdot\bm\nabla U_s)
 \left[1-(4\beta+2\gamma-3-\zeta_2-4\xi)U_e
                 -\tfrac12(\alpha_1-\alpha_3)w^2\right]
 -e^i X_{s,ijk}(\tfrac12\alpha_2w^jw^k-\xi U_e^{jk}).
\end{align*}
这里 $U_s,X_s$ 只由局部源给出，尚应另加与 $r_p$ 成正比的外界潮汐。取点质量 $X_s=Mr$，有 $e^iX_{s,ijk}=-M(\delta_{jk}-e_je_k)/r^2$，因而最后一项同时贡献与 $w^2$ 及 $(\bm e\cdot\bm w)^2$ 有关的两种项。还须注意前面把加速度投影到 $\bm e$ 时用的是已有时钟与尺改正的坐标式，最后换成径向单位标架又产生一个空间度规因子。将所有项按 $M/r_p^2$ 合并，得到本小节的 $G_{\rm L}$。例如对 $w=0$ 且 $\xi=\zeta_i=0$ 的理论，交叉项中 $4\beta+2\gamma-3$ 与三项空间尺度修正相抵，最终是 $4\beta-\gamma-3$。此组合也出现在同一理论的 Nordtvedt 效应中，其物理原因都是内部引力相互作用是否会因外部引力势的均匀部分而改变。

若要把一种此类实验的灵敏度换算给另一种，应先写出两次测值的差。设两次实验处于同一地点与速度，测量轴由 $\bm e$ 转到 $\bm e'$，点源近似下只有二次型 $E_{ij}=\xi U_e^{ij}-\alpha_2w_iw_j/2$ 的无迹部分贡献，即 $\Delta G_{\rm L}/G=E_{ij}(e'^ie'^j-e^ie^j)/c^2$。若始终不转仪器而仅改变所在的外势，才对各向同性项也有灵敏度。由此可知，“$G$ 随地点改变”至少有一个角平均和一个随方向的部分；一次旋转实验不能同时限定这两类的全部系数。

\subsection{自转天体的质心与自旋}

对一团自转物质，应先指定绕哪条世界线取角动量。在最低阶质心系中有 $\bm S_A=\int_A\rho_*\bm y\times\bar{\bm v}\,\d^3x$，其中 $\bm y=\bm x-\bm z_{*A}$。挪动所取的中心而保持同一团物质的总角动量不变，必有一部分角动量在自转与轨道之间挪动。在相对论中尤其如此：判断一颗有宽度的星在不同部位“同一时刻”的能量分布，要先选同时面。设总动量为 $p^\mu$、总角动量为 $J^{\mu\nu}$，在局域近似平直区域内相对于 $z^\mu$ 定义 $S^{\mu\nu}=J^{\mu\nu}-2z^{[\mu}p^{\nu]}$。故换到另一中心 $z^\mu+\delta z^\mu$ 时，$S^{\mu\nu}$ 随之改变 $-2\delta z^{[\mu}p^{\nu]}$。可施加 $S^{\mu\nu}p_\nu=0$ 作为选取中心的条件，在总动量静系便有 $S^{0i}=0$。之后的轨道位置及自旋须配用这一约定。只更换质心而不给轨道角动量作补偿，会使同一观测看来有两种进动量。

取 \S\ref{sec:pn-transport-shell} 的试验自旋，以 PPN 坐标的静止正交标架作瞬时无转动的 Lorentz 推进，指定质心系的空间轴。将同一输运计算用于式 \eqref{eq:exp-ppn-all}，有 $\dot{\bm S}=\bm\Omega\times\bm S$，其中
\[
 \bm\Omega=\frac{1}{c^2}
 \left[(\gamma+\tfrac12)\bm v\times\bm\nabla U
                     -\tfrac12\bm\nabla\times\bm B\right]+O(c^{-4}).
\]
$\bm B$ 是本节已定义的 $g_{0i}$ 势，勿与电磁场混淆。第一项还含有比较沿轨道相继两个质心静系时发生的旋转，故其系数并非单由空间曲率的 $\gamma$ 给出。第二项只取 $B_i$ 的旋度，因此 $\partial_i\dot X$ 对它没有贡献，转动场源的质量流则可留下贡献。若取一个沿质心世界线平行移动的标架当作比较轴，自旋的三个分量反而是常数。后文实验测的是陀螺相对于外部参照方向的变化，必须先说明所取的轴。对质量、自旋均不可忽略的双星，还须把一颗星的质量流及自旋引起的场放回伴星的运动方程，才能同时求自旋与轨道角动量的演化。

\subsection{带自旋的两体：进动与轨道反馈}

只把小陀螺当作探针时，轨道可先于自旋求出。但若二星本身自转快，还须反过来计算自旋怎样改变轨道。以下在相隔很远的物体之间只保留各星的一次自旋及来自不同星的自旋乘积，略去自转导致的自身质量四极矩。暂取几何制。令 $S_a^{ij}=\epsilon^{ijk}S_a^k$，$\sigma_a^{ij}=S_a^{ij}/m_a$，$\bm n=\bm n_{ab}$，并写 $n^{\langle ij\rangle}=n^in^j-\delta^{ij}/3$、$n^{\langle ijk\rangle}=n^in^jn^k-(n^i\delta^{jk}+n^j\delta^{ik}+n^k\delta^{ij})/5$。定义星的代表世界线时，若把自旋的混合分量定为 $S_a^{0j}=-\lambda S_a^{jk}v_a^k$，所选位置相对于原先所取的物质中心便移动 $-\lambda S_a^{jk}v_a^k/m_a$ 到最低非平凡阶。$\lambda=1$ 对应协变的质心条件；$\lambda=0$ 对应用 PPN 坐标的同时面定义原来的物质中心。恢复单位时该位移含 $c^{-2}$，量纲仍为长度。

例如取全守恒理论、$\bm w=0$ 与 $\lambda=1$，则从体积分得到的领先自旋--轨道加速度可写成
\begin{align*}
 (a_{a,{\rm SO}})^j=\frac32\sum_{b\ne a}\frac{m_b}{r_{ab}^3}
  \big\{&n^{\langle jk\rangle}(v_a^p-v_b^p)
          [(2\gamma+2)\sigma_a^{kp}+(2\gamma+2)\sigma_b^{kp}]\\
        &+n^{\langle kp\rangle}(v_a^p-v_b^p)
          [2\gamma\sigma_a^{jk}+(2\gamma+2)\sigma_b^{jk}]\big\}.
\end{align*}
这两个张量的缩并形式不同，不能先把 $\sigma_a+\sigma_b$ 解释为一个总自旋再与 $\bm v$ 叉乘。在半守恒理论中若保留优选参考系参数，大括号内还出现 $\alpha_1 n^{\langle jk\rangle}[(w+v_a)^p\sigma_b^{kp}-(w+v_b)^p\sigma_a^{kp}]/2$ 及 $\alpha_1 n^{\langle kp\rangle}(v_a-v_b)^p\sigma_b^{jk}/2$。这里按半守恒条件取 $\alpha_3=0$；更一般的自加速还要连同前述内部力矩处理。选择协变的代表线，独立坐标速度只通过相对于优选系的速度出现；在全守恒情形则只剩相对速度。用另一质心定义，坐标加速度会多出代表线位移的二阶导，比较脉冲发射位置时也须作同样的位移。

自旋--自旋项则是一个旋转质量流产生的场作用到另一颗星的角动量上。其公共系数为 $4\gamma+4+\alpha_1$，从体积分可写为
\[
 (a_{a,{\rm SS}})^j=-\frac{15}{8}(4\gamma+4+\alpha_1)
    \sum_{b\ne a}\frac{m_b}{r_{ab}^4}\,
    \sigma_a^{kp}\sigma_b^{kq}n^{\langle jpq\rangle}.
\]
自旋掉头一次，SO 及含它的 SS 项均反号；两星自旋一起掉头则 SS 项不变。轨道面或近心点的长期漂移若只量一个方向，未必能将这两种贡献同普通质量四极矩分开。上式不包含因旋转将单颗中子星压扁造成的 $S_a^2$ 项；要计算它须另外给出星的四极矩与自旋的关系。

反过来研究自旋方向时，应先把 PPN 坐标中的 $S_a^i$ 换成随星体瞬时共动的正交标架所测的自旋。将内部可写成全导数的项吸收后，一种取法为
$\bar{\bm S}_a=[1+v_a^2+(2\gamma+1)U_a]\bm S_a
 -\bm v_a(\bm v_a\cdot\bm S_a)/2$。这样做的好处是单纯的进动应保持 $\bar{\bm S}_a^2$。若忽略形变星体可能有的内部自力矩，则 $\dot{\bar{\bm S}}_a=\bm\Omega_a\times\bar{\bm S}_a$，其中三个相互独立的来源为
\begin{align*}
 \bm\Omega_{a,{\rm SO}}&=\frac12\sum_{b\ne a}\frac{m_b}{r_{ab}^2}
 \bm n_{ab}\times[(2\gamma+1)\bm v_a-(2\gamma+2)\bm v_b],\\
 \bm\Omega_{a,{\rm SS}}&=\frac{4\gamma+4+\alpha_1}{8}
  \sum_{b\ne a}\frac{3\bm n_{ab}(\bm n_{ab}\cdot\bm S_b)-\bm S_b}{r_{ab}^3},\\
 \bm\Omega_{a,{\rm F}}&=-\frac{\alpha_1}{4}\sum_{b\ne a}\frac{m_b}{r_{ab}^2}
                  \bm n_{ab}\times(\bm w+\bm v_b).
\end{align*}
这些量是固有自旋相对于用远方物质定向的 PPN 轴、按坐标时计算的转速。换成台站时钟和实际参照方向时还须作相应换算，结果已无上面的 $\lambda$。若保留轴对称但不完全球形的内部结合能张量，优选系或优选位置还可从该张量获得力矩，其形式为 $\epsilon^{jpq}\Omega_a^{pn}[\alpha_2(w+v_a)^q(w+v_a)^n-2\xi\sum_b(m_b/r_{ab})n_{ab}^qn_{ab}^n]$。球对称时 $\Omega_a^{pn}\propto\delta^{pn}$，与反对称符号缩并后为零。由哪些旋量给出轨道上的额外力、哪些只使陀螺的轴转动，至此可以分别计算。下面与卫星和脉冲计时比较时，恢复普通单位。

\section{经典检验}

\subsection{天球上可量的角}

\hyperlink{subsection.2.2.3}{\S\,2.2.3} 已从 Schwarzschild 测地线算得太阳附近的光线偏折。用望远镜只看得到光线到达时的方向，如何知道它原来朝哪里走？日食观测所比的是星与星的角距离。同一片天区在太阳移开后仍可拍摄，用这时的星位作参照，便能从底片中找出有太阳时的相对位移。离日心近的星像向外移得多，离得远的也在移动，两星偏折之差才是实际信号。若底片的焦距略有改变，各星距片心的距离也会同时伸缩。过去因日食期间的温度变化而出现的焦距问题，须借参考星的位置独立标定；多给同一张底片上的星计数，不能消除此项误差。

角距离有与坐标无关的定义。取观者的单位 4-速为 $u$，$u\cdot u=-1$，两束光在接收处的波矢为 $k,k'$。光在观者空间的方向是 $k$ 在该空间的投影除以观者所测频率 $-u\cdot k$。由 $k\cdot k=k'\cdot k'=0$，两方向之间的夹角满足 $\cos\theta=1+(k\cdot k')/[(u\cdot k)(u\cdot k')]$。归一化时两束光的频率消去，观者的速度却留下来。故地球公转引起的光行差须按观测日期改正，同一历元两源在天球上的相对移动也须扣除。

设远处点源与太阳在天空中的未扰动角距为 $\Phi$，地球距太阳 $r_\oplus$。由一阶光线方程积分到观者处，像相对日心向外偏移 $\delta\Phi=(1+\gamma)GM_\odot\cot(\Phi/2)/(c^2r_\oplus)$。日缘处 $b=r_\oplus\sin\Phi$ 约为太阳半径，广义相对论给出熟悉的 $1.75$ 角秒。若光源与太阳相隔 $90^\circ$，偏折仍有约 $4$ 毫角秒。故只要量角精度足够，检验并不限于日食或擦过日缘的光线。若以另一遥远光源作参照，应分别计算两源的位移，再沿两源的连线取差。譬如两颗很近的星在日面同侧时，大部分偏折是共同的，角距改变量可以远小于任一颗星的偏折。

射电观测可在白昼进行。相隔数千公里的两处天线接收同一类星体的信号，按两地原子钟记时，再把两份记录互相关，移动其中一份直至波形最为吻合。这个移动量便给出同一波前抵达两站的时间差。令 $\bm B$ 从第一站指向第二站，$\bm s$ 为从台站指向源的单位矢量，远源的几何时差即 $-\bm B\cdot\bm s/c$。一条基线一次只能量到源方向的一个投影，地球自转却不断改变基线的朝向，多处台站还可组成不同的基线。用这种\textbf{甚长基线干涉}(very-long-baseline interferometry, VLBI) 比较两射电源，便可从时间差还原角距变化。为明白所需精度，取垂直于视线的基线长 $6000\,\mathrm{km}$，$0.1$ 毫角秒的角位移只给出约 $10$ 皮秒的时差。

类星体 3C279 等每年随太阳在天球上的视运动而接近日缘，最易分辨引力引起的位移，但也最受日冕折射干扰。远离日面时信号弱，却有满天的源可用。分析者一并估计台站钟差、对流层延迟、基线和射电源的位置，按太阳的预报位置为各次观测算出引力项，拟合它的公共系数 $(1+\gamma)/2$。例如分析截至 2010 年的 VLBI 资料，得到 $\gamma-1=(-0.8\pm1.2)\times10^{-4}$。近太阳的专门观测和跨度多年的全天观测，纵然都归结为同一个 $\gamma$，限制测量精度的因素并不相同。

星系也能作偏折实验的引力源，但质量的处理与太阳不同。若背景源、前景星系和观者近乎共线，像围成一圈，圈的角半径 $\theta_{\rm E}$ 由视线方向投影的圈内质量决定。记透镜、光源距观者及两者之间的角直径距离为 $D_{\rm l},D_{\rm s},D_{\rm ls}$，对薄的轴对称透镜有 $\theta_{\rm E}^2=2(1+\gamma)GM_{\rm 2D}(<D_{\rm l}\theta_{\rm E})D_{\rm ls}/(c^2D_{\rm l}D_{\rm s})$。单有一个像环只知道质量与 $1+\gamma$ 的乘积。可另从星系吸收线的展宽定出恒星速度弥散，利用恒星所受的 Newton 引力估计同一星系的质量分布，再看它对光的偏折是否相符。Bolton 等在 2006 年分析了十五个椭圆星系透镜，在所取恒星轨道及质量分布模型下，测得 $(1+\gamma)/2$ 与 1 相符，精度约为百分之五。这样可以在远大于太阳系的尺度上比较物质运动与光的传播，代价是必须估计透镜内的暗物质和恒星轨道的方向分布。

\subsubsection{同一块底片上的两颗星}

具体算一算两个光源的角距。记观者指向太阳、两颗星的无扰动单位矢量为 $\bm u,\bm s,\bm s'$，$\bm u\cdot\bm s=\cos\Phi$、$\bm u\cdot\bm s'=\cos\Phi'$、$\bm s\cdot\bm s'=\cos\theta_0$。由太阳造成的星像位移须在星的切平面内，并背离太阳。因而以 $A=(1+\gamma)GM_\odot/(c^2r_\oplus)$ 记共同系数，一阶角位移矢量恰为 $\delta\bm s=A[(\cos\Phi)\bm s-\bm u]/(1-\cos\Phi)$。直接点乘 $\bm s$ 得零，取模得前述 $A\cot(\Phi/2)$，两个条件都满足。把 $\bm s+\delta\bm s,\bm s'+\delta\bm s'$ 代回天球角距，对 $\theta=\theta_0+\delta\theta$ 展开，得
\[
 \delta\theta=\frac{A}{\sin\theta_0}
  \left[
    \frac{\cos\Phi'-\cos\Phi\cos\theta_0}{1-\cos\Phi}
    +\frac{\cos\Phi-\cos\Phi'\cos\theta_0}{1-\cos\Phi'}
  \right].
\]
此式只含三段可在天球上标定的角距。第一项为第一颗星的位移沿两星大圆的投影，第二项为参考星的贡献。若两星分居日心两边，位移增加二者间距；若二者都在同一侧同一大圆上，净改变量取决于谁更靠近日缘。选一颗亮星作参照并不意味着它的理论改正为零。早期对照底片在夜间拍摄，看不到太阳，归算时用的是历表所给日心方向。

缩小视场也有代价。若整个视场只占狭小一角，引力位移在其中可近似写成一个常数平移再加一个线性变形。两张底片独立确定原点、旋转角与放大率，平移可被原点吸收，部分线性项也可被比例尺吸收。要留下引力所特有的随日心距和位置角的变化，必须有足够的星和合适的分布。有云遮住外侧参考星，失去的不仅是几个独立样本，还可能失去标定比例尺的条件。放大率变化给出正比于距片心距离的位移，引力却近似反比于距日心距离；若多颗星恰好处在近似同一日心距，这两种规律将很难区别。

1919 年 Príncipe 地点的十六张底片，只有两张星像可靠，约有五颗星可用，换算到日缘的偏折为 $1.60\pm0.31$ 角秒。Sobral 使用的两套仪器不应视作一套仪器反复读数。较好的八张底片给 $1.98\pm0.12$ 角秒，另一台仪器的像质及比例尺在食甚前后有变化。1979 年用测量机重测 Sobral 底片，这后一组给 $1.55\pm0.34$ 角秒。三组数可以比较，但不能简单按星数相加求总精度。对前两组，除以广义相对论日缘预言即得 $(1+\gamma)/2=0.91\pm0.18$ 和 $1.13\pm0.07$。这里除以二已体现在参数定义中，若改报 $\gamma-1$，误差也须乘二。Cavendish、von Soldner 的微粒偏折与 Einstein 1911 年的等效原理论证都只给约半数的日缘位移。\hyperlink{subsection.2.2.3}{\S\,2.2.3} 已比较它同 Schwarzschild 结果的差别。在此用 $(1+\gamma)/2$ 归一化，正可将各理论的差别同每一仪器的误差放在同一尺度上。

射电实验以时间差作读数，则有另一套标定办法。两台站若有一常数钟差，它给该基线上的所有源加上同样的时差。分别在同一基线上看天空相近的两个源，横向位移的差仍保留，变化缓慢的钟差和相近视线上的大气延迟则可大半消去。若两源方向几乎相同，差异也会太小；若相距太远，两束射线经过大气的路径不再相近。因而要在接近日面、参考源够亮、相对方向合宜三者之间选取观测日期。1995 年对 3C273、3C279 的专门测量得到 $\gamma-1=(-8\pm34)\times10^{-4}$；2009 年加入邻近的两个源后为 $(-2\pm3)\times10^{-4}$。另一条路是用遍布天球的源估计参考架和台站运动。2004 年一项这种分析用了 87 个台站约二百万次观测，得到 $(-1.7\pm4.5)\times10^{-4}$。前者在少数合宜时段对某些参数很敏感，后者逐日提供许多几何方向，因而所得精度不应只按样本数比较。

引力透镜的角分离还涉及另一层识别。1979 年发现的 Q0957+561 由两个相距约 $6$ 角秒的类星体像组成，二者红移都约 $1.41$，光谱十分相似。同一光源通过前景星系形成两个像时，几何路程和引力时延均不同；连续比较亮度起伏应见相同变化先后到达。单见两点光并不能定出透镜质量；连成像结构、光度演变、透镜星系内恒星的速度弥散一起考虑，方可分开“同源多像”和“两源恰好靠近”。球对称薄透镜的圈内投影质量包含星系外围沿视线叠加的物质，不能直接与一个以同一半径截下的三维球质量等同。两项质量估计必须采用同一三维密度分布，然后分别计算天球投影和恒星运动。

\subsection{上合前后的雷达回波}

从地面向行星发一束雷达脉冲，用同一只钟记下发射与回波到达的时刻。若在上合前后观测，地球、太阳与行星接近共线，雷达波的去程与回程都可掠过太阳。太阳引力使往返多费多少时间？

\begin{figure}[htbp]
 \centering
 \begin{tikzpicture}[x=1cm,y=1cm,>=Stealth,font=\small]
  \coordinate (O) at (0,0);
  \coordinate (A) at (-3.8,.92);
  \coordinate (B) at (3.05,.92);
  \coordinate (F) at (0,.92);
  \filldraw[fill=orange!17,draw=black!70] (O) circle(.34);
  \fill (O) circle(.035) node[below=9pt] {日心};
  \draw[thin,black!58] (A)--(O)--(B);
  \draw[thick,-{Stealth[length=2mm]}] (A)--(B);
  \draw[dashed] (O)--(F);
  \draw (0,.78)--(.14,.78)--(.14,.92);
  \node[left] at (0,.55) {$b$};
  \fill (A) circle(.045) node[above=3pt] {发射点 1};
  \fill (B) circle(.045) node[above=3pt] {接收点 2};
  \node[below left] at (-2.02,.47) {$r_1$};
  \node[below right] at (1.55,.46) {$r_2$};
  \node[above] at (-.95,.92) {$\bm k$};
 \end{tikzpicture}
 \caption{上合附近的单程光路。实线箭头是固定两端点的零阶直线，$b$ 是日心到这条直线的距离。取垂足为光路坐标的零点后，势积分中的距离为 $\sqrt{b^2+\ell^2}$。}
 \label{fig:exp-conjunction}
\end{figure}

设 $U=GM_\odot/r$。由前述光行时公式，沿光路每走一段 $\d\ell$，单程即有 $(1+\gamma)U\,\d\ell/c^3$ 的附加时延。在此精度下可沿未偏折的直线积分，因为固定两端点时，直线路长的一阶变分为零。如图~\ref{fig:exp-conjunction}，记发射点和接收点的日心位置为 $\bm x_1,\bm x_2$，$r_i=|\bm x_i|$，两点间距为 $R=|\bm x_2-\bm x_1|$。在连接两点的直线上，由离日心最近处起算有向长度 $\ell$，最近距离为 $b$，于是 $r^2=b^2+\ell^2$。单程坐标光行时为
\begin{align*}
    t_2-t_1
    &=\frac{R}{c}+\frac{(1+\gamma)GM_\odot}{c^3}
       \int_{\ell_1}^{\ell_2}\frac{\d\ell}{\sqrt{b^2+\ell^2}}\\
    &=\frac{R}{c}+\frac{(1+\gamma)GM_\odot}{c^3}
       \ln\frac{r_1+r_2+R}{r_1+r_2-R}.
\end{align*}
等号只写到引力势的一阶。对数项称为\textbf{Shapiro 时延}。后一种写法只用两端点的坐标，亦适合两站不在日心两侧的情形。两端点接近同一视线而分居太阳两侧时，分母很小，测量最为有利。这里的光行时是两个相距甚远的事件之间的坐标时差，任意自由落体观者在光经过自己时测得的光速仍为 $c$。

上合附近去程跨过 $\ell=0$，积分上下限近似为 $-r_\oplus,r_{\rm p}$，其中 $r_{\rm p}$ 为目标至日心距离。略去端点移动，往返时延为 $\Delta t_{\rm S}\simeq 2(1+\gamma)(GM_\odot/c^3)\ln(4r_\oplus r_{\rm p}/b^2)$。取广义相对论的系数以及 $GM_\odot/c^3\simeq4.93\,\mu\mathrm{s}$，若目标距日心也是一个天文单位，信号掠过日缘，得 $\Delta t_{\rm S}\simeq 2.4\times10^{-4}\,\mathrm{s}$。地面雷达若把这个时间全算作平直空间内的路程，会把单程距离高估约 $36\,\mathrm{km}$。

实测时没有一束在平直时空中传播的脉冲可作比较。办法是用远离上合时的测距资料定出轨道，由此算出上合前后改变的几何距离，再寻找随 $b$ 对数变化的部分。具体设发射、转发、接收事件的坐标时各为 $t_1,t_2,t_3$。去程应将地面台站在 $t_1$ 的位置与目标在 $t_2$ 的位置代入单程公式，回程则用目标在 $t_2$、台站在 $t_3$ 的位置。模型还须计入转发器的固定延迟，最后把 $t_3-t_1$ 换成同一地面钟的固有时差，才能与记录相比。早期向金星、水星表面发射雷达波，实际收到的是行星一片表面的回波，表面起伏与自转会改变回波形状。航天器上的转发器有确定的位置，跟踪它可免去这项困难。1976 年 Viking 火星任务已把时延系数测到千分量级。

也可用航天器转发连续无线电波。设地面以频率 $\nu$ 发送，忽略已知的转发倍频比，接收的相位约为 $2\pi\nu[t-T(t)]$，其中 $T$ 是往返光行时。求导即知，扣除轨道运动的 Doppler 频移后，相对频移约为 $y_{\rm S}=-\d\Delta t_{\rm S}/\d t$。在上合附近，$r_\oplus,r_{\rm p}$ 的变化比 $b$ 缓慢，故 $y_{\rm S}\simeq4(1+\gamma)(GM_\odot/c^3)\dot b/b$。光路移近日面时频移为负，离开时转正。恰在最近处，时延最大，频移却过零，这说明测距和测速应当各用什么样的信号去拟合。

2002 年 Cassini 探测器上合时，探测器在距日心约 $8.43$ 个天文单位处，信号一度从距日心 $1.6$ 个太阳半径处经过。这时日冕的自由电子使不同频率的电波产生不同的时延。稀薄等离子体的群时延正比于沿光路积分的电子数，并在领头阶反比于载波频率的平方；引力时延则与频率无关。若两频段的残余频移形如 $y_i=y_{\rm S}+K/\nu_i^2$，作组合 $(\nu_1^2y_1-\nu_2^2y_2)/(\nu_1^2-\nu_2^2)$ 即可消去共同的色散项。实际要区分上下行链路，并计入地面大气和航天器的非引力加速。借 X、Ka 波段定出日冕贡献后，Bertotti 等由这次上合得到 $\gamma-1=(2.1\pm2.3)\times10^{-5}$。

还有一个从数据本身可见的几何特征。以最接近日心的时刻为 $t=0$，短时间内让天空中的光路以近似恒定速度 $v_\perp$ 扫过太阳，便有 $b(t)^2=b_0^2+v_\perp^2t^2$。置 $t_b=b_0/v_\perp$，相对于 $t=0$ 的往返时延变化以及相对频移为
\[
 \Delta t_{\rm S}(t)-\Delta t_{\rm S}(0)
   =-2(1+\gamma)\frac{GM_\odot}{c^3}\ln[1+(t/t_b)^2],
 \qquad
 y_{\rm S}(t)=4(1+\gamma)\frac{GM_\odot}{c^3}\frac{t}{t_b^2+t^2}.
\]
时延变化是偶函数，频移是奇函数，在 $|t|=t_b$ 处频移的绝对值最大。只看上合一侧的短弧段，钟漂和未扣净的航天器加速度都能使测得频率近似线性变化；从两侧连续跟踪，才容易辨出上述双峰形状。离日面最近的几小时内日冕扰动往往最强，太阳还可能遮断链路。最大时延出现在某个时刻，并不要求在该时刻取得一条回波才可测量其系数。从资料的每段起止时间及预报的 $b(t)$ 算出各段平均频率，更切合实际仪器。对积分时段 $[t_1,t_2]$，平均量是 $-[\Delta t_{\rm S}(t_2)-\Delta t_{\rm S}(t_1)]/(t_2-t_1)$，若时段不短，不宜直接取在中点的瞬时导数。

雷达表面测距和相干转发实测的量也不一样。行星的整个可见圆面都受照，首先返来的回波来自靠近雷达的那一小片，后到的来自离它稍远的环带。转动速度不同的环带又产生不同频移。按到达时间和频率分箱，可把一部分地形信息同中心距离分开；但若要从回波前缘反求行星中心，仍须知道该处地貌和行星半径。改用环绕器，将行星地形的困难换成环绕器相对行星的定轨问题。火星表面着陆器的相对位置一天内随火星自转，带来一个可预报的公转以外周期。将这种短周期运动留在模型中，可分别测得着陆点位置和火星质心的轨道；忽略它却会把正常的日运动当成时延的噪声。

在无线电链路中，须区分相位走过的周期数与脉冲包络的传输时间。冷等离子体内群延迟和相位延迟的色散贡献符号相反，不能将测距标定直接作为测速的相位改正。Cassini 的 X 波段和 Ka 波段上行频率约为 $7.2$ 和 $34.3\,\mathrm{GHz}$，这样大的差别使日冕电子柱密度的影响在两组记录中有很大的相对改变量。仅减去两组频移会同时减掉引力信号；先用差值估出色散，再保留共同分量，才有上面的双频组合。所得 $2.3\times10^{-5}$ 是对 $\gamma-1$ 的误差，对时延系数 $(1+\gamma)/2$ 则约为 $1.15\times10^{-5}$。对无质量的 Brans--Dicke 标量，$\beta=1$ 且 $\gamma-1=-1/(\omega_{\rm BD}+2)$，故这种观测特殊的意义在于把正 $\omega_{\rm BD}$ 的下限推到数万量级。若标量有质量，火星与太阳之间的作用、不同时刻无线电经过太阳附近的作用未必有同一个系数。该下限不能原样套给短程的理论。

\subsection{从行星历表中分离近日点进动}\label{sec:pn-ephemeris}

水星近日点每世纪约 $43$ 角秒的额外进动并非直接读自望远镜。把几世纪的天球位置换到固定的参考系，扣除其余行星摄动后才有这个数。雷达和航天器跟踪改变了定轨的做法。观测给出的是某历元的角位置、往返时间或频率比，近日点方向要从用这些资料定出的轨道中算。在参数化理论里尤其要问，改变某个参数能否由另一个参数的改变补偿？

仍记半通径 $p=a(1-e^2)$。忽略水星质量和优选系效应，轨道近于太阳赤道面时，每公转一周的近日点进动中球对称的后 Newton 项与太阳质量四极矩的主项为
\[
    \Delta\varpi\simeq\frac{6\pi GM_\odot}{c^2p}\,
          \frac{2+2\gamma-\beta}{3}
           +3\pi J_2\left(\frac{R_\odot}{p}\right)^2 .
\]
这里 $R_\odot$ 是太阳半径，$J_2$ 是外部势 $U=(GM_\odot/r)[1-J_2(R_\odot/r)^2P_2(\cos\theta)+\cdots]$ 中的无量纲四极矩，$\theta$ 由太阳自转轴起算。位于赤道面的行星感受到的四极势随 $r^{-3}$ 变化，也能使椭圆不闭合。一般倾斜轨道还要把轨道面本身的进动计入，不能把近心点幅角直接当作相对恒星的近日点位移。因而只用一颗行星的一个进动率，至多约束 $2+2\gamma-\beta$ 与 $J_2$ 的一个组合。

水星的球对称修正约为每世纪 $42.98$ 角秒。取由日震学推得的 $J_2\simeq2.2\times10^{-7}$，四极矩项约为每世纪 $0.03$ 角秒。对于早期检验，这个数很小；要把相对论项再量准三个数量级，它已不可忽略。日震学利用太阳的振荡模推断内部自转，给出了不靠水星轨道的 $J_2$。也可同时用其它行星和航天器资料，因为两类进动在每一轨道周期的改变量分别按 $1/p$ 和 $1/p^2$ 改变，而且四极矩会使倾斜轨道的交点进动。光行时独立给出 $\gamma$ 后，对 $\beta$ 的估计也随之改进。

历表的一个解往往含各大行星的初始坐标、速度、质量、小行星的摄动以及台站位置。设所有待估量为 $q_a$，第 $i$ 个观测为 $O_i$，某一组试取参数给出的预报为 $F_i$，则小改正满足 $O_i-F_i\simeq\sum_a(\partial F_i/\partial q_a)\delta q_a$。若改变 $J_2$ 与改变 $\beta$ 在每个历元留下的痕迹几乎成比例，两列偏导便近乎成比例，个别参数的误差由此变大。加入一个新实验的意义正在于它可改变这两列的比例。MESSENGER 在 2011--2015 年绕水星飞行，以航天器为测距目标大大改善了水星的位置资料。包括这类资料的一组历表分析给出 $J_2=(2.4\pm0.2)\times10^{-7}$，且 $\beta$、$\gamma$ 与 1 的差异均与零相容。比较两份历表给出的误差时，还要查清 $J_2$、$\gamma$ 是随数据同时估计，还是已经用外来的测量固定。

\subsubsection{日面形状、进动与定轨误差}

从 \hyperlink{subsection.2.2.2}{\S\,2.2.2} 的计算可见，水星在黄道上的运动从来不只有一个 $43$ 角秒的进动。把岁差等参考系效应扣除，以角秒每世纪计，观测所需的进动约 $574.1$，其中来自金星、地球和木星的摄动分别约为 $277.8,90.0,153.6$，再加上火星约 $2.5$、土星约 $7.3$，计 $531.2$；差约 $42.9$。这里各项经取整，精密历表须以所有天体的非线性方程共同积分，不能用这些取整的数值当作扣除模板。这个数量分配解释了为什么明明测的是水星，改进木星质量和小行星带模型也会改变最后的误差。

太阳的外形能否代替四极矩测量？设太阳光球近似为以角速度 $\omega_\odot$ 旋转的等势面，极半径和赤道半径之差用扁率 $f=(R_{\rm eq}-R_{\rm pol})/R_{\rm eq}$ 表示。分别把极点和赤道的半径代入含 $J_2$ 的引力势，并在赤道加入单位质量离心势 $\omega_\odot^2R_{\rm eq}^2/2$，保留小量的一次项，便有 $f\simeq 3J_2/2+\omega_\odot^2R_{\rm eq}^3/(2GM_\odot)$。可见连在这个刚性转动的简单模型里，所见太阳半径之差也不等于 $J_2$；还要先减掉表面转动的贡献。真实太阳并非刚体，光球边缘还是按一定亮度阈值识别，并非直接以尺量一个物质界面。

Dicke 和 Goldenberg 曾用遮光圆盘挡住太阳像，仅留下很窄的光环，比较在赤道和两极方向上露出的亮度。如果各处表面亮度已知，较亮的一侧意味着该处太阳半径较大。但赤道与两极的温差、光斑以及大气扰动都能改变窄光环的亮度。在 1960 年代，这类观测一度被解释为 $J_2$ 可达 $2.5\times10^{-5}$；按上式的进动换算，对水星会贡献每世纪数角秒。于是即使一个引力理论单独给出的球对称进动偏低，把太阳的贡献加上后也可能符合当时的资料。譬如取 $\omega_{\rm BD}=5$ 的长程标量理论给 $\gamma=6/7$、$\beta=1$，预言的球对称进动约为广义相对论的 $19/21$，小了约每世纪 $4$ 角秒，恰与当时推测的四极矩项相近。这是参数简并在真实实验中的一个例子。

太阳内部的振荡提供了不同的量。太阳表面谱线的 Doppler 变化揭示出许多频率不同的正常模，自转使方位量子数不同的模分裂。用不同深度敏感的模反演角速度剖面，便可检验“内核远比表面转得快，造成很大的重力四极矩”这个解释。对这些资料的分析得到约 $J_2=(2.2\pm0.1)\times10^{-7}$，比上述早期推断小约两个数量级。若以 $2.2\times10^{-7}$ 为标度，水星近日点每世纪进动的近似式可写成 $42.980\,\{(2+2\gamma-\beta)/3 +6.5\times10^{-4} [J_2/(2.2\times10^{-7})]\}$ 角秒。注意四极矩系数位于花括号内，还须乘前面的 $42.980$；日面形状的小不确定量正是这样传到 $\beta$ 的。

其它天体的作用不止是增加水星历表的统计量。火星离太阳较远，轨道要素不同，测得进动对太阳四极矩的相对灵敏度也变了。2011 年一项利用火星轨道器的分析，先采用 Cassini 对 $\gamma$ 的限制，得到 $\beta-1=(0.4\pm2.4)\times10^{-4}$。绕地球的 LAGEOS II 每年约 $3.3$ 角秒的相对论近地点进动也有约百分之二的检验；此处干扰源是地球的多极矩和卫星所受的非引力加速，并非太阳的内部转动。两个精度相近的进动测量若依赖不同场源，其共同约束与各自的误差来源都值得保留。

最后还须说明如何解释“与零相容”。把轨道初值等一组无关参数记为 $q$，待测引力参数记为 $p$。线性化后，若误差权重矩阵为 $W$、两组偏导矩阵为 $H_p,H_q$，允许 $q$ 自由变化时，$p$ 的有效信息矩阵为
\[
 H_p^{\mathsf T}WH_p-
 H_p^{\mathsf T}WH_q(H_q^{\mathsf T}WH_q)^{-1}H_q^{\mathsf T}WH_p .
\]
第二项给出恰能被轨道初值等吸收的那部分灵敏度。若只将已拟合的标准历表减掉，再在固定初值的残差上求 $p$ 的误差，上式第二项就被丢失，实际能分辨的效应会被夸大。独立的日震学数据相当于增加了 $J_2$ 的信息，而非多测了一次水星进动。把 $\gamma$ 固定为 1，则给出以广义相对论光传播为前提的 $\beta$；把 $\gamma$ 连同 Cassini 的测量误差共同估计，所得边界才保留了光行时的不确定量。

\section{强等效原理的检验}

\subsection{地球和月球向太阳下落}\label{sec:pn-sep}

实验室中两块材料的自由落体加速度已经可以比得很准，为何还要测地球和月球？设质量为 $m$、半径为 $R$ 的球自无穷远逐层聚集，每添一层质量 $\d m(r)$ 放出的能量为 $Gm(r)\d m(r)/r$。故其 Newton 引力势能为 $E_{\rm g}=-\int_0^R Gm(r)\d m(r)/r$。均匀球给出 $E_{\rm g}=-3Gm^2/(5R)$，结合能所占的比例随 $Gm/(Rc^2)$ 而增大。一块半径 $10$ 厘米、质量 $1$ 公斤的试验体，这个比例不到 $10^{-26}$；用实际密度分布估算，地球约为 $4.6\times10^{-10}$，月球约为 $0.2\times10^{-10}$。即使两组实验的加速度精度相当，检验引力结合能也应选后者。

若结合能对惯性质量和被动引力质量的贡献不同，地球和月球向太阳的加速度就有差别。对近球形的弱自引力天体，记 $\epsilon_A=E_{{\rm g},A}/(m_Ac^2)<0$，可写 $(m_G/m_I)_A=1+\eta_{\rm N}\epsilon_A$。这里 $m_G$ 指被动引力质量，系数
\[
    \eta_{\rm N}=4\beta-\gamma-3-\tfrac{10}{3}\xi-\alpha_1
                   +\tfrac{2}{3}\alpha_2-\tfrac{2}{3}\zeta_1-\tfrac{1}{3}\zeta_2
\]
称为\textbf{Nordtvedt 参数}。惯性质量可由内能、压强和引力势能各项积分算得，响应外场的被动质量则从合外力中读出，二者取比才有上式的组合。广义相对论给出 $\eta_{\rm N}=0$。若另外假定优选系、优选位置和非守恒参数均为零，就只剩 $4\beta-\gamma-3$。由这类假设推断单一参数的限值时，须把假设一同说明。

取地月质心至太阳的单位矢量为 $\bm N$，日地距离 $R$，太阳在质心处产生的加速度为 $g_\odot\bm N$，其中 $g_\odot=GM_\odot/R^2$。设 $\bm r$ 由地心指向月心，$m=m_\oplus+m_{\rm Moon}$。地月受太阳的加速度相减，把太阳场按 $r/R$ 展开，到第一阶有
\[
    \ddot{\bm r}=-\frac{Gm}{r^3}\bm r
         +\frac{GM_\odot}{R^3}
              [3\bm N(\bm N\cdot\bm r)-\bm r]
         +\delta g_\odot\bm N+\cdots ,\qquad
    \delta=\eta_{\rm N}(\epsilon_{\rm Moon}-\epsilon_\oplus).
\]
为简明，Kepler 项中的 $Gm$ 取由地月轨道测得的值，上式略去 $\delta$ 乘潮汐的小量及其它后 Newton 摄动。右端第二项是寻常潮汐，来自月球比地球更靠近或更远离太阳，正比于 $r$。第三项却在外场均匀时也有，它使整条相对轨道沿太阳方向偏移。由两天体的结合能知 $\delta\simeq4.4\times10^{-10}\eta_{\rm N}$。实验要找的正是太阳加速度中这样小的一个差额。

\subsection{为何在朔望月中出现？}

先让月球在半径 $a$ 的圆轨道上运动，以角频率 $n=\sqrt{Gm/a^3}$ 旋转的径向和切向单位矢量分别记为 $\bm e_r,\bm e_t$。设实际位置为 $(a+x)\bm e_r+y\bm e_t$，其中 $x,y$ 都远小于 $a$。由 $\dot{\bm e}_r=n\bm e_t,\,\dot{\bm e}_t=-n\bm e_r$ 求两次导数，把地月间的引力也展到 $x,y$ 的一次项，得 $\ddot x-2n\dot y-3n^2x=f_r$、$\ddot y+2n\dot x=f_t$。$f_r,f_t$ 是外来摄动力每单位质量在两个方向的分量。这里不能只设一个径向弹簧：切向力会改变角动量，再经 $2n\dot y$ 改变半径。无强迫时方程有频率 $n$ 的自由振荡，这是小离心率的 Kepler 椭圆相对圆轨道的径向起伏。

略去内外轨道的倾角和离心率，设地球绕日角频率为 $n_\odot$，月球与太阳在地心天空中的平黄经差为 $D=(n-n_\odot)t+D_0$。从地球看去，太阳方向与 $\bm e_t$ 的内积为 $-\sin D$，故差分落体力有 $f_r=f\cos D,\,f_t=-f\sin D$，其中 $f=\delta g_\odot$。取受迫解 $x=X\cos D,\,y=Y\sin D$ 并记 $\omega=n-n_\odot$，两方程给出
\begin{align*}
    - (\omega^2+3n^2)X-2n\omega Y&=f,&
    -\omega^2Y-2n\omega X&=-f,\\
    X&=\frac{1+2n/\omega}{n^2-\omega^2}\,f .
\end{align*}
故应找的周期是一个朔望月，约 $29.53$ 日，和月球相对恒星约 $27.32$ 日的轨道周期有别。外轨道慢，强迫频率只比自由振荡频率小 $n_\odot$，分母 $n^2-\omega^2\simeq2nn_\odot$，响应比静态估算的 $f/n^2$ 大得多。用太阳在地球处的 $g_\odot\simeq5.9\times10^{-3}\,\mathrm{m/s^2}$ 代入上式，可得 $X$ 约 $8\eta_{\rm N}$ 米。Will 所用月球轨道模型在计入更高次的太阳摄动后估得距离振幅约为 $9.2\,|\eta_{\rm N}|$ 米；这里的圆轨道估算旨在说明米数和周期的来源。

最大的干扰恰来自太阳。第一阶潮汐力沿半径的投影正比于 $3\cos^2D-1$，沿切向的投影正比于 $\sin 2D$，因而距离有很大的半朔望月项，其振幅达数千公里。记录许多个月便可按频率与 $\cos D$ 项分开。但太阳势的下一阶不再关于地月两端对称，原因是地球和月球到质心的距离相差很大。径向力中出现 $\cos^3D$ 和 $\cos D$，由于 $4\cos^3D=3\cos D+\cos 3D$，此项也含一个朔望月的周期，和 Nordtvedt 项同频。细算此项在距离上的振幅约 $110$ 公里。不能把每一项同频变化都归给新引力，而须由独立测定的质量、距离和轨道频率计算 Newton 潮汐，再从残差中寻找毫米量级的异常。

\subsection{月面反射器与扭秤}

1969 年 Apollo 11 把第一组角反射器留在月面。台站发一束很窄的激光脉冲，角反射器把光沿来路送回，望远镜在约 $2.5$ 秒后收回少数光子。设扣除仪器延迟后的往返钟时为 $T$，最粗略的距离是 $cT/2$，$1$ 厘米距离相当于 $67$ 皮秒往返时间。真正量到的是地表台站到月面反射器的斜距。要算地心到月心的距离，还须知道地球自转使发射站移到了哪里、固体潮怎样改变台站高度，以及月球的天平动怎样改变反射器对地心的距离。地球大气给出的延迟依仰角而变；几个地点、数个反射器及不同月相的观测，有助于把这些影响相互区分。

轨道初值也须跟 $\eta_{\rm N}$ 一起估计。若只测朔日附近，$\cos D$ 总近于 1，常数测距偏差与异常拉伸便很难分开。若每月在不同相位采样，特别是跨过望日前后，二者的作用就不同了。地月潮汐耗散引起的月球长期远离、月球天平动、地月质心绕日运动以及信号本身的相对论时延，均按各自的规律算入理论往返时间。用几十年测距资料的一组分析，得到月距残差中与 $\cos D$ 同相的振幅 $(2.8\pm4.1)\,\mathrm{mm}$。这个量是模型扣除后的振幅拟合值，并非单个光子的测距误差。

还有一个困难。地球比月球富铁，若物质组成也影响落体加速度，单凭测距分不清这种效应与结合能的贡献。可在扭秤两端分别配出模拟地球和月球成分的试样，独立比较它们向太阳的加速度。实验室中的引力结合能极小，因而这里专测组成差异。一组这样的扭秤观测将成分所致差分加速度限制到外场的 $1.4\times10^{-13}$ 量级。将这两种实验合用，Williams 等在 2004 年得到 $\eta_{\rm N}=(4.4\pm4.5)\times10^{-4}$。若有新的长程力且其荷有别于模型采用的成分参数，则也要重新检查模拟试样与真实地月的匹配程度。

这个三体实验还可另换一套天体。毫秒脉冲星 J0337+1715 与一颗白矮星约每 $1.63$ 日互绕一周，外侧另一颗白矮星约每 $327$ 日绕内双星一周。外伴星的作用代替了太阳，中子星和内侧白矮星则代替地球和月球。两内星自身结合能所占比例相差可达十分量级。逐个记录脉冲到达时间，可求内轨道沿视线的位移，看是否随外伴星的方向发生偏移。同样的微小落体差在不同系统里可产生很不相同的位移，须用相应的轨道频率和尺度推算。更要紧的是，中子星内部已超出普通 PPN 假设的弱引力范围，把此类观测直接代入地月的 $\eta_{\rm N}$ 公式并无普适依据。

\subsubsection{潮汐与均匀差分力的谐波}

上面的 $\delta$ 为什么定义为两个质量比之差，而其它系数几乎可按 Newton 值取？设月球为星 1、地球为星 2，记 $u_A=(m_G/m_I)_A$，$\bm r=\bm x_1-\bm x_2$，$\bm X$ 是地月质心的日心位置。两星相对 $\bm X$ 分别位于 $q_2\bm r,-q_1\bm r$，其中 $q_A=m_A/(m_1+m_2)$。若 $\bm g(\bm X)$ 是太阳在质心处的场，则地月相对加速度中的太阳部分可依次写成
\begin{align*}
 \bm a_{\rm ext}={}&(u_1-u_2)\bm g
 +(u_1q_2+u_2q_1)(\bm r\cdot\bm\nabla)\bm g\\
 &+\tfrac12(u_1q_2^2-u_2q_1^2)(\bm r\cdot\bm\nabla)^2\bm g+\cdots .
\end{align*}
第一项在 $u_1=u_2$ 时消失，不用知道内双星的质量分配；第二项在 $u_1=u_2=1$ 时系数等于 1；第三项此时系数为 $(q_2-q_1)/2$。若把地月想象成等质量双星，即使二者绝对间距相同，这项也会消失。这给出一个操作上的区分：均匀差分力来自两星的结合能差，奇次潮汐谐波在 Newton 理论中却由两星相对质心的不对称决定。即使某类理论允许 $\delta\ne0$，在计算大得多的寻常潮汐时，其系数里的微小异常还要再乘 $r/R$，并非每一项都同样重要。

对于点状太阳，前一小节略去的第三项在 $u_1=u_2=1$ 时含沿 $\bm N$ 和 $\bm r$ 的部分。用 $D$ 表示地月向量与朝向太阳的夹角，其径向分量为 $(3/2)(q_2-q_1)g_\odot(r/R)^2(5\cos^3D-3\cos D)$。若换用由太阳指向地月的 $\bm N$，$D$ 与系数的符号须相应换过来。恒等式 $5\cos^3D-3\cos D=(3\cos D+5\cos3D)/4$ 显示，同一个质量差同时产生一倍和三倍朔望角项。因此准确测出三倍频的响应，也有助检验产生同频背景的潮汐模型。相反，只有最先出现的均匀场差时，力的径向、切向投影已恰好是一倍频，不会凭空多出一个独立的三倍频系数。

也可把小量的放大直接看成对月球偏心运动的激发。给一颗圆轨道上的月球在日地方向一个小的持续加速度，在一个月内它时而顺行、时而横向。上一圈留下的离心率矢量转得很慢，下一圈的迫动力与它之间的相位仅因太阳方向移过了一小段而改变，连续数圈的响应可以积累。真正决定结果的是太阳方向的转速与自由径向振荡的失谐。若将太阳方向在每一圈都固定不动，还把自由振荡频率恰取为绕转频率，线性解中会出现随时间增长的项；这不是对真实地月距离的预言。用潮汐修正过的自由运动、实际的朔望频率来求有界周期解，才可与望远镜记录对比。论证所用的相邻落体世界线与潮汐之别，在 \S\ref{sec:tidal} 已用测地偏离方程讨论过。

激光测距还有一种独立的几何误差：反射阵列若略有倾斜，激光脉冲照到阵列两端的光程不同，一簇返回光子在时间上会加宽。月面上的不同反射器朝向和位置各异，使这类误差随观者所见月面姿态改变。1970 年苏联 Lunokhod 1 送到月面的法国反射器，曾因其位置不准而多年难以捕获；2010 年月球轨道器确定月面位置后，重新取得信号，增加了月面一个与原有阵列分开的基准点。激光测距本身变得更精确，并不会自动改正月球的朝向；同一月面上增加地理位置不同的反射器却能把朝向和地月中心距离分开。用这个道理还可理解，研究月球光学天平动和研究 SEP 的两组观测为何会彼此受益。

\subsubsection{星系场与第三颗伴星}

若把一个中子星--白矮星双星整体放在银河的引力场里，外场方向在一圈乃至几年内近乎固定，差分落体将在银河指向上产生受迫离心率。宽轨道毫秒脉冲星多由以往质量转移自旋加速形成，轨道常相当圆，其自然离心率矢量的方向事前并不知道。即使某系统的实测离心率小，也可能是一个较大的受迫分量与自然分量近乎相消。更稳妥的做法是把不同天空位置和轨道朝向的多个双星合在一起。近乎相消对每一个系统均是小概率，因而可在说明取向先验后给出统计上限。对二十余个宽轨道系统的一组分析，给出中子星引力结合能所致的差分质量比绝对值约小于 $4.6\times10^{-3}$。这个上限就相对加速度的数值而言较弱，但其中包含中子星内部的强场效应。

三体脉冲星 J0337+1715 免去了用银河场方向充作近乎不变外场的困难。内轨道由约 $1.44 M_\odot$ 的中子星和 $0.20M_\odot$ 的白矮星组成，外侧白矮星约 $0.41 M_\odot$。2014 年公布的内外轨道离心率分别约为 $6.9\times10^{-4}$ 和 $0.035$，对视线的倾角均接近 $39^\circ$。外伴星在已知的三百多日周期中移过天空，与几乎固定的银河场相比，异常力的相位可由轨道本身定出。实际计时是对中子星位置沿视线的投影敏感，不能把某个物理半径的变化不作投影便一律换成 $\delta r/c$ 的到达时间变化。内星相对质心的位置还须乘另一内星的质量分数。用整套三体方程预报到达时刻便包含这些因子。

在近圆近共面的模型中，取两内星单位静质量结合能相差 $0.1c^2$，并为估量灵敏度试取 $|\eta_{\rm N}|=10^{-4}$，所得内轨道距离极化为数百公里。普通三体潮汐基频和二倍频分量分别可在 $10^4$ 和 $10^5$ 公里量级，自然不能按脉冲计时的钟误差一个数字来估计这项检验。外轨道准确演化的相位、内轨道的强迫和自由偏心、由周期累积的相位漂移都须一起求解。2018 年 Archibald 等报告，在外白矮星的场中，中子星和内白矮星自由落体的相对差异小于 $2.6\times10^{-6}$。这是结合具体系统所得的差分加速度上限；上述 $10^{-4}$ 只是用于演示位移尺度的试取参数，不能把它误读为实测值。

\subsection{搜寻优选方向}

若引力实验依赖实验室相对某一优选系的速度 $\bm w$，双星的轨道也会透露这个速度。譬如 $\alpha_1\ne0$ 时，质量不等的双星即使原在圆轨道上，也能出现受迫离心率，使轨道长轴具有相对于 $\bm w$ 面内投影的优选方向。用指向近心点、大小为离心率的矢量 $\bm e$ 说明尤其方便。一般有 $\bm e(t)=\bm e_{\rm free}(t)+\bm e_{\rm F}$，第一项随近心点长期进动，第二项与 $\bm w$ 的投影有固定关系。某天量到 $e$ 很小，仍不能断言 $\bm e_{\rm F}$ 必小，或许正好与第一项抵消。若两项相消后又过了一段足以让近心点明显进动的时间，固定矢量便再也藏不住。

实际可用脉冲星与白矮星构成的低离心率双星。J1738+0333 的离心率仅约 $3.4\times10^{-7}$，但近心点每年约进动 $1.6^\circ$，跨越十年的观测已可区分离心率的固定和旋转部分。其白矮星伴星还可以用光谱观测，给出独立的视向速度。2012 年 Shao 和 Wex 用这一系统得到 $|\hat\alpha_1|<3.4\times10^{-5}$。帽号提醒读者，该系数控制含中子星的系统与优选系的耦合，其与弱场 PPN 参数的关系要由理论给定。

$\alpha_2\ne0$ 则可使轨道平面绕 $\bm w$ 进动。双星计时精确给出的是投影半长轴 $x=a_{\rm p}\sin i/c$，其中 $i$ 为轨道法线与视线之间的夹角。轨道面的方向改变，即使实际半长轴不变，也给出 $\dot x/x=\cot i\,\dot i$。另一种指向天空固定方向的效应来自 $\xi$：银河系的引力势可影响局域引力的方向依赖，其方向通常与 $\bm w$ 不同。具体作比较时，常取宇宙微波背景各向同性的参考系，太阳系相对它的速度约为 $370\,\mathrm{km/s}$；取指向银心的方向研究 $\xi$。这只是为了把假设变为可检验的角度与相位，宇宙微波背景的存在本身不意味着引力定律选中了这个系。

地面也有次日又可重复的检验。若局域引力强弱与仪器相对银心的取向有关，地球自转就使同一台重力仪扫过许多朝向。设台站铅垂方向为 $\hat{\bm n}(t)$，天空固定的单位矢量为 $\hat{\bm q}$，一个可能的方向项正比于 $(\hat{\bm n}\cdot\hat{\bm q})^2$。内积可写成常数加一个随恒星日转动的正弦，平方后分别有一倍和两倍恒星日频率。理论给出振幅随台站纬度及该方向赤纬的关系。月球和太阳原已产生很强的固体潮，海水涨落还会改变台站附近的质量分布。1976 年 Warburton 和 Goodkind 报告了约十八个月的超导重力仪观测，将附近太平洋的海潮一起拟合，得到 $|\alpha_2|<4\times10^{-4}$、$|\xi|<10^{-3}$。此处须同时比相位和振幅，单从频谱上看到一个接近一日的峰不能辨认其成因。

自转天体亦可接受这种检验。在只保留 $\alpha_2$ 项的近球形模型中，令 $\hat{\bm s}$ 为自转轴，$\omega_{\rm s}$ 是自转角速度，它绕优选系速度方向进动的角速度为 $\bm\Omega_{\rm PF}=-(\alpha_2/2)\omega_{\rm s}(\hat{\bm s}\cdot\bm w)\bm w/c^2$。太阳自转轴与行星总轨道角动量的夹角只有数度，若二者形成时大致平行，则在太阳数十亿年的寿命内不能积累太大进动。不过这个推断须用初始取向。孤立毫秒脉冲星转得快，亦可直接找十多年间脉冲轮廓是否改变：转轴一转，视线便扫过辐射束的另一部分。与双星相比，孤立脉冲星不必扣除伴星引起的测地进动。2013 年对 B1937+21 和 J1744$-$1134 的轮廓分析给出 $|\hat\alpha_2|<1.6\times10^{-9}$、$|\hat\xi|<3.9\times10^{-9}$。这些界限都依赖视线与转轴的几何关系，几颗天空方向不同的脉冲星可减少一次偶然对准的影响。

\subsubsection{优选系怎样进入轨道方程}

把三种扰动并排写在一个方程中，容易看清怎样区分它们。以下只取优选系与优选位置的一阶修正，以 $\bm n=\bm r/r$ 指向星 1，$\Delta=(m_2-m_1)/m$、$\bm v=\dot{\bm r}$，$U_{\rm G}=GM_{\rm G}/R_{\rm G}>0$ 为以集中质量近似银河时的外势。去掉只相当于重新标定两体 $Gm$ 的常数，并先把 $\alpha_3$ 所致自加速放在一旁，外加的相对加速度为
\begin{align*}
 \bm a_{\alpha_1}={}&\frac{Gm}{2c^2r^2}\alpha_1\Delta
       [(\bm w\cdot\bm v)\bm n-(\bm n\cdot\bm v)\bm w] ,\\
 \bm a_{\alpha_2}={}&\frac{Gm}{c^2r^2}\alpha_2
     [\tfrac32(\bm w\cdot\bm n)^2\bm n
                   -(\bm w\cdot\bm n)\bm w],\\
 \bm a_\xi={}&\frac{Gm}{c^2r^2}\xi U_{\rm G}
     [2(\bm n_{\rm G}\cdot\bm n)\bm n_{\rm G}
              -3(\bm n_{\rm G}\cdot\bm n)^2\bm n].
\end{align*}
其中 $\bm n_{\rm G}$ 沿银心方向，在最后一式中把它反号得同一结果。第一式在两质量相等时消失，而后二式不消失。第二、三式的角依赖相同，只是方向和共同因子不同；若这两个天空方向在一个给定轨道上的投影十分接近，限单独一个参数比同时估计二者容易得多。$\bm a_{\alpha_1}$ 中两项的次序应由运动学固定：当 $\bm n\parallel\bm v$ 时，速度的径向部分也可同 $\bm w$ 形成横向加速度。

实际归算以小离心率轨道最为清楚。令 $\hat{\bm h}$ 为轨道法线，$\hat{\bm P}$ 为近心点方向，$\hat{\bm Q}=\hat{\bm h}\times\hat{\bm P}$，例如 $w_P=\bm w\cdot\hat{\bm P}$。把扰动力沿未扰动轨道平均，半长轴首阶无长期增量，离心率在一轨道周期中的改变量为
\[
 \Delta e=-\frac{\pi\alpha_1\Delta}{2c^2}\sqrt{\frac{Gm}{a}}\,w_P
       +\frac{\pi\alpha_2 e}{2c^2}w_Pw_Q
       -\frac{\pi\xi U_{\rm G}e}{c^2}n_{{\rm G}P}n_{{\rm G}Q}
       +O(e^2).
\]
若初始 $e$ 很小，最前一项仍在。这里的 $\hat{\bm P}$ 同时在随正常的近心点进动改变，因而不能把第一项当常数乘几百万圈而得出离心率无限增长。以 $\dot\varpi_0$ 表示正常的平均进动率，受迫离心率的大小在此近似下是 $e_{\rm F}=|\alpha_1\Delta|\,n\sqrt{Gm/a}\,|\bm w_\parallel|/(4c^2|\dot\varpi_0|)$，其中 $\bm w_\parallel$ 为 $\bm w$ 的面内投影，$n$ 是平均运动。可从上式看出分母的来源：在 $\hat{\bm P}$ 转过一角度的周期内，对 $w_P$ 的正弦变化积分便带来 $1/\dot\varpi_0$。更正规的做法是直接积分离心率矢量，避开当 $e\to0$ 时近心点方位的坐标奇性。

对严格的圆轨道，暂忽略 $\alpha_1 e$ 小项，两种方向性力给出的轨道法线每一圈的改变量可写为
\[
 \Delta\hat{\bm h}
  =-\frac{\pi\alpha_2}{c^2}w_h(\bm w\times\hat{\bm h})
    +\frac{2\pi\xi U_{\rm G}}{c^2}n_{{\rm G}h}
                  (\bm n_{\rm G}\times\hat{\bm h}).
\]
无论哪一项，轴与 $\hat{\bm h}$ 平行时叉积为零，轴正好在轨道面内时其法向投影又为零。故最紧的界限并非来自任意一个测得最圆的轨道，轨道方向和视线也很重要。用双星脉冲星的 $x=a_{\rm p}\sin i/c$ 找 $\dot i$ 时，横向自行本身会使观者相对轨道的视线逐年变化，亦产生表观 $\dot x$。已知自行以及升交点在天空中的方向，方可将其扣除。用同一系统的离心率和 $x$，相当于同时观测面内的极化与面外的转动。$\xi$ 虽可出现在全守恒理论中，内双星角动量在银河场中转动并无矛盾：此处内双星本来就不是孤立系统，外轨道可以交换角动量。

$\alpha_3$ 更特殊。只让星 1 有不可忽略的自引力和自转，令 $\epsilon_1=E_{{\rm g},1}/(m_1c^2)<0$、$\bm\omega_1$ 为自转角速度，则随系统速度而变的主项约为 $\bm a_{\rm self}=-(\alpha_3/3)\epsilon_1\,\bm w\times\bm\omega_1$。它并不需要有伴星。取对普通脉冲星适当的 $|\epsilon_1|\sim0.1$、$w\sim300\,\mathrm{km/s}$，利用 $\omega_1=2\pi/P_{\rm s}$，其对自转周期导数的 Doppler 贡献 $P_{\rm s}a_\parallel/c$ 约为 $10^{-4}\alpha_3$ 乘一个与观测方向有关的因子。自转越快时 $a_{\rm self}$ 越大，所取的钟周期却也越短，二者相抵。由普通脉冲星所见很小的 $\dot P_{\rm s}$ 曾得到 $|\hat\alpha_3|\lesssim10^{-10}$ 的数量级限制。

在毫秒脉冲星--白矮星双星内，同一项对两星很不一样，使相对轨道响应近乎固定的差分力 $\bm B$。先对近圆轨道的离心率矢量求导并平均，有 $\langle\dot{\bm e}\rangle_{\bm B}=(3/2)\bm B\times\hat{\bm h}/(na)$。正常近心点进动给 $\dot\varpi_0\hat{\bm h}\times\bm e$。两者合起来给出沿 $\bm B$ 面内投影的受迫分量，大小为 $3|\bm B_\parallel|/(2na|\dot\varpi_0|)$。这与 $\alpha_1$ 产生的受迫离心率具有不同的方向依赖，因为 $\bm B$ 正比于 $\bm w\times\bm\omega_1$。用约二十个宽轨道的毫秒脉冲星系统联合分析，Bell、Damour 和此后的研究把相应的 $|\hat\alpha_3|$ 限制到约 $4\times10^{-20}$。这样的统计界限依赖系统的形成史、未知升交点的方向分布和中子星自引力对系数的改变，不能只读它比前一个数字多了几个零。

\subsubsection{方向、潮汐与自转轴的独立标定}

谈优选系时容易把太阳系在银河内的转速与太阳系相对宇宙微波背景的速度混为一谈。以微波背景偶极给出的 $w_0\simeq368\,\mathrm{km/s}$ 为例，其速度顶点在赤道坐标约为赤经 $168^\circ$、赤纬 $-6.9^\circ$；银心方向却约为赤经 $266^\circ$、赤纬 $-29^\circ$。前者用来投影 $\alpha_i$ 项，后者用来投影 $\xi$ 项。对太阳系绕银河公转，可由圆轨道速度 $v_{\rm G}\simeq230\,\mathrm{km/s}$ 估得 $U_{\rm G}/c^2\sim v_{\rm G}^2/c^2\sim5\times10^{-7}$。银河有延伸很远的暗晕，不能把所有质量当成银心一点，再把整个势的数值都用作可量的方向性。

具体说，影响仪器方向的量是外部质量分布的张量势 $U_{ij}=G\int\rho' s_i s_j\,\d[3]{x'}/s^3$ 的各向异性部分；$C\delta_{ij}$ 与仪器转角无关。由先前的超势恒等式 $U_{ij}=U\delta_{ij}-X_{,ij}$，若暗晕近似球对称，$\bm N$ 为从银心指向太阳的方向，则 $X_{,ij}=(X''-X'/R)N_iN_j+(X'/R)\delta_{ij}$。因此留下来的只是径向导数与切向导数之差。用内核密度近乎常数、外区 $\rho\propto R^{-2}$ 的球形模型，大半径处旋转曲线近于平坦，方向项的系数仍为若干 $10^{-7}$，和集中质量估算在同一量级，但不要求每一圈外层壳都像中心点质量那样贡献。最合宜的做法是带着银河势的不确定度估 $\xi$，而 $\bm w_0$ 的方向和速度可由微波背景更精确地直接给定。

用地球作实验场源时，还要问“引力常数的方向性”怎样进入重力仪。地面仪器通常比较的是自由试块与地表支承的位置，读数约为 $g=G_{\rm L}m/R_\oplus^2$。若场改变了固体地球的平衡形状，$R_\oplus$、附近地层的质量分布及台站的位置也随之改变。寻常天体潮汐既有直接的引潮加速度，又有地球弹性形变的间接改正。对地面重力仪，合起来的固体潮因子约为 $1.16$，海潮质量在站点附近的运动还须另行计算。方向性 $G_{\rm L}$ 的等效体积力取决于内部引力及密度，未必在整个地球内都与太阳或月球的四极引潮势同形。因此不能假定已知 Newton 潮的仪器响应，便连任何新作用的弹性响应也已准确确定。

为看清哪些周期会出现，取台站纬度为 $L$，地球自转频率为 $\Omega_\oplus$，给定天空方向的赤纬为 $\delta$，其在台站的时角记为 $H=\Omega_\oplus t+H_0$。台站铅垂 $\hat{\bm n}$ 与天空方向的内积是 $A+B\cos H$，其中 $A=\sin L\sin\delta$，$B=\cos L\cos\delta$。平方后，一倍频振幅为 $2AB$，两倍频为 $B^2/2$。在赤道台站，一倍频项消失而两倍频仍在；在地理极点，铅垂指向恒定，两个调制项都消失。$H_0$ 对银心和速度顶点又相差约百角度。同一重力仪把资料按两种方向各自的恒星时折叠，或者合用不同纬度的台站，才可能辨出原本在频率上完全重合的两个参数。一个通用的“日振幅上限”并不足以完成这种区分。

地球相对优选系的速度还含公转速度 $\bm v_\oplus$，因而 $w^2=w_0^2+2\bm w_0\cdot\bm v_\oplus+v_\oplus^2$。常数 $w_0^2$ 可吸收到局域 $G$ 的平均标定里，交叉项却有年周期，尺度为 $w_0v_\oplus/c^2\sim10^{-7}$；各向异性的 $(\bm w\cdot\hat{\bm n})^2$ 还会有年频率与恒星日谐波之和、差频率。$G_{\rm L}$ 有年周期也可使地球半径和转动惯量小幅改变，由角动量守恒引起地球自转速率的年变化。对 $\xi$ 则应将各外部质量产生的张量势相加。太阳贡献的势约 $10^{-8}c^2$，比银河小，但太阳方向在一年中扫过天空，给出另一套相位。正常的大气角动量、海潮及地下水质量也随季节变化，地面上各站的环境资料因而是引力检验的一部分。

对长时稳定的自转轴，可由力矩而非下落加速度寻找同类效应。取自引力天体质量四极结构近乎轴对称，转动惯量为 $I$，自转为 $\bm\omega_{\rm s}=\omega_{\rm s}\hat{\bm s}$。若外部方向引起内部各部分引力的微小各向异性，求绕质心力矩时，各向同性压强不贡献；内部动能张量、引力能张量与压强满足的张量维里关系可将结果换成 $I\omega_{\rm s}^2$ 乘方向因子。由 $\alpha_2$ 导致的转轴改变量，在时间 $T$ 内约为 $|\delta\hat{\bm s}|\simeq |\alpha_2|\omega_{\rm s}T(w/c)^2|\sin\beta_w\cos\beta_w|/2$，$\beta_w$ 为速度顶点与自转轴夹角。代入太阳的 $\omega_{\rm s}\simeq2.8\times10^{-6}\,\mathrm{s}^{-1}$，用约 $4\times10^9$ 年积累并要求转轴偏离不得大于观测到的数度，可估得 $|\alpha_2|\lesssim2\times10^{-6}$。若初始倾角本就很大，或演化中存在足够的普通力矩，这一估计便需重做。脉冲轮廓法直接看十余年中的变化而无需倒推十亿年前的初态，但它又须用辐射束的几何模型。两种限制互补的理由就在各自所用的先验不同。

\begin{table}[htbp]\centering\small
\caption{优选方向的几种历史检验。帽号表示含中子星的有效系数。}
\begin{tabular}{@{}lll@{}}\hline
观测与时期 & 约束 & 提供方向信息的量\\ \hline
地月测距（1996--2008） & $\alpha_1=(-0.7\pm1.8)\times10^{-4}$ & 距离与轨道偏心\\
同一类地月资料 & $\alpha_2=(1.8\pm5.0)\times10^{-5}$ & 周期性轨道扰动\\
双星计时（2012--2013） & $|\hat\alpha_2|<1.8\times10^{-4}$ & 轨道面的投影\\
同一组双星 & $|\hat\xi|<3.1\times10^{-4}$ & 轨道面与银心方向\\
孤立脉冲星（2013） & $|\hat\xi|<3.9\times10^{-9}$ & 轮廓与自转进动\\ \hline
\end{tabular}
\end{table}

\subsection{引力常数会改变吗？}

局域 Lorentz 和位置不变性所讲的是非引力实验，SEP 还要求实验室里做 Cavendish 实验也不能测出所处时代。若引力理论含一个随宇宙演化的标量，场的局域边界值会慢慢改变，在实验室以固定的质量、长度、时间标准所量的 $G$ 便可能随之改变。一次实验只能测出某时的值，两次测量之间的漂移却很难靠实验室内的绝对标定分出。行星轨道逐周积累相位，适于找这种缓慢变化。

行星测距实际确定的是 $\mu=GM_\odot$。若 $\mu$ 的变化远慢于公转，近圆轨道的单位质量角动量守恒给出 $r\mu\simeq\text{常数}$，再由 Kepler 定律得 $\dot P/P\simeq-2\dot\mu/\mu$，其中 $P$ 是公转周期。也可从动量矩 $h^2=\mu a(1-e^2)$ 看出，对缓变的各向同性质量参数，在每个周期平均，离心率在此阶不变而 $\dot a/a=-\dot\mu/\mu$。由平均角频率 $n=2\pi/P$ 积分即得一个随时间平方积累的轨道相位差 $\delta\phi\simeq n(\dot\mu/\mu)(t-t_0)^2$。改动初始相位给出常数，改动初始平均运动给出线性项，因而观测跨度本身就帮助分开这些量。

但 $\dot\mu/\mu=\dot G/G+\dot M_\odot/M_\odot$。到了每年 $10^{-13}$ 的精度，必须计入太阳辐射和太阳风带走的质量。对行星历表还要估计小行星质量，否则未计的长期摄动也能被吸入 $\dot\mu$。Konopliv 等在 2011 年用火星轨道器资料得到 $\dot G/G=(0.1\pm1.6)\times10^{-13}\,\mathrm{yr}^{-1}$；Hofmann 等在 2010 年用地月测距得到 $(0.7\pm3.8)\times10^{-13}\,\mathrm{yr}^{-1}$。地月另有较大的潮汐演化，两项研究对质量损失和潮汐耗散的依赖并不相同。

还能回溯到更早的时代。日震学通过今天太阳内部的声速分布间接探查长期演化，原初核合成则借轻元素丰度限制早期宇宙的膨胀率，后者依赖当时的引力强度。将这些结果表示成今天的 $\dot G/G$ 须假定几十亿年来 $G(t)$ 以何种方式改变，与行星近几十年的直接漂移有别。即使是脉冲星计时，也不宜简单套用行星公式。中子星的结合能占比很大，若 $G$ 改变，星的平衡结构、质量及转动惯量都会改变，从而同时影响自转和公转；须先指定引力理论与物质状态方程，才能用观测解释 $\dot G$。

把不同时间尺度的实验放在一起，还可看清 $\dot G$ 所代表的假设。若固有的重力尺度随宇宙演化变化，人们常取 $\dot G/G=\sigma H_0$ 作数量级比较；$H_0$ 的倒数给出宇宙演化时间，$\sigma$ 的符号和大小却要从具体理论求得。以便于比较的 $H_0=70\,\mathrm{km\,s^{-1}\,Mpc^{-1}}$ 为例，即约每年 $7\times10^{-11}$，局域漂移若限制到每年 $10^{-13}$，限制的已是这个尺度的千分量级。在空间平直、以冷物质主导的 Brans--Dicke 宇宙解，$G$ 可随幂律时间变化并给 $\sigma\simeq-1/(2+\omega_{\rm BD})$；有优选系的理论却也可以 $G$ 为常数。例如单位类时矢量若只以其固定模进入局域 Newton 常数，就不能单从有优选方向推出 $\dot G\ne0$。

关于太阳质量的扣除也可定量估算。太阳光度 $L_\odot$ 所代表的静质量损失率是 $-L_\odot/c^2$，尚有太阳中微子以及逸出的太阳风粒子。以前者和已知太阳质量代入，光子和中微子的贡献约为每年 $-0.7\times10^{-13}$ 的相对质量变化；太阳风还可给同一量级的负贡献，其值随时间和活动状态起伏。故历表若精确报告 $GM_\odot$ 的小幅下降，方向恰与已知的质能流失相同，必须先以太阳资料限 $M_\odot(t)$，才能判断 $G$ 是否也有变化。相反，以同一 $G(t)$ 同时拟合行星和地月运动，可借两种中心质量的不同演化有所分辨，但月球加速中的潮汐模型会随同进入误差预算。

\begin{table}[htbp]\centering\small
\caption{不同来源的历史漂移估计。第二列单位为每年 $10^{-13}$。}
\begin{tabular}{@{}lll@{}}\hline
资料与分析年代 & $\dot G/G$ & 换算所需的物理量\\ \hline
火星历表（2011） & $0.1\pm1.6$ & $GM_\odot$ 及其时间变化\\
地月测距（2004） & $4\pm9$ & 地月潮汐摄动\\
脉冲星（2008--2009） & $-7\pm33$ & 星体质量、转动惯量的演化\\
日震学（1998） & $0\pm16$ & 太阳演化模型\\
原初核合成（2004--2005） & $0\pm4$ & $G$ 自早期至今的变化规律\\ \hline
\end{tabular}
\end{table}

最后一行尤其不可直接当作一次本地钟的计时。轻元素在宇宙早期约几分钟内形成，给出的本来是当时与今日的引力强度之比。以可容许的差异在约百分之二十内为例，可以设 $G(t)=G(t_0)(t/t_0)^{-\alpha}$，于是今天 $\dot G/G=-\alpha/t_0$，而远古的比值同时依赖很大的 $\ln(t_0/t)$。若只假定过去短时间内线性变化，却把这个线性式直接外推一百多亿年，统计上相似的“漂移上限”会对应不同历史。表中前三项主要依赖现代运动资料，后两项还依赖从早期演化到目前状态的模型，不能不加条件地加权平均。

\section{后 Newton 引力的其它检验}

\subsection{陀螺的两种进动}\label{sec:pn-gyro-observe}

自由下落的小陀螺不受外力矩时，其自转矢量沿测地线平行移动。若它绕静止球对称天体运动，对照远处恒星，转轴一般每公转一周便略微改变方向，这称为\textbf{测地进动}(geodetic precession)。若中心天体也在转，它的质量流还要通过度规的 $g_{0i}$ 拖动陀螺，这就是\textbf{Lense--Thirring 进动}。前者取决于陀螺的轨道角动量，后者还须有中心天体的自转角动量。这里的自转指有体积物体的转动。即使陀螺和中心星都极接近球形，两种效应仍可存在。地球自转轴在日月引力下的寻常岁差则需要地球的赤道隆起，属于 Newton 四极矩受外力矩的效应。

具体令 $M,\bm J$ 为中心天体质量和自转角动量，$\bm r,\bm v$ 为陀螺的相对位置、速度，$\hat{\bm r}=\bm r/r$。在渐近惯性坐标中以静止恒星为参照，陀螺的物理自转矢量 $\bm S$ 的演化可写成 $\dot{\bm S}=\bm\Omega\times\bm S$。保留质量单极和自转偶极引起的进动，略去优选系速度所致的附加项，得
\begin{align*}
    \bm\Omega_{\rm G}&=\left(\gamma+\tfrac12\right)
          \frac{GM}{c^2r^3}\,\bm r\times\bm v,\\
    \bm\Omega_{\rm LT}&=K\frac{G}{c^2r^3}
            [3\hat{\bm r}(\bm J\cdot\hat{\bm r})-\bm J],
       \qquad K=\frac{4+4\gamma+\alpha_1}{8}.
\end{align*}
第一式既含单位质量产生的空间曲率，也含沿弯曲轨道依次变换局域静止系的效应。第二式有磁偶极场的角分布：到中心相距同样远，沿自转轴与在赤道面上时，进动方向相反，而前者的角速度大小是后者的两倍。广义相对论给出 $K=1$。在已独立测过 $\gamma$ 的条件下，观测拖曳还可检验另一种度规分量所含的参数组合。

实验常追踪转轴许多周，更便于使用平均角速度。记轨道法线单位矢量为 $\hat{\bm h}$，半长轴 $a$、离心率 $e$ 和平均角频率 $n$ 满足 $n^2a^3=GM$。角动量守恒给出 $r^2\d f/\d t=h$，其中 $f$ 为真近点角、$h=\sqrt{GMa(1-e^2)}$。用 $r=a(1-e^2)/(1+e\cos f)$ 对一周平均，便有
\[
    \left\langle r^{-3}\right\rangle=\frac{1}{a^3(1-e^2)^{3/2}},\qquad
    \left\langle\frac{\hat r_i\hat r_j}{r^3}\right\rangle
      =\frac{\delta_{ij}-\hat h_i\hat h_j}{2a^3(1-e^2)^{3/2}} .
\]
后一等式可在以近心点和轨道法线为轴的坐标系中验证：两个面内对角元相等，交叉项的积分为零。故平均的测地进动角速度沿 $\hat{\bm h}$，大小为 $(\gamma+1/2)nGM/[c^2a(1-e^2)]$；每公转一周的转角即 $\pi(2\gamma+1)GM/[c^2a(1-e^2)]$。相同轨道上中心自转所致的平均进动为
\[
    \langle\bm\Omega_{\rm LT}\rangle
       =\frac{KG}{2c^2a^3(1-e^2)^{3/2}}
             [\bm J-3\hat{\bm h}(\bm J\cdot\hat{\bm h})] .
\]
在赤道轨道，$\bm J$ 与轨道法线平行，故两种平均进动共轴，单从转轴方向的变化只量到二者之和。在极轨道，$\bm J\cdot\hat{\bm h}=0$，二者的轴垂直。若陀螺初始转轴在轨道面内，测地进动便使其在面内转动，拖曳造成的平均偏转则垂直于轨道面。此举减小了拖曳信号，却将两个待测量分到了互相垂直的方向。

\subsubsection{质量流怎样拖动一个陀螺}

由 \S\ref{sec:pn-source-moments} 的定常流偶极，在忽略外加优选系速度时，标准 PPN 度规给出 $\bm g_0=(g_{01},g_{02},g_{03})=-2KG(\bm J\times\bm r)/(c^3 r^3)$。这个场随赤道半径和方向改变，即使用匀速旋转的坐标消去一点的表观转动，也不能在整个外部区域同时消去。

从这里既可看轨道也可看陀螺。慢粒子拉氏量中与质量流有关的项为 $mc\,\bm g_0\cdot\bm v$，于是平稳场给粒子的附加加速度是 $c\,\bm v\times(\bm\nabla\times\bm g_0)$。固定能量时把轨道运行方向反过来，这项反号，质量单极矩带来的 Newton 加速度则同向。另一方面，将自转矢量沿轨道平行移动，质量流给出的角速度为 $\bm\Omega_{\rm LT}=-(c/2)\bm\nabla\times\bm g_0$。代入偶极矢量势即得上文的 $[3\hat{\bm r}(\hat{\bm r}\cdot\bm J)-\bm J]/r^3$。同一个旋度关联轨道加速度和自转轴进动，二者却不测量同一个角。特别是陀螺本身的慢速运动另有测地进动，仅用这个旋度不能算出总读数。

另一种比较是让两束光沿同一闭合路径相向传播。对静止的仪器和时间无关的度规，$g_{00}c^2\,\d t^2+2g_{0i}c\,\d t\,\d x^i+g_{ij}\d x^i\d x^j=0$。两束光的几何路径若可视为相同，反向后二次方程根号下的部分不变，线性项反号，两方向的环行坐标时差为 $\delta t=-(2/c)\oint (g_{0i}/g_{00})\d x^i$。将其换成发射点一只钟的固有时，再与旋转仪器自身即有的 Sagnac 时差区分，才能讨论由中心天体引起的拖曳。这里出现环路积分，与 \S\ref{sec:sync-clock} 中的对钟条件相呼应。\hyperlink{subsection.2.5.4}{\S\,2.5.4} 的局域非转动观者则以轴对称时空的角动量来定义；某个坐标角速度、光环行的时差和实验陀螺的自转角属于三种可区分的量。

\subsubsection{质量可比时的自旋与轨道}

陀螺远轻于地球，故地球的运动只参与很小的反冲。双星两星质量可比时，不能再把一颗星的质量设成零，却仍可假定两组自旋的变化远慢于一个轨道周期。令 $\bm S_A$ 为星 $A$ 的自旋，优选系速度项暂予略去。对星 1，一周中由轨道角动量引起的自旋变化为
\[
 \Delta\bm S_1\big|_{\rm G}=
    \pi\left(2\gamma+1+\frac{m_1}{m}\right)
      \frac{Gm_2}{c^2p}\,\hat{\bm h}\times\bm S_1 .
\]
星 2 的式子只需交换指标。这比单陀螺多出的 $m_1/m$ 来自两处：接近星 1 的度规由运动的星 2 产生，而定义星 1 的自旋又须使用它自己的瞬时静系。取 $m_1/m\to0$，回到 $\pi(2\gamma+1)GM/(c^2p)$。若已由其它计时效应估出双星的质量，上式就不必再用自转轴数据调质量。例如就 B1913+16 而言，广义相对论所给的自转进动周期约数百年。视线扫过磁极附近的不同纬度，脉冲轮廓的峰间距和强度便会在十几二十年间有系统地改变。仅凭轮廓的变化速度，仍须靠辐射束形状估真实转角；双脉冲星的磁层掩食另给一个几何约束。

在同一近似下，轨道面也受自旋的反作用。令 $\bm S=\bm S_1+\bm S_2$、$\bm\sigma=(m_2/m_1)\bm S_1+(m_1/m_2)\bm S_2$，置 $\bm C=(2\gamma+1)\bm\sigma+4K\bm S$。仅取对自旋成一次的摄动，半长轴和离心率在一周的长期变化为零，而轨道法线和近心点有
\[
 \Delta\hat{\bm h}=\frac{\pi G}{c^2hp}\,\bm C\times\hat{\bm h},
 \qquad \Delta\varpi=-\frac{2\pi G}{c^2hp}\,\bm C\cdot\hat{\bm h},
 \qquad h=\sqrt{Gmp}.
\]
这里 $\Delta\varpi=\Delta\omega+\cos i\,\Delta\Omega$ 为椭圆在自身轨道面内的转角。若 $\bm C$ 恰与轨道法线同向，椭圆可在固定平面内转；若 $\bm C$ 在轨道面内，直接留下的是法线的岁差。对星体内部“自转中心”选不同的代表世界线，会改变瞬时自旋--轨道加速度的若干项，在一周平均的上述组合中这项自由度消去了。取 $\bm S_1\to0$ 且 $m_1/m_2\to0$，即 $\bm C\to4K\bm J$，便得到下文卫星的拖曳。若两自旋本身也在长期转动，$\bm C$ 不再永远固定，这个平均法只说得清每一圈的变化，应同时积分自旋进动方程。

两颗星都自转时还有自旋乘积项，它不会在仅保留中心自转和无自旋测试粒子的轨道公式中出现。在上述阶次，其方向依赖可由相互作用能
\[
 H_{\rm SS}=\frac{KG}{c^2r^3}
   [\,3(\bm S_1\cdot\hat{\bm r})(\bm S_2\cdot\hat{\bm r})
                              -\bm S_1\cdot\bm S_2\,]
\]
看出。对 $\bm S_1$ 求偏导再叉乘 $\bm S_1$，正给由 $\bm S_2$ 所致的拖曳自旋进动；对位置求负梯度，则给轨道一个按 $r^{-4}$ 变化的相互作用力。作轨道平均，用此前的 $\langle \hat r_i\hat r_j/r^3\rangle$ 即得平均能量正比于 $\bm S_1\cdot\bm S_2-3(\bm S_1\cdot\hat{\bm h})(\bm S_2\cdot\hat{\bm h})$。若两星均缓慢自转，它远小于轨道运动本身造成的测地进动；对两颗快速旋转的紧致天体，应同时计入自旋所诱导的单体四极矩。中心质量不规则形变的 Newton 四极矩、自旋乘积耦合、线性质量流效应是三个不同阶次和来源的修正，不宜都吸收到一个拖曳系数中。

\subsection{让转轴指向远处恒星}\label{sec:pn-gyro-guide}

Gravity Probe B 采用的正是这个安排。2004 年升空的卫星带四只球形石英陀螺，在离地面约 $642$ 公里的近极地圆轨道上运动。转子的直径约 $4$ 厘米，表面覆以铌膜，在 $1.8\,\mathrm K$ 左右处于超导态。旋转的超导体产生沿自转轴的 London 磁矩，只要用周围的超导量子干涉仪测磁矩的指向，就无须在转子上接触一个机械指针。外壳定期绕望远镜视线旋转，许多固连在外壳上的微小力矩在一转中可以抵消。实际上，若转子与外壳表面存在分布不匀的电势，即使两者的总电荷为零也会有静电力矩。转子的惯性主轴又未必与自转轴完全重合，表面电势图样相对外壳缓慢摆动，使这种误差有自己的时间规律。数据归算时须由四只陀螺各自的摆动和卫星转动记录估计其影响。

如何把卫星里的一个磁矩同“远处恒星”相比？星上望远镜指向飞马座 IM 星，陀螺自转轴初始近乎平行于望远镜的光轴。此星既能用光学望远镜跟踪，也有足够强的射电辐射可用地面 VLBI 测出它相对河外类星体的自行。这个改正不能忽略，导星的自行约为每年 $34$ 毫角秒，恰与拖曳效应的量级相同。另外，地球绕日运动也使陀螺有每年约 $20$ 毫角秒的测地进动，地球四极矩产生的进动也须计入。借 VLBI 测出的导星运动，将陀螺相对望远镜的偏转换到类星体参考系，才可与轨道模型预言的漂移相比。

Everitt 等在 2011 年公布的南北、东西两个方向年漂移量的大小分别为 $6602\pm18$ 和 $37.2\pm7.2$ 毫角秒，广义相对论预言各为 $6606$ 和 $39.2$ 毫角秒。为何是同一组陀螺、同时读数，却一个可测到千分之几而另一个只有约五分之一的相对精度？两者的绝对误差同在每年几个到十几个毫角秒，但距离地球这么远时，由地球自转产生的信号已经比绕地公转产生的信号小两个数量级。

用仪器的原始记录再看一次。SQUID 实际输出的是电信号，须标定电信号对一弧度相对偏转的灵敏度。可用地球公转和卫星绕地运动带来的恒星光行差作已知的角度输入。星上望远镜随视位置略微转动，旋转中的陀螺不会跟着光行差转，故望远镜与转轴之间有可预报的周年和轨道周期调制。地球公转的贡献可达约二十角秒，比每年的拖曳漂移大数百倍。正因这个角度大而相位能由轨道给定，可用它标定电子增益，再找扣去周期信号后的微小长期漂移。若把一个未知的增益与一个未知的小漂移同时对一条近似直线拟合，二者几乎不可分辨。此处用光行差解决的，正是仪器刻度和物理信号乘积的简并。

理想转球上只有绕自转轴产生的 London 磁矩，实际铌膜还俘获少量随转子表面转动的磁通，约为 London 信号的百分之一。若刚体惯性主轴与角动量有一小夹角，惯性主轴相对自转的摆动便会通过这部分磁通进入读数。四个转子的摆动周期都是数小时量级，却不相同；受微弱耗散影响，周期和角度也随时间变。残余磁通因而有两面作用。它使磁场读数不能简单当成理想自转轴，但同时给出了转子表面朝向的连续标记，可以由调制重建这段摆动，再估朝向相关的静电力矩。若一概当作无结构噪声舍掉，就失去这一项独立标定。

转子的表面电势斑块与外壳的电势斑块相互作用，使力矩随外壳相对转轴的方向改变。卫星约每 $78$ 秒绕望远镜光轴转一圈，是为了平均掉随壳转动的干扰。若干扰的某个谐波同外壳转动接近共振，平均法却会失效：随着转子的摆动周期变短，穿过一个共振条件的几天内，转轴可突然出现一段额外漂移。它与恒定的相对论进动不同，在磁通记录可推算的特定相位发生，而且四只转子并不同步。用独立的姿态和磁通资料估这类信号，再将各转子的南北、东西漂移分别拟合，四者得到相容的拖曳值；发生共振次数较多的转子给出的误差也较大。把任务划为几个时间段另行归算，是检验误差模型能否在拟合段之外给出正确漂移的办法。

恢复单位后，地球角动量可写成 $J_\oplus=I_\oplus\omega_\oplus$，其中 $I_\oplus\simeq0.3307M_\oplus R_\oplus^2$。这一数值来自地轴的日月岁差与地球引力多极矩的组合，并非由 GP-B 本身反推出的归一化。任务实际选择的导星使初始自转轴与地轴的夹角约 $73^\circ$，而不是便于估算的直角，故极轨道的平均拖曳还须投影一个 $\sin73^\circ$。在 $r$ 稍大于 $7000\,\mathrm{km}$ 的轨道上，以球形地球算测地进动约每年 $6.6$ 角秒，再计轨道、太阳外场与地球形状等项后，才是实测应比的两维向量。将平行移动所定的进动与日常机械陀螺仪的力矩偏转分开，依靠的是这些互相独立的标定和方向信息。

\subsection{让轨道平面作陀螺}

换一个尺度，地月系统本身也可充作陀螺。把月球绕地球的轨道角动量当作它的“自转”，用 $M=M_\odot$ 及地月质心的日心轨道要素代入测地进动公式，便得月球轨道面绕日心轨道法线的缓慢进动。这也称\textbf{de Sitter 进动}。取 $a$ 为一个天文单位、$e\simeq0.017$，广义相对论给出的每年进动角为 $3\pi GM_\odot/[c^2a(1-e^2)]\simeq19.2$ 毫角秒，亦即每世纪约 $1.92$ 角秒。太阳对月球轨道的 Newton 潮汐本来就使交点显著退行，故应从完整的月球历表中估计额外的进动。Williams 等在 2004 年报告的地月激光测距检验精度约为 $0.6\%$。它与 Nordtvedt 检验共用月面反射器，却测量不同的量：一个是轨道面的长期转动，一个是在朔望角的固定相位上来回伸缩的距离。

若要单独看地球自转对卫星\emph{轨道面}的拖曳，则不必把一个转子带上天，只须不断测它的轨道。以地球赤道为参考平面，令 $i$ 为轨道倾角，$\Omega$ 为升交点经度，$\omega$ 为近地点幅角。地球自转角动量 $J_\oplus$ 对卫星轨道的作用在一周期平均后给出
\[
    \dot\Omega_{\rm LT}=\frac{2KGJ_\oplus}{c^2a^3(1-e^2)^{3/2}},\qquad
    \dot\omega_{\rm LT}=-3\cos i\,\dot\Omega_{\rm LT}.
\]
两式作用的是卫星环绕地球的角动量矢量与轨道椭圆，而上一小节两式作用的是卫星内的陀螺轴。设地球引力四极矩为 $J_2$，赤道半径为 $R_\oplus$。同一升交点在 Newton 力学中已有 $\dot\Omega_{J_2}=-(3/2)nJ_2(R_\oplus/a)^2\cos i/(1-e^2)^2$ 的进动，它通常远大于拖曳。幸而拖曳项不依赖倾角，四极矩项正比于 $\cos i$。若两颗卫星 $a,e$ 相同而 $i_1+i_2=\pi$，直接将两升交点的进动率相加，四极矩影响就消掉，拖曳却加倍。

LAGEOS 卫星是覆有角反射器的致密球体，阻力和辐射压的单位质量效应较小。地面以激光测得各台站至卫星的距离，再由这些距离反求几圈乃至几年的轨道。LAGEOS I、II 的轨道半长轴均约 $1.22\times10^4\,\mathrm{km}$，倾角却分别为 $109.84^\circ$ 和 $52.64^\circ$，并不严格互补。按各自轨道算出四极矩系数 $C_i=\partial\dot\Omega_i/\partial J_2$，便可用 $\dot\Omega_1-(C_1/C_2)\dot\Omega_2$ 消除这一项。LAGEOS I 的拖曳约为每年 $30.7$ 毫角秒；要用这种方法识别它，还须知道地球更高的偶数阶质量矩及其随时间的变化。GRACE 的双星测距提供地球引力场的独立模型，LARES 的第三条轨道则多给一个系数组合。2016 年的三颗卫星联合分析报告了约 $5\%$ 的检验精度，其中对地球质量多极矩误差的估计影响很大。

同一种平行移动还能在双星中转得更快。双脉冲星 J0737$-$3039 的轨道几乎侧向地球，A 星的脉冲经过 B 星磁层时会被周期性遮蔽。B 星自转轴的取向不同，磁层穿过视线的位置也不同。比较多年的掩食光变，结合 B 星的自转调制，就可估计 B 星自转轴绕轨道角动量的进动率。Breton 等在 2008 年用磁层模型得到每年 $4.7\pm0.7$ 度，双星质量、轨道周期和离心率在广义相对论中预言约 $5.07$ 度。此处两颗星质量相当，需用二体自转进动公式；把上文单个小陀螺的式子直接套给一颗中子星并不合适。

\subsubsection{多颗卫星怎样消去四极矩？}

用实际轨道可以将前述消元写成一次数值计算。取两颗 LAGEOS 的半长轴为 $12257$、$12168\,\mathrm{km}$，离心率为 $0.0045$、$0.0135$。Newton 四极矩项在升交点进动中的系数 $C_i$ 正比于 $-a_i^{-7/2}(1-e_i^2)^{-2}\cos i_i$。因第一颗逆行、第二颗顺行，四极矩使两颗卫星的升交点向相反方向走。由所给轨道可算出 $-C_1/C_2\simeq0.545$，故可测的组合应取 $\dot\Omega_1+0.545\dot\Omega_2$，并非两交点之差。对广义相对论，在只保留自转偶极时，这两颗星的拖曳率各为每年约 $30.7$、$31.4$ 毫角秒，组合给每年约 $47.8$ 毫角秒。相比之下，第一颗卫星由地球四极矩产生的交点进动约每年 $125^\circ$，两者相差七个数量级。角度单位没有写清，便无法从数值的大小判断这项消元需要何等精度。

更精确地说，观测不是在两条已知的完美椭圆上逐日读出 $\Omega$。用地面激光往返时间拟合每天或若干天的短轨道弧段，估出各弧的升交点以及可能的一些经验加速度，再将长时间内的升交点残差与地球引力模型作比。设余下的长期线性趋势为 $b_i$，$L_i$ 为理论拖曳率，$C_{ij}=\partial\dot\Omega_i/\partial J_j$，最简的模型是 $b_i=L_i\lambda+C_{i2}\delta J_2+C_{i4}\delta J_4+\cdots$，其中 $\lambda=1$ 表示广义相对论预言成立。绝不能由于对某一条轨道拟合得很好，就断言 $\lambda$ 自动已定；沿单轨道改变地球的偶极拖曳与改变高阶质量矩都可以调整轨道面。如果定轨时用自由的跨轨推力吸收所有长期残差，$\lambda$ 的一部分也会同时被吸收，须把待测参数和这类经验项的相关性算进去。

LARES 的轨道倾角约 $69.5^\circ$，很接近 LAGEOS I 倾角的补角，但半长轴约 $7821\,\mathrm{km}$，远小于后二者，所以不能仅因倾角接近就说四极矩已相消。第三颗卫星使独立交点观测从两个增为三个。以 $\bm C_2=(C_{12},C_{22},C_{32})$、$\bm C_4=(C_{14},C_{24},C_{34})$ 记两个系数列，在两列不平行时，可取 $\bm k$ 正比于 $\bm C_2\times\bm C_4$。于是 $\bm k\cdot\bm b/(\bm k\cdot\bm L)$ 消去 $\delta J_2$ 与 $\delta J_4$，给出 $\lambda$ 的估计。前提是 $\bm k\cdot\bm L$ 不太小；若它也接近零，恰在消去 Newton 项时连拖曳信号也大部抵消，误差就要放大。这是选择不同轨道半径和倾角的另一个理由。更高质量矩仍有贡献，特别是低轨的 LARES 对高阶多极矩更灵敏，同时也较易受稀薄大气等非引力加速度影响。

质量相同而形状不同的卫星受辐射压和阻力很不相同。LAGEOS 选球形、相对小的面积质量比、被动的角反射器，目的就是少给未知面力留下机会。LAGEOS II 的离心率虽小但非零，因此还可观测近地点；过去有人尝试以其近地点进动率作第三个量，以便只有两颗卫星也可消除 $J_4$。然而近圆轨道近地点的误差随 $1/e$ 放大，日照、地球反照以及卫星表面不同时刻重辐射形成的微弱热推力又会改变近地点。与直接在两交点解中多添一个近地点读数相比，GRACE 测量两星相对距离变化建立的地球重力模型，是另一种来源的先验。2011 年公布的 LAGEOS 双星分析给出约百分之十的检验，2016 年加入 LARES 后，分析组报告的精度为约百分之五。两次结果还用了不同的地球引力模型；比较它们的误差不能全归因于增加了一颗卫星。

类似的量级竞争在水星处也有。太阳自转给水星近心点的附加进动约为每世纪若干 $10^{-3}$ 角秒，小于太阳质量四极矩已经很小的进动。若历表只拟合 $J_2$ 和球对称 PPN 项，而把这个自转贡献固定为零，达到相应精度后，其一部分就会被吸入 $J_2$ 的估计。地球激光卫星和绕日水星用的都是 $a^{-3}(1-e^2)^{-3/2}$ 的质量流效应，究竟宜测哪一颗，取决于场源角动量、计时精度和竞争的质量矩三个条件。

\subsection{从作用反作用到质心运动}

扭秤还能区分引力质量的主动、被动两面。记物体产生场的主动质量为 $m_{\rm A}$，响应场的被动质量为 $m_{\rm P}$。设 $\bm n$ 由物体 2 指向物体 1，两者相距 $r$，它们所受的 Newton 引力之和为 $\bm F_1+\bm F_2=\frac{G}{r^2} (m_{{\rm A}1}m_{{\rm P}2}-m_{{\rm A}2}m_{{\rm P}1})\bm n$。自由落体实验比较的是 $m_{\rm P}/m_I$，上式却只在不同材料的 $m_{\rm A}/m_{\rm P}$ 相等时自动为零。若各物体的这个比值均为同一常数，还可以把常数吸入 $G$ 的定义。要检验动量守恒，必须能独立地比较两种物质作为场源和作为受力者的作用。

1968 年，Kreuzer 使一根富氟的聚四氟乙烯圆柱悬浮在富溴的液体中。调节液体密度，直至圆柱既不升也不沉，这时圆柱与其排开的等体积液体有相等的被动引力质量。再于容器附近放一副扭秤，让圆柱在容器内来回移动。对秤而言，容器中的一个区域交替由圆柱、液体占据，几何外形没有改变。若二者主动引力质量有别，圆柱移近与移远时秤所受力矩便不相同。实验未见相应的周期性力矩，将两种材料主动、被动引力质量之比的相对差异限制在 $5\times10^{-5}$ 以内。故浮力给出被动质量相等的条件，容器外的扭秤专门比较相同体积内物质产生的引力。

这如何写成 PPN 约束？将许多带电粒子组成的中性系统作平均，其内部动能、电磁结合能和压强会分别出现在引力质量中。若微观模型的平均能与流体的 PPN 表述相接，对电磁束缚系统可得 $m_{\rm A}/m_{\rm P}=1+\zeta_3 E_{\rm e}/(2m_{\rm P}c^2)$，其中 $E_{\rm e}$ 为静电能。不同核素的质子数不同，单位质量所占的核内静电能也不同。将实验中的成分和半经验核质量公式代入，Kreuzer 实验相当于 $|\zeta_3|<0.06$ 的约束。在同样的微观与流体描述相容条件下，各种能量和压强的引力贡献还满足 $6\zeta_4=3\alpha_3+2\zeta_1-3\zeta_3$。这最后一个等式是对模型的理论条件，在解释实验时不能把它当成另一项独立的观测。

月球居然也能作检验用的“两块试样”。月幔较富铁，月壳较富铝，两者质心没有重合。若铝与铁主动、被动质量比不同，幔壳对彼此的引力就会有一点不平衡，整个月球的质心将有异常加速度。只有当这个力严格沿地月连线时才只是径向扰动；在月球内部结构模型中，质心偏移方向有偏离地月连线的分量，故亦会给轨道以切向加速。寻常地月潮汐已在长期改变月球的平均运动，须用卫星资料研究潮汐耗散，并与激光测距给出的月球加速度比较。Bartlett 和 van Buren 在 1986 年得到两种材料的主动、被动质量比之差小于约 $4\times10^{-12}$，在上述静电能模型下对应 $|\zeta_3|<10^{-8}$。这个天文限制的精度还取决于月球结构和潮汐模型。

同一问题也可在双星上问。即使取 $\alpha_3=0$，若 $\zeta_2\ne0$，偏心轨道上质量不等的两星仍可能使质心沿近心点方向获得平均加速度。设 $m=m_1+m_2$，$\eta=m_1m_2/m^2$，$\Delta=(m_2-m_1)/m$，$\hat{\bm p}$ 指向星 1 的近心点。在弱场 PPN 方程中保留 $\zeta_2$ 的一阶贡献并作轨道平均，有
\[
    \langle\bm a_{\rm CM}\rangle
       =-\frac{\zeta_2}{2}\frac{G^2m^2}{c^2a^3}\,
           \frac{\eta\Delta e}{(1-e^2)^{3/2}}\hat{\bm p}.
\]
这个形式可由对称性查验。圆轨道没有特定近心点方向，故结果必须随 $e$ 消失；交换质量相等的两星也不应选出一个净力方向。在非圆且不等质量时，近心点本身会进动，质心加速度的视线分量因而随时间改变。若 $a_\parallel$ 取远离观者为正，附加的 Doppler 项为 $(\dot P_{\rm s}/P_{\rm s})_{\rm kin}\simeq a_\parallel/c$，进动又给 $\ddot P_{\rm s}$ 一项。与单独量 $\dot P_{\rm s}$ 相比，高阶导数及其相对近心点的相位有助于同脉冲星固有的缓慢自转变慢分开。1992 年对 B1913+16 自转周期二阶导数应用这一计算，给出 $|\zeta_2|<4\times10^{-5}$；这里与其它借助脉冲星的 PPN 界限一样，需假定强自引力未给该项系数以显著不同的修正。


\subsubsection{哪一些守恒参数能被单独测量？}

试估一下富氟、富溴材料何以能测 $\zeta_3$。半经验核质量公式的 Coulomb 项近似给 $E_{\rm e}/(mc^2)\simeq7.6\times10^{-4}Z(Z-1)A^{-4/3}$，这里 $Z,A$ 各为核的质子数与质量数。以 $^{19}\mathrm F$ 和质量数近 $80$ 的溴为例，两者单位质量所占的核静电能不同。由主动、被动质量比之差等于 $\zeta_3/2$ 乘两个静电能分数之差，便知所用的是原子核内的高能结合，而非化学反应时每分子释放的那一点能量。圆柱与溶液都不是纯元素，实算要对所有核素按质量分数加权。这说明哪一种材料合适，还取决于要检验的异常耦合；单以两块样品的化学性质差异最大作为选择准则并不充分。

月球的材料检验中，面对地球的平原玄武岩与背面高地在成分上有别。历史模型用富铁月幔和富铝月壳来拟合，二者中心偏离的尺度约 $10$ 公里，偏移方向据当时的模型又与地月连线成约 $14^\circ$ 的角。若月球处处严格球对称，即使不同材料的主动、被动质量比不等，各方向的自力也会彼此抵消；只有质心偏移的方向和不对称程度已由引力场、地形等资料定出，才能预报异常自力的切向投影。持续的切向力会改变轨道能量，留下平均运动的长期加速。地月潮汐耗散也会改变这一量，但同时改变地球自转，并可借海潮和卫星测量求出。因此“月球实验”事实上同时用到了天体成分分布、轨道相位和独立的潮汐资料，最后给出的是在该结构模型下的界限。

对 $\zeta_2$ 的双星计时，可进一步看看哪些系统最有利。每轨平均的质心加速度正比于 $\eta\Delta e/(1-e^2)^{3/2}$。近等质量减少 $\Delta$，一颗星极轻又使 $\eta$ 变小；最有利的是质量有可观差异且轨道明显偏心的双星。将加速度沿视线投影，取升交点作角度原点，便有 $a_\parallel\propto\sin i\sin\omega$，其中 $\omega$ 为近心点幅角。若轨道面和平均加速度大小在观测期间几乎不变，脉冲钟的额外二阶导数就近似正比于 $\sin i\cos\omega\,\dot\omega$。因此在某些方位即使 $a_\parallel$ 很大，二阶导数也可暂时很小；跨越可观测的近心点进动弧段，又有助于摆脱这一几何盲区。一般孤立系统受外部银河场的共同加速也影响一阶钟漂，但不会天然以已知近心点角为相位给出这种周期项。

至于 $\zeta_1$ 和 $\zeta_4$，本节这些实验没有一个能在其余系数全部自由时将二者逐一读出。$\zeta_1$ 参与地月测距中的 Nordtvedt 组合，可与独立的 $\beta,\gamma,\xi,\alpha_i,\zeta_2$ 界限一起推断。$\zeta_4$ 衡量压强产生引力的方式，若接受前述微观动能和宏观流体压强相容的理论关系，则可经 $\alpha_3,\zeta_1,\zeta_3$ 间接限制。给十个参数各列一个数时，应分清哪一些来自自由拟合、哪一些是把其它参数先定为零、哪一些还用了额外的理论关系。不同实验的误差再小，也不能把没有独立测到的方向变成一个直接读数。

\section{紧致天体的结构与运动}

设两颗中子星相距 $a$，其中一颗质量为 $m_A$、半径为 $R_A$。取 $m_A=1.4M_\odot$、$R_A=12$\,km，其表面的紧致度 $Gm_A/(R_Ac^2)$ 约为 $0.17$。双星相隔很远时，轨道上的 $G(m_1+m_2)/(ac^2)$ 却可以很小。要计算此时的运动，哪里可以作后 Newton 展开？

\subsection{一颗星有几种质量？}

先单独看一颗星。对静止球对称的中子星，给定压强随能量密度变化的状态方程和中心密度，便可按\hyperlink{subsection.2.4.1}{\S\,2.4.1} 的静力平衡方程向外积分。压强降至零处给出星体半径 $R$，其外的 Schwarzschild 场给出质量 $M$。改变中心密度即得一条质量与半径的关系曲线。若某种状态方程的稳定支连一颗已测质量约为 $2M_\odot$ 的中子星也容不下，便可在广义相对论下排除它。反过来要用同一颗星检验引力理论，就须把状态方程的不确定度计入。

这里的质量可从远处一颗伴星的轨道测出。若把星体拆成一个个重子，移到无穷远再去称量，结果是多少？以下构造星体时暂用几何制，以 $\varepsilon$ 表示含静质量的总能量密度，以 $n$ 表示流体共动系所测的重子数密度。冷物质沿同一状态方程有 $\d(\varepsilon/n)=-p\,\d(1/n)$，故 $\d\varepsilon=(\varepsilon+p)\d n/n$。已知 $p(\varepsilon)$，并在低密度处约定每个重子的参考静质量 $m_b$，就可由此算出 $n(\varepsilon)$。若 $r$ 为面积半径，$m(r)$ 为球内的引力质量，径向的固有长度为 $\d r/\sqrt{1-2m(r)/r}$。于是星体总重子数与引力质量分别为
\[
 N=4\pi\int_0^R\frac{n(r)r^2\,\d r}{\sqrt{1-2m(r)/r}},
 \qquad M=m(R)=4\pi\int_0^R\varepsilon(r)r^2\,\d r .
\]
只有前式按空间切片的固有体积积分。后式出自径向场方程，用它构造的质量已经计入结合能。在上述重子质量约定下，若 $Nm_b>M$，把星体还原为分散的物质须补给能量 $(Nm_b-M)c^2$。因此变动引力强弱而保持同一团物质时，应当固定 $N$。若反而固定中心密度，变化后的平衡星一般连重子数都不同。

再问，一次天文测量到底给出星体的哪个量？远处所见的圆轨道周期给出 $GM$。若能识别星面谱线的固有频率和红移 $z_s$，对静止球星又有 $1+z_s=(1-2GM/Rc^2)^{-1/2}$，它单独只给出紧致度。两者合用才有 $R$。用热辐射估计半径也有类似区别。设整颗静止星都以同一温度 $T_s$ 按黑体辐射，且半径在光子球以外。每个光子的能量和单位坐标时发出的光子数各少一个红移因子，故 $L_\infty=L_s/(1+z_s)^2$，而谱温度 $T_\infty=T_s/(1+z_s)$。在已知距离下由 $L_\infty=4\pi R_\infty^2\sigma T_\infty^4$ 反算，所得是 $R_\infty=R(1+z_s)$。实际还要判断大气如何改变光谱、发射区是否覆盖整个星面。可见即使在广义相对论内，“测得半径”也须交代用的是什么信号。

\subsection{从一族平衡解到可观测量}

为辨清状态方程和引力理论各起什么作用，可先作两个具体模型。第一种取 $\varepsilon$ 为常数，则广义相对论中 $m(r)=Mr^3/R^3$。将这个已知函数代回静力方程，自 $p(R)=0$ 向内积分，得到
\[
 \frac{p(r)}{\varepsilon}=
 \frac{\sqrt{1-2Mr^2/R^3}-\sqrt{1-2M/R}}
      {3\sqrt{1-2M/R}-\sqrt{1-2Mr^2/R^3}} .
\]
$M/R\ll1$ 时，中心压强约为 $\varepsilon M/(2R)$，与常密度 Newton 球相符。若继续压缩而保持均匀，中心压强却在 $2M/R=8/9$ 处先发散，尚未达到 $R=2M$。这个解便于检查数值积分的符号、星面位置和坐标，但它要求物质在压强改变时保持密度不变，不能据此预测中子星的真实最大质量。特别是改变引力场方程后，上述显式解一般随之失效。

第二种取冷物质的多方模型 $p=K\rho_b^\Gamma$，$\rho_b=m_b n$ 为重子静质量密度。同一状态方程的总能量密度还须由热力学关系求出。代入 $\d(\varepsilon/\rho_b)=p\,\d\rho_b/\rho_b^2$ 并选取稀薄极限，得 $\varepsilon=\rho_b+p/(\Gamma-1)$。这样 $c_s^2=\d p/\d\varepsilon=\Gamma p/(\varepsilon+p)$ 给出局域声速平方（几何制）。用一条多方关系向很高密度外推时要检查 $0\le c_s^2\le1$，尤其不能把本就违反因果条件的分支拿去约束引力。若给出的是真实核物质的表，用热力学一致的插值同样重要；只分别插值 $p,\varepsilon,n$ 而破坏微分关系，算得的重子数或下文的质量导数会产生假偏差。

实际积分时，星面附近压强很小时，径向步长稍大就会算出负压。可换为焓变量 $H(p)=\int_0^p\d p'/(\varepsilon+p')$，此时静力方程成为 $H'=-\Phi'$，星面总在 $H=0$。对上面的冷多方体有 $H=\ln[(\varepsilon+p)/\rho_b]$。在中心 $m(r)=4\pi\varepsilon_c r^3/3+O(r^5)$，故 $H(r)=H_c-2\pi(\varepsilon_c/3+p_c)r^2+O(r^4)$。可在一个很小却非零的 $r$ 以这两个展开式起步，避开方程中形式上的 $m/r^2$。把中心焓一点点增大，记录 $(M,R,N)$，即得一条平衡序列。沿冷的无自转单参数序列，第一处质量极大通常是径向基模失稳的起点。若含强一级相变，可能有第二个稳定支，这时仅有“同质量”并不能唯一选定中心密度，还须检查径向稳定性。

太阳系实验若已定出局域 $G$，能否直接拿它代入星体内部？这取决于理论中还有哪些长度尺度。在仅以一个常数 $G$ 控制引力、物态微观尺度独立于 $G$ 的一族模型中，固定中心密度后，由静力与质量方程可令 $r\propto G^{-1/2}$，$m,N\propto G^{-3/2}$。原因是势强度 $Gm/r$ 及每重子的热力学性质均不变，只有体积在改变。对质量、重子数作全微分，再消去随 $G$ 改变的中心密度，得
\[
 -\left(\frac{\p\ln m}{\p\ln G}\right)_{\!N}
 =\frac32\left[1-\left(\frac{\p\ln m}{\p\ln N}\right)_{\!G}\right].
\]
右端的斜率可在同一固定 $G$ 的平衡序列上求。$\p m/\p N$ 有每增加一颗重子所增加的总能量的意义，与把一颗重子送入星内的结合能有关。对于很弱的自引力物体，$m$ 几乎正比于 $N$，这个导数自然很小。如果理论还有有限作用程、独立的曲率长度，或改变 $G$ 的同时也大幅改变无量纲耦合，上述标度即被破坏，必须回到完整的边值问题。

除了质量半径曲线，还能测转动惯量。对缓慢匀速转动的星体，自转角速度 $\Omega_s$ 的一次项只扰动度规的 $t\phi$ 分量，可写作角线元中的 $\d\phi-\omega(r)\d t$。$\omega(r)$ 是局部惯性轴相对远方的拖曳角速度，外部真空解在广义相对论中给出 $\omega=2J/r^3$。因几何在改变自转方向时质量和形状不变，球形背景上的一次扰动足以求 $J$ 和 $I=J/\Omega_s$，而自转扁率要从二次扰动求。给定一条状态方程，先按前法选定质量对应的背景星，再令中心拖曳解正则，并在星面接上外解；外部 $1/r^3$ 项的系数就确定 $I$。这样可把一次独立的 $I$ 测量同 $M,R$ 并列，减少仅用一条质量半径关系时引力与物态的退化。

给出一个带优选参考系的实例。对 Einstein-以太理论中相对背景矢量静止的球星，若还取矢量的径向分量为零，它的单位模长条件固定时间分量 $K^t=\exp(-\Phi)$。这种取法比静态球对称本身多一个假定。此时径向方程只用到组合 $c_{14}=c_1+c_4$，在保持弱场 $\alpha_1=\alpha_2=0$ 的模型中也可改变最大质量。一组使用较硬强子物态的计算给出，$c_{14}$ 从零升到 $1.9$ 时，序列最大质量从约 $2.4M_\odot$ 降到约 $1.8M_\odot$；换为中等硬度的物态，相应两端约为 $2.2M_\odot$ 与 $1.7M_\odot$。这是固定模型的数值比较。一次可靠的 $2M_\odot$ 测量便能排除其中部分「物态、耦合」组合，却不能不经物态选择直接把四个 $c_i$ 逐个定下来。此处仅借这些历史模型说明检验方法。

张量-矢量-标量理论的情况又多一步。辅助度规与物质度规既有共形因子，又含优选矢量方向。求内部方程时用前者的面积半径，计算谱线红移时却须用物质度规的 $g_{tt}$；两处若误用了同一个度规函数，会在场很弱时似乎符合 Newton 极限，而在星面强场的红移上出错。原有模型选取 $\Gamma=2.46$ 的多方物态，单在广义相对论中所得最大质量就约为 $1.65M_\odot$。据此比较耦合变化的趋势是可行的，用它给高质量脉冲星作绝对的理论界限却不合适。另一例是将标量同 Gauss--Bonnet 密度耦合。恒定系数时该密度在四维只给边界项，而系数依赖于标量时，源中出现曲率平方，原本同紧致度、不同尺度的星也可能给不同效应。与之相反，Pontryagin 密度在静态球对称背景因宇称而消失，要追寻它的主要效应应让星自转。决定实验灵敏程度的既有紧致度，也有耦合的对称性。

\subsection{标量如何进入星体的边值问题}

可把星体平衡条件分出与理论有关和无关的两步。在物质所耦合的度规中写 $g_{tt}=-\exp(2\Phi)$。静止流体的 4-速沿时间 Killing 矢量，径向 Euler 方程便是 $p'=-(\varepsilon+p)\Phi'$。给出度规即可用此式求压力梯度，这一步只用到了物质能动张量的协变守恒。决定 $\Phi'$ 与 $\varepsilon,p$ 之间关系的才是场方程。在广义相对论中，质量函数与压强即足以闭合方程；若引力还有标量自由度，还须同时解一个径向微分方程及其边界条件。

举一个可实际积分的例子。取无质量、单标量且具有正动能的标量-张量理论。为使方程简明，取 $G_*=c=1$，用辅助度规 $g^*_{\mu\nu}$ 写引力作用量，使其中的两项正比于 $R_*-2(\nabla_*\varphi)^2$。物质所测的度规仍为 $g_{\mu\nu}=A^2(\varphi)g^*_{\mu\nu}$。记 $\alpha(\varphi)=\d\ln A/\d\varphi$，$\varepsilon,p,n$ 始终是流体用自己的钟尺测得的量。由度规变分、标量变分分别得
\[
 G^*_{\mu\nu}=8\pi T^*_{\mu\nu}
       +2\p_\mu\varphi\,\p_\nu\varphi-g^*_{\mu\nu}(\nabla_*\varphi)^2,
 \qquad
 \square_*\varphi=-4\pi\alpha(\varphi)T_*,
 \qquad T_*=A^4(-\varepsilon+3p).
\]
特别是标量的物质源含 $\varepsilon-3p$。把中子星近似为一个加大了密度的尘埃球，会漏掉压强对这项源的贡献。

取辅助度规为 $\d s_*^2=-\exp(2\nu)\d t^2+B^{-1}\d r^2+r^2\d\Omega^2$，其中 $B=1-2m_*(r)/r$，记 $\psi=\d\varphi/\d r$。代入时间、径向场方程和上述平衡条件得到一组一阶方程
\begin{align*}
 m_*'&=4\pi r^2 A^4\varepsilon+\tfrac12r^2B\psi^2,
 &\nu'&=\frac{m_*+4\pi r^3A^4p}{r^2B}+\tfrac12r\psi^2,\\
 p'&=-(\varepsilon+p)(\nu'+\alpha\psi),
 &\varphi'&=\psi .
\end{align*}
最后一个未知导数由 $(r^2\e^\nu\sqrt B\,\psi)'=4\pi r^2\e^\nu A^4\alpha(\varepsilon-3p)/\sqrt B$ 给出。压强方程中的 $\alpha\psi$ 正是 $\Phi=\nu+\ln A$ 的第二项所致。在 $A=1,\psi=0$ 的广义相对论极限，以上各式还原为前章的静力方程。改变引力理论后，星体外面也可能储存标量场能量，故不能一到 $p=0$ 就以 $m_*(r)$ 的值作为远处测出的质量。

如何给定初值？中心处 $m_*(0)=0$、$\psi(0)=0$，以保证球心无特定径向，另给中心压强 $p_c$。还可猜一个 $\varphi_c$，从中心向外积分，遇到 $p(r_s)=0$ 后去掉物质源，继续积分真空方程。若无穷远的 $\varphi$ 未达到指定的 $\varphi_0$，就改动 $\varphi_c$ 重算。最后定时间坐标的单位，使无穷远的物理钟以单位速率走动。所得星面物理半径是 $A(\varphi_s)r_s$，总重子数是 $4\pi\int_0^{r_s}A^3n r^2\,\d r/\sqrt B$。改变 $p_c$ 得一族星体，比较不同 $\varphi_0$ 的结果时，可在各族中插值找出相同 $N$ 的星。一个边界条件在中心，另一个却在远处，这正是星体能对外界引力环境作出反应的原因。

若还有标量引力场 $\phi$，求星体的平衡解须给定外部标量场的边界值 $\phi_0$。自由落体坐标不能消去这个值，伴星又可使它发生改变。局域引力强度随之变化，因而即使重子数 $N_A$ 不变，结合能和总质量 $m_A(\phi_0)$ 也可改变。记 $s_A=(\partial\ln m_A/\partial\ln\phi_0)_{N_A}$ 为星体对标量场的\textbf{敏感度}。求同一重子数在两个邻近背景值下的平衡解，即可估计 $s_A$。

这里须分清 $\phi$ 与 $\varphi$。前者是物质度规表述中乘在曲率标量前的场，因而 $\phi\propto A^{-2}(\varphi)$。若用辅助度规计时，同一星体的有效质量为 $m_A^*=A\,m_A$，这可由 $m_A\,\d s=m_A^*\,\d s_*$ 直接读出。于是 $\alpha_A=(\partial\ln m_A^*/\partial\varphi_0)_{N_A}=\alpha_0(1-2s_A)$，其中 $\alpha_0=\alpha(\varphi_0)$。这一微分关系适用于光滑的平衡分支。实际遇到 $\alpha_0$ 很小、$\phi$ 作为参数不便时，直接求星外 $\varphi$ 的 $1/r$ 系数更稳妥。在渐近辅助坐标中，取出的标量场为 $\varphi=\varphi_0-G_*m_A^*\alpha_A/(rc^2)+\cdots$。物质总质量和标量荷由不同的远场系数给出，测出其中一个还不足以得知另一个。

某些标量-张量理论中，即使弱场物质与标量的耦合很小，中子星内部也能维持很强的标量场，此称\textbf{自发标量化}(spontaneous scalarization)。这个可能性已隐含在径向方程中。将耦合近似写为 $\alpha(\varphi)=\alpha_0+\beta_0(\varphi-\varphi_0)+\cdots$。令 $\alpha_0=0$，$\varphi=\varphi_0$ 是一个解。在此解上加小扰动 $\delta\varphi$，标量方程的线性项相当于有效质量平方 $m_{\rm eff}^2=-4\pi\beta_0 T_*$。若星内 $T_*<0$ 而 $\beta_0$ 又充分负，$m_{\rm eff}^2$ 可以为负。有限大小的星有梯度能代价，故符号改变还不足以断定失稳；须解满足中心正则、外部衰减的扰动，判断有无随时间增长的模式。若出现这种模式，原来的常标量分支便不能描述稳定星，完整非线性方程可给出具有有限标量荷的另一分支。临界点依赖质量和状态方程。太阳系实验所容许的弱耦合，未必也能通过双星实验。

\subsection{黑洞的外场与双星的运动}

把中子星换成黑洞，外部解的内边界由星面变为视界。仍以上述无标量势的理论为例，在无物质的静态外部有 $\square_*\varphi=0$。球对称方程可直接积分为 $r^2\e^\nu\sqrt B\,\varphi'=C$。若 $C\ne0$，则 $(\nabla_*\varphi)^2=C^2/(r^4\e^{2\nu})$ 会在规则静态视界处发散。故正则性要求 $C=0$，给定无穷远常数后外部标量就处处等于它。此时外部度规满足真空 Einstein 方程。更一般的定态无毛证明也依赖于场继承相应的时空对称性、边界条件和正的标量动能。若有外部物质直接提供标量源、无穷远的背景随时间变，或曲率本身也可激发标量，这个积分论证均须重做。因而在某一理论中测纯双黑洞波形，可能对强场标量耦合仍不敏感；中子星可提供黑洞所缺的物质源。

再把伴星放回来。设 $r$ 为到第一颗星中心的距离。在 $R_A\ll r\ll a$ 的区域内，该星自身的场已弱，伴星的场又几乎均匀。可在此将星体附近的强场解与双星之间的后 Newton 解匹配。若忽略自转和潮汐形变，广义相对论中球形星体的外场只留下质量 $m_A$ 这个参量。于是计算低阶轨道运动时，结合能虽已计入质量，却无须另给星体的密度分布。这说明前文的后 Newton 运动方程何以也适用于相隔很远的中子星。

匹配还有另一种写法。在双星之间用有效点粒子作用量 $S_A=-c\int m_A^*(\varphi)\,\d s_*$ 代替星内方程，所有关于球星结构的信息暂存于函数 $m_A^*(\varphi)$。取 $u_*^\mu u^*_\mu=-1$ 并用几何制，对世界线作变分得 $u_*^\nu\nabla^*_\nu u_*^\mu =-\alpha_A(g_*^{\mu\nu}+u_*^\mu u_*^\nu)\p_\nu\varphi$。投影算符保证力与 4-速正交。对一个无自引力的质点有 $m_A^*=A m_A$，这条方程就是同一质点在物质度规中的测地线方程。对一颗中子星，右端的斜率却可偏离 $\alpha_0$。将伴星产生的标量远场代入，标量加速度含有 $\alpha_A\alpha_B$，再加上张量引力，即得两体的最低阶耦合。在物质钟尺的单位中它是 $G_{AB}=G_*A_0^2(1+\alpha_A\alpha_B)$，其中 $A_0=A(\varphi_0)$。扭秤上两颗弱自引力试样对应的则是 $G_{\rm N}=G_*A_0^2(1+\alpha_0^2)$。若只测一个双星的周期与间距，由 Kepler 定律得出的是 $G_{AB}(m_A+m_B)$。把这个量直接叫作“已知的两星质量”，再用来检验同一条 Kepler 定律，将永远得不出不符合的结果。

下一阶还会怎样？伴星使本星所处的 $\varphi$ 随距离改变，展开 $\ln m_A^*(\varphi)$ 就会遇到二阶系数 $\beta_A=\partial\alpha_A/\partial\varphi_0$。在偏心轨道上，这项响应随着两星远近而改变。测量近日点进动涉及的是相应的保守运动，测量轨道衰减还须知道标量场带走了多少能量。两种观测即使都来自一条轨道，也并非测同一个参数。自转的脉冲星又多一项可观测量。若外力矩可略，$\Omega_s I$ 近乎守恒，伴星造成的转动惯量变化给出 $\delta\Omega_s/\Omega_s=-\delta I/I$。它可在脉冲走时中产生与轨道同步的项，故强场标量理论的计时公式中，有时还要用到 $I$ 对背景场的敏感度。

匹配时若不再略去外场梯度，首先出现潮汐四极矩。以外部潮汐张量 $\mathcal E_{ij}$ 的符号约定为 $Q_{ij}=-\lambda_A\mathcal E_{ij}$，$\lambda_A$ 度量星体易于形变的程度。量纲上 $\lambda_A$ 正比于 $R_A^5/G$，比例系数由内部解与外部受扰解在星面的匹配决定。因 $\mathcal E$ 的量级是 $Gm_B/a^3$，诱导四极矩的相互作用能量级为 $\lambda_A(Gm_B/a^3)^2$，与两体 Newton 势能相比小约 $(m_B/m_A)(R_A/a)^5$。这给出了两个不同的近似条件。远距双星的计时可把每颗星压缩为质量和自旋，而双中子星临近并合时，波形相位已经可能分辨不同的潮汐形变能力。

\subsection{黑洞边界与试验粒子所见的外场}

上述静态论证还有一个不必指定面积半径的写法。以 $\delta\varphi=\varphi-\varphi_0$ 乘真空方程 $\square_*\varphi=0$，在外视界、远方大球以及两片由定态对称联系的空间切片围成的区域 $\mathcal V$ 积分。分部积分后有
\[
 \int_{\mathcal V}(\nabla_*\varphi)^2\sqrt{-g_*}\,\d[4]{x}
   =\int_{\partial\mathcal V}\delta\varphi\,\nabla_*^\mu\varphi\,\d\Sigma_\mu .
\]
两个时间边界给出的量相等、取向相反。若 $\delta\varphi=O(1/r)$，大球面积虽为 $O(r^2)$，其积分仍随 $1/r$ 消失。视界的面元沿其生成元，而标量若继承定态与轴对称性，便沿生成元不变。余下的体积分为零。对静态外部，标量的梯度纯为空间方向；对满足通常圆性的定态轴对称外部，同时去掉两对称方向后，梯度也在类空的二维平面内。由非负积分才得梯度处处为零。各边界一项项考察，可以看出哪些改变真的会破坏证明。例如在无穷远处强加一个随宇宙时间变的标量，两个时间切片的抵消就得重新检查；仅说“黑洞引力强”并不会破坏这一步。

若某种理论只能数值给出 $g_{tt},g_{rr}$，怎样同观测或另一个解定量比较？对物质度规采用面积半径，写 $\d s^2=-A(r)\d t^2+B_r(r)\d r^2+r^2\d\Omega^2$，在无穷远将 $t$ 按物理钟归一化。赤道面测地线的能量和角动量为 $E=A\dot t$、$L=r^2\dot\phi$。对类时圆轨道，径向方程和它的一阶径向导数给出
\[
 \Omega^2=\frac{A'}{2r},\qquad
 E^2=\frac{2A^2}{2A-rA'},\qquad
 L^2=\frac{r^3A'}{2A-rA'} .
\]
因此无论径向坐标如何重标，远方所测的 $M\Omega$ 都可比较。光子圆轨道在 $2A-rA'=0$ 处，若类时圆轨道自外向内第一次失稳，边界上 $\d L^2/\d r=0$，等价于 $3AA'+rAA''-2rA'^2=0$。检查 Schwarzschild 情形，分别给出 $r=3M$ 和 $6M$。还可测粒子在最内稳定轨道上发出的光的频移 $(k_\mu u^\mu)_{\rm em}/(k_\mu u^\mu)_{\rm obs}$，它涉及发射体的轨道运动和沿不同方向离去的光，不仅是 $\sqrt{A(r)}$ 一个函数。比较视界周长、光子轨道频率、最内稳定轨道频率、谱线红端，分别检查到度规的不同性质。

在带单位类时矢量的理论中，须再判定不同扰动的因果边界。若某模式相对以太的特征速率为 $c_I$，几何光学中其径向特征面未必与物质光锥重合。一个数值解在物质光的视界正则，还应把附加场延拓到其本征视界，检查是否另遇奇点。专取与优选系相对静止的孤立黑洞，并在数值上选定一组稳定参数，算出的 $M\Omega_{\rm ISCO}$ 等量可与 Schwarzschild 的值逐项作比。有的早期 Einstein-以太模型中，沿固定耦合方向变化，由 $c_+=c_1+c_3\simeq0.6$ 的几个至十几个百分点，到了 $c_+\simeq0.9$ 可达数十个百分点。这里改变的是指定模型的参数；一个具体黑洞观测要能区分这种差异，还须知道质量、自旋、吸积流乃至所用谱发射模型的误差。若只看一个轨道频率，把质量重新标定便能隐藏部分差异。

\subsection{改进的 EIH 拉氏量}

轨道上有时没有一种适用于所有天体的“强场 PPN 参数”。设一白矮星和一中子星都在同一遥远源 $C$ 的场中下落。上一节的标量模型里，源产生的弱场含 $m_C\alpha_C$；两星受标量加速度则分别正比于 $\alpha_A\alpha_C$、$\alpha_B\alpha_C$。自身场虽强，轨道展开的变量仍可取 $v^2$ 与 $m/r$，但每个系数须注明对应哪个星体。为此可把弱场的多体 EIH 拉氏量推广为一族强场拉氏量。以下只考虑半守恒理论，略去 Whitehead 势、潮汐与自转，使用以弱自引力试样实测的 $G_{\rm N}=c=1$，并在可能存在的优选系内写式。对每颗星取其惯性质量 $m_a$，把与 $v_a^2$ 成正比的自能修正都吸收到它里面，略去常数后有
\begin{align*}
 L={}&\sum_a m_a\left[\frac{v_a^2}{2}
              +\frac{1+\mathcal A_a}{8}v_a^4\right]
 +\frac12\sum_{a\ne b}\frac{m_am_b}{r_{ab}}
        \left(\mathcal G_{ab}+K_{ab}\right)\\
 &-\frac12\sum_a\sum_{b\ne a}\sum_{c\ne a}
      \mathcal D_{abc}\frac{m_am_bm_c}{r_{ab}r_{ac}},\\
 K_{ab}={}&3\mathcal B_{ab}v_a^2
 -\frac12\left(\mathcal G_{ab}+6\mathcal B_{(ab)}
                   +\mathcal C_{ab}\right)\bm v_a\cdot\bm v_b\\
 &-\frac12(\mathcal G_{ab}+\mathcal E_{ab})
       (\bm n_{ab}\cdot\bm v_a)(\bm n_{ab}\cdot\bm v_b).
\end{align*}
这里 $\bm n_{ab}=(\bm x_a-\bm x_b)/r_{ab}$，$(ab)$ 为两标对称化，双重求和遍及有序对，最后一项允许 $b=c$。由交换同一对物体可知 $\mathcal G,\mathcal C,\mathcal E$ 对称，$\mathcal D_{abc}=\mathcal D_{acb}$；$\mathcal B_{ab}$ 一般不与 $\mathcal B_{ba}$ 相等。一个经常出错的简化是只写三颗不同物体的三体势。即使只有两颗星，也有 $\mathcal D_{122},\mathcal D_{211}$，它们正是球外势在这一级的非线性项。

这些系数各怎样观察？最低阶沿轨道给的是 $\mathcal G_{ab}$，如 $\ddot{\bm x}_a=-\sum_{b\ne a}\mathcal G_{ab}m_b\bm n_{ab}/r_{ab}^2$。$K_{ab}$ 是同一物质对因相对运动产生的引力变化，$\mathcal A_a$ 则改变单体动量与速度的高阶关系。分别在 $r$ 不同、$v$ 不同的轨道段测同一对星，才有机会从运动中分开 $K_{ab}$ 与三体势贡献。在广义相对论中取 $\mathcal G_{ab}=\mathcal B_{ab}=\mathcal D_{abc}=1$、其余为零，可直接还原为前面的 EIH 拉氏量。在满足上述假定的弱自引力 PPN 极限则有 $\mathcal G_{ab}=1$、$3\mathcal B_{ab}=2\gamma+1$、$\mathcal D_{abc}=2\beta-1$、$\mathcal C_{ab}=\alpha_1-\alpha_2$、$\mathcal E_{ab}=\alpha_2$。最后两个组合已说明，改进并不限于把 $G$ 换成一套有效常数。

若实验室随双星质心以 $\bm w$ 在优选系中运动，直接把 $\bm v_a$ 换成 $\bm v_a+\bm w$ 是否足够？不够。在同一实验室时刻量两颗星的位置，对优选系却非同时；钟的速率和沿运动方向的距离也有同阶修正。应对作用量 $L\,\d t$ 作低速 Lorentz 变换，再减去不影响轨道的全时间导数。结果中与 $\bm w$ 有关的项由 $\mathcal A_a,\mathcal B_{[ab]},\mathcal C_{ab},\mathcal E_{ab}$ 控制。若它们均为零，可选质心静系而不留对宇宙速度的依赖；弱场参数 $\alpha_1,\alpha_2$ 在太阳系中消失仍不足以断言这四组强场系数也都消失。观测受迫离心率或轨道法线进动所给的，正是这些系数按组分质量加权后的组合。

\subsection{强自引力双星相对优选系的运动}

上面的拉氏量已足以算出每颗星的坐标加速度，但须先把优选系中的作用量变到双星的观测系。以下仍取 $G_{\rm N}=c=1$，$m_a$ 是在 $v_a^2/2$ 前标定的惯性质量，$g=\mathcal G_{12}$。令 $\bm n=(\bm x_1-\bm x_2)/r$，$\bm v=\bm v_1-\bm v_2$，$\bm w$ 是观测系相对优选系的速度。在一个轨道周期内可把 $\bm w$ 看作常量。置 $B_\pm=(\mathcal B_{12}\pm\mathcal B_{21})/2$、$C=\mathcal C_{12}$、$E=\mathcal E_{12}$，并记 $v_{an}=\bm v_a\cdot\bm n$、$w_n=\bm w\cdot\bm n$。这里 $E$ 是配对的结构系数。在运动系去掉常数和全时间导数后，将 $L$ 写成 $L_0+L_1$，其中 $L_0=\sum_a m_av_a^2/2+gm_1m_2/r$。单体和配对的 1PN 部分分别为
\begin{align*}
 L_1={}&\sum_a f_a+\frac{m_1m_2}{r}\,S-\frac{m_1m_2}{r^2}\,d,\qquad
 d=\tfrac12(m_2\mathcal D_{122}+m_1\mathcal D_{211}),\\
 f_a={}&\frac{m_a(1+\mathcal A_a)}{8}v_a^4+
   \frac{m_a\mathcal A_a}{4}
       [v_a^2 w^2+2v_a^2(\bm v_a\cdot\bm w)+2(\bm v_a\cdot\bm w)^2],\\
 S={}&\tfrac12(K_{12}+K_{21})-\tfrac12 C w^2
     +[3B_-\bm v-\tfrac12 C(\bm v_1+\bm v_2)]\cdot\bm w
     -\tfrac12 E w_n(w_n+v_{1n}+v_{2n}) .
\end{align*}
$K_{ab}$ 是前面在优选系给出的速度二次式，此处将它按 $\bm v_1,\bm v_2$ 求值。$S$ 的最后三项和 $f_a$ 中的 $w$ 项已经计入同时面、时间间隔、距离在两个参考系之间的变换。若 $\mathcal A_a=B_-=C=E=0$，则即使 $g,B_+,d$ 随两颗星的种类而变，所有 $\bm w$ 项仍同时消失。

要分别算星 1、星 2 的加速度，可用这个较短的求导式。以 $e_1=+1,e_2=-1$ 表示对两个位置微分的符号，$\dot r=\bm n\cdot\bm v$，$P_{ij}=\delta_{ij}-n_i n_j$，$S_{,\bm n}$ 表示先把 $\bm n$ 的三个分量当作独立变量求梯度。对每颗星有
\begin{align*}
 \bm a_a^{\rm N}&=-e_a\frac{gm_1m_2}{m_ar^2}\bm n,\\
 m_a(\bm a_a-\bm a_a^{\rm N})
   &= e_a\frac{m_1m_2}{r^2}
       \left[-S\bm n+P S_{,\bm n}+\frac{2d}{r}\bm n\right]
       -D_0\left(f_{a,\bm v_a}+\frac{m_1m_2}{r}S_{,\bm v_a}\right),\\
 D_0&=\dot r\,\p_r+
      \frac{\bm v-\dot r\bm n}{r}\cdot\p_{\bm n}
       +\sum_{b=1}^2\bm a_b^{\rm N}\cdot\p_{\bm v_b}.
\end{align*}
式中 $P S_{,\bm n}$ 为投影后的三维矢量。函数 $S,f_a,d$ 全已显式给定，这里不用解含未知加速度的隐方程。例如 $S$ 中的 $-Ew_n^2/2$ 对位置的梯度，除 $1/r$ 的径向导数以外还有 $-Ew_n P\bm w/r$；只对距离求导就会漏掉使轨道平面旋转的力。相反，单体的 $v_a^4$ 不含位置，须在 $D_0 f_{a,\bm v_a}$ 中以 Newton 引力代入方能影响运动。由此式单独对每颗星积分，可直接预报其脉冲发射位置；只对相对加速度积分时，若另需两颗星在天球上的位置，还要指定质心运动。

\subsubsection{相对加速度中实际能分辨的组合}

令 $m=m_1+m_2$、$q_a=m_a/m$、$\eta=q_1q_2$、$\Delta=q_2-q_1$。在已经是一阶修正的量中可代 $\bm v_1=q_2\bm v$、$\bm v_2=-q_1\bm v$。置 $D=q_2\mathcal D_{122}+q_1\mathcal D_{211}$ 及 $A^{(k)}=q_2^k\mathcal A_1-(-q_1)^k\mathcal A_2$。先让 $\bm w=0$，相对加速度便是
\begin{align*}
 \bm a_{\rm L}={}&-\frac{gm}{r^2}\bm n+
     \frac{m}{r^2}\left[
       \bm n\left(H_1v^2+H_2\dot r^2+H_3\frac m r\right)
                          +H_4\dot r\,\bm v\right],\\
 H_1={}&\tfrac12[g(1-6\eta)-3B_+-3\Delta B_-
                                 -\eta(C+2E)+gA^{(3)}],\\
 H_2={}&\tfrac32\eta(g+E),\qquad
 H_3=g[2\eta g+3B_++\eta(C+E)+3\Delta B_-]+D,\\
 H_4={}&g(1-2\eta)+3B_++3\Delta B_-+\eta C+gA^{(3)} .
\end{align*}
这里 $H_i$ 只用作四项运动学结构的系数。在广义相对论中 $H_1=-1-3\eta$、$H_2=3\eta/2$、$H_3=4+2\eta$、$H_4=4-2\eta$。仅使优选系效应消失而仍允许标量荷时，不应预先用这四个数；它们会由该对双星的 $g,B_+,D$ 改变。若用近心点进动检验，沿 Newton 椭圆平均后给 $\Delta\varpi_{\rm L}=[\pi m/p](-H_3/g-2H_1+2H_4)$，其中 $p=a(1-e^2)$；$H_2$ 给周期扰动，但不单独出现在这次长期进动里。

在 $\bm w\ne0$ 的系中定出两种轨道组合
\[
 \widehat\alpha_1=\Delta(C+E)-6B_- -2gA^{(2)},\qquad
 \widehat\alpha_2=E-gA^{(1)} .
\]
对同一系统在这段观测期间为常数的 $w^2$ 修正可收入实测的 Kepler 质量。若 $g$ 坚持取优选系中定义的数，上一组 Euler 方程还给各向同性的一项 $[m/(2r^2)](C+gA^{(1)})w^2\bm n$。将它计入拟合的 $gm$ 之后，余下的方向依赖为
\begin{align*}
 \bm a_{\rm PF}=\frac{m}{r^2}\bigg\{&
 \bm n\big[(\tfrac12\widehat\alpha_1+2gA^{(2)})(\bm w\cdot\bm v)
       +\tfrac32(\widehat\alpha_2+gA^{(1)})w_n^2\big]\\
 &-\bm w\big[\tfrac12\widehat\alpha_1\dot r+\widehat\alpha_2 w_n\big]
       +gA^{(2)}\bm v w_n\bigg\}.
\end{align*}
两颗星即使在弱引力近似下满足同一套 PPN 系数，第一项的归一仍须留意。让自能趋零，$g\to1$、$B_-\to0$、$A^{(k)}\to0$，则 $\widehat\alpha_2\to\alpha_2$，而此处的 $\widehat\alpha_1\to\Delta\alpha_1$，已含质量差。有些计时分析把 $\Delta$ 提到所报参数之外，两种限值应先换到同一约定；对质量近乎相等的双星尤其如此。将星的标签互换后，$B_-,\Delta,A^{(2)}$ 与 $\widehat\alpha_1$ 都变号，$\widehat\alpha_2$ 不变，恰与 $\bm n,\bm v$ 反向时所需的相对加速度一致。若两星在质量和内部结构上完全相同，所有这些反对称组合都为零。

为何近圆双星常只报告这两个带帽参数，而无法分别给 $\mathcal A_a,\mathcal B_{ab},\mathcal C_{ab},\mathcal E_{ab}$？置 $n_b=\sqrt{gm/a^3}$、$v_{\rm K}=n_ba$，令 $\hat{\bm h}$ 为轨道法线，$\bm w_\parallel$ 为 $\bm w$ 的面内投影。在一个周期内用 Newton 圆轨道平均，离心率矢量 $\bm e$ 的一阶方程给
\[
 \langle\dot{\bm e}\rangle_{w,\,e=0}
    =-\frac{m}{4a^2}\widehat\alpha_1\bm w_\parallel,\qquad
 \Delta\hat{\bm h}\big|_{e=0}
    =-\frac{\pi\widehat\alpha_2}{g}
      (\bm w\cdot\hat{\bm h})(\bm w\times\hat{\bm h}).
\]
第一式中的 $A^{(2)}$ 单独出现在瞬时的径向、切向加速度，但在圆轨道平均时两者恰相抵。与普通的 $\dot\varpi_0\hat{\bm h}\times\bm e$ 合在一起，强迫离心率的大小为 $m|\widehat\alpha_1\bm w_\parallel|/(4a^2|\dot\varpi_0|)$。第二式所用的是加速度的法向投影，平行于 $\bm n$ 或 $\bm v$ 的项均不参与。于是观测近心点缓慢转动时离心率矢量是否保持一个固定分量，可限制前一种组合；观测脉冲星投影半长轴或轨道倾角是否漂移，可限制后一种。质量比在这里只宜由两个组分的轨道尺度比或独立的光谱速度等资料确定。仅有周期和相对半长轴给出的是 $gm$，不能从中同时独立读出 $q_1,q_2$。恢复普通单位时用 $G_{\rm N}m$ 代替几何单位的质量因子，全部 1PN 加速度修正以及上面两项世俗摄动另乘 $c^{-2}$。

\subsection{标量双星的加速度和进动}

现在把前面的标量边值解与这个拉氏量接起来。记 $\zeta=(4+2\omega_0)^{-1}$，$\lambda=\phi_0\omega'_0/[(3+2\omega_0)(4+2\omega_0)]$，并取 $s'_a=\p^2\ln m_a/\p(\ln\phi_0)^2$；对物质星均固定重子数。$\lambda$ 在本小节是无量纲耦合。在 $G_{\rm N}=1$ 的单位下，球对称星体间有
\begin{align*}
 \mathcal G_{ab}&=1-2\zeta(s_a+s_b-2s_as_b),
 &3\mathcal B_{ab}&=\mathcal G_{ab}+2(1-\zeta),\\
 \mathcal D_{abc}&=\mathcal G_{ab}\mathcal G_{ac}
  +2\zeta(1-2s_b)(1-2s_c)
      [\lambda(1-2s_a)+2\zeta s'_a].
\end{align*}
其余 $\mathcal A,\mathcal C,\mathcal E$ 为零。第一式与前文 $G_{ab}/G_{\rm N}=(1+\alpha_a\alpha_b)/(1+\alpha_0^2)$ 相同，只须代入 $\alpha_0^2=1/(3+2\omega_0)$ 和 $\alpha_a=\alpha_0(1-2s_a)$ 即可核对。若 $s_a,s_b$ 都很小，最低阶相互作用近似能写成主质量乘被动质量；两星都紧致时 $s_as_b$ 不小，单独给每颗星指定一个引力质量就不能普遍描述所有配对。这是仅凭地面不同材料的自由落体实验难以看出的效应。

取 $\bm x=\bm x_1-\bm x_2$，$\bm v=\dot{\bm x}$，$\bm n=\bm x/r$、$\dot r=\bm n\cdot\bm v$、$\Delta=(m_2-m_1)/m$，并简写 $g=\mathcal G_{12}$。在质心系内只须在 1PN 项里用 Newton 关系 $\bm v_1=(m_2/m)\bm v$、$\bm v_2=-(m_1/m)\bm v$。直接对拉氏量变分再以低阶加速度消去右端的加速度，得到
\begin{align*}
 \ddot{\bm x}&=-\frac{gm}{r^2}\big[(1-A_{\rm PN})\bm n
                              -B_{\rm PN}\dot r\,\bm v\big],\\
 A_{\rm PN}&=-(1+3\eta+\bar\gamma)v^2+\tfrac32\eta\dot r^2
       +2(2+\eta+\bar\gamma+\bar\beta_++\Delta\bar\beta_-)\frac{gm}{r},\\
 B_{\rm PN}&=2(2-\eta+\bar\gamma).
\end{align*}
以上 $c=1$，恢复单位时，整个相对加速度的修正项各带 $1/c^2$。带横线的量是针对这对星的结构参数修正，可用场变量的一阶、二阶响应直接表示：
\[
 \bar\gamma=-\frac{2\alpha_1\alpha_2}{1+\alpha_1\alpha_2},\qquad
 \bar\beta_1=\frac{\beta_1\alpha_2^2}{2(1+\alpha_1\alpha_2)^2},
 \qquad \bar\beta_2=\frac{\beta_2\alpha_1^2}{2(1+\alpha_1\alpha_2)^2},
 \qquad \bar\beta_\pm=\tfrac12(\bar\beta_1\pm\bar\beta_2).
\]
这里 $\alpha_a,\beta_a$ 是前述对 $\varphi_0$ 的导数；勿与优选系 PPN 参数 $\alpha_1,\alpha_2$ 混用。在两星均为弱自引力试样时 $\bar\gamma=\gamma-1$、$\bar\beta_1=\bar\beta_2=\beta-1$。令标量荷全为零，因质量比而来的 $3\eta v^2$、$\eta\dot r^2$ 等普通 1PN 项仍在，恰与前面的广义相对论二体式相合。

轨道实验更便于用近心点的进动角，而不直接使用坐标加速度。把修正加速度投影到径向、切向，沿最低阶椭圆 $r=p/(1+e\cos f)$ 作一周平均，其中 $p=a(1-e^2)$、$f$ 为真近点角。$\dot r^2$ 的部分与径向速度项、$\dot r\bm v$ 的切向分量必须同时保留；若先把后者略掉，连广义相对论极限的进动角也不对。以实测的 $n_b=2\pi/P_b$ 消去 $a$，可写成
\[
 \dot\omega
 =\frac{3n_b}{1-e^2}(gmn_b)^{2/3}
       \left[1+\frac{2\bar\gamma-\bar\beta_+-\Delta\bar\beta_-}{3}\right].
\]
系数只涉及这对星的保守运动。进动测得的是 $gm$ 同方括号的组合，哪怕 $\dot\omega$ 本身极精确，也不能先假定方括号为一再用所得的总质量去检验同一标量理论。换一个观测，如由白矮星谱线与脉冲星走时测质量比，才可再添独立条件。若两星都是无标量荷的定态黑洞，$s_a=1/2$、$g=1-\zeta$，$\bar\gamma,\bar\beta_\pm$ 均为零，此阶运动中仅余一个不可独立辨认的质量重标。若只有一星是这种黑洞，另一星可带标量荷，保守运动在此阶仍很接近广义相对论，下一节的偶极辐射却未被这个条件排除。

\subsection{矢量背景下的敏感度}

有优选时间方向的理论须换一种求敏感度的实验。设背景矢量为 $K^\mu$，满足 $K_\mu K^\mu=-1$，与质心 4-速 $u^\mu$ 可构成不变量 $\Gamma=-K_\mu u^\mu$。一颗相对于以太静止的球星恒有 $\Gamma=1$；只改变静态中心密度无法算出质量对 $\Gamma$ 的导数。应固定该星重子数，令它以很小的速度相对以太匀速运动。在星体的静系，远处 $K^i$ 不再是零，求出静态解的一阶绕流扰动并匹配其外场，就能定出运动诱发的结构变化。因球星没有一个内禀优选的空间方向，孤立质量的标量响应从 $v^2$ 开始。若写 $m(\Gamma)=m_0[1+\sigma(\Gamma-1)+\cdots]$，由于 $\Gamma-1=v^2/2+O(v^4)$，平直背景作用量 $-m(\Gamma)\sqrt{1-v^2}$ 的二阶项变成 $m_0(1-\sigma)v^2/2$。因此实验按动量与速度所定的惯性质量和原来的静止质量参数已有不同。

四次项还会用到 $m''(\Gamma)$。把二次项的系数重新定为惯性质量后，四次项中剩下的修正就进入上述 $\mathcal A_a$；两星间运动诱发的矢量偶极场则进入 $\mathcal B,\mathcal C,\mathcal E$。例如用 $\sigma_a$ 表示上述未重标质量的导数，在一种常用归一化中有效两体耦合为 $G_{\rm N}/[(1-\sigma_a)(1-\sigma_b)]$。这些量均须由运动星模型给定。在一个质量较轻的实验样品上把 $\alpha_1,\alpha_2$ 调到零，主要抹去了弱场优选系信号，未必同时抹去一对强自引力星体的信号。若对星体质量和状态方程从不建模，仅把地面 $|\alpha_1|$ 的界限套到双中子星上，便跳过了最值得检验的物理。

最后可把两体运动放进一个层级三体系统。设内双星为 1、2，远处第三星为 3，内外间距分别为 $r,R$。在 $\bm x_1-\bm x_2$ 的加速度中，第三星的共同加速度的一部分若不能相消，会留下 $(\mathcal G_{13}-\mathcal G_{23})m_3\hat{\bm R}/R^2$，另有通常的潮汐梯度项 $O(m_3r/R^3)$；$\hat{\bm R}$ 指向第三星。两项的区别是前者在 $r\to0$ 仍在，后者来自两质心本来所在的位置不同。内轨道在固定方向上的极化又随外星的轨道相位转动。因此同一计时序列中，从内轨道尺度、相位与外轨道频率的组合辨认这两项，便是把地月--太阳实验的思想移入强自引力环境。下面用三体脉冲星说明这种测法。

\subsection{由运动星的解填写矢量理论的 EIH 系数}

设 $\sigma_a=(\p\ln m_a^{(0)}/\p\ln\Gamma)_1$，$\sigma'_a=(\p^2\ln m_a^{(0)}/\p(\ln\Gamma)^2)_1$，下标 1 表示在 $\Gamma=1$ 求导，$m_a^{(0)}(\Gamma)$ 是进入单星作用量的质量函数。重子数固定。记 $\Psi=\Gamma-1$，则在 $\Psi^2$ 前的质量系数为 $(\sigma_a^2+\sigma'_a-\sigma_a)/2$。给定质心相对于背景矢量的小速度，这个项连同钟慢因子产生四次动能。由二次动能定标后的惯性质量为 $m_a=m_a^{(0)}(1)(1-\sigma_a)$。把这一定义代回改进的 EIH 拉氏量，设 $g_{ab}=1/[(1-\sigma_a)(1-\sigma_b)]$，便有
\begin{align*}
 \mathcal G_{ab}&=g_{ab},&
 \mathcal B_{ab}&=g_{ab}(1-\sigma_a),\\
 \mathcal D_{abc}&=g_{ab}g_{ac}(1-\sigma_a),&
 \mathcal A_a&=\sigma_a-\frac{\sigma'_a}{1-\sigma_a},\\
 \mathcal C_{ab}&=g_{ab}
   [\alpha_1-\alpha_2+3(\sigma_a+\sigma_b)-\mathcal Q_{ab}-\mathcal R_{ab}],\\
 \mathcal E_{ab}&=g_{ab}[\alpha_2+\mathcal Q_{ab}-\mathcal R_{ab}].
\end{align*}
前面的拉氏量已取 $G_{\rm N}=c=1$。在质量用 kg、距离用 m 的单位中，两体 Newton 势的常数为 $G_{\rm N}g_{ab}$，三体 1PN 势的常数为 $G_{\rm N}^2\mathcal D_{abc}/c^2$，四次动能带 $1/c^2$。$\alpha_1,\alpha_2$ 在这组式中都是弱自引力极限的 PPN 参数。$\mathcal B_{ab}-\mathcal B_{ba}=g_{ab}(\sigma_b-\sigma_a)$，所以即使 $\alpha_1=\alpha_2=0$，不对称双星的优选系组合仍可非零。

对 Einstein-以太理论令 $c_\pm=c_1\pm c_3$（故 $c_+=c_{13}$）。剩下两个可由运动星解与弱场解共同读出的系数为
\begin{align*}
 \mathcal Q_{ab}&=
 \frac{(2-c_{14})(\alpha_1-2\alpha_2)}{2(2c_+-c_{14})}
      (\sigma_a+\sigma_b)
       +\frac{2-c_{14}}{c_{123}}\sigma_a\sigma_b,\\
 \mathcal R_{ab}&=\frac{8+\alpha_1}{4c_1}
       [c_-(\sigma_a+\sigma_b)+(1-c_-)\sigma_a\sigma_b].
\end{align*}
实际上，利用前文该理论的 $\alpha_2$，$\mathcal Q_{ab}$ 的第一项系数可直接写成 $(c_++c_{14}+3c_2)/c_{123}$。这样在 $2c_+=c_{14}$ 而物理模式未退化时尤为方便。若把实测 $\alpha_1,\alpha_2$ 先随意置零再做这个约化，则会把同趋于零的分子分母误算成零。

若优选矢量须正交于一族时间叶面，低能叶层理论可由固定 $c_{14}=\alpha_K$、$c_+=\beta_K$、$c_2=\lambda_K$ 并令扭转代价趋无穷得到。例如令 $c_1=-L,c_3=\beta_K+L,c_4=\alpha_K+L$，最后在保守系数中取 $L\to\infty$ 的代数极限。由前式即得
\begin{align*}
 \mathcal Q_{ab}^{K}&=
 \frac{(\alpha_K+\beta_K+3\lambda_K)(\sigma_a+\sigma_b)
                    +(2-\alpha_K)\sigma_a\sigma_b}{\beta_K+\lambda_K},\\
 \mathcal R_{ab}^{K}&=\tfrac12(8+\alpha_1)(\sigma_a+\sigma_b-\sigma_a\sigma_b).
\end{align*}
前面六组 EIH 系数的映射保持上述形式，$\alpha_{1,2}$ 则改用叶层理论的值，$G_{\rm N}=2G_{\rm bare}/(2-\alpha_K)$。这里取的是作用量系数的受约束极限；不能在计算辐射时又保留已经由叶层条件消去的独立横向矢量模式。对实际算例，可先求两颗星的 $\sigma_a,\sigma'_a$，再构造 $\mathcal G,\mathcal B,\ldots$ 变分出轨道。尤其不能把同一根质量--半径曲线在零速度处的质量斜率代作这里的运动响应。

\section{引力辐射}

\subsection{由源区到波区}\label{sec:pn-matching}

一对相距很远的中子星，附近的单星场很强，绕转却慢，传到地球上的波更弱。这三处不宜用同一近似处理。设双星间距为 $a$，角频率为 $\Omega$，记 $\ell=c/\Omega$。轨道缓慢时 $a/\ell\sim v/c\ll1$。在距双星质心 $r\ll\ell$ 处，一个光行时还不足以使源显著改变，推迟积分可按时间导数逐次展开。在 $r\gg a$ 处，双星的角大小很小，可按源的多极矩展开，但要保留推迟时刻 $u=t-r/c$。两种展开在 $a\ll r\ll\ell$ 都有意义，可令两展开代表同一个场。如此定下近区解中的齐次项、远区波中的多极系数。还须在远处规定无入射辐射；否则除了双星自己发的波，解中还可有从外面打来的任意自由波。

先在平直背景上看推迟积分的两种展开有何区别。在源区对 $S(t-|\bm x-\bm x'|/c,\bm x')$ 求 Taylor 展开，各项含同一坐标时的 $S,\dot S,\ddot S,\ldots$，但这个展开对 $r\to\infty$ 并不一致。在远处令 $\bm N=\bm x/r$，有 $|\bm x-\bm x'|=r-\bm N\cdot\bm x'+O(a^2/r)$。于是 $1/r$ 辐射项用到的是 $S(u+\bm N\cdot\bm x'/c,\bm x')$。此时保留了到观者的整段传播时间，只按源自身的光行时对第二项展开。对每个源元相隔的推迟作何种处理，决定了所得是瞬时的近区势，还是一个会向外传播的波形。

测试粒子怎样从远区场读出 TT 分量、质量单极与偶极为何不给出最低阶辐射，可见 \hyperlink{subsection.1.2.4}{\S\,1.2.4}。用该处的有效能动张量 \eqref{eq:gw-effective-stress} 对远方球面积分，便可由波形求得能流。对由自身引力束缚的源，用来构造辐射矩的是含物质和引力场的总源 $\tau^{\mu\nu}$。在谐和坐标中它满足 $\p_\nu\tau^{\mu\nu}=0$，由此可将最低阶的空间应力积分换成质量四极矩的二阶时间导数；管外场应力在更高阶的贡献仍须与波区匹配。

同一个守恒恒等式也说明为何要谨慎使用波的“能量损失”。在近区作坐标变换可改变瞬时机械能与场能的分配，远处积分所得的能通量则可同轨道量的长期变化比较。设在许多周转期间轨道尺度只略有改变，取满足 $P_b\ll\Delta t\ll a/|\dot a|$ 的平均时段。近区中随绕转振荡、有界的储能项求导后平均趋于零，因而有 $\langle\dot E_{\rm orb}\rangle=-\langle\mathcal P_\infty\rangle$，这里略去星体吸收。也可从近区解直接求两体的辐射反作用力，再作同样的平均。最低阶的两条途径给出相同的长期衰减，用远处的辐射反推轨道演化才有根据。若一周之内轨道已显著改变，必须重算运动，不能沿用这一平均步骤。

运动与辐射之所以长期被当作两个问题研究，与求近似解的办法有关。1938 年的 EIH 计算只在天体以外解真空场方程。包围一颗星作积分，边界条件要求星际弱场同星旁近球形场可相接，由此即可反求星的运动；这条途径擅长给出保守运动，却不能从一个只算到 1PN 的近区解判断波带走多少能量。另一条途径以流体质元的守恒方程为出发点，先算有限大小天体的力，再对整颗星积分。质量固定、半径远小于间距，并可略去潮汐时，两种算法对轨道应当相合。这个交叉检查尤其重要：取点质量极限太早会遇到发散的自场，只在形式发散项上选不同丢弃规则，可能制造出一个实际不存在的力。

脉冲星使同一问题可以接受实验检查。所测的近心点时刻在许多年内发生世俗漂移，要求既算出近区的坐标无关轨道预言，又要说明依据远区波动求得的能流为何等于该系统失去的能量。数学上，先作慢运动展开再令观者半径趋于无穷远，与先作辐射极限再展开，并非同一极限。物理上，一条只适用于太阳系每个物质元的弱场展开，不能不加说明地套在已有强自引力的中子星内部。近远区匹配、星体的有效参数与从反作用功独立检查能量平衡，分别回答了这三项不同的困难。

\subsection{松弛场方程、任意分界球与辐射反作用}

引力由天体产生，天体又受这同一个场加速。在第一章求给定源的线性波时，可以把物质运动看成外部已知量；若还要求轨道能量随波的发射而减少，就必须把产生波的运动和波作用于源的运动放回同一道方程中。广义相对论中可取 $\mathfrak h^{\mu\nu}=\eta^{\mu\nu}-\sqrt{-g}\,g^{\mu\nu}$，谐和条件为 $\p_\nu\mathfrak h^{\mu\nu}=0$。在暂取 $G=c=1$ 的单位中，场方程精确改写为
\[
 \square_\eta\mathfrak h^{\mu\nu}=-16\pi\tau^{\mu\nu},\qquad
 \tau^{\mu\nu}=(-g)T^{\mu\nu}+\frac1{16\pi}\Lambda^{\mu\nu}
                [\mathfrak h,\p\mathfrak h,\p^2\mathfrak h].
\]
$\Lambda$ 至少从场的二次项开始，还可包含一个场乘另一个场的二阶导数。式中左边的平直波算符是一种求解工具；总场满足的仍是原来的非线性方程。和谐和条件相容的必要条件是 $\p_\nu\tau^{\mu\nu}=0$，其物质部分恰含运动方程。若任意把两团粒子放上一条加速轨迹，只用 $\sum m v^iv^j$ 当作给定的辐射源，这个必要条件就不再成立。

考察领头四极波究竟须把场迭代几次。在第一次迭代中只保留物质项，从质量密度求出近区 Newton 势 $U$。第二次迭代的空间场源含 $T^{ij}$，也含 $\p_iU\p_jU/(4\pi)$ 一类项。场梯度平方在两体附近的积分量级为 $m_1m_2/a$，而 $\int T^{ij}\d[3]{x}$ 的量级为 $\mu v^2$，依束缚轨道关系二者同阶。更直接地看，二阶时间导数 $\ddot I^{ij}$ 除了速度平方还含 $m_Ax_A^{(i}a_A^{j)}$；所需加速度正来自另一质量的 Newton 场。前文用总源的守恒关系化为 $\ddot I^{ij}$，已经包含这个结果。它允许最后把领头质量矩写成物质密度的积分，并不允许在求解之前同时删掉场应力、又把物质留在一条加速轨道上。

推迟积分还隐藏着一个次序问题。用半径 $\mathcal R$ 的世界管围住双星，取 $a\ll\mathcal R$ 且 $\mathcal R$ 与一段源变化时间的光行距离同量级。对观察点的过去光锥，把积分分成位于管内和管外的两部分。第一部分的源在有限球内，可按 $\bm N\cdot\bm x'$ 展开得到具有明确截断的多极矩；管外则已有向外传播的场，波的应力也是新一轮迭代的源。以场点为中心的推迟距离与以源为中心的球壳半径一般不同，不能在管外统一把推迟时刻替换为 $t-r/c$。两块积分各自依赖 $\mathcal R$，它们的和却必须与这道任意分界无关。将分界球向外移动一个薄层，前一块增添的原积分与后一块失去的原积分完全相同，因而在同一近似阶上也应抵消。这可用作高阶多极计算的一项具体检查。

例如最低阶静态场随 $1/r'$ 下降，梯度平方随 $1/r'^4$ 下降。若对遍及所有空间的场应力仍照紧支集物质源的办法，先乘 $r'^l$ 再积分，一些高阶矩的径向积分便在无穷远发散。引入分界并非赋予引力场一个物质表面，而是使管外的推迟积分在该保留波动性时仍保持原状。近区展开若留下一项 $\ln\mathcal R$，只有管外传播项的相应对数也并入后才可谈物理对数。由波被静态总质量散射所留下的尾波，正是这种非局域传播的一个例子。反过来，从远区把展开继续到源内也会出现 $1/r^{l+1}$ 型的形式奇点；这些只是在源外求得的多极解，不能拿来代替星内的正则解。

能否不用远区辐射直接看见轨道衰减？在一种常用的近区规范中，四极辐射反作用可表示为与 $x^ix^jQ^{(5)}_{ij}$ 成正比的势。恢复单位且沿用正 Newton 势的记号，取
\[
 U_{\rm react}(t,\bm x)
   =-\frac{G}{5c^5}x^ix^jQ_{ij}^{(5)}(t),\qquad
 a_{\rm react}^i=-\frac{2G}{5c^5}x^jQ_{ij}^{(5)} .
\]
上标括号表示对时间求这么多次导数。在最低阶 $\Omega^2\sim Gm/a^3$、$Q\sim\mu a^2$，故该加速度比引力约小 $\eta(v/c)^5$；其中有奇数次时间导数，反转时间会改变符号，这正是 2.5PN 的耗散项。将它对所有质点作功，利用 $Q$ 无迹而对称，得 $\dot E_{\rm mech}=-G\dot Q_{ij}Q_{ij}^{(5)}/(5c^5)$。这与远区能流并非每一时刻都相等。两次分部积分给
\[
 \dot Q_{ij}Q_{ij}^{(5)}
   =Q_{ij}^{(3)}Q_{ij}^{(3)}
      +\frac{\d}{\d t}
       \left(\dot Q_{ij}Q_{ij}^{(4)}-\ddot Q_{ij}Q_{ij}^{(3)}\right).
\]
最后一项是近场储存的、有界的周期项。作轨道平均后消去它，便与从远方波幅算出的四极能流严格相合。换近区坐标时势与这个全导数均能改变，平均周期减少量却不变。在散射或并合这样的非周期过程中若舍去全导数，还需检查所选起止切片上的储能是否同样可忽略。

在非广义相对论的度规理论中也可分别对传播本征模式求推迟解。不同模式若有不同波速 $v_I$，给定同一观察事件时应在各自的过去特征锥上求源，并在约 $v_I/\Omega$ 的尺度区分该模式的近区和辐射区。若某额外场有质量，频率低于静质量对应频率的分量只作空间衰减，不能把在近区测到的 Yukawa 场直接换成一个同频的远区出射波。因而静态轨道受新场扰动、探测器直接测到新模式、双星因新模式而损失能量，是三个相关但并不等价的问题。

\subsection{双星的能量、角动量与啁啾}

四极辐射功率怎样反映在钟的读数上？考虑质量为 $m_1,m_2$、间距为 $a$ 的圆轨道双星，记总质量 $m=m_1+m_2$、约化质量 $\mu=m_1m_2/m$。在最低阶，轨道能量为 $E=-G\mu m/(2a)$，四极辐射功率为 $\mathcal P=32G^4\mu^2m^3/(5c^5a^5)$。由 $\dot E=-\mathcal P$ 得 $\dot a<0$。轨道越小，绕转越快。双星每绕行一周，四极矩重复两次，故主辐射频率 $f$ 是轨道频率的两倍。代入 Kepler 定律，得到 $\dot f=(96/5)\pi^{8/3}(G\mathcal M/c^3)^{5/3}f^{11/3}$，其中 $\mathcal M=\mu^{3/5}m^{2/5}$ 称\textbf{啁啾质量}(chirp mass)。故记录波峰愈来愈密的程度，即可测得这个质量组合。若源有宇宙学红移 $z$，用观者的钟测得的是 $(1+z)\mathcal M$。

积分即得 $t_c-t=(5/256)(G\mathcal M/c^3)^{-5/3}(\pi f)^{-8/3}$，这里的 $t_c$ 是从低频外推的时间常数，靠近并合时须用别的解。频率从 $f_1$ 升到 $f_2$，观者可数到的周期数为 $\int_{f_1}^{f_2}f\,\d f/\dot f$。它随 $f_1^{-5/3}-f_2^{-5/3}$ 而变，故在低频段尤其容易积累相位。若有微小的额外能量损失，只要能持续数许多周期，最后的波峰到达时刻仍可显著偏离原来的预报。

偏心轨道与圆轨道又有何不同？此时 $a$ 表示半长轴，令 $e$ 为离心率。星体在近心点附近转得最快，四极矩的三阶导也最大。因此用平均间距代入圆轨道公式会低估辐射。取偏近点角 $u_b$，有 $r=a(1-e\cos u_b)$、$\Omega\,\d t=(1-e\cos u_b)\d u_b$。按此权重对一圈平均，能量通量与角动量通量分别为
\begin{align*}
 \langle\mathcal P\rangle&=
 \frac{32}{5}\frac{G^4\mu^2m^3}{c^5a^5}
 \frac{1+73e^2/24+37e^4/96}{(1-e^2)^{7/2}},\\
 -\langle\dot J\rangle&=
 \frac{32}{5}\frac{G^{7/2}\mu^2m^{5/2}}{c^5a^{7/2}}
 \frac{1+7e^2/8}{(1-e^2)^2}.
\end{align*}
先用 $E=-G\mu m/(2a)$ 求 $\dot a$，再对 $J=\mu\sqrt{Gma(1-e^2)}$ 取对数导数，可把两条损失率分别换成 $a,e$ 的变化率。结果为
\[
 \left\langle\dot e\right\rangle=
 -\frac{304}{15}\frac{G^3\mu m^2}{c^5a^4}
 \frac{e(1+121e^2/304)}{(1-e^2)^{5/2}} .
\]
圆轨道的 $e=0$ 得以保持，小离心率时两式相除给出 $\d\ln e/\d\ln a=19/12$，故仅受这种耗散的双星在旋近中趋于圆化。偏心脉冲星却常离并合还很远，它在观测期间可保有相当大的 $e$。计时实验中应保留完整的离心率因子。

四极矩本身给的是波的方向图。设轨道角动量与视线夹角为 $i$，观者到源的距离为 $D$。取相位零点和偏振基，使圆轨道的张量波写为 $h_+=H(1+\cos^2i)\cos\Phi/2$、$h_\times=H\cos i\,\sin\Phi$，其中 $H=4(G\mathcal M)^{5/3}(\pi f)^{2/3}/(c^4D)$。沿轨道轴看，两种偏振等幅且相差四分之一周期，故为圆偏振；沿轨道面看，只有一个线偏振。若只测到一个振幅，小质量而近的源同大质量而远的源不易区分。相位先给出啁啾质量，比较不同朝向仪器的偏振振幅还可限制 $i$，这才使距离测量成为可能。推广到宇宙学距离，振幅式中用光度距离和红移后的质量。单靠最低阶啁啾无法分开红移与源系质量。

\subsection{能流之外的矩与通量}

波带走的角动量为何还能传递方向信息？取远处球面上穿出的通量为正，用质量无迹四极矩写广义相对论的最低阶总角动量通量，有 $\mathcal F_{J_i}=(2G/5c^5)\epsilon_{ijk}\langle\ddot Q_{ja}Q^{(3)}_{ka}\rangle$。它涉及两个方向分量的相对相位。如果沿固定轴作直线往返运动，四极张量虽可改变并产生能流，但其主轴不转，$\ddot Q$ 与 $Q^{(3)}$ 可在同一基下对角化，反对称缩并为零。圆轨道则有持续转动的主轴，角动量通量沿轨道法线。同样的能量损失率，未必意味着同样的轨道面或离心率演化。

总线动量通量则相当于将每个方向的能量通量乘 $\bm N/c$ 再在全天积分。单独的质量四极图样在 $\bm N\mapsto-\bm N$ 下为偶函数，乘 $\bm N$ 后积分为零，故不能从最简单的四极功率算出双黑洞遗留体的反冲速度。产生首个非零项需让具有相反宇称的角向图样干涉。例如用通常的对称无迹质量矩 $I_{ij}=Q_{ij}$、质量八极矩 $I_{ijk}$ 以及流四极矩 $J_{ij}$，最低阶可写为
\[
 \mathcal F_{P_i}=\frac{G}{c^7}
 \left[\frac{2}{63}I^{(4)}_{ijk}I^{(3)}_{jk}
       +\frac{16}{45}\epsilon_{ijk}I^{(3)}_{ja}J^{(3)}_{ka}\right] .
\]
矩的领头源分别为 $\sum_A m_Ax_A^{\langle i}x_A^{j}x_A^{k\rangle}$ 与 $\sum_A m_A x_A^{\langle i}(\bm x_A\times\bm v_A)^{j\rangle}$，尖括号表示对称去迹。对于等质量无自旋的双星，交换两星使瞬时源绕质心反演不变，这两种奇宇称矩都消失；质量不等或自旋排列破坏对称时，辐射才可有净动量。反冲是将矢量通量随时间积分，轨道各相位有时相消，故仅从某一时刻的峰值不能估出最后的反冲。并合处慢速近似已失准，应由数值相对论提取辐射多极矩再积分。

高阶波形中“源矩”与远方所读出的“辐射矩”也要分清。以张量质量四极为例，领头辐射矩是 $I_{ij}^{(2)}(u)$。若把近区、外区解匹配到下一非线性传播阶，还会出现形如
\[
 U_{ij}^{\rm rad}(u)=I_{ij}^{(2)}(u)
 +\frac{2GM}{c^3}\int_0^\infty I_{ij}^{(4)}(u-\tau)
        \left[\ln(\tau/\tau_0)+C\right]\d\tau+\cdots
\]
的尾项。$\tau_0$ 是为定义辐射时间而引入的参考尺度，$C$ 随约定改变，单独的这两个常数不是观测量。与直接到达的 $I_{ij}^{(2)}(u)$ 不同，积分确实取遍源的过去。对近于单频的源，两个多出的时间导数经一次时间积分后给出量级 $\Omega$，故尾项相对大小是 $GM\Omega/c^3\sim(v/c)^3$。这给前面无自旋能流的半整数阶一个传播上的解释。求高精度旋近相位时还要保留它同直接波的干涉，不能把尾项仅理解为并合后才开始出现的一段微弱余波。

\subsection{谐波、角动量和偏心率的观测窗口}

若轨道是圆形，最低阶四极矩只含 $2\Omega$ 的时间依赖。若是椭圆，$x=a(\cos u_b-e)$、$y=a\sqrt{1-e^2}\sin u_b$ 仍是偏近点角的简单函数，观者的钟却按平近点角 $\ell=u_b-e\sin u_b$ 均匀走动。把同一个 $x^ix^j$ 展成 $\exp(\mathrm{i}n\ell)$，就有全部整数次谐波。给定一条 Kepler 轨道，不必另解波动方程即可算出它在各个频率上发了多少能量。可对 $Q_{ij}$ 的各谐波系数分别求三阶时间导数，按一周平均时不同 $n$ 正交，然后用 $J_k(z)=(2\pi)^{-1}\int_0^{2\pi}\cos(ku-z\sin u)\,\d u$ 写出结果。

具体令 $J_k=J_k(ne)$，第 $n$ 次谐波的平均功率为 $\mathcal P_n=\mathcal P_{\rm circ}(a)\,g_n(e)$。这里圆轨道比较基准与原椭圆有相同半长轴。将依偏近点角积分的 Jacobian $1-e\cos u_b$ 一并计入，得
\begin{align*}
 A_n&=J_{n-2}-2eJ_{n-1}+\tfrac2nJ_n+2eJ_{n+1}-J_{n+2},\\
 B_n&=J_{n-2}-2J_n+J_{n+2},\\
 g_n(e)&=\frac{n^4}{32}
 \left[A_n^2+(1-e^2)B_n^2+\frac{4}{3n^2}J_n^2\right].
\end{align*}
圆轨道上 $J_0(0)=1$，其余非零阶 Bessel 函数在零处均为零，故 $g_2=1$，其余各项均为零。任意 $e$ 的 $\sum_{n\ge1}g_n$ 恰等于前述能流的离心率因子。以 B1913+16 的 $e=0.617134$ 作数值检查，总功率约为同半长轴圆轨道的 $11.86$ 倍，而第二谐波只约占总功率的 $0.55\%$，第六、第七谐波各约占 $12\%$。轨道相位仍可由近心点准确定出，若只在两倍轨道频率搜寻波，所利用的却不是它所发的主要功率。这是轨道计时和波形搜寻面对同一个天体仍会采用不同策略的一个例子。

也可从能量与角动量的频谱看圆化。若按 $\exp[\mathrm i(m\phi-\omega t)]$ 展开辐射，单色模式的轴向角动量通量与能通量之比为 $m/\omega$。圆轨道的主模 $m=2$、$\omega=2\Omega$，所以 $\mathcal P=-\Omega\dot J$。圆轨道序列的 $E=-G\mu m/(2a)$、$J=\mu\sqrt{Gma}$ 也本来满足 $\d E/\d J=\Omega$，从而辐射可在一系列缓慢收缩的圆轨道上连续过渡。椭圆轨道则同时发出若干不同比值的模式，不能只拿一个 $m/\omega$ 去反推 $\dot J$；必须分别积分能流和角动量流，才能确定 $\dot e$。这也解释为何测到一个 $\dot P_b$ 还没有独立测到辐射的角动量通量。

如果在演化的若干阶段只能测到不同的 $(a,e)$，还可消去演化所需的质量和时间。在给定四极衰减率中直接取 $\dot a/\dot e$ 并积分，得守恒的轨迹参数
\[
 C_0=a(1-e^2)e^{-12/19}
          \left(1+\frac{121}{304}e^2\right)^{-870/2299},
 \qquad 0<e<1 .
\]
这条关系只在轨道平均、四极辐射主导、无外体激发离心率时成立。若密集星团的近距离相遇持续给系统角动量，在某一频段测得大的离心率便不宜用四极圆化倒推更早时期的 $e$。相反，对可以认为已孤立的双星，已知 $a_0,e_0$ 后将上式给出的 $a(e)$ 代回 $\dot e$，积分即可估计从现时到下一给定频率所需的时间。圆轨道则单独积分 $\dot a$ 即得四极外推时间 $5c^5a_0^4/(256G^3\mu m^2)$；它的 $a_0^4$ 依赖说明，为何轨道稍宽的脉冲星能在相当长的望远镜观测期内只改变很小一点轨道周期。

\subsection{张量以外的辐射}

在一些标量-张量理论中，还须计入偶极辐射。记星体 $A$ 的标量荷为 $q_A$，则双星的标量偶极矩为 $\bm D_\varphi=\sum_A q_A\bm x_A$。在 Newton 质心系，$\bm D_\varphi=\mu(q_1/m_1-q_2/m_2)(\bm x_1-\bm x_2)$。两星的荷质比相等时此项消失；比值相差越大，随轨道转动的偶极矩就越大。其辐射功率正比于 $\lvert\ddot{\bm D}_\varphi\rvert^2$，比四极项少两个 $v/c$ 的幂次。在前述无质量标量理论中，荷质比正比于 $1-2s_A$。若无外来物质且远处标量场趋于常数，定态黑洞的外部标量场可处处为常数，黑洞不带标量荷；两颗这样的黑洞就没有标量偶极辐射。中子星与白矮星的结合能占比相差很大，这种配对更适合寻找偶极辐射。

可将“少两个幂次”算得更具体。暂令 $A_0=1$，并取上一节的场归一化，即 $q_A=\alpha_A m_A$。远区标量展开的一阶项为 $-G_*\,\bm N\cdot\dot{\bm D}_\varphi/(c^3D)$。标量能流正比于 $\dot\varphi^2$，利用 $\int N_iN_j\,\d\Omega=4\pi\delta_{ij}/3$，求得最低阶偶极功率 $\mathcal P_D=G_*|\ddot{\bm D}_\varphi|^2/(3c^3)$。对圆轨道，$\lvert\ddot{\bm D}_\varphi\rvert=\mu|\alpha_1-\alpha_2|a\Omega^2$。与同一轨道上的张量四极功率 $32G_*\mu^2a^4\Omega^6/(5c^5)$ 比较，有 $\mathcal P_D/\mathcal P_T=(5/96)(\alpha_1-\alpha_2)^2(c/v)^2$，其中 $v=a\Omega$。在这一节，$\alpha_A$ 的脚标标记星体，须与 PPN 的优选系参数分清。双星转得慢会使两项辐射都变弱，但会使偶极占总能流的比重增大，因而较宽的不对称双星仍有检验价值。

标量能否产生单极辐射？若低阶标量荷可视为常数，其总和不变，单极项就不随时间变。但引力结合能、动能或物质的应力若改变有效标量荷，较高阶完全可能出现随时间变化的单极项。即使偶极因两星对称而消失，标量四极也未被这个对称条件排除。换作有优选矢量场的理论，波的本征模式还可含矢量场的横向振动，且各模式可有不同波速。要从观测排除这类理论，须先在给定理论中算清源究竟能激发哪种模式、耦合有多强，再算探测器对该模式的响应。仅从场方程数一数自由度，并不能预言某次事件的六个偏振振幅都很大。

\subsection{标量波形与一般双星能流的参数}

把探测到的振幅和轨道损失率混为一谈，在标量理论中尤其容易出错。线性近似下，辅助度规的张量扰动记为 $h^{*{\rm TT}}_{ij}$，物质度规另有共形扰动 $2\alpha_0\delta\varphi\,\eta_{\mu\nu}$。取 $A_0=1$，两部分不同的角向投影分别给出张量潮汐和横向呼吸潮汐。前面求得的标量偶极远场已正比于 $\alpha_1-\alpha_2$；一台由普通物质制成的探测器要从它读到呼吸振幅，还要再乘与 $\alpha_0$ 成正比的钟尺响应。反之，标量波自身的能量由其动能项给出，轨道的阻尼随 $|\alpha_1-\alpha_2|^2$ 变化。哪怕源中的标量化很强，若远处弱自引力物质的 $\alpha_0$ 很小，在地球直接看到呼吸模式仍可能比在双星时间尺度中察觉额外能量损失困难。

角向分布和谐波也不同。圆轨道的偶极矩沿两星连线，每周转一圈，故 $\bm N\cdot\dot{\bm D}_\varphi$ 的投影正比于 $\sin i\sin(\Omega t+\phi_0)$。即使轨道完全为圆，标量波也可以在轨道基频处出现，而张量主模仍在两倍基频。若沿轨道轴看，偶极的视向投影消失，但它仍向其它方向输送能量，观者仅凭当前视线没看到标量波，不能据此认定双星没有额外阻尼。对称星体的标量四极波则可在 $2\Omega$ 处同张量波共存，其呼吸极化有不同的仪器响应。因两星引力势和动能而变化的标量单极，圆轨道在这一近似下是常数，偏心轨道却可随距离起伏。比较频率、天空方向和偏振，可把几种本来落在同一条总能流中的效应拆开。

为了看清相同 PN 阶都要保留什么，可在辐射区把标量的推迟解写成
\[
 \delta\varphi=-\frac{G_*}{c^2D}
 \left[I_s+\frac{N_i}{c}\dot I_s^i
       +\frac{N_iN_j}{2c^2}\ddot I_s^{ij}
       +\frac{N_iN_jN_k}{6c^3}I_s^{ijk(3)}+\cdots\right]_u .
\]
这里用普通对称矩，迹稍后才分解为不可约矩。在最低阶，$I_s^i$ 为有效标量荷的偶极矩；若要更高精度，每个矩本身也要由迭代场源求出修正。标量的辐射功率正比于 $D^2\int(\p_t\delta\varphi)^2\d\Omega$。单极与无迹四极的角积分为零，偶极与任意偶数次的 $N_i$ 项相乘也因角向奇偶性积分为零。四个 $N_i$ 的积分则由三种 $\delta_{ij}\delta_{kl}$ 配对构成，所以未去迹的三阶矩和偶极相乘可保留迹项。若一开始已改成对称无迹基，这份贡献便被归入修正的辐射偶极，不能再加一次。

从相对阶次看，设张量四极振幅为零阶，强场标量偶极振幅可比它低半阶。因此要把功率算到普通四极的精度，不能只留「最低偶极平方加最低四极平方」。偶极的低半阶与高半阶修正相乘已是零阶，作用在它上面的轨道加速度也须用到 1PN。若天体可产生标量单极，还应计入单极能流。将这些项在同一精度整理后，一类无自转双星的瞬时能量损失可写成
\[
 -\dot E=\frac{\eta^2}{g}\left(\frac{gm}{r}\right)^4
 \left[\frac{8}{15}\big(\kappa_1 v^2-\kappa_2\dot r^2\big)
        +\frac13\kappa_D S^2\right].
\]
此处回到 $G_{\rm N}=c=1$。$g=\mathcal G_{12}$，$S$ 度量两星的荷质差，$\kappa_1,\kappa_2,\kappa_D$ 是在这个归一化下定义的能流系数。广义相对论给 $\kappa_1=12,\kappa_2=11,\kappa_D=0$。对圆轨道 $\dot r=0$，只能测到 $\kappa_1$ 和偶极项的组合；椭圆轨道的径向运动使 $\kappa_2$ 也进入周平均。要把这些系数推回一种理论的基本耦合，须同时代入其 $g$、平衡解的敏感度以及实际采用的质量定义。

一般标量双星的两个零阶系数仍可从其平衡解算出。沿用上文 $g=\mathcal G_{12}$ 和 $\Delta=(m_2-m_1)/m$，令 $u_a=1-2s_a$、$L_a=\lambda u_a+2\zeta s'_a$；这里的 $s_a$ 是对标量背景的敏感度。取 $S_-=(s_2-s_1)/\sqrt g$、$S_+=(1-s_1-s_2)/\sqrt g$。若以普通单位计算，上式左端为 J/s，右端整体乘以 $c^5/G_{\rm N}$，其中 $gm/r$ 换成 $G_{\rm N}gm/(rc^2)$，$v,\dot r$ 分别换为 $v/c,\dot r/c$。此处 $\bar\gamma,\bar\beta_\pm$ 与保守轨道所用的量相同，但通量还要用两个不同的线性组合
\[
 X=\frac{S_-\bar\beta_++S_+\bar\beta_-}{\bar\gamma}
   =\frac{u_1L_2-u_2L_1}{4g^{3/2}},\qquad
 Y=\frac{S_+\bar\beta_++S_-\bar\beta_-}{\bar\gamma}
   =-\frac{u_1L_2+u_2L_1}{4g^{3/2}} .
\]
等号右端由 $\bar\gamma=-2\zeta u_1u_2/g$、$\bar\beta_1=\zeta u_2^2 L_1/g^2$ 及交换 1、2 所得。它们也规定了 $\bar\gamma=0$ 时的连续取值。对一般光滑平衡分支，非自旋双星至四极通量阶的结果为
\begin{align*}
 \kappa_1={}&12+5\bar\gamma
  -5\zeta S_-^2(3+\bar\gamma+2\bar\beta_++2\Delta\bar\beta_-)
  +10\zeta S_-(X+\Delta Y),\\
 \kappa_2={}&11+\tfrac{45}{4}\bar\gamma-40\bar\beta_+
  -5\zeta S_-^2(17+6\bar\gamma+\eta+8\bar\beta_++8\Delta\bar\beta_-)\\
 &+30\zeta S_-(3X-\Delta Y)-120\zeta Y^2 .
\end{align*}
最后一项也能由标量单极的平方产生。$X,Y$ 中出现二阶结构导数，是因为对偶极波再作时间求导时，有效荷会随伴星标量场改变；只把弱场 $\gamma,\beta$ 代入四极矩做不到这一点。在对称双星中 $S_-=0$，与荷质差相乘的项消失，$Y^2$ 却不一定消失，特别可影响偏心轨道的径向能流。若两星均取广义相对论的极限，则 $\zeta,\bar\gamma,\bar\beta_\pm\to0$，得回 $12,11$。

对标量-张量理论，可取 $S=(s_2-s_1)/\sqrt g$、$\kappa_D=4\zeta$。若 $s_1=s_2=1/2$，$g=1-\zeta$、偶极源为零，余下的张量能流和波形可用轨道所量到的重标质量换回广义相对论形式。对一颗黑洞、一颗中子星，相同重标使保守轨道在低阶仍与广义相对论相像，能流中却有 $\kappa_D S^2=Q$，$Q=\zeta(1-2s_{\rm NS})^2/(1-\zeta)$。这种配对的强场信息因而可集中在一个非负组合里。在此理论及该归一化下，零阶通量系数还给出 $\kappa_1=12-15Q/4$、$\kappa_2=11-5Q(17+\eta)/4$。可见一旦把相位算到次于偶极的阶，逐项只加一个偶极功率而保持其它辐射系数全不变，未必仍是自洽的理论计算。

若背景矢量也可辐射，通常先在远区将线性方程对角化。设某实本征模 $q_I$ 的二次作用量密度为 $K_I[\dot q_I^2-c_I^2(\nabla q_I)^2]/2$，这里吸收了归一化常数。乘波动方程以 $\dot q_I$ 再分部积分，可得能流 $\bm S_I=-K_I c_I^2\dot q_I\,\bm\nabla q_I$。远处出射波 $q_I=D^{-1}F_I(t-D/c_I,\bm N)$ 的径向通量即 $K_I c_I(\dot F_I)^2/D^2$。故知某个模式传播很快仍不足以算出损失，还须知道 $K_I$ 的符号和大小、源能否激发 $F_I$。正动能各模向外携能使束缚双星失去能量。历史上的 Rosen 双度规模型曾给出携负能的辐射，用同一定义有 $\kappa_1=-21/2,\kappa_2=-23/2,\kappa_D=-20/3$。对圆轨道甚至 $\dot E$ 的符号都与广义相对论相反，轨道应逐渐变大。脉冲星的近心点时刻不断提前，单凭损失的方向已构成强约束。对有正能矢量模式的理论，实际还须由运动星的敏感度计算矢量振幅，即使弱场静态质量给得很准也代替不了这一步。

\subsection{以太与叶层理论中的可测能流系数}

怎样用一个带优选时间方向的具体模型预言 $\dot P_b$？须把单星的运动敏感度、各辐射模式的归一化和两体轨道的质量单位配成一套。以下取无自旋，写 $\chi=c_{14}$、$d=2c_1-c_+c_-$，以 $w_2,w_1,w_0$ 表示张量、横矢量、标量模式的相速除以 $c$，它们的平方见下文。$G_{\rm ae}=G_{\rm N}(1-\chi/2)$ 为场作用量中的常数，而本对星的轨道耦合是 $G_{12}=G_{\rm N}g$。有些文献把单星的原始敏感度定义为 $-\p\ln m_a^{(0)}/\p\ln\Gamma$，与前文的 $\sigma_a$ 异号。为联用其波区结果，记 $t_a=-\sigma_a/(1-\sigma_a)$。于是 $g=(1-t_1)(1-t_2)$，而 $m_a=m_a^{(0)}(1)(1-\sigma_a)$ 仍是此处惯性质量。若把 $t_a$ 当作未经换算的 $\sigma_a$，即使只求最前面的偶极项，也会在强敏感度情况下混淆质量和荷。

定义 $\Sigma=(m_2t_1+m_1t_2)/m$、$\delta_t=t_1-t_2$ 与 $Z=(\alpha_1-2\alpha_2)(1-c_+)/[3(2c_+-\chi)]$。利用弱场关系还可写成 $Z=2(1-c_+)(c_++\chi+3c_2)/[3c_{123}(2-\chi)]$。把不同模式的贡献写成平方，可少引入六个单独的常数。以下 $z$ 只是多项式的无量纲自变量：
\begin{align*}
 A(z)={}&\frac1{w_2}
  +\frac{2\chi}{w_1}\left(\frac{c_+}{d}
                  +\frac{z}{2\chi w_1^2}\right)^2
  +\frac{3\chi}{2(2-\chi)w_0}
       \left(Z-1-\frac{2z}{3\chi w_0^2}\right)^2,\\
 B(z)={}&\frac{\chi}{4(2-\chi)w_0}
       \left(Z-\frac{2z}{3\chi w_0^2}\right)^2,\\
 C_D={}&\frac{4}{3\chi}
       \left[\frac1{w_1^3}+\frac{1}{(2-\chi)w_0^3}\right],
 \qquad D_V=\frac{1}{6\chi w_1^5}.
\end{align*}
$A$ 中三项依次来自张量、矢量、标量模。$B$ 在质心静止时只对有径向运动的源起作用，$C_D$ 是两星运动荷质差平方的偶极系数。取稳定、正能的传播分支，$\chi$ 及 $2-\chi$ 为正，能流因子也显出相应的正号。耦合趋零时，涉及 $1/\chi$ 的式子必须连同运动敏感度取极限，不能只保留形式上发散的系数。

以质心静止在优选系中的轨道为例，记 $\bm u=\bm v/c$、$u_r=\bm u\cdot\bm n$。保留领头偶极及已知的四极阶项，在普通单位中向外的能流与周期衰减为
\begin{align*}
 {\cal P}_{\rm out}&=\frac{G_{12}^{\,2}G_{\rm ae}\mu^2m^2}{c^3 r^4}\,{\cal F},\\
 {\cal F}&=\tfrac8{15}A(\Sigma)(12u^2-11u_r^2)
                   +4B(\Sigma)u_r^2+C_D\delta_t^2,\\
 \frac{\dot P_b}{P_b}&=-\frac{3a}{G_{12}\mu m}
                              \langle{\cal P}_{\rm out}\rangle .
\end{align*}
最后一式由 $E=-G_{12}\mu m/(2a)$ 和 Kepler 定律推出，平均按实际偏心轨道的一周求。含 $\bm u$ 的项相对偶极多两次 $1/c$，$G_{\rm ae}$ 来自波的作用量，两个 $G_{12}$ 来自将源的加速度换成轨道变量。对圆轨道，取 $u_r=0$、$u=(G_{12}mn_b/c^3)^{1/3}$，同质量、同观测 $n_b$ 的广义相对论预言相比，所得
\[
 \frac{\dot P_b}{\dot P_b^{\rm GR}}
  =\frac{G_{\rm ae}}{G_{\rm N}}g^{2/3}
    \left[A(\Sigma)+\frac5{32}C_D\,\frac{\delta_t^2}{u^2}\right].
\]
输入是理论耦合、两星在指定物态上的运动敏感度与质量。对相同运动敏感度的双星，$\delta_t=0$ 而 $\Sigma$ 仍在，故没有偶极不保证普通四极阻尼恰等于广义相对论。混合双星则常因 $u\ll1$ 而放大 $\delta_t^2/u^2$。独立的保守轨道实验还可涉及 $\sigma'_a$，在以上领头辐射系数中却不必另给这个二阶导数。

若质心相对优选系有速度 $c\bm W$，不可把它混入两星的相对速度。仍取 $\bm n$ 从 2 指向 1，令 $W_n=\bm W\cdot\bm n$。写 $A_3=A''/2$、$B_3=B''/2$，导数是上面两个二次多项式对 $z$ 的导数。在上述精度下，$\mathcal F$ 再加
\begin{align*}
 {\cal F}_{W}={}&\delta_t^2
 \left[\left(\tfrac{18}{5}A_3+2D_V\right)W^2
       +\left(\tfrac65 A_3+36B_3-2D_V\right)W_n^2\right]\\
 &+\delta_t\left[12 B'(\Sigma)W_n u_r
          +\tfrac85 A'(\Sigma)(3\bm W\cdot\bm u-2W_nu_r)\right].
\end{align*}
线性含 $\bm W$ 的项在圆轨道一周平均时消失，有些二次项却不消失。对偏心双星和希望保留方向信息的分析，应先带方向作平均。以上近似保留了可控制的领头偶极、主要四极及所示质心速度项。若要把偶极本身也推进相对 1PN，还会有依赖敏感度的源修正和反作用项；不能据此声称已求全任意强场双星该精度的矢量波形。这些能流系数按 Yagi 等的修订稿换算了敏感度的符号，见 \href{https://arxiv.org/abs/1311.7144v4}{arXiv:1311.7144v4} 式 (107)--(124)、(192)--(193)。

在叶层理论中取 $\chi=\alpha_K$、$c_+=\beta_K$、$c_2=\lambda_K$，$\mathcal Q_{ab}^{K},\mathcal R_{ab}^{K}$ 所用的弱场参数也随之更换。传播速度分别为 $w_2^2=1/(1-\beta_K)$ 和 $w_0^2=(2-\alpha_K)(\beta_K+\lambda_K)/[\alpha_K(1-\beta_K)(2+\beta_K+3\lambda_K)]$。保留 $A(z)$ 的第一、三项及整个 $B(z)$，取 $C_D=4/[3\alpha_K(2-\alpha_K)w_0^3]$、$D_V=0$，便得叶层理论的能流。这个取法消去横矢量波而保留标量偶极；将所有非张量辐射一起置零，会丢掉不对称脉冲星在该理论中最敏感的观测量。

\subsection{从六种潮汐形变到干涉仪读数}

探测器怎样识别波的偏振？自由质点的相对加速度取决于潮汐张量 $R_{0i0j}$。在广义相对论中，平面引力波只给出加、叉两种横向形变。某些度规理论还容许横截面各方向同时伸缩的标量模式、沿传播方向的伸缩及两种矢量模式。取来波沿局域正交标架的 $z$ 轴，记 $-\tfrac12\ddot H_{ij}=c^2 R_{\hat0 i\hat0 j}$。相邻自由质点若原先相对静止，扣除无波背景后有 $\delta\xi^i=H_{ij}\xi_0^j/2$，积分常数由波到来前的状态选定。于是所有可测形变可排列为
\[
 H_{ij}=
 \begin{pmatrix}
 h_{\rm b}+h_+&h_\times&h_{\rm x}\\
 h_\times&h_{\rm b}-h_+&h_{\rm y}\\
 h_{\rm x}&h_{\rm y}&h_{\rm l}
 \end{pmatrix}.
\]
把一圈质点放在横截面上，$h_{\rm b}$ 使圆保持为圆而交替扩缩，故称呼吸模式。$h_+,h_\times$ 则使圆变成椭圆，且轴向相差 $45^\circ$。要看出 $h_{\rm x},h_{\rm y}$，还须知道横向相邻的质点有无纵向相对运动。沿传播方向放置的质点之间的距离则能测到 $h_{\rm l}$。转动横截面的两根坐标轴，两个张量分量按两倍转角彼此混合，两个矢量分量按一倍转角混合，两标量分量各自不变。这里分类的是潮汐力的方向图；一个理论也可能使两项具有固定振幅比，从而只有一个独立的传播自由度。

例如长程标量在物质度规中造成共形扰动。在线性、局域平直近似下，取 $h_{\mu\nu}^{(s)}=b(t,z)\eta_{\mu\nu}$，则 $c^2 R_{\hat0 x\hat0 x}=c^2 R_{\hat0 y\hat0 y}=-\ddot b/2$，而 $c^2 R_{\hat0 z\hat0 z}=(c^2\p_z^2 b-\ddot b)/2$。若标量以 $c$ 沿 $z$ 传播，后者为零，探测器只看到呼吸模式。若标量有质量，满足 $\omega^2=c^2 k^2+\omega_0^2$，对单频波纵向与横向潮汐的振幅比为 $\omega_0^2/\omega^2$。此时观测到纵向响应并不意味着另有一个独立的纵向场，须连同频率依赖一并辨认。

长波极限下，等臂干涉仪的读数可写为 $h=(\delta L_1-\delta L_2)/L=F_+h_++F_\times h_\times$，这里先只写广义相对论的两项，$F_+,F_\times$ 由来波方向、偏振基及两臂朝向决定。设单位臂矢为 $\bm u,\bm v$，则 $D^{ij}=(u^iu^j-v^iv^j)/2$，任一偏振基张量 $e^A_{ij}$ 的方向因子为 $F_A=D^{ij}e^A_{ij}$。若波垂直入射，加偏振的两主轴恰沿两臂，则一臂伸长时另一臂缩短，差分读数最大。两臂若各自转过 $45^\circ$，它们对该偏振的投影相同，差分信号反而消失。等臂差分还有一个与方向无关的限制。取 $e^{\rm b}_{ij}=\delta_{ij}-N_iN_j$、$e^{\rm l}_{ij}=N_iN_j$，因 $D^{ij}\delta_{ij}=0$，总有 $F_{\rm b}=-F_{\rm l}$。许多台同类型仪器也不能仅凭长波下的差分振幅分开这两个标量。

实际用激光比较的是两臂光的往返相位。把镜间距离当作同一时刻的长度，要求臂长远小于波长；若一个光行时中场变化已显著，就须沿光的出入两段积分度规扰动。这使仪器的传递函数随频率、来波方向而变。给定某台仪器的方向因子和噪声，不能因某次事件在它的记录上很弱就断定波本身很弱；来波可能恰从它不灵敏的方向进入。不同台站的定时先给出传播方向的信息。若平面波沿 $\bm N$ 传播，两个台站的位置差为 $\Delta\bm x$，到达时刻差便是 $\bm N\cdot\Delta\bm x/c_g$。一个时差只限定天空上一条圆环，再添一条基线可缩小源区；已知宿主的位置则可反过来限制 $c_g$。

一台仪器只测得一个偏振组合。取向不同的仪器记录同一来波，对准到达时刻后比较各自振幅，便可区分不同的偏振假设。若假定只有两个张量偏振，把各台 $F_+$、$F_\times$ 看作两个台站空间中的列矢量，信号矢量必须落在它们张成的平面内。三台方向合适的仪器便可取一组权重，使两个张量贡献在加权和中同时消失。若其中仍有随同一波形的相干残差，需检验方向估计、仪器响应和额外偏振。若三台几乎平行，两列的有效独立性很差，这种检验自然较弱。例如 GW170814 同时被两台 LIGO 和一台 Virgo 干涉仪记录，在纯张量、纯矢量和纯标量三种假设中，实测的相对振幅更符合张量偏振。这项比较测试的是三个特定假设，尚不能给出任意混合的六个分量。

还可用脉冲星与地球构成很长的测量基线。设自地球指向脉冲星的单位矢为 $\bm p$、距离为 $L_p$。引力波沿射电信号的路径改变其到达频率，长时间累计又成为到达时差。对沿 $\bm N$ 传播的张量波，脉冲周期的相对变化可写成
\[
 z(t)=\frac{p^ip^j}{2(1+\bm N\cdot\bm p)}
 \left[h_{ij}(t)-h_{ij}\!\left(t-\frac{L_p}{c}(1+\bm N\cdot\bm p)\right)\right],
 \qquad R(t)=\int^t z(t')\,\d t'.
\]
括号第一项是在地球接收时的场，第二项是在脉冲发射处的场，两个场的自变量都已换算成波在地球处的相位。许多不同方向的脉冲星共用第一项，却通常不共用第二项。若有统计各向同性的张量波背景，多对脉冲星计时残差间的相关性应随两条视线的夹角而变。地面基准钟若有误差，各方向反而产生同相的残差；太阳系质心位置若有误差，在各视线留下的是位移矢量的投影。比较多对脉冲星而不只寻一条异常时间序列，就能分开这些来源。这里把脉冲星当作接收低频引力波的远端钟，下面则把它当作自身双星运动的示踪体。

\subsection{波速、零标架与偏振的分类}

上面的几个 $h_A$ 是在一组选定的标架中积分潮汐所得的分量。若只改变横截面内两根轴的方向，实验室对一个张量波究竟叫“加”还是“叉”会有分歧，但都同意它是横向四极形变。若允许观者在保持波矢不变的同时再作 Lorentz 变换，哪些关于额外偏振的陈述还能不依赖观者？为说明这一点，先只考虑沿 $z$ 轴的类光平面波，用局域正交标架的单位矢 $e_0,e_x,e_y,e_z$ 组成
\[
 \ell=\frac{e_0+e_z}{\sqrt2},\qquad
 k=\frac{e_0-e_z}{\sqrt2},\qquad
 m=\frac{e_x+\mathrm i e_y}{\sqrt2},\qquad
 \bar m=\frac{e_x-\mathrm i e_y}{\sqrt2}.
\]
它们均为零矢量，非零的内积只有 $\ell\cdot k=-1$、$m\cdot\bar m=1$。这里 $k$ 是与波矢反向的基矢，波矢本身沿 $\ell$。波各分量只随 $t-z/c$ 改变，因而沿波面切向的导数为零。将这一点代入线性微分 Bianchi 恒等式，辐射阶的 Riemann 分量可约化为 $R_{kpkq}$ 及同对称性联系的分量，其中 $p,q$ 各取 $\ell,m,\bar m$。实曲率的独立量又只剩六个实数，恰好与实验室的六个潮汐分量一样多。

用这种零标架取曲率分量的办法称 Newman--Penrose 表述。对这里的类光波，约定 $\Psi_2=-R_{k\ell k\ell}/6$、$\Psi_3=-R_{k\ell k\bar m}/2$、$\Psi_4=-R_{k\bar m k\bar m}$、$\Phi_{22}=-R_{kmk\bar m}$。前一个和最后一个为实数，中间两个各给两实分量。含两个横向 $\bar m$ 的 $\Psi_4$ 在绕 $z$ 旋转偏振基时乘以 $\exp(-2\mathrm i\theta)$，含一个的 $\Psi_3$ 乘以 $\exp(-\mathrm i\theta)$，其余两者不变。代回正交基，$\operatorname{Re}\Psi_4,\operatorname{Im}\Psi_4$ 分别量出两张量潮汐，$\Psi_3$ 量出两个混合纵横潮汐，$\Phi_{22}$ 量出横向迹，$\Psi_2$ 量出纵向潮汐。符号和数因子依零矢归一化而变，有无相应曲率模式则可由足够多方向的质点相对加速度直接检验。

保持一个给定零波矢的 Lorentz 变换所成的群同二维 Euclid 群 $E(2)$ 同构。除了绕传播方向旋转，还可作保持 $\ell$ 不动的零旋转，例如对复参数 $b$ 取 $m'=m+b\ell$、$k'=k+\bar b\,m+b\,\bar m+|b|^2\ell$。内积不变，却改变其余基矢的投影。把变换后的基代入曲率便见，$\Psi_2$ 不变，$\Psi_3$ 可获得正比于 $\Psi_2$ 的项，$\Psi_4$ 又可从原来的 $\Psi_3,\Psi_2$ 获得贡献；横向迹也可受较低螺旋度的分量混入。故在某一观者处测得纯纵向潮汐，不意味着所有标准观者都会把同一波描述为纯纵向。可靠的分类须先指出从哪一级起有不可变为零的分量。

\begin{center}
\begin{tabular}{lll}
\hline
类别 & 分类所需的非零分量 & 必为零的分量\\ \hline
$\mathrm{II}_6$ & $\Psi_2$ & 无\\
$\mathrm{III}_5$ & $\Psi_3$ & $\Psi_2$\\
$\mathrm N_3$ & $\Psi_4,\Phi_{22}$ & $\Psi_2,\Psi_3$\\
$\mathrm N_2$ & $\Psi_4$ & $\Psi_2,\Psi_3,\Phi_{22}$\\
$\mathrm O_1$ & $\Phi_{22}$ & $\Psi_2,\Psi_3,\Psi_4$\\
$\mathrm O_0$ & 无 & 全部\\ \hline
\end{tabular}
\end{center}
下标是这类波在标准观者间最多可呈现的实潮汐分量数，并非某一理论必须以独立初值激发的自由度数。用一个理论所允许的一般波给它分类，一次特定事件因源对称性欠激发某模式并不矛盾。广义相对论中真空 $R_{\mu\nu}=0$，逐项消掉 $\Phi_{22},\Psi_2,\Psi_3$，剩下 $\mathrm N_2$。线性无质量标量波满足 $\square\delta\varphi=0$，在物质度规中却有 $R_{\mu\nu}\propto\p_\mu\p_\nu\delta\varphi$；其 $kk$ 分量可非零，故与通常张量波叠加可给 $\mathrm N_3$。对于严格非类光的波矢，小群已变，上表不能原封不动使用。波速极接近 $c$ 时可以近似按此表辨认主项，能实际运用的始终是局域曲率和仪器响应。

波速本身还怎样从场方程读出？取扰动的波长远小于背景变化尺度，保留方程中对扰动的二阶导数，设一切模式都正比于 $\exp[\mathrm{i}(\bm k\cdot\bm x-\omega t)]$。系数矩阵有非零解需要其在消去规范与约束后的行列式为零；不同因子就给出不同的特征锥。以单位类时矢量理论为例，在平坦、矢量静止的背景取 $c=1$，写 $c_{13}=c_1+c_3$、$c_{14}=c_1+c_4$、$c_{123}=c_1+c_2+c_3$。线性张量、横矢量和纵向标量模式的平方速度分别为
\begin{align*}
 v_T^2&=\frac{1}{1-c_{13}},&
 v_V^2&=\frac{2c_1-c_1^2+c_3^2}{2c_{14}(1-c_{13})},\\
 v_S^2&=\frac{c_{123}(2-c_{14})}
    {c_{14}(1-c_{13})(2+c_{13}+3c_2)} .
\end{align*}
表达式以相应模式有非退化动能为前提；分母趋零时应回原方程检查约束和强耦合，不能一概说有可观测的无穷波速。要求平方速度为正只保证小扰动有实传播频率，正能量还要另查二次作用量。高能粒子若快过某个与之耦合的引力模式，运动学上可向该模发出类似 Cherenkov 的辐射，长程传播而未显著失能的粒子可提供另一类界限。这个推论取决于耦合和粒子能量，而张量波与电磁波直接比较到达时刻，首先限制的只是 $v_T$。同一实验给出的 $c_{13}$ 约束并不能代替对其它两个锥的观测。

\subsection{把潮汐力转换成读数}

最直接的探测器未必是两块彼此分离的反射镜。设一根两端近乎自由的圆柱棒有纵向本征振动，以 $q(t)$ 表示某模的幅度，$L_{\rm eff}$ 为其位移图样同潮汐力沿棒积分所得的有效长度。若波长远大于棒长，取棒轴为 $\bm u$、投影应变 $h_\parallel=u^iu^jH_{ij}$，线性受迫方程为
\[
 \ddot q+\frac{\omega_0}{Q_{\rm m}}\dot q+\omega_0^2q
        =\frac{L_{\rm eff}}{2}\ddot h_\parallel .
\]
近共振的持续单频波，振幅可提高约 $Q_{\rm m}$；响应带宽却约为 $\omega_0/Q_{\rm m}$，达到稳态又需约 $Q_{\rm m}/\omega_0$ 的时间。持续时间远短于此的脉冲即使频率合适，也只能激发一段由仪器自由衰减的振铃。热运动同样激发这个本征模。两座远隔的棒在正确光行时范围内出现一致的能量激发，可排除多数单台热涨落，却仍要以错开时间的记录估计偶然符合率。用的是实体的本征振动，故信号的带宽、能量和到达时刻要按机械响应反推。

激光干涉仪则用电磁场作位移的读出。设等臂长 $L$、激光波长 $\lambda_{\rm L}$，忽略光学腔时，一次往返的相位差变化为 $\delta\phi_{\rm L}=4\pi(\delta L_1-\delta L_2)/\lambda_{\rm L}$。这项测量不要求各处先放上同步的钟，两束光从同一分光镜出发并在同一点重新比较。地面长臂仪器以悬挂的反射镜近似自由质点：高于悬挂本征频率时，镜间的相对位移主要反映潮汐；很低频处地面运动和悬挂的回复力已不能略。两臂光学腔可使光多次往返，把小的相位差积累起来，但腔的储光时间也限制仪器追随快速应变的能力。若不把实测输出除以这一有频率依赖的响应函数，峰值出现的时间和相位都可能出现人为偏差。

自由镜在 TT 坐标中的坐标位置可保持不变，是否与“臂长改变”矛盾？用同一坐标沿未扰动光路算往返时刻即可检验。例如垂直入射时，沿一臂的光速由零线元给出 $c\,\d t\simeq[1+h_\parallel(t)/2]\d\ell$。对出发时刻为 $t_0$ 的光，附加的往返时间近似为 $\tfrac12\int_{t_0}^{t_0+2L/c}h_\parallel(t')\,\d t'$。低频极限是 $Lh_\parallel/c$，恰等于往返两次穿越增量 $Lh_\parallel/2$ 的光行时。对单频波，积分多出 $\sin(\omega L/c)/(\omega L/c)$ 与中点相位；光行期间先拉长又缩短的贡献可部分抵消。有限臂长的仪器不宜把所有模式一律用一个与频率无关的方向因子概括。

将天空位置与偏振基也写入响应，有助于理解为何把 Virgo 加入两台 LIGO 后就能作新的比较。取等长直角两臂沿 $x,y$，令 $\bm N$ 在仪器系的极角、方位角为 $\theta,\phi$，用球面基 $\bm e_\theta,\bm e_\phi$ 作横向偏振轴。在前面采用的呼吸、纵向模式归一化下，逐一将 $D^{ij}$ 与偏振张量缩并，有
\begin{align*}
 F_+&=\tfrac12(1+\cos^2\theta)\cos2\phi,&
 F_\times&=-\cos\theta\sin2\phi,\\
 F_{\rm x}&=\sin\theta\cos\theta\cos2\phi,&
 F_{\rm y}&=-\sin\theta\sin2\phi,\\
 F_{\rm b}&=-\tfrac12\sin^2\theta\cos2\phi,&
 F_{\rm l}&=-F_{\rm b}.
\end{align*}
若把偏振基再转 $\psi$，加叉两列按 $2\psi$ 旋转，两个矢量模式则按 $\psi$ 旋转。同一台仪器随地球自转，天空中的同一方向就扫过这些图样的不同部分。对持续很久的窄带源，可利用这种日调制补充偏振信息；对只有零点几秒的并合，地球来不及转动，几何多样性主要来自不同地点仪器的取向。

对已知方向的短脉冲，还可显式构造\textbf{零流}(null stream)。将各台记录校正到统一的到达时刻和相位，以噪声谱白化，记三个读数为列矢量 $\bm d$，张量响应列为 $\bm f_+,\bm f_\times$。若二者不平行，$N=(\bm f_+\times\bm f_\times)\cdot\bm d$ 不含任何形状的纯张量波。逐频率应用时，复数响应宜用矩阵的 Hermite 转置处理。对多台站可用 $P_{\rm null}=1-F(F^\dagger F)^{-1}F^\dagger$ 投影至零空间，这要求保留的响应列独立；近线性相关时需用奇异值分解处理。若零流的能量恰和噪声水平相符，说明记录可由两个张量模式及噪声解释。若要进一步分辨额外的一个模式，还应检查它的响应列在这个零空间是否有非零投影；落在原来两列张成平面内的分量无法由这种实验单独排除。

基线比地面臂长更长时还有别的读出方法。航天器的双程无线电链路以地面原子钟发射信号，航天器作相干转发，地面再测瞬时频移。地球、航天器和先后两个传播方向都采样同一波在不同事件的值；往返光行时长到与波周期相近时，频移成为几个错开时刻的应变之差。空间激光干涉的各臂若有不同光行时，同一个激光的频率噪声就不再自动消去。已知各臂光行时后，可把若干单程读数依各自延迟平移、相加减，使同一个发射事件的频率噪声在两条合成路径等次数出现，这就是时间延迟干涉的出发点。与它相比，前述脉冲星计时不掌握远端激光相位，而依靠多个天然的长时钟以及它们之间的角向相关。

例如对统计稳定、各向同性、非偏振的张量背景，若两颗不同脉冲星足够远且距离不相干，脉冲星端的交叉项通常平均消失。地球端保留的相关正比于两个方向因子的乘积对全天积分。令 $x=(1-\cos\Theta)/2$，$\Theta$ 是两条视线的夹角，取通常归一化后，角关联为 $\Gamma(\Theta)=\tfrac12-\tfrac14x+\tfrac32x\ln x$。$\Theta\to0$ 的不同脉冲星对相关为 $1/2$；同一颗脉冲星自相关还另含自身的远端项。此关系的非单调形状与钟误差在全天给出的常数相关不同，也与历表误差主要产生的偶极相关不同。对非各向同性背景、单个明亮源或标量/矢量偏振，则要重新作角积分。再降到接近宇宙尺度的频率，张量扰动还可在最后散射附近给光子分布产生四极各向异性，并经 Thomson 散射转为宇宙微波背景偏振。实验上须从尘埃偏振、透镜将偶宇称偏振转换出的奇宇称分量中分离原初信号，这是另一种依靠电磁光子探测引力波的方法。

\section{强场、相对论性引力的动力学检验}

\subsection{双星脉冲星的计时}

一颗自转的中子星若以射电波束周期性地扫过地球，就给了我们一座随轨道运动的钟。把各脉冲的到达时刻换算至太阳系质心，扣除自转所给的脉冲序列，剩下的时差中，最大的一项来自脉冲星在轨道上的远近变化。由此可测出轨道周期 $P_b$、离心率 $e$、投影半长轴 $a_1\sin i$ 等 Kepler 根数。长期计时还能分辨近心点进动率 $\dot\omega$、脉冲钟因引力红移和运动而产生的周期性时差振幅 $\gamma_{\rm E}$，以及轨道周期变化率 $\dot P_b$。这些量称为\textbf{后 Kepler 参数}，其中 $\gamma_{\rm E}$ 具有时间量纲。

怎样从一串到达时刻定出这么多参数？单个射电脉冲的形状可有相当大的涨落。观者把一定时间内的许多脉冲按自转周期折叠，和稳定的平均轮廓对准，由此定出某个脉冲相位到达天线的时刻；再设法在两次观测之间也数准转过的整圈数。设 $\tau$ 是脉冲星的固有时，则脉冲编号可写为 $N=N_0+\nu_0(\tau-\tau_0)+\dot\nu_0(\tau-\tau_0)^2/2+\cdots$。常数和缓慢自转变慢的项由整段资料拟合。射电波穿过星际等离子体还多一个近似正比于射电频率平方倒数的群延迟，比较多个射电频段便能将它和几乎不随射电频率改变的轨道效应分开。地球公转、自转也改变接收位置，故还要把地面记录的时刻按地球历表和太阳系内光传播换算到太阳系质心。若源方向估错，地球公转的投影会在残差中留下一年周期项，反过来可改善源的位置估计。

以下记双星质心系内的发射坐标时为 $t_e$，把观测时刻扣去源距离产生的常数，记为 $t_{\rm arr}$。从星上的钟读数到观者的钟读数，可分为两个问题：何时发出的信号，以及信号从哪个地点出发、沿怎样的路径到达。采用离开观者的位移取正的约定，把最低几项写成 $t_{\rm arr}=\tau+\Delta_{\rm E}+\Delta_{\rm R}+\Delta_{\rm S}+\text{常数}$。各轨道函数在发射时计算，实际拟合时可迭代求 $t_e$；若全部改在到达时刻求值，展开中就要多出光行时乘轨道速度的修正。

先去掉引力红移，单看几何。记 $n_b=2\pi/P_b$，偏近点角 $u_b$ 由 $u_b-e\sin u_b=n_b(t_e-T_0)$ 确定，$T_0$ 是一次近心点时刻。沿轨道长、短轴的质心坐标分别正比于 $\cos u_b-e$、$\sqrt{1-e^2}\sin u_b$。将两者绕近心点角 $\omega$ 旋转，再投影到视线，得到\textbf{R\o mer 时延}
\[
 \Delta_{\rm R}=x\,F(u_b),\qquad
 F(u_b)=\sin\omega(\cos u_b-e)+\sqrt{1-e^2}\cos\omega\sin u_b,
 \qquad x=a_1\sin i/c .
\]
可见一颗脉冲星的计时只给出 $x$。与 Kepler 定律合用，可测的质量函数为 $(n_b^2/G)(cx)^3=(m_2\sin i)^3/(m_1+m_2)^2$。单凭它不能分别给出两质量和倾角。若轨道近乎圆形，近心点本身也难定位，直接拟合偏心矢量的两分量通常比拟合很小的 $e$ 与不定的 $\omega$ 承担更少的退化。

再来看星上的钟。设 $r=a(1-e\cos u_b)$ 为两星距离，脉冲星的质心速度为 $v_1$。扣去自身静态场所给的常量走时率后，$\d\tau/\d t_e=1-[Gm_2/r+v_1^2/2]/c^2+\cdots$。前一项为伴星的引力红移，后一项为运动钟慢。用 $v_1=(m_2/m)v_{\rm rel}$ 和 $v_{\rm rel}^2=Gm(2/r-1/a)$，两效应随时间改变的部分都正比于 $1/r$。由于 $\d t_e=(1-e\cos u_b)\d u_b/n_b$，积分 $1/r$ 给出 $u_b/(a n_b)$。再由 Kepler 方程减去匀速部分，留下随 $\sin u_b$ 改变的\textbf{Einstein 时延} $\Delta_{\rm E}=\gamma_{\rm E}\sin u_b$，其中
\[
 \gamma_{\rm E}=
 \frac{e}{n_b}\left(\frac{Gmn_b}{c^3}\right)^{2/3}
 \frac{m_2}{m}\left(1+\frac{m_2}{m}\right).
\]
这说明一个严格的圆轨道为何难测此量。红移和钟慢都仍存在，但整圈相同的走时率可重新定义为脉冲星的自转频率，观者不能单独量出它。

最后算传播路径。取伴星为弱场的瞬时中心，将其引起的光行时沿未扰动路径积分，在略去与轨道相位无关的项后，得 $\Delta_{\rm S}=-2r_s\ln[1-e\cos u_b-sF(u_b)]$，其中 $r_s=Gm_2/c^3$、$s=\sin i$。这里 $r_s$ 有时间量纲，和上节星体表面的径向坐标无关。若 $i$ 接近直角，脉冲在上合时贴近伴星，括号变小，时延出现尖峰；其宽度与最大高度不同程度地依赖于 $s,r_s$。小倾角近圆轨道的峰则很平缓，把对数按 $s$ 展开，最低次项与几何时延的正弦变化很相似，不易分离。由此可知，只给一颗伴星一个较大的质量，未必足以产生能从数据中辨认的 Shapiro 信号。

\subsection{换引力理论时计时公式的哪一项会变？}

计时有一层分工。首先用一套带有自由振幅、形状、进动率等参数的公式拟合到达时刻，这时无需先说脉冲星由哪条状态方程支撑；其次才由某种引力理论把拟合所得的后 Kepler 参数解释为质量、结构和耦合。在第二步用了广义相对论，并不意味着每一个已经拟合的时延都依赖广义相对论。但若新理论产生原模型根本没有的周期项，或显著改变两个已用参数的相关性，最好重新拟合脉冲到时，不能只把沿用原理论定出的质量换个名字。

为了区分这些量，仍取 $G_{\rm N}=c=1$。记伴星的远场 $g_{00}=-1+2\alpha_2^*m_2/r+\cdots$、$g_{ij}=(1+2\gamma_2^*m_2/r)\delta_{ij}+\cdots$，并记两体轨道的耦合为 $g=\mathcal G_{12}$。$\alpha_2^*$ 决定脉冲钟在伴星场中变慢，$g$ 决定它以什么速度走过这条轨道，二者原则上须分别保留。若局域引力强度还随伴星势按 $\delta G_{\rm loc}/G_{\rm loc}=-\eta_2^*m_2/r$ 改变，取 $\kappa_I=-\p\ln I_1/\p\ln G_{\rm loc}$ 为固定重子数、自转角动量时的转动惯量响应。在缓慢变化下 $\delta\nu/\nu=-\delta I/I=\kappa_I\delta G_{\rm loc}/G_{\rm loc}$。它在自转相位中引入一个 $\int\delta\nu\,\d t$，同红移时延的那一项具有相同的 $\sin u_b$。沿用前面给的时延符号，结果为
\begin{align*}
 \gamma_{\rm E}&=
  e m_2(gmn_b)^{-1/3}
       \big[\alpha_2^*+g m_2/m+\kappa_I\eta_2^*\big],\\
 r_s&=\tfrac12(\alpha_2^*+\gamma_2^*)m_2,\qquad
 s=\sin i .
\end{align*}
Shapiro 系数用到的是伴星的时间势和空间势之和，这是由光走零测地线决定的。两个形状参数都来自信号的几何投影时，改变引力理论仍会经质量函数、轨道倾角间接改变所画的质量约束带。

例如在上节的无质量标量理论中，有 $\alpha_2^*=1-2\zeta s_2$、$\gamma_2^*=1-2\zeta+2\zeta s_2$。单看走时率，伴星结合能已在 $\alpha_2^*$ 中出现；两者相加却给 $r_s=(1-\zeta)m_2$，在此阶伴星敏感度恰好抵消。此消去不能推到所有计时量上。对于近心点进动率，要用带 $\bar\gamma,\bar\beta_\pm$ 的二体式；对于周期变化率，要加上椭圆轨道平均后的偶极、标量与张量通量。因此即使总质量相同，两颗 $\dot\omega$ 相等的候选模型在其余参数的平面中也可分开。三个后 Kepler 量对广义相对论的两个质量已给过多约束，对含额外耦合和未知敏感度的理论，却还要数清这些额外未知量。

实际拟合怎样知道某一项真的被测到了？以近圆轨道的 Shapiro 信号为例，令轨道相位为 $\Phi_b$ 并近似 $F=\sin\Phi_b$。$-2r_s\ln(1-s\sin\Phi_b)$ 的第一谐波部分可混入 $x\sin\Phi_b$，第二谐波也可与小离心率的几何项混合。若只有前两项在噪声以上，只调 $x,e,\omega$ 已可得到外观良好的拟合。倾角足够大时，第三及更高谐波能长成明显的尖峰，此时单独的传播时延才易于识别。类似地，$\gamma_{\rm E}\sin u_b$ 与几何时延的某一正弦项在单圈中有相关性；若 $\omega$ 在多年间显著进动，轨道几何的正弦、余弦系数一起改变，而 $\gamma_{\rm E}$ 的红移项按另一规律改变，长期跟踪就可分离它们。

若脉冲自转轴有缓慢进动，射电束的视线截面也在变。此时两种时刻漂移须分清。一种是星体作为钟真的改变走时或轨道改变，另一种是轮廓形状改变而观者拿固定模板对准了不同的束相位。不同射电频率、食与非食的时段、总强度和偏振轮廓，都可帮助检查这种变化。地球位置或星际色散的拟合有误也会给出周期结构，但一项多随地球公转，一项多随射电频率的平方倒数，双星几何时延则随其自身轨道相位。不是看见残差中有一个正弦就足以称为 Einstein 时延，还要看它相对于哪个时钟、哪一条轨道有固定相位。

\subsection{近心点的时刻与质量平面}

以 Hulse--Taylor 双星脉冲星 B1913+16 为例。给定 $P_b,e$ 后，广义相对论的 $\dot\omega$ 决定总质量，$\gamma_{\rm E}$ 再给出两颗星各自的质量。至此，预言第三个量 $\dot P_b$ 已无可调的质量参数。具体地，$\dot\omega=3n_b(Gmn_b/c^3)^{2/3}/(1-e^2)$，给定观测值后，$m_1$--$m_2$ 平面上的允许区是一条近似直线 $m_1+m_2=\text{常数}$。与之相交的 $\gamma_{\rm E}$ 带是弯曲的，两者测量误差给出交叠区的宽度。按该系统已有的一组计时解，轨道周期约 $7.75$ 小时，$e\simeq0.617$，近心点每年进动约 $4.23^\circ$。两带交会给出约 $1.438M_\odot$ 和 $1.390M_\odot$。这些数是以广义相对论解释两个后 Kepler 参数所得的质量。若将理论改换成带标量荷的理论，各条带都应依新理论重画。

把四极辐射功率按一个椭圆轨道周期平均，结合 $\dot P_b/P_b=3\dot a/(2a)$，得
\[
    \dot P_b=-\frac{192\pi}{5}
    \left(\frac{2\pi G\mathcal M}{P_bc^3}\right)^{5/3}
    \frac{1+73e^2/24+37e^4/96}{(1-e^2)^{7/2}} .
\]
代入上述参数，预言 $\dot P_b^{\rm GR}\simeq-2.40\times10^{-12}$。这相当于一年只少约 $76$ 微秒。若每次只拿相邻两次近心点的间隔相减，细小的周期变化很难看出。实际可一直数经过了多少周。设起始周期为 $P_{b0}$，在观测时间内 $\dot P_b$ 几乎不变，则经过 $k$ 周的时刻相对于固定周期预报的差约为 $\dot P_bP_{b0}k^2/2$。由于 $\dot P_b<0$，第 $k$ 次近心点提前，这项时差按观测年数的平方积累，几十年内已可达数十秒。

但太阳系与双星并非在同一个惯性系中静止。若沿视线远离观者的速度为 $v_\parallel$，则一阶 Doppler 效应使观测周期多一个 $1+v_\parallel/c$ 的因子。对两地距离 $D$ 求两次时间导数，并区分视向、横向，得到
\[
 \left(\frac{\dot P_b}{P_b}\right)_{\rm obs}
 =
 \left(\frac{\dot P_b}{P_b}\right)_{\rm int}
 +\frac{(\bm a_{\rm bin}-\bm a_\odot)\cdot\hat{\bm n}}{c}
 +\frac{\mu_{\rm pm}^2D}{c}.
\]
这里 $\hat{\bm n}$ 从太阳系质心指向双星，$\mu_{\rm pm}$ 是天球自行的角速度，最后一项即 $v_\perp^2/(cD)$。即使源沿直线匀速横过视线，到观者的距离也有正的二阶导数，这叫作 \textbf{Shklovskii 效应}。它总使周期看起来增长，银河引力的投影项却可取两种符号。脉冲自转周期也经受同一相对改正。一个小数位该不该用于引力实验，取决于对距离、自行和银河势的了解能否支持它。

扣除银河运动等贡献后，2016 年的 B1913+16 计时结果给出 $\dot P_b^{\rm int}/\dot P_b^{\rm GR}=0.9983\pm0.0016$。这项累积的时差给出了双星辐射能量的间接证据。反过来，若先接受广义相对论的辐射预言，也可把周期残差用于估计双星的距离或银河加速度。这两种用法依赖的先验不同，不能既取计时所得的距离修正同一条 $\dot P_b$，又把修正后的符合程度当作独立检验。

图~\ref{fig:exp-b1913-mass} 可直接用上述三式画出。横轴取脉冲星质量，纵轴取伴星质量。先由 $\dot\omega$ 定出总质量 $2.82836M_\odot$，再以 $\gamma_{\rm E}$ 定出两星约为 $1.4379M_\odot$ 和 $1.3905M_\odot$。将此质量对代入四极公式，得 $\dot P_b^{\rm GR}=-2.40263\times10^{-12}$。图上第三条线取自独立的轨道衰减测量。应比较的是三项约束的允许区域，不能要求三条中心线精确交于一点。

\begin{figure}[htbp]
    \centering
    % P_b=0.322997448918 d, e=0.6171340, GM_sun/c^3=4.925490947 us.
    % For each m1 solve gamma_E(m1,m2)=4.307(4) ms and
    % -Pbdot(m1,m2)=2.398(4)e-12 for m2. Shading uses +/-2 errors.
    \begin{tikzpicture}
        \begin{axis}[
            width=.82\linewidth,height=5.5cm,
            xmin=1.422,xmax=1.454,ymin=1.374,ymax=1.406,
            xtick={1.425,1.435,1.445},ytick={1.38,1.39,1.40},
            scaled ticks=false,
            xticklabel style={/pgf/number format/fixed,/pgf/number format/precision=3},
            yticklabel style={/pgf/number format/fixed,/pgf/number format/precision=2,
                             /pgf/number format/fixed zerofill},
            tick label style={font=\footnotesize},tick align=outside,
            xlabel={脉冲星 $m_1/M_\odot$},ylabel={伴星 $m_2/M_\odot$},
            label style={font=\small},
            legend style={at={(.5,1.02)},anchor=south,draw=none,font=\footnotesize,
                          inner sep=1pt,column sep=3pt},legend columns=3,
        ]
            \addplot[forget plot,draw=none,fill=blue!14] coordinates {
                (1.422,1.3970876) (1.426,1.3931847) (1.430,1.3893079) (1.434,1.3854570) (1.438,1.3816317) (1.442,1.3778317) (1.446,1.3740568) (1.450,1.3703068) (1.454,1.3665814)
                (1.454,1.3775035) (1.450,1.3812647) (1.446,1.3850507) (1.442,1.3888619) (1.438,1.3926984) (1.434,1.3965605) (1.430,1.4004485) (1.426,1.4043626) (1.422,1.4083031)
            } \closedcycle;
            \addplot[forget plot,draw=none,fill=orange!23,fill opacity=.75] coordinates {
                (1.422,1.3827638) (1.426,1.3840674) (1.430,1.3853690) (1.434,1.3866688) (1.438,1.3879666) (1.442,1.3892626) (1.446,1.3905567) (1.450,1.3918489) (1.454,1.3931393)
                (1.454,1.3983011) (1.450,1.3970069) (1.446,1.3957109) (1.442,1.3944131) (1.438,1.3931133) (1.434,1.3918117) (1.430,1.3905082) (1.426,1.3892027) (1.422,1.3878954)
            } \closedcycle;
            \addplot[black,thick,no marks] coordinates {(1.422,1.4063603) (1.454,1.3743603)};
            \addlegendentry{$\dot\omega$}
            \addplot[orange!65!black,thick,dashed,no marks] coordinates {
                (1.422,1.3853294) (1.426,1.3866348) (1.430,1.3879383) (1.434,1.3892400) (1.438,1.3905397) (1.442,1.3918376) (1.446,1.3931336) (1.450,1.3944277) (1.454,1.3957199)
            };
            \addlegendentry{$\gamma_{\rm E}$}
            \addplot[blue!65!black,thick,dashdotted,no marks] coordinates {
                (1.422,1.4026924) (1.426,1.3987707) (1.430,1.3948752) (1.434,1.3910058) (1.438,1.3871621) (1.442,1.3833439) (1.446,1.3795509) (1.450,1.3757829) (1.454,1.3720396)
            };
            \addlegendentry{$\dot P_b^{\rm int}$}
        \end{axis}
    \end{tikzpicture}
    \caption{B1913+16 在广义相对论中的质量约束，放大至两质量各约 $0.03M_\odot$ 的范围。橙、蓝带分别由 $\gamma_{\rm E}=4.307(4)$\,ms 与 $\dot P_b^{\rm int}=-2.398(4)\times10^{-12}$ 的中心值上下各取两倍标准误差得到。$\dot\omega$ 的误差带窄于图中线宽。}
    \label{fig:exp-b1913-mass}
\end{figure}

双脉冲星 J0737$-$3039A/B 的两颗星都曾测到脉冲，故可直接比较两条轨道的大小。用光行时表示两条轨道的投影半长轴，记 $x_A=a_A\sin i/c$、$x_B=a_B\sin i/c$。由质心条件，到 1PN 精度有 $m_A/m_B=x_B/x_A$。倾角 $i$ 在比值中消去。此系统的视线又几乎落在轨道面上，脉冲从伴星附近经过时的 Shapiro 时延尤其显著。在广义相对论中，时延的幅度尺度为 $Gm_{\rm c}/c^3$，形状参数为 $\sin i$，其中 $m_{\rm c}$ 是伴星质量。把这些独立的约束画在两星质量的平面上，各条允许带应交于同一区域；若不相交，再调整质量也不能使理论与计时相容。

双脉冲星还有一只可看的“钟针”。A 的射电束在经过 B 的磁层时会受遮掩，食的细致结构随 B 的自转相位而变。如果 B 的自转轴随轨道场发生测地进动，食的图样在多年间也会逐渐改变。由这种长期变化可估计轴的进动率，所用的是脉冲轮廓和遮掩几何，与以往返光行时或轨道衰减测得的参数有不同的误差来源。源的进动还会使波束渐渐扫离视线。看不到它不必意味着星停止自转，原来用来观测它的几何条件已经改变。

\subsection{为什么还需要不对称双星？}

两颗质量相近的中子星都有很强的自引力，若它们的标量荷质比也相近，则偶极辐射被差的平方压低。在这类双星中看到精确的四极阻尼，仍可容许两星各自带有相当的标量荷。若伴星是自引力较弱的白矮星，荷质比的差异可以大得多。对前述无质量的长程标量，暂取 $A_0=1$，并忽略更高阶的荷变化。由 $\ddot{\bm D}_\varphi=\mu(\alpha_1-\alpha_2)\ddot{\bm x}$ 及最低阶反平方吸引力，可知偶极能流的轨道平均涉及 $\langle 1/r^4\rangle$。用偏近点角积分，得到 $\langle 1/r^4\rangle=(1+e^2/2)/[a^4(1-e^2)^{5/2}]$。于是偶极对周期衰减的贡献为
\[
 \dot P_b^{D}=-\frac{4\pi^2G_*\mu}{P_b c^3}
       (\alpha_1-\alpha_2)^2
       \frac{1+e^2/2}{(1-e^2)^{5/2}} .
\]
推导中取 $E=-G_{12}\mu m/(2a)$、$n_b^2a^3=G_{12}m$；同一个保守力常数在消去 $a$ 后抵消。圆轨道极限下，偶极衰减随 $P_b^{-1}$，四极衰减随 $P_b^{-5/3}$。若有若干轨道周期、组分质量均不同的系统，能检验的就不只是某一颗星的总损失率。

以低离心率的中子星--白矮星 J1738+0333 为例，其轨道周期约 $8.5$ 小时。可见伴星够亮，光谱线的轨道 Doppler 振幅与脉冲计时所得的质心速度振幅一比，就给出质量比。再由白矮星的大气分析及质量--半径关系估计其质量，便得到白矮星约 $0.181M_\odot$、脉冲星约 $1.46M_\odot$。这条定质量的途径不必先假定四极辐射率。2012 年的分析扣去自行和银河加速度后给出 $\dot P_b^{\rm int}=(-25.9\pm3.2)\times10^{-15}$，与按两质量预测的四极衰减相容。将残值及质量的不确定度代入上式，即可限制 $|\alpha_{\rm NS}-\alpha_{\rm WD}|$。若标量场有质量、伴星受照后显著失去物质，或另有进入同一残值的衰减机制，这种直接换算就须改动。

不同质量的中子星还有另一层用处。标量化若有星体紧致度门槛，一颗普通质量的中子星过关，不能代替一颗高质量中子星的检验。J0348+0432 也有白矮星伴星，光谱与计时合用所得的脉冲星质量为 $2.01\pm0.04M_\odot$（2013 年）。它的可取之处正在于试探了另一段质量范围。实际用这一族双星限制标量理论时，可对每个候选状态方程、每组背景耦合先解星体边值问题，给出 $\alpha_A(m_A)$，再用各系统的质量和轨道计算 $\dot P_b$。不同系统可排除曲线上的不同部分。把质量、状态方程、耦合全都忘掉，声称一次计时给出了一个与理论无关的“标量荷上限”，会失去这种检验的物理含义。

\subsection{把几种双星的记录放在一起}

下面选取几套已完成长期计时的历史结果作比较。同一张表中取同一去除运动学效应的 $\dot P_b^{\rm int}$，不把原始到时拟合出的 $\dot P_b^{\rm obs}$ 混在其中。最后一列按 $P_b/|\dot P_b^{\rm int}|$ 换算为周期发生显著相对变化的时间尺度；它并非用 $\dot P_b$ 不变的假定外推出并合时间。数字只取足以分辨各组的位数。
\begin{center}
{\small
\begin{tabular}{lcccc}
\hline
系统 & $P_b$/小时 & $e$ & $\dot P_b^{\rm int}/10^{-12}$ &
 $P_b/|\dot P_b^{\rm int}|$/年\\ \hline
B1913+16       &7.75&0.6171&$-2.398(4)$ & $3.7\times10^8$\\
J0737$-$3039A/B&2.45&0.0878&$-1.25(2)$  & $2.2\times10^8$\\
J1738+0333     &8.51&$3.4\times10^{-7}$&$-0.0259(32)$&$3.8\times10^{10}$\\
J1141$-$6545   &4.74&0.1719&$-0.401(25)$&$1.3\times10^9$\\ \hline
\end{tabular}}
\end{center}
前两行是中子星对，后两行含白矮星，括号表示末位误差。B1913+16 与 J1738+0333 的周期相似，$\dot P_b$ 却差约两个数量级。前者既有两颗较重的星，又有颇大的离心率；就算没有任何额外辐射，普通四极功率也应比后者大。反之，在较小的总衰减背景中寻找中子星与白矮星的偶极，虽绝对计时困难，若已控制好距离、自行和伴星质量，仍能对两星标量响应的差异给出很强的界限。

双脉冲星还可作几个不用另外观测的量纲检查。取 $x_A\simeq1.415$\,s、$x_B\simeq1.516$\,s，则质量比约为 $1.071$。取 $\dot\omega\simeq16.90^\circ/\mathrm{yr}$，用广义相对论进动式定出总质量，再与比值联立，得到 $m_A\simeq1.34M_\odot$、$m_B\simeq1.25M_\odot$。相应 A 发出的脉冲在 B 场中的 Shapiro 尺度应约为 $1.25GM_\odot/c^3\simeq6.2$ 微秒。虽然这个尺度只有微秒，$\sin i$ 接近一时乘上的对数却可显著放大实际上合时延。用质量比而不用 $\dot P_b$ 定质量后，实验还另有 $\gamma_{\rm E}$、Shapiro 形状与幅度、引力辐射阻尼可检验。若实验中的多项量有共用的历表或定时误差，统计上仍应保留协方差，不能把每条交于同一区域的带都当作一次完全独立的显著性。

J1141$-$6545 的轨道并不很圆，近心点每年进动约 $5.31^\circ$。可用 $\gamma_{\rm E}\simeq0.77$\,ms 与 $\dot\omega$ 给出质量，再检验周期衰减。由此得到的脉冲星、白矮星质量约 $1.27M_\odot,1.02M_\odot$，与 J1738+0333 的低质量白矮星明显不同。同为中子星--白矮星，稳定光谱是否可得、轨道是否偏心，决定了用来定质量的两个实验可以完全不同。J0348+0432 的内禀衰减在所引早期资料中只同四极预言符合至约两成，精度低于前两例，却把中子星的质量推到了约 $2M_\odot$。若某模型在较轻星全部未标量化、在很窄的高质量区间突然带上很大标量荷，这条误差较大的约束也能比许多低质量星更有排除力。

并非所有显著的 $\dot P_b^{\rm obs}$ 都适合检验引力。对银河系内 B1534+12，一度限制阻尼精度的是距离所涉的运动学改正；若双星在球状星团中，如 B2127+11C，还可有难以准确估计的当地星团视向加速度。在把视向二阶距离导数换算为周期漂移时，若只知道天空上的投影半径，某个天体位于星团中心的近侧还是远侧尚未确定，视向力甚至有两种符号。按某个星团模型给一个中心改正值却不传播其不确定度，会使所谓检验过于精确。与此形成对照，层级三体 J0337+1715 的外伴星运动在脉冲到时中本就可以跟踪，因而外场强度和方向能由三体整体拟合定出，引力梯度与不等自由下落可同时进入残差模型。

能否把强场检验平面上的窄带直接同太阳系 $\gamma-1$ 界限排在一起？在对数耦合 $\ln A=\alpha_0(\varphi-\varphi_0)+\beta_0(\varphi-\varphi_0)^2/2+\cdots$ 的模型族内，可以选择 $(\alpha_0,\beta_0)$ 为共同坐标。弱场光传播约束 $\gamma-1=-2\alpha_0^2/(1+\alpha_0^2)$，主要排除较大的 $|\alpha_0|$。某个脉冲星的排除线则要对每个参数点解星体结构，求 $\alpha_{\rm NS}(m)$，再对两质量和可能的状态方程处理测量误差。固定 $\beta_0=0$ 得 Brans--Dicke 型截面；若 $\alpha_0$ 很小且两中子星质量相近，两者的偶极差也小。在 $\beta_0$ 足够负的一侧，星体可自发标量化，双星约束便能穿入弱场所容许的区域。把理论平面按两种不同的方向切开，才可理解“太阳系比双星强”或相反说法为何都需要指明比较的是哪一类耦合。

\subsection{波形、传播与外部几何}

若能收到双星本身发出的引力波，便可逐个周期比较实测相位与理论波形。旋近时，最低阶相位给出啁啾质量，高阶项对质量比、自旋的依赖各不相同。用若干项确定源参数后，便可检验其余项的系数。这里为何特别注意相位？设频域中相邻两个波形的振幅相同，只有相位差 $\delta\Psi(f)$。作相关时，各频率的贡献就多一个 $\exp(\mathrm{i}\delta\Psi)$ 的因子。若 $\delta\Psi$ 随频率累积到几弧度，不同频段互相抵消，即使各自的能谱都和信号一样，也不能同实测波形重合。实际以噪声的单边功率谱 $S_n(f)$ 作权，两个波形的内积为 $(h_1|h_2)=4\operatorname{Re}\int_0^\infty\tilde h_1^*\tilde h_2\,\d f/S_n$。一段频带中可数的周期虽多，也必须该处仪器的噪声够小，才能为检验提供信息。

在广义相对论中举一个可算的例子。对准圆、无自旋轨道令 $\eta=\mu/m$，以可测频率构成 $y=(\pi Gmf/c^3)^{1/3}$，其最低阶等于 $v/c$。将轨道能量与辐射能流都写为 $y$ 的函数，再作 $\dot f=-\mathcal P/(\d E/\d f)$，得到
\[
 \dot f=\dot f_{\rm N}\!
 \left[1-\left(\frac{743}{336}+\frac{11}{4}\eta\right)y^2
             +4\pi y^3+O(y^4)\right].
\]
$y^2$ 项中第一次出现独立于啁啾质量的质量比组合。$y^3$ 的 $4\pi$ 项在无自旋时也存在，来自四极波在总质量所生曲率上的散射。直接传播的波和被背景曲率散射的波行时不同，后者使当前的波含有源在较早时刻的贡献，此即辐射的尾项。若有自旋，同一阶还可有自旋与轨道角动量耦合的贡献。一次测准了 $y^3$ 项并不能在未知自旋下独自证明尾项，须同时利用别的阶次或进动造成的调制。

再以 $\dot\Phi=2\pi f$ 从频率演化求相位。在驻相近似下，Fourier 积分在 $f(t)=f$ 的时刻附近相消最少，故取相位约定 $\Psi(f)=2\pi f t(f)-\Phi[t(f)]-\pi/4$。对它一次求导给出 $2\pi t(f)$，二次求导给出 $2\pi/\dot f$。将上式的倒数逐项积分两次，便有
\[
 \Psi(f)=2\pi f t_c-\Phi_c-\frac{\pi}{4}
 +\frac{3}{128\eta y^5}
 \left[1+\left(\frac{3715}{756}+\frac{55}{9}\eta\right)y^2
             -16\pi y^3+O(y^4)\right].
\]
$t_c,\Phi_c$ 是时刻和初相的积分常数。最低阶的 $\eta y^5$ 只依赖 $\mathcal M$；第二、第三项则给出可相互校验的不同频率幂。一个办法是在某个系数前引入 $1+\delta_k$，和质量、自旋一起拟合，然后问 $\delta_k=0$ 是否落在允许区内。若同时把所有系数都当作自由数，有限频带内相邻幂次往往很相似，误差会放大。因而报告界限时应说明只放开了哪一项，哪些系数仍由广义相对论给定。上节的偶极通量若存在，相对领头的四极通量含 $y^{-2}$，于是 $\Psi$ 的首项修正正比于 $f^{-7/3}$。它在低频比普通的质量比修正更明显，可与脉冲星计时所给的偶极限制比较。

对于双黑洞，还可分别分析并合前后的信号。旋近段给出双星的质量和自旋，借助数值相对论可预言遗留黑洞的质量和自旋。若它趋于 Kerr 黑洞，振铃的固有频率与衰减时间也由这两个量决定，因而可用并合后的信号独立核对。更明确地，对足够晚的一段信号可试用 $h(t)=\sum_k A_k e^{-(t-t_0)/\tau_k}\cos[\omega_k(t-t_0)+\phi_k]$。初始振幅随碰撞方式而变，准正规模复频率却由遗留黑洞的 $M_f,J_f$ 决定。在几何制中，任一模都有 $M_f\omega_k=F_k(J_f/M_f^2)$、$\tau_k/M_f=T_k(J_f/M_f^2)$。若能分辨一个模的频率与寿命，原则上可定质量、自旋；再测出另一模就是额外检验。不同模未必都被强烈激发，若用同一个短暂的峰值便同时声称测出了多个模，须先检查资料能否把它们分开。

2015 年的双黑洞事件 GW150914 在地面探测器的敏感频带内仅约 $0.2$ 秒，却已可见先加快、再达到峰值、随后衰减的过程。将信号分为低频和高频两段，可各自估计遗留黑洞的质量和自旋，比较两组允许区域；这一步会用到并合的数值波形。若将来只用晚时振铃作上述多模检验，依赖的主要是单黑洞微扰理论及其边界条件。两种检验对能分辨的信号长度、噪声以及开始分析的时刻要求不同。系统若含中子星，临近并合还可发生潮汐形变或瓦解，不能拿双黑洞的模板残差直接当作引力理论偏离。

\subsection{保守能量、通量与频域模板}\label{sec:pn-wave-energy}

旋近资料能测出的不是某个坐标半径的势能。为避免把坐标选择当成新物理，先用从波测出的 $\Omega=\pi f$ 构成 $x=(Gm\Omega/c^3)^{2/3}=y^2$。对无自旋圆轨道，可用径向运动方程消去坐标半径，将广义相对论的束缚能和能流各写到相对 2PN：
\begin{align*}
 E(x)={}&-\tfrac12\mu c^2 x
 \left[1-\left(\tfrac34+\tfrac{\eta}{12}\right)x
 -\left(\tfrac{27}{8}-\tfrac{19}{8}\eta+\tfrac{\eta^2}{24}\right)x^2+\cdots\right],\\
 {\cal P}(x)={}&\frac{32c^5}{5G}\eta^2 x^5
 \left[1-\left(\tfrac{1247}{336}+\tfrac{35}{12}\eta\right)x
      +4\pi x^{3/2}\right.\\[-2pt]
 &\hspace{35mm}\left.
      +\left(-\tfrac{44711}{9072}+\tfrac{9271}{504}\eta
               +\tfrac{65}{18}\eta^2\right)x^2+\cdots\right].
\end{align*}
作一项简短检查。$\d E/\d x$ 的 1PN 系数是方括号中能量系数的两倍，故 $-\mathcal P/(\d E/\d x)$ 的 1PN 系数等于 $-1247/336-35\eta/12+3/2+\eta/6=-743/336-11\eta/4$，正是前述 $\dot f$ 的系数。若直接把谐和坐标的 $Gm/(rc^2)$ 放在 $x$ 的位置，连最前面的 $x^5$ 也会向一阶系数作贡献。不同规范发表的能流看起来系数不同，须先辨明所展开的变量。

将上式再代入相位的二阶导数，原有 $\Psi$ 的方括号在 $y^4$ 处多出 $15293365/508032+(27145/504)\eta+(3085/72)\eta^2$。这样一套无自旋系数全由两个质量确定，时间零点和相位零点另算。若自旋均平行或反平行于轨道角动量，最低阶的自旋--轨道修正可用
\[
 \beta_{\rm SO}=\sum_{A=1}^2
 \left[\frac{113}{12}\left(\frac{m_A}{m}\right)^2
          +\frac{25}{4}\eta\right]\hat{\bm L}\cdot\bm\chi_A,
 \qquad \bm\chi_A=\frac{c\,\bm J_A}{Gm_A^2}
\]
表示。此时 $y^3$ 的相位系数从 $-16\pi$ 变为 $-16\pi+4\beta_{\rm SO}$。2PN 还会有两自旋相互耦合及依天体自转四极矩而定的项。若自旋倾斜，轨道面会进动，张量波的倾角因子也随时间改变。将所有这类效应都压成一个常数可用于某些短而信噪比较低的信号，绝不是任意取向双星的普适相位公式。

靠近并合的中子星还有一种能帮助区分物态的相位。记 $\Lambda_A=c^{10}\lambda_A/(G^4m_A^5)$ 为无量纲潮汐形变能力，其中 $\lambda_A$ 用前节的 $Q_{ij}=-\lambda_A\mathcal E_{ij}$ 约定。两颗非自旋星的领头组合是
\[
 \widetilde\Lambda=\frac{16}{13m^5}
 \left[(m_1+12m_2)m_1^4\Lambda_1+
       (m_2+12m_1)m_2^4\Lambda_2\right].
\]
在准圆绝热近似下，频域相位方括号含 $-(39/2)\widetilde\Lambda y^{10}$。这里出现十次幂，是由于诱导潮汐与 Newton 引力相比小 $O[(R/a)^5]$；精确数因子同时包括束缚能和辐射四极矩所受的修正。低频时这项远小于点质量项，高频时相位的灵敏度才可能使它显现。若只报末端相位比点质量预报提前，却未把潮汐、自旋和有限带宽同时拟合，就不能把差额唯一归为引力场方程的修正。

\subsection{信号和噪声在何处可分}

设本地记录为 $d=h(\bm\theta)+n$，其中噪声在分析小段内平稳、近似 Gaussian，只有其功率谱已知。对参数 $\bm\theta$ 的对数似然等于 $-(d-h|d-h)/2$ 加一常数。若模板振幅可自由缩放，对它求极值，得标准化的相关 $(d|h)/\sqrt{(h|h)}$；对已知真实波形，其期望的最优信噪比满足 $\rho^2=(h|h)$。因此多个噪声独立的探测器通常给 $\rho_{\rm net}^2=\sum_I\rho_I^2$，但各台模板相位和到达时刻要按同一源方向而关联，不能各自从噪声中挑一个最高峰便视作同一事件。噪声有突发非 Gaussian 成分时，经验的背景和台站间符合检验仍不可少。

对准圆旋近的实信号 $h(t)\sim H(t)\cos\Phi(t)$，驻相点附近把相位保留到二次。Gaussian 型积分的振幅正比于 $H[t(f)]/\sqrt{\dot f}$。由 $H\propto\mathcal M^{5/3}f^{2/3}/D$ 和 $\dot f\propto\mathcal M^{5/3} f^{11/3}$ 得 $|\tilde h|\propto\mathcal M^{5/6}f^{-7/6}/D$。于是每对数频带对 $\rho^2$ 的贡献正比于 $f^{-4/3}/S_n(f)$。在理想白噪声仪器中低频段尤其重要，实际地面台低频噪声增长很快；若只按总周期数而不乘噪声权重，会高估对低频偶极的灵敏度。质量大的双星虽然振幅大，却更早离开慢速旋近段，高质量源适合检验并合，并不必然最适合测每一个 PN 系数。

在拟合的极大值附近展开 $\delta h\simeq\delta\theta^a\p_a h$，似然的二次项给 $\Gamma_{ab}=(\p_a h|\p_b h)$。若信噪比够高、模板光滑、后验在参数边界内近似单峰，协方差近似为 $\Gamma^{-1}$。这个式子说明某个系数“可以精确测量”的条件。若两个导数只差一个常数倍，$\Gamma$ 便有近零特征值，再大一个相同形状的信号也难以区分这两个参数。例如频域相位中一个任意的 $f$ 项可完全吸收到到达时刻，一个常数项可吸收到初相；领头 $f^{-5/3}$ 项若对每一源完全自由，又可重新定义啁啾质量。应检查扣除已知参数张成的子空间后，新项还剩多少独立信息。低信噪比和明显不对称的误差区需直接计算似然而非将上述局部椭球硬作置信区。

可对标量偶极作一次小量计算。若耦合小到整个观测频带内均有 $b y^{-2}\ll1$，写圆轨道的 $\dot f=\dot f_{\rm N}(1+b y^{-2}+\cdots)$，$b$ 在固定归一化下正比于 $(\alpha_1-\alpha_2)^2$。将倒数展开，$\Psi''$ 中的修正正比于 $-b f^{-13/3}$。对频率两次积分，得到 $\Psi$ 的领头方括号多出 $-(4/7)b y^{-2}$。$4/7$ 是两次幂函数积分的系数比，不是一个新的耦合。若耦合大到低频已由偶极主导，则这种先围绕四极展开的小参数近似失效，应直接由两项通量之和积分。另一方面，色散也会改变 $\Psi''$，但它随传播距离而累积，且下文有质量传播给出的幂次为 $f^{-1}$。从一颗低信噪比源拟合多种幂次容易退化；比较不同距离、不同紧致度和不同低频截止的若干源，才可能把产生和传播区分开来。

在理论波形已知的频段内，可以依次放开 PN 系数检验 $\delta_k=0$；若要检验没有完整并合波形的理论，也可先把观测波形写为 $\tilde h_{\rm GR}(1+a y^p)\exp(\mathrm i b y^q)$。每一个 $a,b,p,q$ 的界限都以该形状假定为前提。固定某个辐射机制会在振幅、相位、偏振间给出关系，比四个数都自由变化更容易受约束，但不能把只在旋近弱耦合区成立的幂律任意延伸到振铃。若一套理论只给出了色散、尚无与之对应的并合解，报告传播假定下的界限比把它说成该理论的全程波形检验更明确。

\subsection{群速度、色散与发射时差}

传播速度可借同一源的两种信号检验。2017 年的双中子星事件 GW170817 中，伽马暴的起始时刻比引力波并合时刻晚 $1.74\pm0.05$\,s，而源距约 $40$\,Mpc。设距离为 $D$、引力波速度为 $c_g$。在小红移近似下，到达时差包含源处发射时差和传播时差 $D(1/c-1/c_g)$。若伽马暴在并合后 $0$ 至 $10$ 秒内发出，采用距离估计的保守下限，所得双侧界为 $-3\times10^{-15}<(c_g-c)/c<7\times10^{-16}$。之所以界限不对称，正是因为假定的发射延迟取值不对称。若伽马暴本来较晚发出，抵达时却只相差一点，引力波本身便可以比光稍慢。把两者的测得时差直接除以光行时而不给发射模型，只能作数量级估计。

若源不发可见电磁信号，仍可考察一个啁啾内部各频率的相对到达时间。设传播满足 $E^2=p^2c^2+m_g^2c^4$，且 $m_gc^2\ll h_{\rm P} f$，其中 $h_{\rm P}$ 为 Planck 常数。波包的群速为 $v_g=\d E/\d p=pc^2/E$，故 $v_g/c\simeq1-c^2/(2\lambda_g^2 f^2)$，$\lambda_g=h_{\rm P}/(m_gc)$。近邻源不同频率的到达时差除发射时差以外，又有 $\Delta t_{\rm disp} \simeq \frac{Dc}{2\lambda_g^2} \left(\frac{1}{f_1^2}-\frac{1}{f_2^2}\right)$。低频的波较晚到。旋近时本来是低频先发，色散会压缩从低频到高频的观测间隔。由 $\d\Psi/\d f=2\pi t(f)$ 可见这种传播给频域相位添一个正比于 $f^{-1}$ 的项；它与源处偶极能流所给的 $f^{-7/3}$ 不同。对遥远源，传播途中物理频率也随宇宙膨胀而变，式中的 $D$ 要换成沿路径带有相应红移权重的积分，不能直接代入光度距离。

这类实验怎样与场方程相接？若假定只改变传播关系，可以把一个色散系数加入广义相对论源波形中拟合。2016 年发表的 GW150914 分析按此假设给出 $\lambda_g$ 约大于 $10^{13}$\,km 的界限。若采用完整的有质量引力理论，则静态势、双星能流、偏振内容也可能变化，要用该理论预言的整段波形重新比较。在带有优选系的矢量理论中，传播速度还可随本征模式而异。一次将到达的主导张量波同伽马暴相比较，直接限制的是该传播路径和频段内的张量模速度；对未被源有效激发的另一种模式，还须另找实验。

若源远到传播频率随宇宙膨胀显著改变，低红移式的哪一项需要换掉？在平坦的 Friedmann 背景中归一化 $a(t_0)=1$，频率为 $f$ 的波在红移 $z'$ 处的物理频率为 $(1+z')f$。相同共动距离满足 $\int v_g\,\d t/a(t)=\text{常数}$。将 $1-v_g/c$ 的领头色散项沿未扰动光路积分，并用 $\d t/a=-\d z'/H(z')$，可得同发射事件中两个频率分量的观测到达时差
\[
 \Delta t_{\rm disp}
  =\frac{c^2}{2\lambda_g^2}
     \left(\frac1{f_1^2}-\frac1{f_2^2}\right)
     \int_0^{z_s}\frac{\d z'}{(1+z')^2H(z')}.
\]
若源先后发出，还要加 $(1+z_s)\Delta t_{\rm em}$。低红移时积分近似为 $D/c$，正好还原前式。积分中多出的 $(1+z')^{-2}$ 抑制高红移贡献，原因是同一分量在早期频率较高，传播速度更接近 $c$。观者还可用波形振幅估计光度距离，或用宿主星系的电磁红移估计 $z_s$；两者不是上式的色散积分本身。若混用，在低红移例子中不显著，到较远双星却会系统地改变质量界限。

这种推导也适用于更一般的传播关系。例如在一套具体单位下写 $E^2=p^2c^2+m_g^2c^4+A p^\alpha$。只要最后两项对 $p^2c^2$ 都小，群速的质量修正随 $p^{-2}$，附加项的修正则正比于 $(\alpha-1)A p^{\alpha-2}$。$\alpha=2$ 在此阶给频率无关的速度改动，用单一短暂的引力波信号常会吸收到源的到达时刻；与同源电磁辐射或已知天空位置的多站时差相比却可显现。$\alpha=1$ 的附加项在能量中给常数，不能从该阶群速度差直接量出。若同时让源相位任意逐频改变，波形内部也无法唯一分辨是传播改变还是源先后多发了一些相位，因此要说明对产生机制作了哪些假定。

\subsection{沿黑洞附近的轨道测几何}

轨道本身也可用来测黑洞的外部几何。在几何制中，中性的 Kerr 黑洞满足 $M_2=-J^2/M$。若用黑洞附近多颗恒星的轨道分别测定质量 $M$、角动量 $J$ 及四极矩 $M_2$，这三个数便应满足上述关系。难处在于周围星团的摄动也会使轨道进动，须从测得的进动中分离出来。

先看一条可直接分辨的恒星轨道。连续成像给天球上的角半长轴 $\theta_a$、椭圆的投影和周期。若黑洞到地球距离为 $D$，由天体测量得的是 $M/D^3=4\pi^2\theta_a^3/(GP_b^2)$，单靠改进测角无法分开 $D$ 与 $M$。再测恒星的视向谱线速度，就引入了一个线速度尺度，结合轨道倾角可解出物理半长轴和距离。银河中心恒星 S2 的周期约为 16 年，这种联合轨道测量正好给出一个例子。近心点处 $v$ 和 $GM/r$ 同时达到大值，谱线除一阶视向 Doppler 效应，还含 $v^2/(2c^2)+GM/(rc^2)$。用靠近中心的一次观测约束组合还不够，须有远离中心时的轨道与谱线作为基准。

若进一步比较每转的进动，记 $p=a(1-e^2)$，$\chi=cJ/(GM^2)$。球对称质量引起的近心点进动正比于 $GM/(pc^2)$，并不转动轨道面。自旋拖曳可转动轨道面，每转量级为 $\chi[GM/(pc^2)]^{3/2}$。若四极矩服从 Kerr 关系，它的贡献又约为 $\chi^2[GM/(pc^2)]^2$，并依赖轨道法线和自旋轴的夹角。于是第二条方向不同、半通径不同的轨道特别有用：两轨道所见的是同一黑洞的质量和自旋方向，拖曳及四极矩的投影却不同。若周围有恒星或气体，其潮汐摄动一般也随轨道距离和方向而变，须将外部质量分布与黑洞多极矩同时估计。

极端质量比双星提供了不同的测轨办法。质量很小的紧致星绕大黑洞运行时，第一近似下沿背景测地线运动。一般偏心、有倾角的 Kerr 束缚轨道有径向振荡、纬向振荡和方位绕转三种基本频率。一个轨道产生的辐射谱就可含 $k\Omega_r+l\Omega_\theta+m\Omega_\phi$ 的许多组合，其中 $k,l,m$ 为整数。由一连串谱线可推回基本频率，比单独测一个接近圆轨道的周期多了两种几何信息。随能量、角动量被辐射带走，轨道缓慢穿过不同半径，同一套质量、自旋和多极矩还必须预言频率怎样共同改变。计算这种长时相位既须精确的背景几何，又须算小天体自身场对轨道的摄动。忽视后者积累的相移，会把轨道动力学的误差错认成中心黑洞的额外四极矩。

吸积盘的谱线给出的信息又不同。静止的发射体本有确定谱线能量，盘上不同半径、不同方位的发射体却有不同引力红移和视向速度。到达观者的谱是所有可见发射面经过偏折、红移和光行时效应后的叠加。一条谱线若有长红翼，可以检验发射区是否深入强引力场；但盘倾角、内缘半径、激发铁谱线的照明分布也会改变谱形。准周期振荡若来自轨道本征频率，可给出另一个尺度，前提是已正确识别振荡模式。恒星轨道、波形与光谱各自从不同的物理量求几何，联合使用时应分别保留它们依赖的天体模型。

可用一组真实轨道估计所需精度。下列数值来自早期对银河中心恒星的 Kepler 轨道拟合，只用作比较轨道尺度；近日点距一列由前两列作 $r_p=a(1-e)$ 所得。
\begin{center}
\begin{tabular}{lrrrr}
\hline
恒星 & $a$/AU & $e$ & $P_b$/年 & $r_p$/AU\\ \hline
S2   & $1020$ & $0.880$ & $15.8$ & $122$\\
S4   & $2470$ & $0.406$ & $59.5$ & $1467$\\
S14  & $2125$ & $0.963$ & $47.3$ & $79$\\
S102 & $812$  & $0.680$ & $11.5$ & $260$\\ \hline
\end{tabular}
\end{center}
S14 的周期比 S2 长约三倍，近心点距离却更短，故一个周期内最强的相对论效应可能更大。若只按周期短选目标容易漏掉高离心率轨道。另一方面，轨道近心点处速度大，停留时间短，若实际观测恰错过这段，名义上很强的信号也未化成很强的测量。对参考质量 $4.4\times10^6M_\odot$，$GM/c^2$ 约为 $0.043$\,AU。S2 的 $p\simeq230$\,AU，故球对称近心点进动约为 $6\pi GM/(pc^2)\simeq3.5\times10^{-3}$ 弧度每转。紧邻一个质量很大的黑洞，并不意味着这条恒星轨道已经位于视界附近。

怎样从轨道的天球投影分离自旋和四极矩？取 $\hat{\bm h}$ 沿轨道角动量，自旋轴为 $\hat{\bm J}$，夹角为 $\alpha$。在轨道面内从升交点方向量自旋投影的方位角 $\beta$，轨道面对天球面的倾角为 $i$。定义一转的自旋与四极矩修正尺度 $A_J=4\pi J/\sqrt{Mp^3}$、$A_Q=3\pi M_2/(Mp^2)$（几何制）。圆环平均的轨道力矩使法线绕 $\hat{\bm J}$ 进动，由球面几何可得
\[
 \Delta i=\sin\alpha\cos\beta\,(A_J-A_Q\cos\alpha),\qquad
 \sin i\,\Delta\Omega_{\rm node}
 =\sin\alpha\sin\beta\,(A_J-A_Q\cos\alpha).
\]
两个变化量的比值先确定 $\tan\beta$，且共同因子消失。若两颗星的轨道面不同，两套角度关系可用来定位黑洞的自旋方向，再利用它们各自的 $p,\cos\alpha$ 分开 $J$ 和 $M_2$。若所有星都在同一平面或其轨道面正好使某个投影为零，再高的角分辨率也可能只给某个组合。用距离 $D$ 观测时，每转在天球上的位移量级为 $(a/D)\Delta i$，故须同时权衡近心距离、角半长轴和几年内可见多少轨道弧段。

电磁成像探测的又是射线。对上一节的无自转外度规，光子圆轨道 $r_{\rm ph}$ 确定临界冲量参数 $b_{\rm ph}=r_{\rm ph}/\sqrt{A(r_{\rm ph})}$。在远处看来，临界曲线的角尺度约为 $b_{\rm ph}/D$。基于很远处恒星轨道所给的 $M/D$，即可比较无量纲的 $b_{\rm ph}/M$。由图像直接测到的往往是某种发光等离子体环的亮度峰或光线吸收后的暗区，峰位还随自转、发射分布、吸收及其时间变化而动。因而“量出一个环的直径”与“已测出多极矩”之间还隔着光线追迹和辐射传输。两种黑洞外场如果给出近似相同的阴影角直径，也仍可在不同倾角轨道的进动或准正规模寿命上分开。

\subsection{用宇宙演化比较引力理论}

早期宇宙的轻元素丰度取决于核反应速率与膨胀率之比。核反应速率可由实验确定，故氦、氘等的丰度能检验引力场方程预言的早期膨胀率。举例说，中子与质子间的弱相互作用在高温下维持热平衡。温度下降，反应速率 $\Gamma_{\rm w}$ 降至膨胀率 $H$ 以下，数目比例便逐渐偏离平衡值。若相同温度的宇宙膨胀更快，冻结会更早发生。此时中子与质子质量差所受的 Boltzmann 压低较弱，又较少中子能在形成原子核前衰变，因而通常会形成更多的 $^4\mathrm{He}$。辐射占优时，广义相对论给 $H^2=8\pi G\rho_{\rm rad}/3$，这里 $\rho_{\rm rad}$ 是能量密度除以 $c^2$。若把修改引力暂近似为一个有效引力常数，改变 $G$ 与改变辐射密度都能改变 $H$。故须带入相对论性粒子数目、重子丰度和反应速率的信息，不能仅从一种元素读出当时的 $G$。

较晚时期则可将距离与红移给出的膨胀史，同引力透镜和物质聚集给出的结构增长史相比较。两种引力理论可以有相同的尺度因子 $a(t)$，只要调节均匀物质能量的演化即可。这样按标准烛光或标准尺度量出的距离可一样，密度扰动却不必相同。在线性标量扰动下取 $\d s^2=-(1+2\Psi)c^2\d t^2+a^2(1-2\Phi)\d\bm x^2$，其中 $\Psi,\Phi$ 是相对此宇宙学背景定义的无量纲势。非相对论物质受 $\bm\nabla\Psi$ 加速，光的偏折则由 $\bm\nabla_\perp(\Psi+\Phi)$ 沿视线的积分决定。因此星系聚集与视向速度主要检验物质如何响应 $\Psi$，弱透镜还会看到 $\Phi$。将两类观测分红移、尺度比较，可检验在给定物质应力和引力定律下两势的关系。这里的背景自身正在演化，并非太阳系的渐近平直背景，不能把两个宇宙学势的比值直接等同于标准 PPN 参数 $\gamma$。
核合成的灵敏度还可作一个数量估计。若在冻结的短时间内有效相对论自由度不变，辐射时代 $\rho_{\rm rad}\propto T^4$，故 $H\propto G_{\rm eff}^{1/2}T^2$。弱相互作用维持中子与质子转换的速率在这一阶段约随 $T^5$ 变化，令 $\Gamma_{\rm w}(T_f)=H(T_f)$ 即得 $T_f\propto G_{\rm eff}^{1/6}$。冻结前 $n/p$ 大致为 $\exp[-(m_n-m_p)c^2/(k_{\rm B}T_f)]$。升高 $G_{\rm eff}$ 时，$T_f$ 升高而指数压低减轻，冻结后到氘能稳定生存前的时间又缩短，两项效应都增加可进入氦核的中子。若最后几乎每个幸存中子都在 $^4$He 中，近似氦质量分数 $Y_p=2(n/p)/[1+(n/p)]$。氘丰度则还对重子密度和燃烧反应率敏感。用这些关系可解释不同轻元素提供的约束方向为何不同，作精确比较仍要联立完整核反应网络，不能把粗略的 $1/6$ 幂当作元素丰度的精确灵敏度。

用今天的实验室 $G$ 或 $\dot G/G$ 对这一时代作线性外推一般靠不住。若引力中有标量，其宇宙学背景 $\varphi(t)$ 所受的物质源正比于 $T=-\varepsilon+3p$；理想辐射满足 $p=\varepsilon/3$，普通非相对论物质却有 $T\simeq-\varepsilon$。从辐射主导转为物质主导，标量背景的演化方程即使耦合函数一成不变也可改变。再靠近粒子变为非相对论性或正负粒子湮灭的时期，混合气体的迹未必为零。因此某些模型早期偏离较大、晚期却演化到弱耦合点。核合成观测、背景微波辐射和太阳系实验定的是几个不同时空环境下的物理量，要在一条明确给出的背景解上比较。

较晚时期如何实际组合观测？在尺度远短于 Hubble 半径、普通冷物质的压强可略、星系上小尺度非线性演化也可暂略时，把密度写为 $\rho_m=\bar\rho_m(1+\delta)$。线性连续性方程与 Euler 方程联立，得 $\ddot\delta+2H\dot\delta=c^2\nabla^2\Psi/a^2$。若将静力近似下的场方程参数化为
\[
 -k^2\Psi=\frac{4\pi G a^2\bar\rho_m}{c^2}\,
               \mu(k,a)\delta,\qquad
 \Phi=\eta_{\rm slip}(k,a)\Psi ,
\]
则增长方程变为 $\ddot\delta+2H\dot\delta=4\pi G\bar\rho_m\mu\delta$。弱透镜测量的 $(\Phi+\Psi)$ 主要限制 $\Sigma=\mu(1+\eta_{\rm slip})/2$。举例说，某种模型增加对慢粒子的吸引，使星系聚集变快；若又减小 $\Phi/\Psi$，透镜强度未必按同一倍数增大。将背景距离测出的 $H(a)$ 先放入增长方程，再交叉比较星系速度和透镜，可把膨胀阻尼和有效引力的变化分开。

在这个写法中，$k$ 为共动波数，$\mu,\eta_{\rm slip}$ 是待检验的函数而非太阳系 PPN 常数。即使在广义相对论中，自由传播粒子的各向异性应力也能使 $\Phi$、$\Psi$ 不等；星系数密度同物质密度之间还可能有随尺度变化的偏差。可利用星系红移空间畸变估计速度散度，以星系与背景星系剪切的交叉相关约束透镜势，再改变红移分层考察随时间的变化。若进入接近 Hubble 半径的尺度，时间导数、辐射和规范选择中的扰动均不可随意丢弃；两函数的准静态近似应换回具体理论完整的线性扰动方程。这样作比较，实验给出的才是哪些距离尺度、哪一时代的引力发生了可分辨的偏离。
